% ─────────────────────────────────────────────────────────────────────
% TWIN-LOCATION NOTICE. This file is one of the shared homes for the same
% results; edit every listed home or the drift checker fails.
%   Standalone papers : docs/harbor-research/tex/paper1.tex (S2-S5)
%                        docs/harbor-research/tex/paper6.tex (S(sole-roles))
%   Chapter           : book/legible-swarm.tex (this file)
%   Index ids         : R1, R2, R3, R4, R15
%   Check             : python3 scripts/harbor-research/check_library_index.py
%                        python3 scripts/harbor-research/check_propagated_corrections.py
%                        python3 scripts/harbor-research/check_citations.py
%   System document   : docs/harbor-research/LIBRARY-SYSTEM.md
% ─────────────────────────────────────────────────────────────────────
\documentclass[11pt,a4paper]{article}

% ─── Packages ────────────────────────────────────────────────────────────────
\usepackage[utf8]{inputenc}
\usepackage[T1]{fontenc}
\usepackage{lmodern}
\usepackage{microtype}
\usepackage{geometry}
\usepackage{amsmath,amssymb,amsthm}
\usepackage{listings}
\usepackage{xcolor}
\usepackage{booktabs}
% For \needspace at the two subsections that were landing in the last
% two lines of a page. The Book's own preamble already loads it; the
% generator strips this block, so adding it here only reaches the
% standalone twin, which would otherwise fail on an undefined command.
\usepackage{needspace}
\usepackage{tabularx}
\usepackage{xltabular}
\usepackage{graphicx}
\usepackage{fancyhdr}
\usepackage{tikz}
\usepackage{pgfplots}
\pgfplotsset{compat=1.18}
\usetikzlibrary{arrows.meta,positioning,fit,backgrounds,calc,shapes.geometric,shapes.symbols,decorations.pathmorphing,decorations.pathreplacing}
\usepackage{float}
\usepackage{caption}
\usepackage{enumitem}
\usepackage[backref=page]{hyperref}
\captionsetup{font=small,labelfont=bf}

% ─── Geometry ────────────────────────────────────────────────────────────────
\geometry{margin=2.5cm}

% ─── Colors (the hh-* harbor palette, verbatim from house style) ──────────────
\definecolor{codebg}{RGB}{248,248,248}
\definecolor{codeframe}{RGB}{200,200,200}
\definecolor{hhsand}{HTML}{E9DCC4}
\definecolor{hhsanddeep}{HTML}{D8C7A6}
\definecolor{hhebony}{HTML}{121212}
\definecolor{hhink}{HTML}{1B1712}
\definecolor{hhcobalt}{HTML}{003FB8}
\definecolor{hhamber}{HTML}{6B4500}
\definecolor{hhteal}{HTML}{00564C}
\definecolor{hhpaper}{HTML}{FBF7EF}
\definecolor{hhgray}{HTML}{5C5650}
% The Book's semantic color system: palette v2, light values. One hue, one meaning.
% Source of record for every hex: website-v2/src/styles/tokens.semantic.css (light
% theme); website-v2/scripts/check-figure-palette.mjs fails the build if this file
% and the tokens disagree. Twin copy under book/figures/ (byte-identical).
%
% Separation: contrast says an ink is READABLE on the cream; it says nothing about
% whether two inks are TELLABLE APART, and until now only the first was checked.
% check-figure-palette-separation.mjs measures all 21 pairs in OKLab with protanopia
% and deuteranopia simulated. indigo and violet are re-stepped (deeper, on their own
% hues, both gaining contrast) because they were the pairs with room: cobalt/indigo
% was dE 9.1 and is 15.1; cobalt/violet collapsed to 4.8 under protanopia and is 8.6.
% Three pairs CANNOT be fixed -- teal/health, gold/health and teal/gold have ceilings
% of 12.2, 13.9 and 14.3 against a floor of 15, because sRGB has no chroma to spend
% on the 112-182 degree arc this dark. Only ONE of gold, health and teal may carry a
% series by colour alone; the others need a label, a gap, or a texture.
%
% Discipline: story colors are accents (rules, fills, marks, display text); body
% pages stay cream and ink. Amber is stripes, dots, and display only, never small
% text (3.71:1 on cream). Violet and gold are AA as text; use them at display size
% or on rules and fills. Line weights: 1px texture, 1.5px linework, 2px enclosure.
\definecolor{pdcobalt}{HTML}{003FB8}      % kernel / truth            --brand-primary
\definecolor{pdteal}{HTML}{006B5F}        % legibility                --brand-accent
\definecolor{pdhealth}{HTML}{1F7A4D}      % ready / coordinated       --story-health
\definecolor{pdindigo}{HTML}{312865}      % protocol / federation     --story-indigo
\definecolor{pdviolet}{HTML}{822586}      % identity / continuity     --story-violet
\definecolor{pdrust}{HTML}{7A4514}        % reputation / trust earned --story-rust
\definecolor{pdgold}{HTML}{666A00}        % economy / value           --story-gold
\definecolor{pdslate}{HTML}{384856}       % foundation / bedrock      --story-slate
\definecolor{pderror}{HTML}{BF2F2F}       % breach / correction       --status-error
\definecolor{pdamber}{HTML}{A66F00}       % warning (display only)    --status-warning
\definecolor{pdlime}{HTML}{CAD900}        % highlight fill, ink text  --chart-yellow
\definecolor{pdink}{HTML}{121212}         % ink                       --text-primary
\definecolor{pdinkmuted}{HTML}{403B34}    % links, secondary text     --text-secondary
\definecolor{pdcream}{HTML}{F2EEE6}       % page ground               --surface-base
\definecolor{pdcreamraised}{HTML}{F7F3EB} % panels                    --surface-raised
\definecolor{pdcreamstrong}{HTML}{E9E2D5} % inset wells               --surface-strong
% Print inks for the Swiss edition of the Book (SWISS-BRIEF, B); print only.
\definecolor{pdswissblue}{HTML}{001489}
\definecolor{pdswissviolet}{HTML}{582C83}
\definecolor{pdswissred}{HTML}{DA291C}
% And the maritime edition's fourth ink: antique gold, the trim of the engraved
% plates carried onto the page (opener rules, part-title ornament, fillets round
% framed matter). Flat stand-in for a metallic pass; print only.
\definecolor{pdmaritimegold}{HTML}{805A14}

\newcommand{\pdchapterprefix}{ls}% this chapter's key in book/textbook.json
% GENERATED FILE. Do not edit by hand.
% Source of record: book/textbook.json (chapter order, numbers, parts, edition).
% Regenerate canonical map and site data: node scripts/generate-mega-whitepaper.mjs --sync-shared
% Drift check (CI): node --test scripts/generate-mega-whitepaper.test.mjs
% Canonical map: book/figures/pd-textbook-map.tex.
%
% Every standalone chapter preamble defines \pdchapterprefix before inputting
% this file; the Book redefines it at every \pdchapter. Everything else derives.

\providecommand{\pdeditiontitle}{Governing AI Agents: The Harbor, the Person, and the Economy}
\providecommand{\pdeditionsubtitle}{A Textbook of Concurrency, Identity, and Accountable Autonomous Work}
\providecommand{\pdeditionversion}{4.0 (Textbook Edition)}
\providecommand{\pdeditiondate}{5 November 2026}
\providecommand{\pdeditionclaim}{For autonomous work to remain accountable, each action needs an authority, a record, and someone responsible for its consequences.}
\providecommand{\pdbooktitle}{\emph{\pdeditiontitle}}
\providecommand{\pdchaptercount}{8}
\providecommand{\pdchaptercountword}{eight}
\providecommand{\pdpartcount}{4}
\providecommand{\pdchapterprefix}{none}
\expandafter\providecommand\csname pdchapternumberofnone\endcsname{0}
\expandafter\providecommand\csname pdchaptertitleofnone\endcsname{\pdeditiontitle}
\expandafter\providecommand\csname pdchaptercolorofnone\endcsname{pdcobalt}
\expandafter\providecommand\csname pdchapternumberofswk\endcsname{1}
\expandafter\providecommand\csname pdchaptertitleofswk\endcsname{The Single-Writer Kernel}
\expandafter\providecommand\csname pdchaptercolorofswk\endcsname{pdcobalt}
\expandafter\providecommand\csname pdchapterformerofswk\endcsname{II}
\expandafter\providecommand\csname pdchapterpartindexofswk\endcsname{1}
\expandafter\providecommand\csname pdchapterpartlabelofswk\endcsname{Part I \textperiodcentered\ Ground Truth}
\expandafter\providecommand\csname pdchapterquestionofswk\endcsname{Where can a rule be made real?}
\expandafter\providecommand\csname pdchapteronelineofswk\endcsname{A single program makes every change to a local database, one at a time; the chapter states which failures that record survives.}
\expandafter\providecommand\csname pdchapterepigraphofswk\endcsname{A transaction commits at the instant its commit record appears on disk.}
\expandafter\providecommand\csname pdchapterepigraphsourceofswk\endcsname{Jim Gray et al., The Recovery Manager of the System R Database Manager, 1981, p. 232}
\expandafter\providecommand\csname pdchapternumberofanchor\endcsname{2}
\expandafter\providecommand\csname pdchaptertitleofanchor\endcsname{The Anchor Protocol}
\expandafter\providecommand\csname pdchaptercolorofanchor\endcsname{pdcobalt}
\expandafter\providecommand\csname pdchapterformerofanchor\endcsname{V}
\expandafter\providecommand\csname pdchapterpartindexofanchor\endcsname{1}
\expandafter\providecommand\csname pdchapterpartlabelofanchor\endcsname{Part I \textperiodcentered\ Ground Truth}
\expandafter\providecommand\csname pdchapterquestionofanchor\endcsname{How may authority travel without silently growing?}
\expandafter\providecommand\csname pdchapteronelineofanchor\endcsname{When an agent passes a permission on, the new permission can only be narrower than the one it came from, at every step, and the checker verifies the whole chain.}
\expandafter\providecommand\csname pdchapterepigraphofanchor\endcsname{Every program and every user of the system should operate using the least set of privileges necessary to complete the job.}
\expandafter\providecommand\csname pdchapterepigraphsourceofanchor\endcsname{Jerome Saltzer and Michael Schroeder, The Protection of Information in Computer Systems, 1975}
\expandafter\providecommand\csname pdchapternumberofsealed\endcsname{3}
\expandafter\providecommand\csname pdchaptertitleofsealed\endcsname{The Sealed Harbor}
\expandafter\providecommand\csname pdchaptercolorofsealed\endcsname{pdcobalt}
\expandafter\providecommand\csname pdchapterformerofsealed\endcsname{}
\expandafter\providecommand\csname pdchapterpartindexofsealed\endcsname{1}
\expandafter\providecommand\csname pdchapterpartlabelofsealed\endcsname{Part I \textperiodcentered\ Ground Truth}
\expandafter\providecommand\csname pdchapterquestionofsealed\endcsname{What can a sealed agent leak, and how would we know?}
\expandafter\providecommand\csname pdchapteronelineofsealed\endcsname{A job run in a sealed environment can release results only through named checks, and all its releases together may disclose only a fixed amount of information, under the model's stated assumptions about channels.}
\expandafter\providecommand\csname pdchapterepigraphofsealed\endcsname{This note explores the problem of confining a program during its execution so that it cannot transmit information to any other program except its caller.}
\expandafter\providecommand\csname pdchapterepigraphsourceofsealed\endcsname{Butler W. Lampson, A Note on the Confinement Problem, 1973}
\expandafter\providecommand\csname pdchapternumberofls\endcsname{4}
\expandafter\providecommand\csname pdchaptertitleofls\endcsname{The Legible Swarm}
\expandafter\providecommand\csname pdchaptercolorofls\endcsname{pdteal}
\expandafter\providecommand\csname pdchapterformerofls\endcsname{I}
\expandafter\providecommand\csname pdchapterpartindexofls\endcsname{2}
\expandafter\providecommand\csname pdchapterpartlabelofls\endcsname{Part II \textperiodcentered\ The Cost of Seeing}
\expandafter\providecommand\csname pdchapterquestionofls\endcsname{What must remain visible, and what does seeing cost?}
\expandafter\providecommand\csname pdchapteronelineofls\endcsname{Useful supervision needs summaries that keep what the operator must decide on and link back to evidence the operator can inspect.}
\expandafter\providecommand\csname pdchapterepigraphofls\endcsname{In an information-rich world, the wealth of information means a dearth of something else: a scarcity of whatever it is that information consumes. What information consumes is rather obvious: it consumes the attention of its recipients.}
\expandafter\providecommand\csname pdchapterepigraphsourceofls\endcsname{Herbert A. Simon, Designing Organizations for an Information-Rich World, 1971}
\expandafter\providecommand\csname pdchapternumberofstp\endcsname{5}
\expandafter\providecommand\csname pdchaptertitleofstp\endcsname{From Spawn to Person}
\expandafter\providecommand\csname pdchaptercolorofstp\endcsname{pdviolet}
\expandafter\providecommand\csname pdchapterformerofstp\endcsname{III}
\expandafter\providecommand\csname pdchapterpartindexofstp\endcsname{3}
\expandafter\providecommand\csname pdchapterpartlabelofstp\endcsname{Part III \textperiodcentered\ What Survives the Restart}
\expandafter\providecommand\csname pdchapterquestionofstp\endcsname{What remains long enough to owe anything?}
\expandafter\providecommand\csname pdchapteronelineofstp\endcsname{An AI worker process starts, does a job and exits. Continuity links each replacement process, through a checked succession record, to one continuing account in a shared ledger, which shows what the worker owes, what it may do, which of its results an independent checker confirmed and which earlier processes it descends from. Reputation is what that record of confirmed results is worth to the next buyer.}
\expandafter\providecommand\csname pdchapterepigraphofstp\endcsname{We are but whirlpools in a river of ever-flowing water. We are not stuff that abides, but patterns that perpetuate themselves.}
\expandafter\providecommand\csname pdchapterepigraphsourceofstp\endcsname{Norbert Wiener, The Human Use of Human Beings: Cybernetics and Society, revised edition, 1954}
\expandafter\providecommand\csname pdchapternumberofhe\endcsname{6}
\expandafter\providecommand\csname pdchaptertitleofhe\endcsname{The Harbor Economy}
\expandafter\providecommand\csname pdchaptercolorofhe\endcsname{pdgold}
\expandafter\providecommand\csname pdchapterformerofhe\endcsname{IV}
\expandafter\providecommand\csname pdchapterpartindexofhe\endcsname{4}
\expandafter\providecommand\csname pdchapterpartlabelofhe\endcsname{Part IV \textperiodcentered\ Trade Between Strangers}
\expandafter\providecommand\csname pdchapterquestionofhe\endcsname{What makes cooperation rational without making loss unbounded?}
\expandafter\providecommand\csname pdchapteronelineofhe\endcsname{Labor, rented assets and licensed skills carry different risks, but all of them settle through one shared ledger in which no transfer creates or destroys any unit of payment.}
\expandafter\providecommand\csname pdchapterepigraphofhe\endcsname{The main reason why it is profitable to establish a firm would seem to be that there is a cost of using the price mechanism.}
\expandafter\providecommand\csname pdchapterepigraphsourceofhe\endcsname{Ronald H. Coase, The Nature of the Firm, 1937}
\expandafter\providecommand\csname pdchapternumberofbonded\endcsname{7}
\expandafter\providecommand\csname pdchaptertitleofbonded\endcsname{The Bonded Commons}
\expandafter\providecommand\csname pdchaptercolorofbonded\endcsname{pdgold}
\expandafter\providecommand\csname pdchapterformerofbonded\endcsname{VI}
\expandafter\providecommand\csname pdchapterpartindexofbonded\endcsname{4}
\expandafter\providecommand\csname pdchapterpartlabelofbonded\endcsname{Part IV \textperiodcentered\ Trade Between Strangers}
\expandafter\providecommand\csname pdchapterquestionofbonded\endcsname{Who may decide, on what evidence, with what consequence?}
\expandafter\providecommand\csname pdchapteronelineofbonded\endcsname{Some doubt about an agent's work remains even after access checks. The chapter handles it with shared rules: providers post deposits (bonds), signed records show who did what, audits sample the results, and either side can appeal a verdict. A model checker confirms, on a bounded model, that the deposit ledger never creates or loses money.}
\expandafter\providecommand\csname pdchapterepigraphofbonded\endcsname{The power to constrain an adversary may depend on the power to bind oneself.}
\expandafter\providecommand\csname pdchapterepigraphsourceofbonded\endcsname{Thomas C. Schelling, The Strategy of Conflict, 1960}
\expandafter\providecommand\csname pdchapternumberoffh\endcsname{8}
\expandafter\providecommand\csname pdchaptertitleoffh\endcsname{The Federated Harbor}
\expandafter\providecommand\csname pdchaptercoloroffh\endcsname{pdgold}
\expandafter\providecommand\csname pdchapterformeroffh\endcsname{VII}
\expandafter\providecommand\csname pdchapterpartindexoffh\endcsname{4}
\expandafter\providecommand\csname pdchapterpartlabeloffh\endcsname{Part IV \textperiodcentered\ Trade Between Strangers}
\expandafter\providecommand\csname pdchapterquestionoffh\endcsname{What can cross a trust boundary without moving the authority itself?}
\expandafter\providecommand\csname pdchapteronelineoffh\endcsname{Machines run by owners who do not trust each other can pass each other signed records, but not the power to approve actions on each other's systems. The chapter works out which information can spread from machine to machine, which decisions still need one designated approver, and why the receiving machine must issue its own permission for its own resources instead of narrowing the sender's.}
\expandafter\providecommand\csname pdchapterepigraphoffh\endcsname{We reject: kings, presidents and voting. We believe in: rough consensus and running code.}
\expandafter\providecommand\csname pdchapterepigraphsourceoffh\endcsname{David D. Clark, A Cloudy Crystal Ball: Visions of the Future, 1992}
\providecommand{\pdchapternumber}{\csname pdchapternumberof\pdchapterprefix\endcsname}
\providecommand{\pdchaptertitle}{\csname pdchaptertitleof\pdchapterprefix\endcsname}
\providecommand{\pdchaptercolor}{\csname pdchaptercolorof\pdchapterprefix\endcsname}
\providecommand{\pdchapterpartindex}{\csname pdchapterpartindexof\pdchapterprefix\endcsname}
\providecommand{\pdchapterpartlabel}{\csname pdchapterpartlabelof\pdchapterprefix\endcsname}
\providecommand{\pdchapterquestion}{\csname pdchapterquestionof\pdchapterprefix\endcsname}
\providecommand{\pdchapteroneline}{\csname pdchapteronelineof\pdchapterprefix\endcsname}
\providecommand{\pdchapterepigraph}{\csname pdchapterepigraphof\pdchapterprefix\endcsname}
\providecommand{\pdchapterepigraphsource}{\csname pdchapterepigraphsourceof\pdchapterprefix\endcsname}

% Parts, for the weighted four-hue rule: numeral, hue, the label that ends the
% part (the next part, or the appendices), and the chapter count used as the
% weight until page references exist. \pdweightrule (Book preamble) consumes it.
\providecommand{\pdfirstpartnumeral}{I}
\providecommand{\pdweightsegmentslist}{%
  \pdweightsegment{I}{pdcobalt}{part:II}{3}%
  \pdweightsegment{II}{pdteal}{part:III}{1}%
  \pdweightsegment{III}{pdviolet}{part:IV}{1}%
  \pdweightsegment{IV}{pdgold}{book:appendices}{3}%
}
\expandafter\providecommand\csname pdpartindexofI\endcsname{1}
\expandafter\providecommand\csname pdpartindexofII\endcsname{2}
\expandafter\providecommand\csname pdpartindexofIII\endcsname{3}
\expandafter\providecommand\csname pdpartindexofIV\endcsname{4}
\expandafter\providecommand\csname pdparttitleofI\endcsname{Ground Truth}
\expandafter\providecommand\csname pdparttitleofII\endcsname{The Cost of Seeing}
\expandafter\providecommand\csname pdparttitleofIII\endcsname{What Survives the Restart}
\expandafter\providecommand\csname pdparttitleofIV\endcsname{Trade Between Strangers}

% Locator map: every part and chapter of the Book, the current chapter marked.
% \pdchaplink is provided by figures/pd-hyperlinks.tex (a live link inside the
% Book, plain text in a standalone chapter).
\providecommand{\pdmapchapter}[3]{%
  \ifnum#1=\pdchapternumber\relax
    \textbf{\pdchaplink{#2}{#1.~#3}}~\textnormal{(this chapter)}%
  \else
    \pdchaplink{#2}{#1.~#3}%
  \fi}
\providecommand{\pdtextbookmap}{%
  \begin{tabular}{@{}l@{\quad}p{\dimexpr\linewidth-12em\relax}@{}}
  \textsc{\textbf{Part I}}\ Ground Truth & \raggedright \pdmapchapter{1}{swk}{The Single-Writer Kernel} \textperiodcentered\ \pdmapchapter{2}{anchor}{The Anchor Protocol} \textperiodcentered\ \pdmapchapter{3}{sealed}{The Sealed Harbor}\tabularnewline[2pt]
  \textsc{\textbf{Part II}}\ The Cost of Seeing & \raggedright \pdmapchapter{4}{ls}{The Legible Swarm}\tabularnewline[2pt]
  \textsc{\textbf{Part III}}\ What Survives the Restart & \raggedright \pdmapchapter{5}{stp}{From Spawn to Person}\tabularnewline[2pt]
  \textsc{\textbf{Part IV}}\ Trade Between Strangers & \raggedright \pdmapchapter{6}{he}{The Harbor Economy} \textperiodcentered\ \pdmapchapter{7}{bonded}{The Bonded Commons} \textperiodcentered\ \pdmapchapter{8}{fh}{The Federated Harbor}\tabularnewline[2pt]
  \end{tabular}}

% Shared pedagogic apparatus for the Book and every standalone chapter: boxed
% claims, honest boundaries, worked examples, retrieval prompts, numbered
% exercises, and their deferred solutions.
% Twin copy under book/figures/ (byte-identical; checked by
% scripts/generate-mega-whitepaper.test.mjs).
%
% \input this immediately after figures/pd-textbook-map (so \pdchapterprefix
% and \pdchapternumber already resolve) and before figures/pd-hyperlinks.
%
% Page grammar (Wave 12): a content kind is signalled by typography and the
% margin, never by a tinted rectangle. Claims: a hairline above, a small-caps
% run-in head, a short hairline below. Worked examples: two hairlines and an
% italic head. Boundaries: a plain ink bar at the left, no fill. Recitation:
% margin prompts in the Book, a compact list in a standalone chapter.
% Exercises: a numbered hanging paragraph with kind and stars in small caps,
% the solution page in the margin. Nothing here paints a background.
%
% \ifpdmargincolumn is false by default (standalone A4 chapters have no margin
% column); the Book's preamble sets \pdmargincolumntrue after its geometry.
% mdframed is used only for the bars (leftline) and hairlines (top/bottom),
% with the DEFAULT frame method -- never framemethod=tikz.
%
% Key-value parsing for \pdexercise uses pgfkeys: pgf ships as part of tikz,
% which every preamble already loads, so this adds no new dependency (keyval
% would have meant one). The closed kind/rating vocabularies are enforced with
% the LaTeX kernel's own \in@/\ifin@ comma-list test -- always available,
% no package of its own -- wrapped in an \edef that expands both sides first:
% \in@ does not expand its own arguments, so hunting for a *macro* (rather
% than literal text) requires pre-expanding it, or every kind reads as absent.
\makeatletter
\@ifpackageloaded{mdframed}{}{\usepackage{mdframed}}
\@ifpackageloaded{answers}{}{\usepackage{answers}}
\@ifpackageloaded{marginnote}{}{\usepackage{marginnote}}
\@ifpackageloaded{environ}{}{\usepackage{environ}}
\@ifpackageloaded{etoolbox}{}{\usepackage{etoolbox}}
\@ifpackageloaded{listings}{}{\usepackage{listings}}
\@ifpackageloaded{caption}{}{\usepackage{caption}}
\@ifpackageloaded{tcolorbox}{}{\usepackage[skins,breakable,listings]{tcolorbox}}

% ---------------------------------------------------------------------------
% Shared helpers: the margin-column switch, the margin head, the margin glyph.
% ---------------------------------------------------------------------------
% The Book defines \pdcurrentchaptercolor (per part) AFTER this file; a
% standalone chapter never does. Resolve it lazily, falling back to cobalt.
\newcommand{\pd@chaptercolor}{\ifdefined\pdcurrentchaptercolor\pdcurrentchaptercolor\else pdcobalt\fi}
% The eight standalone chapter PDFs this false branch exists for are retired
% as of this change (the chapters no longer build or publish an A4 edition of
% themselves). Every \ifpdmargincolumn false branch below is now dead weight
% kept alive only by the chapters' .tex sources still being independently
% compilable; do NOT delete it in this change -- the margin apparatus work is
% in flight on another branch and this would collide with it. Removing the
% false branch (and, eventually, the conditional itself) is that wave's job.
\newif\ifpdmargincolumn
\pdmargincolumnfalse
% The Book overrides this with an unnumbered, baseline-aligned margin aside.
% Keep chapter sources independently compilable without duplicating prose.
\providecommand{\pdsourceaside}[1]{\footnote{#1}}
% An intentional body boundary: flush pending figures before a new margin
% discussion. The generator strips legacy literal chapter-opening breaks,
% but preserves this explicit editorial instruction. Never use it in frames.
\newcommand{\pdmarginpagebreak}{\clearpage}
\renewcommand*{\marginfont}{\footnotesize\raggedright\color{pdinkmuted}}
% A small-caps head: set in the margin when the Book has a margin column,
% otherwise run in at the start of the paragraph.
% A head stays a \marginnote, and the reason is worth keeping. It was set as
% a \marginpar for one build (2026-09-08), so the page builder would keep the
% Recall block clear of it, and the build died with "Float(s) lost" at the end
% of chapter 1. \ifinner is true inside a minipage or mdframed only BETWEEN
% paragraphs; the moment a paragraph starts it is ordinary horizontal mode,
% \ifinner is false, and a \marginpar issued there -- a \keyidea inside a
% pdexample -- has no page to be shipped to. The heads are one line each and
% are never what runs into something; the Recall block was, and it is now a
% \marginpar issued in vertical mode, where the test is honest.
% Each head records the page and the page-height at which it was issued, so a
% Recall block raised later on the same page can stop short of it. A head is
% issued at the start of its paragraph, where \pagetotal is the page above the
% paragraph -- the head's own position, to within a line.
% Margin occupancy is tracked both ways. A head records where it was issued
% so a Recall block raised later on the page stops short of it; a Recall block
% records where its foot came to rest so a head issued later on the same page
% (the next section's first Key idea) steps down below it instead of printing
% inside it. \pagetotal at a head lags by the current paragraph, which puts
% every error on the safe side: the block is capped a little low, the head
% steps a little far.
\def\pd@lastheadpage{0}\newlength{\pd@lastheadbottom}
\def\pd@blockpage{0}\newlength{\pd@blockbottom}\newlength{\pd@headshift}
\newsavebox{\pd@headbox}\newsavebox{\pd@ptrbox}\newsavebox{\pd@sidebox}
% Run the page builder so \pagetotal counts the paragraph that just ended.
% Without it \pagetotal is the page above the LAST paragraph the builder saw,
% which under-reads by however much prose has since been set -- and every
% device here decides where it may sit from that number. A Boundary head read
% a stale \pagetotal and recorded its foot a paragraph too high, and the Recall
% block that followed cleared that phantom and printed on the head itself
% (0.1 pt apart on p. 253 of the technical edition). \penalty\@M contributes
% without opening a breakpoint. Outer vertical mode only: inside a box the
% penalty lands on the box's own list and settles nothing.
\newcommand{\pd@settle}{\ifvmode\ifinner\else\penalty\@M\vskip\z@\fi\fi}
% ---------------------------------------------------------------------------
% Figure-band tracking, for vertical CENTERING of a margin device against the
% main-column figure it accompanies (2026-09). \pd@figtop/\pd@figbot/\pd@figpage
% are declared here but SET in the Book preamble's \pgfsys@typesetpicturebox
% hook (coordination-papers-mega-volume-preamble.tex), not by an
% \AtBeginEnvironment{figure}/\AtEndEnvironment{figure} pair on \pagetotal --
% that was the first attempt and it does not work: for a \begin{figure}[H]
% float, the `float' package defers inserting the box into the main vertical
% list past \end{figure} itself, so \pagetotal read there under-reads by the
% whole figure and \pd@figbot always came out equal to \pd@figtop (measured:
% every firing in the built Book). The pgf hook that already intercepts every
% TikZ picture on its way into the page knows the box's own \ht+\dp directly,
% which is exact regardless of page-builder timing, so that is where the span
% is recorded now. \pd@figpage guards against reading a stale span left over
% from a figure several pages back: \pd@placemargin only trusts the span when
% \pd@figpage still equals the numeric page counter. Only one span is kept -- the most
% recently drawn picture on this page -- which is deliberately the figure a
% margin device issued shortly afterward is almost always about; a page with
% two unrelated figures and one margin device between them is not a case this
% file has seen.
\newlength{\pd@figtop}\newlength{\pd@figbot}\def\pd@figpage{0}
% Does OFFSET (already sitting in \pd@blockoff) sit inside the room this page
% actually has for it, without moving? \pd@reqmin/\pd@reqmax are computed by
% \pd@marginbounds below, from exactly the three collision facts \pd@placemargin
% already knew about (page top, the last margin head, the last margin block) --
% this is a read of the same facts, not a new rule. A candidate that fits
% clean is preferred over one merely clamped INTO range, because a clamped
% value is a collision this page already had; a candidate the bounds do not
% touch is one that was never in anyone's way.
\newlength{\pd@reqmin}\newlength{\pd@reqmax}\newif\ifpd@havemax\newif\ifpd@fits
\newcommand{\pd@marginbounds}[1]{%
  \setlength{\pd@reqmin}{-\pagetotal}%
  \ifnum\pd@lastheadpage=\value{page}\relax
    \ifdim\pd@reqmin<\dimexpr\pd@lastheadbottom-\pagetotal\relax
      \setlength{\pd@reqmin}{\dimexpr\pd@lastheadbottom-\pagetotal\relax}\fi\fi
  \ifnum\pd@blockpage=\value{page}\relax
    \ifdim\pd@reqmin<\dimexpr\pd@blockbottom-\pagetotal+2pt\relax
      \setlength{\pd@reqmin}{\dimexpr\pd@blockbottom-\pagetotal+2pt\relax}\fi\fi
  \ifdim\pagegoal<\maxdimen
    \pd@havemaxtrue
    \setlength{\pd@reqmax}{\dimexpr\pagegoal-\pagetotal-\ht#1-\dp#1\relax}%
  \else\pd@havemaxfalse\fi}
\newcommand{\pd@tryoffset}{%
  \pd@fitstrue
  \ifdim\pd@blockoff<\pd@reqmin\pd@fitsfalse\fi
  \ifpd@fits\ifpd@havemax\ifdim\pd@blockoff>\pd@reqmax\pd@fitsfalse\fi\fi\fi}
\newcommand{\pd@marginhead}[1]{%
  \ifpdmargincolumn
    % Settle before the first \pagetotal read, unlike this macro's earlier
    % form. Every OTHER margin device in this file settles first (the
    % comment two macros up even says why: "\pagetotal is the page above
    % the LAST paragraph the builder saw"), but this one read \pagetotal
    % raw at the head's own issuing point and called the resulting slop
    % "to within a line" -- acceptable when a block follows a head closely,
    % not when many lines of prose come between them and a margin block
    % later clamps to \pd@lastheadbottom expecting it exact. Measured: a
    % \pdboundary head ("Where this stops") 14 lines above its section's
    % \end{pdrecitation}, recorded 2.35 pt short, so the Recall block's
    % clamp (aimed at a 3 pt gap under the head) landed 2.35 pt INSIDE it
    % instead -- built and measured overlapping, PDF p. 244.
    \pd@settle
    \setlength{\pd@headshift}{0pt}%
    \ifnum\pd@blockpage=\value{page}\relax
      \ifdim\pagetotal<\pd@blockbottom
        \setlength{\pd@headshift}{\dimexpr\pd@blockbottom-\pagetotal+3pt\relax}\fi
    \fi
    % And around the last HEAD on this page, which this macro recorded below
    % and then never read. Blocks stepped around heads, heads stepped around
    % blocks, and two heads walked straight through each other: a Pitfall head
    % printed 1.0 pt from a Recall head at the top of p. 306 of the maritime
    % edition, inside the 10.4 pt leading the collision check measures. The
    % rule is the same one \pd@placemargin already obeys three lines at a
    % time -- never above what is already in the column -- so it is applied
    % here in the same shape, after the block cap so whichever foot sits
    % lower wins.
    \ifnum\pd@lastheadpage=\value{page}\relax
      \ifdim\dimexpr\pagetotal+\pd@headshift\relax<\pd@lastheadbottom
        \setlength{\pd@headshift}{\dimexpr\pd@lastheadbottom-\pagetotal+3pt\relax}\fi
    \fi
    % A twin of the head, set at the margin's width and font, measured but
    % never shipped: a head of three or four words wraps to two lines there,
    % and a block raised later on the page has to clear the head's real foot.
    % Assuming one line put "RECALL" 0.1 pt under "WHERE THIS STOPS" in the
    % Swiss and technical editions, whose pagination differs from the
    % maritime edition's by a page.
    \sbox{\pd@headbox}{\begin{minipage}{\marginparwidth}\marginfont\scshape #1\end{minipage}}%
    \marginnote{\scshape #1}[\pd@headshift]%
    \xdef\pd@lastheadpage{\number\value{page}}%
    % +8pt, not +3pt: \pd@headbox measures the head's own text, but the
    % marginnote package sets every note's content behind a \strut (its
    % \marginfont\raggedrightmarginnote\strut...) whose height+depth come
    % from \marginfont's normal font metrics, not from \scshape "Where this
    % stops" alone -- the strut is measurably taller, so \ht+\dp here
    % under-reads the head's true rendered bottom. +3pt covered one 0.1pt
    % case (comment above); it did not cover this one -- built and measured
    % a Recall block clamped to +3pt still overlapping "Where this stops" by
    % 2.35 pt (PDF p. 244, pdboundary fourteen lines above its section's
    % \end{pdrecitation}). +8pt re-measured clear on that page and does not
    % reopen the "1.3 pt under" case the \newpage reservation above already
    % handles (a different mechanism -- a migrated page, not a short strut).
    \global\pd@lastheadbottom=\dimexpr\pagetotal+\pd@headshift
      +\ht\pd@headbox+\dp\pd@headbox+8pt\relax
  \else{\scshape #1.}\enspace\fi}
% A small square in the chapter's hue, in the margin, marking a worked example
% so a reader flipping through finds every one.
\newlength{\pd@blockoff}
% Place a measured margin block so its LAST line sits on the line that issues
% it -- the block occupies the page above its anchor rather than hanging below
% it, which is what keeps a note issued near the foot on the paper at all.
% Three caps, in order: never above the page top; never above the last margin
% head on this page (measured to that head's own foot); never above the last
% block on this page. Then the block records where its foot came to rest, so a
% head issued later steps down below it. \pagetotal is the shared clock for all
% of this, which is why every device that can settles the page builder first.
% Every margin device goes through here, including the three the margin
% apparatus added (\pdcite, \pdprov, \pdprovedon): a device that places its
% own note without asking what is already in the column is how the Book got
% seven collisions and a portrait below the paper, and there are now several
% hundred short notes in the column where there were dozens.
% Place a margin block with its TOP on the issuing line rather than its foot.
% Same occupancy bookkeeping as \pd@placemargin -- it must, or a later block
% would be told the column is free where this one sits -- but it never raises,
% so it cannot leave the page above. Clamped both ways: the offset is never
% negative (nothing above the text block) and, where \pagegoal is a real
% dimension, never so large that the block's foot passes the bottom of the
% text block. \pagegoal is \maxdimen on an empty page, which is room for
% anything; a block already taller than its page is left alone rather than
% shoved, because moving it cannot make it fit.
\newlength{\pd@blockfloor}
\newcommand{\pd@placemargintop}[1]{%
  \setlength{\pd@blockoff}{0pt}%
  \ifnum\pd@lastheadpage=\value{page}\relax
    \ifdim\pd@blockoff<\dimexpr\pd@lastheadbottom-\pagetotal\relax
      \setlength{\pd@blockoff}{\dimexpr\pd@lastheadbottom-\pagetotal\relax}\fi
  \fi
  \ifnum\pd@blockpage=\value{page}\relax
    \ifdim\pd@blockoff<\dimexpr\pd@blockbottom-\pagetotal+2pt\relax
      \setlength{\pd@blockoff}{\dimexpr\pd@blockbottom-\pagetotal+2pt\relax}\fi
  \fi
  \ifdim\pd@blockoff<0pt \setlength{\pd@blockoff}{0pt}\fi
  % The floor the two caps above just established: whatever is already in the
  % column on this page ends here, and nothing may be placed above it.
  \setlength{\pd@blockfloor}{\pd@blockoff}%
  % Pull the block up if its foot would pass the bottom of the text block --
  % but never above that floor. Reducing the offset past it is how the
  % foot-anchored placement used to undo its own downward pushes: a clamp
  % applied last re-collides the block with the head or the note it was just
  % moved clear of, which is the defect this file's other comment records as
  % 132 collisions on one merge. A block too tall to satisfy both is left at
  % the floor and overruns the foot, where page_overflow.py's foot-intrusion
  % and off-the-paper checks will report it -- a visible, measured overrun is
  % worth more than a silent overlap.
  \ifdim\pagegoal<\maxdimen
    \ifdim\dimexpr\pagetotal+\pd@blockoff+\ht#1+\dp#1\relax>\pagegoal
      \setlength{\pd@blockoff}{\dimexpr\pagegoal-\pagetotal-\ht#1-\dp#1\relax}%
      \ifdim\pd@blockoff<\pd@blockfloor \setlength{\pd@blockoff}{\pd@blockfloor}\fi
    \fi
  \fi
  \marginnote{\usebox#1}[\pd@blockoff]%
  \xdef\pd@blockpage{\number\value{page}}%
  \global\pd@blockbottom=\dimexpr\pagetotal+\pd@blockoff+\ht#1+\dp#1\relax}

% Default placement is now CENTERED on the vertical span of the most recently
% closed figure on this page, not anchored to the line that issues the note --
% a Recall block or a gloss discussing a figure reads naturally beside its
% middle, not hanging off one edge. Center only WHEN IT FITS CLEAN: tried
% against \pd@marginbounds's read of the same three collision facts the old
% single anchor already obeyed (page top, the last margin head, the last
% margin block), with no clamping. A center that does not fit clean is not
% clamped into place -- clamping a center is neither center nor top nor
% bottom, it is just wherever the clamp landed, and the point of asking for an
% aligned center was to get one honestly or not at all. Top-of-figure is tried
% next, then bottom-of-figure; the first of the three that needs no clamp
% wins. If none does, this falls through unchanged to the pre-existing
% issuing-line anchor and its own clamps below -- the same fallback this file
% always had, now reached only when no figure-relative placement clears.
\newif\ifpd@figcentered
\newcommand{\pd@placemargin}[1]{%
  \pd@figcenteredfalse
  \ifnum\pd@figpage=\value{page}\relax\ifdim\pd@figbot>\z@
    \pd@marginbounds{#1}%
    \setlength{\pd@blockoff}{\dimexpr(\pd@figtop+\pd@figbot)/2-\pagetotal-(\ht#1+\dp#1)/2\relax}%
    \pd@tryoffset
    \ifpd@fits\pd@figcenteredtrue\else
      \setlength{\pd@blockoff}{\dimexpr\pd@figtop-\pagetotal\relax}%
      \pd@tryoffset
      \ifpd@fits\pd@figcenteredtrue\else
        \setlength{\pd@blockoff}{\dimexpr\pd@figbot-\pagetotal-\ht#1-\dp#1\relax}%
        \pd@tryoffset
        \ifpd@fits\pd@figcenteredtrue\fi
      \fi
    \fi
  \fi\fi
  \ifpd@figcentered
    \marginnote{\usebox#1}[\pd@blockoff]%
    \xdef\pd@blockpage{\number\value{page}}%
    \global\pd@blockbottom=\dimexpr\pagetotal+\pd@blockoff+\ht#1+\dp#1\relax
  \else
  \setlength{\pd@blockoff}{\dimexpr\baselineskip-\ht#1-\dp#1\relax}%
  % A note never starts below the line that issues it -- clamped here, BEFORE
  % the occupancy caps, not after. Clamping last undid every downward push:
  % two \pdcite calls in one paragraph read the same \pagetotal, so the second
  % was told to start below the first's foot and then clamped straight back
  % onto it (132 collisions on the merge, all of them pairs of citations on
  % one line). Below the anchor is exactly where a second note belongs.
  \ifdim\pd@blockoff>0pt \setlength{\pd@blockoff}{0pt}\fi
  \ifdim\pd@blockoff<-\pagetotal
    \setlength{\pd@blockoff}{-\pagetotal}%
    % Headroom at the column top, and the reason is the one hole left in the
    % occupancy model. Every cap below is keyed on the numeric page counter, which is the
    % page the device was ISSUED on -- and a head issued in the last line of
    % a page whose paragraph then breaks prints at the top of the NEXT page
    % while having recorded the page before. That is exactly the Recall-on-
    % Pitfall collision at the top of p. 306 of the maritime edition: the
    % Pitfall head began its paragraph at the foot of p. 305, migrated with
    % the break, and the Recall block clamped flush to p. 306's column top
    % and landed 1.0 pt from a head no cap could see. TeX cannot tell us the
    % shipout page here, so the fix is not a better cap but a reservation:
    % when the top clamp binds, leave two lines for whatever may have come
    % over the break. Two, not one, because a head of three or four words
    % wraps to two lines at the margin's width. Only when there is room --
    % \pagegoal is \maxdimen on an empty page, which is room enough for
    % anything, and a block already too tall for its page is left where the
    % clamp put it rather than pushed off the paper.
    \ifdim\dimexpr\pagetotal+\pd@blockoff+2\baselineskip+\ht#1+\dp#1\relax>\pagegoal
      \ifdim\dimexpr\pagetotal+\pd@blockoff+\baselineskip+\ht#1+\dp#1\relax>\pagegoal
      \else\setlength{\pd@blockoff}{\dimexpr\pd@blockoff+\baselineskip\relax}\fi
    \else\setlength{\pd@blockoff}{\dimexpr\pd@blockoff+2\baselineskip\relax}\fi
  \fi
  \ifnum\pd@lastheadpage=\value{page}\relax
    \ifdim\pd@blockoff<\dimexpr\pd@lastheadbottom-\pagetotal\relax
      \setlength{\pd@blockoff}{\dimexpr\pd@lastheadbottom-\pagetotal\relax}\fi
  \fi
  \ifnum\pd@blockpage=\value{page}\relax
    \ifdim\pd@blockoff<\dimexpr\pd@blockbottom-\pagetotal+2pt\relax
      \setlength{\pd@blockoff}{\dimexpr\pd@blockbottom-\pagetotal+2pt\relax}\fi
  \fi
  \marginnote{\usebox#1}[\pd@blockoff]%
  \xdef\pd@blockpage{\number\value{page}}%
  \global\pd@blockbottom=\dimexpr\pagetotal+\pd@blockoff+\ht#1+\dp#1\relax
  \fi}
% The same placement for a note the page can do without. A Recall block, a
% margin figure and a gloss carry content that is nowhere else, so they are
% always placed. A short-form citation, a provenance note and a discharge
% pointer are second copies: \pdcite has already emitted the \cite, so the
% work is in the text and in the bibliography whatever the margin does. There
% are several hundred of them and on a dense page the column has no room for
% one beside every line that cites -- arithmetic cannot make space that is not
% there. So when placing one would push its foot past the text block's foot,
% it is left out rather than printed over its neighbour or below the paper.
% Nothing is lost that a reader could otherwise have read.
\newcommand{\pd@placemarginopt}[1]{%
  \setlength{\pd@blockoff}{\dimexpr\baselineskip-\ht#1-\dp#1\relax}%
  \ifdim\pd@blockoff>0pt \setlength{\pd@blockoff}{0pt}\fi
  \ifdim\pd@blockoff<-\pagetotal \setlength{\pd@blockoff}{-\pagetotal}\fi
  \ifnum\pd@lastheadpage=\value{page}\relax
    \ifdim\pd@blockoff<\dimexpr\pd@lastheadbottom-\pagetotal\relax
      \setlength{\pd@blockoff}{\dimexpr\pd@lastheadbottom-\pagetotal\relax}\fi
  \fi
  \ifnum\pd@blockpage=\value{page}\relax
    \ifdim\pd@blockoff<\dimexpr\pd@blockbottom-\pagetotal+2pt\relax
      \setlength{\pd@blockoff}{\dimexpr\pd@blockbottom-\pagetotal+2pt\relax}\fi
  \fi
  % Room left on this page, measured to the block's foot, with two lines of
  % slack. \pagegoal is \maxdimen while the page is empty, which is room
  % enough for anything. The slack is there because \pagetotal under-reads
  % mid-paragraph, which is where most of these are called: without it the
  % guard believed there was room and put an Exercises pointer 8.2 pt off the
  % foot of p. 309 of the Swiss edition -- a page that only exists in the CI
  % build's pagination, which is why the local builds were clean and the check
  % running over every edition in CI is what found it.
  \ifdim\dimexpr\pagetotal+\pd@blockoff+\ht#1+\dp#1+2\baselineskip\relax>\pagegoal
  \else
    \marginnote{\usebox#1}[\pd@blockoff]%
    \xdef\pd@blockpage{\number\value{page}}%
    \global\pd@blockbottom=\dimexpr\pagetotal+\pd@blockoff+\ht#1+\dp#1\relax
  \fi}
\newcommand{\pd@marginglyph}{%
  \ifpdmargincolumn\marginnote{\textcolor{\pd@chaptercolor}{\rule[-0.5pt]{4.5pt}{4.5pt}}}\fi}

% \pdgloss{Term}{one-line definition} -- a defined term set bold where it is
% defined, in the running text, the way the Book already sets any term of art
% at its point of definition; the definition itself is carried as a margin
% note rather than left inline, so a reader can recall it from the margin
% without re-parsing the sentence. In the margin column (Book) this reuses
% the small-caps head convention \pd@marginhead sets for \keyidea, \pitfall
% and \xrefbox's own heads, but folds head and body into ONE \marginnote call
% -- calling \pd@marginhead and a separate \marginnote back to back would
% place two independent margin boxes at the same source line and risk them
% colliding, exactly what \pdmarginfigure avoids by building one combined box
% for its image and caption. The margin text renders at \marginfont, the
% same footnotesize sidenote size every other margin note in this file uses.
% Outside the margin column (standalone chapters) the definition folds back
% into the sentence as a parenthetical, so no content is lost. No box, no
% rule -- a bare labelled aside like every other device here. No trailing
% \ignorespaces: unlike \pdboundary/\pdexample (environments whose content
% starts on the next source line), \pdgloss is called mid-sentence and the
% genuine word space that follows it in the prose must survive.
% The note is measured and RAISED so its last line sits on the line that
% issues it, for the same reason the Recall block is: a gloss called in the
% last line or two of a page hung its definition below the paper, where the
% log says nothing and the reader has nothing (Lamport's gloss, 81 pt off the
% foot of p. 68 in the technical edition). Raised, the note occupies the page
% above its anchor, which is page the anchor is on by construction; the only
% cap it needs is the page top. \pagetotal is read mid-sentence here and so
% lags by the lines of the current paragraph -- it under-reads, so the top cap
% fires a little early and the note is raised a little less than it could be,
% which is the harmless direction.
\newsavebox{\pd@glossbox}
\newcommand{\pdgloss}[2]{%
  % The margin box is placed BEFORE the term: the placement machinery
  % consumes the token that follows it, and when the term came first that
  % token was the word space after the term ("fencing epochof its role",
  % found the first time the macro was used, 2026-09-11). Ending the macro
  % on the term leaves the prose's own space where it was.
  \ifpdmargincolumn
    \sbox{\pd@glossbox}{\begin{minipage}[b]{\marginparwidth}\marginfont
      {\scshape #1}\par\nobreak\vspace{1pt plus 1pt}#2\end{minipage}}%
    \pd@placemargin{\pd@glossbox}%
    {\bfseries #1}%
  \else
    {\bfseries #1}\space{\itshape (#2)}%
  \fi}

% Marginalia (Wave 12, 12.2): a duotone portrait or title page in the margin
% column at the point where the person's idea is decisive, with a
% one-sentence caption. Book only: the standalone chapters have no margin
% column and cannot see the plates; each plate has a licence sidecar and is
% listed on the Image credits page. \pdmarginfigure{slug}{caption}
\newsavebox{\pd@mfbox}
\newcommand{\pdmarginfigure}[2]{%
  \ifpdmargincolumn
    \IfFileExists{plates/marginalia/#1.jpg}{%
      \sbox{\pd@mfbox}{\begin{minipage}[b]{\marginparwidth}
        \noindent\includegraphics[width=\marginparwidth]{plates/marginalia/#1.jpg}\par
        \vspace{3pt}{\raggedright\scriptsize\itshape #2\par}\end{minipage}}%
      % This was the file's last \marginpar, on the belief that the page
      % builder moves one up when it would run off the foot. It does not: it
      % moves marginpars DOWN to clear each other and never up, so Lamport's
      % portrait and its caption printed 81 pt below the paper on p. 68 of the
      % technical edition, where no warning and no reader could reach them.
      % Placed like every other margin block instead: measured, raised to sit
      % on its own line, capped by what is already in the column.
      % Top-anchored, unlike every other margin block, and that is the whole
      % fix for a portrait that printed off the top of the paper. The others
      % raise themselves so their LAST line sits on the line that issued them,
      % which is right for a note you read beside its sentence. A figure is
      % tall, and the raise is computed from \pagetotal -- true only for the
      % page TeX is filling now. Ostrom's portrait was issued right after
      % \paragraph{Commons governance.}, whose paragraph then broke to the
      % next page: the offset was measured against the foot of p. 327 and
      % applied at the top of p. 328, putting 88 pt of a 117 pt image above
      % the paper, over the running head, on a page two-thirds empty. Anchored
      % at the top it can only extend downward, where \pd@placemargin's foot
      % and occupancy caps already apply, so no stale \pagetotal can lift it
      % off the page.
      \pd@placemargintop{\pd@mfbox}%
    }{}%
  \fi}

% Every Book caption belongs in the OUTER margin, including long captions.
% Capture it separately and attach it at the exhibit's top. The Book's
% pd-margin-layout override allocates the margin independently at shipout.
% Wide artwork must be recomposed; captions never get a band below the art.
\newsavebox{\pd@captionbox}
\newsavebox{\pd@floatcaptions}
\newif\ifpd@captionfloat
\newif\ifpd@captionwide
\let\pd@inlinecaption\caption
% A preceding Recall block can extend below its issuing line. It is outside
% TeX's vertical list, so finish its occupied band before admitting a float;
% otherwise the new caption would share that same margin space.
\newcommand{\pd@clearcaptionband}{%
  \par\pd@settle
  \ifnum\pd@blockpage=\value{page}\relax
    \ifdim\pd@blockbottom>\pagetotal
      \ifdim\dimexpr\pd@blockbottom+8pt\relax>\pagegoal
        \newpage
      \else\vskip\dimexpr\pd@blockbottom-\pagetotal+8pt\relax\fi
      \pd@settle
    \fi
  \fi}
\newcommand{\pd@begincaptionfloat}{%
  \ifpdmargincolumn
    \pd@clearcaptionband
    \pd@captionfloattrue
    \global\pd@captionwidefalse
    \global\setbox\pd@floatcaptions=\box\voidb@x
  \fi}
\AtBeginEnvironment{figure}{\pd@begincaptionfloat}
\AtBeginEnvironment{figure*}{\pd@begincaptionfloat}
\AtBeginEnvironment{table}{\pd@begincaptionfloat}
\AtBeginEnvironment{table*}{\pd@begincaptionfloat}
\newcommand{\pd@storecaption}[1]{%
  \sbox{\pd@captionbox}{\begin{minipage}{\marginparwidth}\marginfont
    #1\par\end{minipage}}%
  \ifpd@captionfloat
    \global\setbox\pd@floatcaptions=\vbox{%
      \ifvoid\pd@floatcaptions\else
        \unvbox\pd@floatcaptions\vskip6pt
      \fi
      \hbox{\usebox{\pd@captionbox}}}%
  \else
    % Non-floating captionof: preserve the outer-margin contract. The Book's
    % page audit must still check the caller's reserved space.
    \pd@placemargintop{\pd@captionbox}%
  \fi}
\newcommand{\pdmargincaption}[2][]{%
  \ifpdmargincolumn
    \ifx\@captype\@undefined
      \PackageError{pd-pedagogy}{\string\pdmargincaption needs a caption type}%
        {Use figure, table, or captionof.}%
    \else
      \refstepcounter{\@captype}%
      \ifpd@captionfloat
        \xdef\pd@exhibitnumber{\@nameuse{the\@captype}}%
        \xdef\pd@exhibitkind{\@captype}%
      \fi
      \def\@currentlabelname{#2}%
      \def\pd@captionshort{#1}%
      \ifx\pd@captionshort\@empty\def\pd@captionshort{#2}\fi
      \edef\pd@captionext{\@nameuse{ext@\@captype}}%
      \addcontentsline{\pd@captionext}{\@captype}{%
        \protect\numberline{\@nameuse{the\@captype}}{\pd@captionshort}}%
      \pd@storecaption{{\bfseries\@nameuse{fnum@\@captype}.}\par\nobreak #2}%
      \typeout{PD-MARGIN-CAPTION: \@captype\space \@nameuse{the\@captype}}%
    \fi
  \else
    \def\pd@captionshort{#1}%
    \ifx\pd@captionshort\@empty\pd@inlinecaption{#2}\else\pd@inlinecaption[#1]{#2}\fi
  \fi}
% Route by caption TYPE, not the innermost environment: a caption inside a
% minipage or center within a figure is still a figure caption.
\newif\ifpd@routecaption
\newcommand{\pd@captiontype}{%
  \pd@routecaptionfalse
  \ifpdmargincolumn\ifx\@captype\@undefined\else
    \edef\pd@captionenv{\@captype}%
    \def\pd@captiontarget{figure}\ifx\pd@captionenv\pd@captiontarget\pd@routecaptiontrue\fi
    \def\pd@captiontarget{table}\ifx\pd@captionenv\pd@captiontarget\pd@routecaptiontrue\fi
  \fi\fi}
\newcommand{\pd@captiondispatch}[2]{%
  \pd@captiontype
  \ifpd@routecaption
    \pdmargincaption[#1]{#2}%
  \else
    \def\pd@captionshort{#1}%
    \ifx\pd@captionshort\@empty\pd@inlinecaption{#2}\else\pd@inlinecaption[#1]{#2}\fi
  \fi}
\newcommand{\pd@captionstar}[1]{%
  \pd@captiontype
  \ifpd@routecaption\pd@storecaption{#1}\else\pd@inlinecaption*{#1}\fi}
\newcommand{\pd@captionrouter}{\@ifstar\pd@captionstar{\@ifnextchar[\pd@captionrouter@opt\pd@captionrouter@plain}}
\long\def\pd@captionrouter@opt[#1]#2{\pd@captiondispatch{#1}{#2}}
\long\def\pd@captionrouter@plain#1{\pd@captiondispatch{}{#1}}
% Called INSIDE the finished float's box, so strict page checking measures
% the shipped leaf, not the page on which TeX encountered the source.
\newcommand{\pd@floatcaptionedge}{%
  \checkoddpage
  \ifoddpage
    \rlap{\hskip\dimexpr\textwidth+\marginparsep\relax
      \vtop{\vskip0pt\unvbox\pd@floatcaptions}}%
  \else
    \llap{\vtop{\vskip0pt\unvbox\pd@floatcaptions}\hskip\marginparsep}%
  \fi}
\newcommand{\pd@attachfloatcaption}{%
  \ifpdmargincolumn\ifpd@captionfloat\ifvoid\pd@floatcaptions\else
    \ifdim\wd\@currbox>\dimexpr\textwidth+1pt\relax
      \PackageError{pd-pedagogy}{Exhibit exceeds the captioned text column}%
        {Recompose the exhibit; do not shrink labels or put its caption below.}%
    \fi
    \global\setbox\@currbox=\vbox{%
      \hbox to\textwidth{\pd@floatcaptionedge
        \vtop{\vskip0pt\unvbox\@currbox}\hss}\vskip0pt}%
    \ifdim\dimexpr\ht\@currbox+\dp\@currbox\relax>\textheight
      \PackageError{pd-pedagogy}{Exhibit and margin caption exceed a Book page:
        \@captype\space\@nameuse{the\@captype},
        height \the\ht\@currbox}%
        {Shorten the caption or split the exhibit; never shrink its type.}%
    \fi
  \fi\fi\fi}
\AtBeginDocument{%
  \let\caption\pd@captionrouter
  \let\pd@endfloatbox\@endfloatbox
  \def\@endfloatbox{\pd@endfloatbox\pd@attachfloatcaption}}

% Page-breaking tables replace \caption internally, so the ordinary float
% router cannot reach them. Keep longtable's own counters, list entries and
% repeated column heads; change only the first caption cell's visual placement.
% Reserve the measured caption height from the previous pass before the
% table starts. Like page references, this converges with the Book build; a
% fixed fraction of the page cannot protect an arbitrarily long caption.
% Subsequent pages retain just column heads.
\newcounter{pd@captiontable}
\newcommand{\pdtablecaptionheight}[2]{%
  \expandafter\gdef\csname pd@tablecaptionheight@#1\endcsname{#2}}
\newcommand{\pd@begincaptiontable}{%
  \ifpdmargincolumn
    \pd@clearcaptionband
    \stepcounter{pd@captiontable}%
    \edef\pd@tablecaptionkey{\arabic{pd@captiontable}}%
    \ifcsname pd@tablecaptionheight@\pd@tablecaptionkey\endcsname
      \Needspace{\dimexpr\csname pd@tablecaptionheight@\pd@tablecaptionkey\endcsname+24pt\relax}%
    \else\Needspace{.4\textheight}\fi
  \fi}
\AtBeginEnvironment{xltabular}{\pd@begincaptiontable}
\AtBeginEnvironment{longtable}{%
  \ifx\pd@tablecaptionkey\@undefined\pd@begincaptiontable\fi}
\AtBeginDocument{%
  \let\pd@LTmakecaption\LT@makecaption
  \long\def\LT@makecaption#1#2#3{%
    \ifpdmargincolumn
      \LT@mcol\LT@cols{@{}l@{}}{%
        \sbox{\pd@captionbox}{\begin{minipage}{\marginparwidth}\marginfont
          #1{{\bfseries#2.}\endgraf\nobreak}#3\endgraf\end{minipage}}%
        \ifdim\dimexpr\ht\pd@captionbox+\dp\pd@captionbox+24pt\relax>\textheight
          \PackageError{pd-pedagogy}{Page-breaking table caption exceeds a Book page}%
            {Shorten the caption; do not move it inline or shrink its type.}%
        \fi
        \immediate\write\@auxout{\string\pdtablecaptionheight
          {\pd@tablecaptionkey}{\the\dimexpr\ht\pd@captionbox+\dp\pd@captionbox\relax}}%
        \smash{\hbox to0pt{%
          \checkoddpage
          \ifoddpage\hskip\dimexpr\textwidth+\marginparsep\relax
          \else\hskip-\dimexpr\marginparwidth+\marginparsep\relax\fi
          \vtop{\vskip0pt\hbox{\usebox{\pd@captionbox}}}\hss}}}%
    \else\pd@LTmakecaption#1{#2}{#3}\fi}}

% listings has its own caption renderer and can break across pages. Preserve
% its counters/list entries, but force the caption to the first page's margin.
% Reserve the measured caption before the listing begins; after a short listing
% clear its remaining occupied margin band before admitting more exhibits.
\newsavebox{\pd@listingcaptionbox}
\newcommand{\pd@makelistingcaption}[2]{%
  \sbox{\pd@listingcaptionbox}{\begin{minipage}{\marginparwidth}\marginfont
    {\bfseries#1.}\par\nobreak #2\par\end{minipage}}%
  \pd@settle
  \xdef\pd@blockpage{\number\value{page}}%
  \global\pd@blockbottom=\dimexpr\pagetotal+\ht\pd@listingcaptionbox+\dp\pd@listingcaptionbox+8pt\relax
  \nointerlineskip\hbox to\textwidth{%
    \smash{\checkoddpage
      \ifoddpage
        \rlap{\hskip\dimexpr\textwidth+\marginparsep\relax
          \vtop{\vskip0pt\hbox{\usebox{\pd@listingcaptionbox}}}}%
      \else
        \llap{\vtop{\vskip0pt\hbox{\usebox{\pd@listingcaptionbox}}}\hskip\marginparsep}%
      \fi}\hfil}\nobreak}
\AtBeginDocument{%
  \let\pd@lstMakeCaption\lst@MakeCaption
  \let\pd@lstmakecaption\lst@makecaption
  \def\lst@MakeCaption#1{%
    \ifpdmargincolumn
      \ifx#1t%
        \ifpd@captionfloat\else\ifx\lst@caption\@empty\else
          \pd@clearcaptionband
          \sbox{\pd@listingcaptionbox}{\begin{minipage}{\marginparwidth}\marginfont
            {\bfseries Listing.}\par\nobreak\lst@caption\par\end{minipage}}%
          \Needspace{\dimexpr\ht\pd@listingcaptionbox+\dp\pd@listingcaptionbox+24pt\relax}%
        \fi\fi
      \fi
      \def\lst@captionpos{t}%
    \fi
    \pd@lstMakeCaption{#1}}%
  \long\def\lst@makecaption#1#2{%
    \ifpdmargincolumn
      \ifpd@captionfloat
        \pd@storecaption{{\bfseries#1.}\par\nobreak #2}%
      \else\pd@makelistingcaption{#1}{#2}\fi
    \else\pd@lstmakecaption{#1}{#2}\fi}%
  \lst@AddToHook{DeInit}{\ifpdmargincolumn\ifpd@captionfloat\else\pd@clearcaptionband\fi\fi}}


% ---------------------------------------------------------------------------
% Marginalia (Wave 16): citations, provenance, and proved-on pointers at the
% point of use. Three GENERATED data tables carry the per-chapter content
% these three macros look up; the macros themselves are hand-written, here,
% once. \input immediately: \pdcite and \pdprovedon need their tables loaded
% before any chapter body calls them, and no chapter calls them before this
% file returns.
%
%   figures/pd-cite-shortforms.tex   -- scripts/harbor-research/build_cite_shortforms.py
%   figures/pd-discharges.tex        -- scripts/harbor-research/build_discharge_pointers.py
%
% Both are generated, committed, twinned the same way this file is; do not
% hand-edit either.
% ---------------------------------------------------------------------------

% \pdciteshort{key}{short form} -- one row of the citation short-form table,
% stored in a \csname so a chapter's own \bibitem keys (hyphens, digits, and
% all) are safe control-sequence-name fragments. \pd@citeshort@one looks a
% key up and falls back to printing the bare key when the table has no row
% for it (an UNPARSED \bibitem, or a key the Book's citation collation did
% not alias -- see mega-volume-cite-aliases.tex, generated by
% scripts/generate-mega-whitepaper.mjs once per Book build).
\newcommand{\pdciteshort}[2]{\expandafter\gdef\csname pd@citeshort@#1\endcsname{#2}}
\newcommand{\pd@citeshort@one}[1]{%
  \ifcsname pd@citeshort@#1\endcsname\csname pd@citeshort@#1\endcsname\else#1\fi}
% GENERATED by scripts/harbor-research/build_cite_shortforms.py from the
% eight Book chapters' own \bibitem entries -- do not hand-edit. Regenerate
% with: python3 scripts/harbor-research/build_cite_shortforms.py
%
% One \pdciteshort{key}{short form} per parseable \bibitem key, read by
% \pdcite (figures/pd-pedagogy.tex) for its Book-only margin note. A key with
% no row here falls back to printing its own bibliography key.
% Twin copy: book/figures/pd-cite-shortforms.tex and
% book/figures/pd-cite-shortforms.tex are kept
% byte-identical by this script (both preambles \input it via
% figures/pd-pedagogy.tex).

\pdciteshort{aacp11}{Alvim et al. 2011, \textit{On the relation between\ldots}}
\pdciteshort{acp2025}{OpenAI \& Stripe 2025, \textit{Agentic Commerce Protocol}}
\pdciteshort{activitypub2018}{Lemmer-Webber et al. 2018, \textit{ActivityPub}}
\pdciteshort{agache2020firecracker}{Agache et al. 2020}
\pdciteshort{agarwal2019market}{Agarwal et al. 2019, \textit{A Marketplace for\ldots}}
\pdciteshort{agha1986actors}{Agha 1986, \textit{Actors: A Model of Concurrent\ldots}}
\pdciteshort{akata2023repeated}{Akata et al. 2023, \textit{Playing Repeated Games\ldots}}
\pdciteshort{akerlof1970}{Akerlof 1970, \textit{The Market for ``Lemons'':\ldots}}
\pdciteshort{allen1990}{Allen 1990, \textit{Probability, Statistics, and\ldots}}
\pdciteshort{ameloot2013}{Ameloot et al. 2013}
\pdciteshort{ananthanarayanan2019keeping}{Ananthanarayanan et al. 2019}
\pdciteshort{anderson1972}{Anderson 1972, \textit{Computer Security\ldots}}
\pdciteshort{anthropic2025}{Anthropic 2025, \textit{Effective Context\ldots}}
\pdciteshort{anthropic2026hooks}{Anthropic 2026, \textit{Hooks reference}}
\pdciteshort{anthropic2026memory}{Anthropic 2026, \textit{How Claude Remembers Your\ldots}}
\pdciteshort{anthropic2026skills}{Anthropic 2026, \textit{Agent Skills: overview.\ldots}}
\pdciteshort{ap2mandates2026}{AP2 2026, \textit{Mandates}}
\pdciteshort{aps1990}{Abreu et al. 1990, \textit{Toward a Theory of\ldots}}
\pdciteshort{arbitrum-fraud-proofs}{Offchain Labs, \textit{Interactive Fraud Proofs}}
\pdciteshort{armstrong2006}{Armstrong 2006, \textit{Competition in Two-Sided\ldots}}
\pdciteshort{arrow1951social}{Arrow 1951, \textit{Social Choice and Individual\ldots}}
\pdciteshort{atiasaligrama2012}{Atia \& Saligrama 2012, \textit{Boolean compressed\ldots}}
\pdciteshort{aumann1974subjectivity}{Aumann 1974, \textit{Subjectivity and Correlation\ldots}}
\pdciteshort{austin1962how}{Austin 1962, \textit{How to Do Things with Words}}
\pdciteshort{aws2022kani}{AWS Automated Reasoning Group 2022, \textit{Kani Rust Verifier}}
\pdciteshort{aws2022s2n}{s2n-tls 2022, \textit{An Implementation of the\ldots}}
\pdciteshort{azm10}{Askarov et al. 2010, \textit{Predictive black-box\ldots}}
\pdciteshort{bainbridge1983}{Bainbridge 1983, \textit{Ironies of Automation}}
\pdciteshort{barbosa2021sok}{Barbosa et al. 2021, \textit{SoK: Computer-Aided\ldots}}
\pdciteshort{basin2013enforceable}{Basin et al. 2013, \textit{Enforceable Security\ldots}}
\pdciteshort{basin2018formal}{Basin et al. 2018, \textit{A Formal Analysis of 5G\ldots}}
\pdciteshort{bayer1993timestamping}{Bayer et al. 1993, \textit{Improving the\ldots}}
\pdciteshort{bdr04}{Barthe et al. 2004, \textit{Secure information\ldots}}
\pdciteshort{becker1968}{Becker 1968, \textit{Crime and punishment: An\ldots}}
\pdciteshort{beckerstigler1974}{Becker \& Stigler 1974, \textit{Law Enforcement,\ldots}}
\pdciteshort{behrouz2026sleep}{Behrouz et al. 2026, \textit{Language Models Need\ldots}}
\pdciteshort{bernheim1986common}{Bernheim \& Whinston 1986, \textit{Common Agency}}
\pdciteshort{bernstein2012high}{Bernstein et al. 2012}
\pdciteshort{bernstein2014orleans}{Bernstein et al. 2014}
\pdciteshort{bhargavan2017everest}{Bhargavan et al. 2017, \textit{Everest: Towards a\ldots}}
\pdciteshort{bhatt2025ctrlz}{Bhatt et al. 2025, \textit{Ctrl-Z: Controlling AI\ldots}}
\pdciteshort{birgisson2014macaroons}{Birgisson et al. 2014, \textit{Macaroons: Cookies\ldots}}
\pdciteshort{biscuit2026}{Biscuit 2026, \textit{Biscuit: a delegated,\ldots}}
\pdciteshort{blackwell1953}{Blackwell 1953, \textit{Equivalent Comparisons of\ldots}}
\pdciteshort{blanchet2007cryptoverif}{Blanchet 2007, \textit{CryptoVerif: A\ldots}}
\pdciteshort{blanchet2016modeling}{Blanchet 2016, \textit{Modeling and Verifying\ldots}}
\pdciteshort{blanchet2017proverif}{Blanchet et al. 2022, \textit{ProVerif with\ldots}}
\pdciteshort{bordini2007agentspeak}{Bordini et al. 2007}
\pdciteshort{bradleyterry1952}{Bradley \& Terry 1952, \textit{Rank Analysis of\ldots}}
\pdciteshort{bratman1987intention}{Bratman 1987, \textit{Intention, Plans, and\ldots}}
\pdciteshort{brown2022privacy}{Brown et al. 2022, \textit{What does it mean for a\ldots}}
\pdciteshort{brumley2011remote}{Brumley \& Tuveri 2011, \textit{Remote Timing\ldots}}
\pdciteshort{brun2011proactive}{Brun et al. 2011, \textit{Proactive Detection of\ldots}}
\pdciteshort{buchman2016tendermint}{Buchman 2016, \textit{Tendermint: Byzantine Fault\ldots}}
\pdciteshort{burrows2006chubby}{Burrows 2006, \textit{The Chubby Lock Service for\ldots}}
\pdciteshort{buterin2014ethereum}{Buterin 2014, \textit{A Next-Generation Smart\ldots}}
\pdciteshort{buterin2017casper}{Buterin \& Griffith 2017, \textit{Casper the\ldots}}
\pdciteshort{butler1736identity}{Butler 1736, \textit{Of Personal Identity}}
\pdciteshort{cai2025}{Cai et al. 2025, \textit{Are You Getting What You\ldots}}
\pdciteshort{campbell1979}{Campbell 1979, \textit{Assessing the Impact of\ldots}}
\pdciteshort{carlini2019secret}{Carlini et al. 2019, \textit{~Erlingsson, J.~Kos,\ldots}}
\pdciteshort{carlini2021extracting}{Carlini et al. 2020}
\pdciteshort{cciso15408}{ISO/IEC 15408-1:2022, \textit{Evaluation criteria}}
\pdciteshort{cerruti2026eal}{Cerruti et al. 2026, \textit{Agent Memory Is a\ldots}}
\pdciteshort{chalkias2020taming}{Chalkias et al. 2020, \textit{Taming the Many\ldots}}
\pdciteshort{chandratoueg1996}{Chandra \& Toueg 1996, \textit{Unreliable Failure\ldots}}
\pdciteshort{charrier2014bigbrother}{Charrier et al. 2014, \textit{Big Brother Logic:\ldots}}
\pdciteshort{chatterjee1983}{Chatterjee \& Samuelson 1983}
\pdciteshort{chen2023gptrationality}{Chen et al. 2023, \textit{The Emergence of\ldots}}
\pdciteshort{chen2024agentpoison}{Chen et al. 2024, \textit{AgentPoison: Red-teaming\ldots}}
\pdciteshort{chengfriedman2005}{Cheng \& Friedman 2005}
\pdciteshort{cheraghchi2013}{Cheraghchi 2013, \textit{Noise-resilient group\ldots}}
\pdciteshort{chhikara2025mem0}{Chhikara et al. 2025, \textit{Mem0: Building\ldots}}
\pdciteshort{chiu2025}{Chiu et al. 2025}
\pdciteshort{chiu2025strategic}{Chiu et al. 2025}
\pdciteshort{chudnov2018continuous}{Chudnov et al. 2018, \textit{Continuous Formal\ldots}}
\pdciteshort{clarke1971multipart}{Clarke 1971, \textit{Multipart pricing of public\ldots}}
\pdciteshort{cohenlevesque1990}{Cohen \& Levesque 1990, \textit{Intention Is\ldots}}
\pdciteshort{consul}{HashiCorp, \textit{Consul Service Discovery and DNS Interface}}
\pdciteshort{conway1968}{Conway 1968, \textit{How do Committees Invent?}}
\pdciteshort{cooper2024unlearning}{Cooper et al. 2024, \textit{Machine unlearning\ldots}}
\pdciteshort{corbet2023merge65}{Corbet 2023, \textit{The first half of the 6.5\ldots}}
\pdciteshort{cremers2017comprehensive}{Cremers et al. 2017, \textit{A Comprehensive\ldots}}
\pdciteshort{critch2017servant}{Critch \& Russell 2017, \textit{Servant of Many\ldots}}
\pdciteshort{crosby2009tamper}{Crosby \& Wallach 2009, \textit{Efficient Data\ldots}}
\pdciteshort{ctgvb2020}{Tosatto et al. 2020, \textit{Business Process Full\ldots}}
\pdciteshort{ctgvb2021}{Tosatto et al. 2021, \textit{Proving Regulatory\ldots}}
\pdciteshort{cvach2012}{Cvach 2012, \textit{Monitor Alarm Fatigue: An\ldots}}
\pdciteshort{cve20159235}{CVE-2015-9235 2015, \textit{Algorithm Confusion in\ldots}}
\pdciteshort{cve20181000531}{CVE-2018-1000531 2018, \textit{JWT ``none''\ldots}}
\pdciteshort{cve202348238}{CVE-2023-48238 2023, \textit{Algorithm confusion\ldots}}
\pdciteshort{dagv1979}{d'Aspremont \& G\'erard-Varet 1979}
\pdciteshort{dai2024}{Dai et al. 2024, \textit{Artificial Leviathan:\ldots}}
\pdciteshort{dao2022flashattention}{Dao et al. 2022, \textit{FlashAttention: Fast and\ldots}}
\pdciteshort{dartmouth1819}{U.S. Supreme Court 1819, \textit{Trustees of Dartmouth College v.\ Woodward}}
\pdciteshort{das2026selfpropagation}{Das et al. 2026}
\pdciteshort{dastani2017normenforcement}{Dastani et al. 2017, \textit{Norm Enforcement as\ldots}}
\pdciteshort{demers1987}{Demers et al. 1987, \textit{Epidemic Algorithms\ldots}}
\pdciteshort{demers1987epidemic}{Demers et al. 1987, \textit{Epidemic Algorithms\ldots}}
\pdciteshort{denning77}{Denning \& Denning 1977, \textit{Certification of\ldots}}
\pdciteshort{dennis1966programming}{Dennis \& Van Horn 1966}
\pdciteshort{dg84}{Dowling \& Gallier 1984}
\pdciteshort{dickens1989}{Dickens et al. 1989, \textit{Employee Crime and\ldots}}
\pdciteshort{dls1988}{Dwork et al. 1988, \textit{Consensus in the\ldots}}
\pdciteshort{dmp91}{Dechter et al. 1991, \textit{Temporal Constraint\ldots}}
\pdciteshort{doerr2017rumor}{Doerr \& Kostrygin 2017, \textit{Randomized Rumor\ldots}}
\pdciteshort{dolev1983security}{Dolev \& Yao 1983, \textit{On the Security of\ldots}}
\pdciteshort{dong2025minja}{Dong et al. 2025, \textit{Memory injection attacks\ldots}}
\pdciteshort{dong2026harmful}{Dong et al. 2026, \textit{Agent Skills Can Be\ldots}}
\pdciteshort{dorfman}{Dorfman 1943, \textit{The Detection of Defective\ldots}}
\pdciteshort{dorigo1999}{Dorigo \& Di Caro 1999, \textit{The Ant Colony\ldots}}
\pdciteshort{douceur2002}{Douceur 2002, \textit{The Sybil Attack}}
\pdciteshort{douceur2002sybil}{Douceur 2002, \textit{The Sybil Attack}}
\pdciteshort{drv}{Dwork et al. 2010, \textit{Boosting and\ldots}}
\pdciteshort{duan2024mia}{Duan et al. 2024, \textit{Do membership inference\ldots}}
\pdciteshort{dudik2011dr}{Dud\'{\i}k et al. 2011, \textit{Doubly Robust\ldots}}
\pdciteshort{duhwang}{Du \& Hwang 2000, \textit{Combinatorial Group\ldots}}
\pdciteshort{dwork2006dp}{Dwork 2006, \textit{Differential privacy}}
\pdciteshort{eckhardt1985}{Eckhardt \& Lee 1985, \textit{A Theoretical Basis\ldots}}
\pdciteshort{ellison1999spki}{Ellison et al. 1999, \textit{SPKI Certificate\ldots}}
\pdciteshort{elo1978}{Elo 1978, \textit{The Rating of Chessplayers, Past\ldots}}
\pdciteshort{endsley1995}{Endsley 1995, \textit{Toward a Theory of Situation\ldots}}
\pdciteshort{endsleykiris1995}{Endsley \& Kiris 1995, \textit{The Out-of-the-Loop\ldots}}
\pdciteshort{epw67}{Esary et al. 1967, \textit{Association of random\ldots}}
\pdciteshort{erlang1917}{Erlang 1917, \textit{Solution of some problems in\ldots}}
\pdciteshort{ethereum2spec}{Foundation 2020, \textit{Ethereum 2.0 Phase 0 --\ldots}}
\pdciteshort{expansion}{Si et al. 2013, \textit{Lossy Compression of\ldots}}
\pdciteshort{fan2014cuckoo}{Fan et al. 2014, \textit{Cuckoo Filter:\ldots}}
\pdciteshort{farina2024correlated}{Farina 2024, \textit{Setting and Equilibria: The Correlated Equilibrium}}
\pdciteshort{feldman2020renyi}{Feldman \& Zrnic 2020, \textit{Individual privacy\ldots}}
\pdciteshort{fema2017nims}{FEMA 2017, \textit{National Incident Management\ldots}}
\pdciteshort{fickinger2020multi}{Fickinger et al. 2020}
\pdciteshort{fipa00023}{FIPA 2002, \textit{Agent Management}}
\pdciteshort{fipa2002acl}{FIPA 2002, \textit{FIPA ACL Communicative Act\ldots}}
\pdciteshort{flp1985}{Fischer et al. 1985, \textit{Impossibility of\ldots}}
\pdciteshort{fong2019}{Fong \& Spivak 2019, \textit{Seven Sketches in\ldots}}
\pdciteshort{friedmanresnick2001}{Friedman \& Resnick 2001, \textit{The Social Cost\ldots}}
\pdciteshort{frieze1985shortest}{Frieze \& Grimmett 1985, \textit{The Shortest-Path\ldots}}
\pdciteshort{fudenberglevine1989}{Fudenberg \& Levine 1989, \textit{Reputation and\ldots}}
\pdciteshort{garcia1987sagas}{Garcia-Molina \& Salem 1987, \textit{Sagas}}
\pdciteshort{garciamolina1987}{Garcia-Molina \& Salem 1987, \textit{Sagas}}
\pdciteshort{garciamolina1987sagas}{Garcia-Molina \& Salem 1987, \textit{Sagas}}
\pdciteshort{gelernter1985}{Gelernter 1985, \textit{Generative Communication\ldots}}
\pdciteshort{gibbard1973manipulation}{Gibbard 1973, \textit{Manipulation of voting\ldots}}
\pdciteshort{githubcodeowners}{GitHub 2026, \textit{About Code Owners}}
\pdciteshort{goes2020ibc}{Goes 2020, \textit{The Interblockchain\ldots}}
\pdciteshort{goguen-meseguer-82}{Goguen \& Meseguer 1982, \textit{Security policies\ldots}}
\pdciteshort{goguen-meseguer-84}{Goguen \& Meseguer 1984, \textit{Unwinding and\ldots}}
\pdciteshort{goldberg1991floating}{Goldberg 1991, \textit{What every computer\ldots}}
\pdciteshort{goldwasser2008}{Goldwasser et al. 2008}
\pdciteshort{gordonroth2019locke}{Gordon-Roth 2025, \textit{Locke on Personal\ldots}}
\pdciteshort{grasse1959}{Grass\'e 1959, \textit{La reconstruction du nid et\ldots}}
\pdciteshort{gray1993transaction}{Gray \& Reuter 1993}
\pdciteshort{grayreuter1992}{Gray \& Reuter 1993}
\pdciteshort{greenblatt2024control}{Greenblatt et al. 2024, \textit{AI Control:\ldots}}
\pdciteshort{greenporter1984}{Green \& Porter 1984}
\pdciteshort{greenswets1966}{Green \& Swets 1966, \textit{Signal Detection\ldots}}
\pdciteshort{greshake2023indirect}{Greshake et al. 2023, \textit{Not what you've\ldots}}
\pdciteshort{grice1941identity}{Grice 1941, \textit{Personal Identity}}
\pdciteshort{grossman1986}{Grossman \& Hart 1986, \textit{The Costs and\ldots}}
\pdciteshort{grossmanhart1986}{Grossman \& Hart 1986, \textit{The Costs and\ldots}}
\pdciteshort{groves1973incentives}{Groves 1973, \textit{Incentives in teams}}
\pdciteshort{gu2025caching}{Gu et al. 2025, \textit{Auditing prompt caching in\ldots}}
\pdciteshort{gui2025hypertree}{Gui et al. 2025, \textit{HyperTree Planning:\ldots}}
\pdciteshort{gupta2010}{Gupta et al. 2010, \textit{On the\ldots}}
\pdciteshort{haberstornetta1991}{Haber \& Stornetta 1991, \textit{How to Time-Stamp\ldots}}
\pdciteshort{haeberlen2007peerreview}{Haeberlen et al. 2007}
\pdciteshort{hafner2023dreamerv3}{Hafner et al. 2023, \textit{Mastering Diverse\ldots}}
\pdciteshort{halpernmoses1990}{Halpern \& Moses 1990, \textit{Knowledge and Common Knowledge in a Distributed Environment}}
\pdciteshort{hamlen2006}{Hamlen et al. 2006, \textit{Computability Classes\ldots}}
\pdciteshort{hansengrist2019}{Hansen \& Ghrist 2019, \textit{Toward a spectral\ldots}}
\pdciteshort{hardy1988confused}{Hardy 1988, \textit{The Confused Deputy (or why\ldots}}
\pdciteshort{hayek1945}{Hayek 1945, \textit{The Use of Knowledge in\ldots}}
\pdciteshort{headscale}{Brackeva et al. 2020}
\pdciteshort{hedberg2024openidfed}{Hedberg (ed.) et al. 2024, \textit{OpenID Federation 1.0}}
\pdciteshort{heilman2015}{Heilman et al. 2015, \textit{Eclipse Attacks on\ldots}}
\pdciteshort{helland2012idempotence}{Helland 2012, \textit{Idempotence Is Not a Medical\ldots}}
\pdciteshort{hellerstein2010}{Hellerstein 2010, \textit{The Declarative\ldots}}
\pdciteshort{hellerstein2020}{Hellerstein \& Alvaro 2020, \textit{Keeping CALM:\ldots}}
\pdciteshort{herbrich2007trueskill}{Herbrich et al. 2007, \textit{TrueSkill: A\ldots}}
\pdciteshort{herlihy2018}{Herlihy 2018, \textit{Atomic Cross-Chain Swaps}}
\pdciteshort{herlihy2018atomic}{Herlihy 2018, \textit{Atomic Cross-Chain Swaps}}
\pdciteshort{herlihywing1990}{Herlihy \& Wing 1990, \textit{Linearizability: A\ldots}}
\pdciteshort{hewitt1973actor}{Hewitt et al. 1973, \textit{A universal modular\ldots}}
\pdciteshort{hickman2019krakoa}{Hickman 2019, \textit{House of X / Powers of X}}
\pdciteshort{hirschman1970}{Hirschman 1970, \textit{Exit, Voice, and Loyalty:\ldots}}
\pdciteshort{hls2022}{Herlihy et al. 2022, \textit{Cross-Chain Deals and\ldots}}
\pdciteshort{hoare1978csp}{Hoare 1978, \textit{Communicating sequential\ldots}}
\pdciteshort{hobbes1651}{Hobbes 1651, \textit{Leviathan, or The Matter,\ldots}}
\pdciteshort{hobbes1651leviathan}{Hobbes 1651, \textit{Leviathan, or The Matter,\ldots}}
\pdciteshort{holmstrom1982}{Holmstr\"om 1982, \textit{Moral hazard in teams}}
\pdciteshort{hong2024metagpt}{Hong et al. 2024, \textit{MetaGPT: Meta\ldots}}
\pdciteshort{horwitz1989integrating}{Horwitz et al. 1989}
\pdciteshort{huang2025ans}{Huang et al. 2025, \textit{Agent Name Service\ldots}}
\pdciteshort{hubinger2024sleeper}{Hubinger et al. 2024, \textit{Sleeper Agents:\ldots}}
\pdciteshort{hume1748}{Hume 1748, \textit{Of the Original Contract}}
\pdciteshort{hwang}{Hwang 1972, \textit{A Method for Detecting All\ldots}}
\pdciteshort{irving2018}{Irving et al. 2018, \textit{AI Safety via Debate}}
\pdciteshort{isa182}{ANSI/ISA 2016, \textit{ANSI/ISA-18.2: Management\ldots}}
\pdciteshort{izraelevitz2016}{Izraelevitz et al. 2016}
\pdciteshort{jasonfailureimplementation}{Jason developers, rev.\ \texttt{4b98ddb1f9c8}, \textit{Goal-failure handling}}
\pdciteshort{jasonfaq}{Jason developers, accessed 2026, \textit{Jason FAQ}}
\pdciteshort{jonessergot1993}{Jones \& Sergot 1993, \textit{On the\ldots}}
\pdciteshort{jonessergot1993regimentation}{Jones \& Sergot 1993, \textit{On the\ldots}}
\pdciteshort{josang2002beta}{Jøsang \& Ismail 2002, \textit{The Beta Reputation\ldots}}
\pdciteshort{josefsson2017edwards}{Josefsson \& Liusvaara 2017}
\pdciteshort{jpbb04}{Jung et al. 2004, \textit{Fast portscan detection\ldots}}
\pdciteshort{kamvar2003}{Kamvar et al. 2003, \textit{The EigenTrust\ldots}}
\pdciteshort{kamvar2003eigentrust}{Kamvar et al. 2003, \textit{The EigenTrust\ldots}}
\pdciteshort{kandorimatsushima1998}{Kandori \& Matsushima 1998}
\pdciteshort{karp2000rumor}{Karp et al. 2000, \textit{Randomized Rumor\ldots}}
\pdciteshort{kasi2013cassandra}{Kasi \& Sarma 2013, \textit{Cassandra: Proactive\ldots}}
\pdciteshort{kemmerer83}{Kemmerer 1983, \textit{Shared resource matrix\ldots}}
\pdciteshort{kendall1953}{Kendall 1953, \textit{Stochastic processes\ldots}}
\pdciteshort{kephartchess2003}{Kephart \& Chess 2003, \textit{The Vision of\ldots}}
\pdciteshort{kiefer1955}{Kiefer \& Wolfowitz 1955, \textit{On the theory of\ldots}}
\pdciteshort{kingman1961}{Kingman 1961, \textit{The single server queue in\ldots}}
\pdciteshort{klein1998}{Klein 1998, \textit{Sources of Power: How People\ldots}}
\pdciteshort{klein2009sel4}{Klein et al. 2009, \textit{seL4: Formal\ldots}}
\pdciteshort{kleppmann2016locking}{Kleppmann 2016, \textit{How to do distributed\ldots}}
\pdciteshort{kleppmann2017}{Kleppmann 2017, \textit{Designing Data-Intensive\ldots}}
\pdciteshort{kleppmann2019interleaving}{Kleppmann et al. 2019}
\pdciteshort{knight1986}{Knight \& Leveson 1986, \textit{An Experimental\ldots}}
\pdciteshort{kobeissi2017automated}{Kobeissi et al. 2017}
\pdciteshort{kofmanlawarree1993}{Kofman \& Lawarrée 1993, \textit{Collusion in\ldots}}
\pdciteshort{kohli2026ninejudges}{Kohli 2026, \textit{Nine Judges, Two Effective\ldots}}
\pdciteshort{komenda2016distributed}{Komenda et al. 2016, \textit{On a Distributed\ldots}}
\pdciteshort{krepswilson1982}{Kreps \& Wilson 1982, \textit{Reputation and\ldots}}
\pdciteshort{kwon2023vllm}{Kwon et al. 2023, \textit{Efficient Memory\ldots}}
\pdciteshort{lamport1978}{Lamport 1978, \textit{Time, Clocks, and the\ldots}}
\pdciteshort{lamport1998part}{Lamport 1998, \textit{The Part-Time Parliament}}
\pdciteshort{lamport2002specifying}{Lamport 2002, \textit{Specifying Systems: The\ldots}}
\pdciteshort{lampson1971}{Lampson 1971, \textit{Protection}}
\pdciteshort{lattimore2020bandit}{Lattimore \& Szepesv\'ari 2020}
\pdciteshort{laurie2013ct}{Laurie et al. 2013}
\pdciteshort{laurie2014}{Laurie 2014, \textit{Certificate Transparency}}
\pdciteshort{leesee2004}{Lee \& See 2004, \textit{Trust in Automation:\ldots}}
\pdciteshort{leroy2009}{Leroy 2009, \textit{Why Is It So Hard to Do My\ldots}}
\pdciteshort{lesaege2021kleros}{Lesaege et al. 2021, \textit{Kleros Long Paper\ldots}}
\pdciteshort{levien-advogato}{Levien \& Aiken 1999}
\pdciteshort{lewontin1974}{Lewontin 2006, \textit{The Analysis of Variance\ldots}}
\pdciteshort{li2004sundr}{Li et al. 2004, \textit{Secure Untrusted Data\ldots}}
\pdciteshort{li2026skillsbench}{Li et al. 2026, \textit{SkillsBench: Benchmarking\ldots}}
\pdciteshort{liang2022helm}{Liang et al. 2023, \textit{Holistic Evaluation of\ldots}}
\pdciteshort{ligatti2005edit}{Ligatti et al. 2005, \textit{Edit Automata:\ldots}}
\pdciteshort{lin-wonham}{Lin \& Wonham 1988, \textit{On observability of\ldots}}
\pdciteshort{lin2025sleeptime}{Lin et al. 2025, \textit{Sleep-time Compute:\ldots}}
\pdciteshort{littlewood1989}{Littlewood \& Miller 1989}
\pdciteshort{litzkow1997condor}{Litzkow et al. 1997, \textit{Checkpoint and\ldots}}
\pdciteshort{liu2014}{Liu \& Skrzypacz 2014, \textit{Limited Records and\ldots}}
\pdciteshort{liu2023agentbench}{Liu et al. 2024, \textit{AgentBench: Evaluating\ldots}}
\pdciteshort{liu2024}{Liu et al. 2024, \textit{Lost in the Middle: How\ldots}}
\pdciteshort{liu2024lost}{Liu et al. 2024, \textit{Lost in the Middle: How\ldots}}
\pdciteshort{liuskrzypacz2014}{Liu \& Skrzypacz 2014, \textit{Limited Records and\ldots}}
\pdciteshort{locke1689}{Locke 1689, \textit{Two Treatises of Government}}
\pdciteshort{locke1689identity}{Locke 1694, \textit{An Essay Concerning Human\ldots}}
\pdciteshort{locke1689two}{Locke 1689, \textit{Two Treatises of Government}}
\pdciteshort{lowe1997hierarchy}{Lowe 1997, \textit{A Hierarchy of Authentication\ldots}}
\pdciteshort{lukas2023pii}{Lukas et al. 2023, \textit{Analyzing leakage of\ldots}}
\pdciteshort{lw88}{Lin \& Wonham 1988, \textit{On observability of\ldots}}
\pdciteshort{lykouris2018}{Lykouris \& Vassilvitskii 2021}
\pdciteshort{macaroons2014}{Birgisson et al. 2014, \textit{Macaroons: Cookies\ldots}}
\pdciteshort{mackworth1948}{Mackworth 1948, \textit{The Breakdown of Vigilance\ldots}}
\pdciteshort{mailath-samuelson}{Mailath \& Samuelson 2001, \textit{Who Wants a\ldots}}
\pdciteshort{maler2008venn}{Maler \& Reed 2008, \textit{The Venn of Identity:\ldots}}
\pdciteshort{matrixspec}{The Matrix.org Foundation, \textit{The Matrix Specification}}
\pdciteshort{mazieres2002byzantine}{Mazi\`eres \& Shasha 2002, \textit{Building Secure\ldots}}
\pdciteshort{mclean2015critical}{McLean 2015, \textit{Critical Vulnerabilities in\ldots}}
\pdciteshort{mcmahan2017userdp}{McMahan et al. 2017}
\pdciteshort{mediratta2026parallact}{Mediratta et al. 2026}
\pdciteshort{meier2013tamarin}{Meier et al. 2013, \textit{The TAMARIN Prover for\ldots}}
\pdciteshort{merkle1987}{Merkle 1987, \textit{A Digital Signature Based on\ldots}}
\pdciteshort{meyer1992}{Meyer 1992, \textit{Applying ``Design by\ldots}}
\pdciteshort{miller1956}{Miller 1956, \textit{The Magical Number Seven,\ldots}}
\pdciteshort{miller2006robust}{Miller 2006, \textit{Robust Composition: Towards a\ldots}}
\pdciteshort{miller2006robustcomposition}{Miller 2006, \textit{Robust Composition: Towards a\ldots}}
\pdciteshort{mireshghallah2023confaide}{Mireshghallah et al. 2023, \textit{Can LLMs keep a\ldots}}
\pdciteshort{mohan1992aries}{Mohan et al. 1992, \textit{ARIES: A Transaction\ldots}}
\pdciteshort{morris2023vec2text}{Morris et al. 2023, \textit{Text embeddings reveal\ldots}}
\pdciteshort{mrv1999}{Micali et al. 1999, \textit{Verifiable Random\ldots}}
\pdciteshort{myers-liskov}{Myers \& Liskov 1997, \textit{A decentralized\ldots}}
\pdciteshort{myerson1981}{Myerson 1981, \textit{Optimal Auction Design}}
\pdciteshort{myerson1981optimal}{Myerson 1981, \textit{Optimal Auction Design}}
\pdciteshort{myerson1983}{Myerson \& Satterthwaite 1983}
\pdciteshort{myerson1983efficient}{Myerson \& Satterthwaite 1983}
\pdciteshort{needham1978using}{Needham \& Schroeder 1978}
\pdciteshort{needham2025evalaware}{Needham et al. 2025, \textit{Large Language Models\ldots}}
\pdciteshort{newman2022sigstore}{Newman et al. 2022, \textit{Sigstore: Software\ldots}}
\pdciteshort{nisan2007}{Nisan et al. 2007, \textit{Algorithmic Game Theory}}
\pdciteshort{nisan2007agt}{Nisan et al. 2007, \textit{Algorithmic Game Theory}}
\pdciteshort{nissenbaum2004ci}{Nissenbaum 2004, \textit{Privacy as contextual\ldots}}
\pdciteshort{nrtv2007}{Nisan et al. 2007, \textit{Tardos \& V. Vazirani\ldots}}
\pdciteshort{nygard2018}{Nygard 2018, \textit{Release It! Design and Deploy\ldots}}
\pdciteshort{olson2007}{Olson 2007, \textit{What Are We? A Study in\ldots}}
\pdciteshort{olson2024defeasible}{Olson et al. 2024, \textit{A Defeasible Deontic\ldots}}
\pdciteshort{ongaro2014search}{Ongaro \& Ousterhout 2014, \textit{In Search of an\ldots}}
\pdciteshort{ostrom1990}{Ostrom 1990, \textit{Governing the Commons: The\ldots}}
\pdciteshort{ostrom1990governing}{Ostrom 1990, \textit{Governing the Commons: The\ldots}}
\pdciteshort{ostrom2010polycentric}{Ostrom 2010, \textit{Beyond Markets and States:\ldots}}
\pdciteshort{ouyang2022instruct}{Ouyang et al. 2022, \textit{Training language\ldots}}
\pdciteshort{packer2023memgpt}{Packer et al. 2023, \textit{MemGPT: Towards LLMs\ldots}}
\pdciteshort{panda2025audit}{Panda et al. 2025, \textit{Privacy auditing of\ldots}}
\pdciteshort{pappu2024rlhf}{Pappu et al. 2024, \textit{Measuring memorization\ldots}}
\pdciteshort{parasuraman2000}{Parasuraman et al. 2000, \textit{A Model for Types\ldots}}
\pdciteshort{parfit1984}{Parfit 1984, \textit{Reasons and Persons}}
\pdciteshort{park2023}{Park et al. 2023, \textit{Generative Agents:\ldots}}
\pdciteshort{park2023generative}{Park et al. 2023, \textit{Generative Agents:\ldots}}
\pdciteshort{patel2026sixsigma}{Patel et al. 2026, \textit{The Six Sigma Agent:\ldots}}
\pdciteshort{pateman1979}{Pateman 1979, \textit{The Problem of Political\ldots}}
\pdciteshort{peterson2020augur}{Peterson et al. 2020, \textit{Augur: a\ldots}}
\pdciteshort{pillai2014allfilesystems}{Pillai et al. 2014, \textit{All File Systems Are\ldots}}
\pdciteshort{pirolli1999}{Pirolli \& Card 1999, \textit{Information Foraging}}
\pdciteshort{pittel1987rumor}{Pittel 1987, \textit{On Spreading a Rumor}}
\pdciteshort{pkube21}{Luo et al. 2021, \textit{Privacy budget scheduling}}
\pdciteshort{polanyi1966}{Polanyi 1966, \textit{The Tacit Dimension}}
\pdciteshort{rafailov2023dpo}{Rafailov et al. 2023, \textit{Direct Preference\ldots}}
\pdciteshort{ramadge1987supervisory}{Ramadge \& Wonham 1987}
\pdciteshort{rao1991modeling}{Rao \& Georgeff 1991, \textit{Modeling Rational\ldots}}
\pdciteshort{rao1995bdi}{Rao \& Georgeff 1995, \textit{BDI agents: From\ldots}}
\pdciteshort{rao2010}{Rao 2010, \textit{A Big Little Idea Called\ldots}}
\pdciteshort{rasmussen2025zep}{Rasmussen et al. 2025, \textit{Zep: A Temporal\ldots}}
\pdciteshort{rats}{Birkholz et al. 2023, \textit{Remote ATtestation\ldots}}
\pdciteshort{rdc}{Nguyen et al. 2026, \textit{Rate-Distortion-Classification Representation Theory}}
\pdciteshort{reid1785memory}{Reid 2017, \textit{Essays on the Intellectual\ldots}}
\pdciteshort{resnick2006}{Resnick et al. 2006, \textit{The Value of\ldots}}
\pdciteshort{resnick2006value}{Resnick et al. 2006, \textit{The Value of\ldots}}
\pdciteshort{rfc7009revocation}{Lodderstedt et al. 2013, \textit{OAuth 2.0 Token\ldots}}
\pdciteshort{rfc8725jwtbcp}{Sheffer et al. 2020, \textit{JSON Web Token Best\ldots}}
\pdciteshort{rfc9449dpop}{Fett et al. 2023, \textit{OAuth 2.0 Demonstrating\ldots}}
\pdciteshort{rivestlampson1996sdsi}{Rivest \& Lampson 1996, \textit{SDSI --- A Simple\ldots}}
\pdciteshort{rochet2003}{Rochet \& Tirole 2003}
\pdciteshort{rochet2006}{Rochet \& Tirole 2006, \textit{Two-Sided Markets:\ldots}}
\pdciteshort{rochettirole2003}{Rochet \& Tirole 2003}
\pdciteshort{rodl1985}{Rödl 1985, \textit{On a Packing and Covering\ldots}}
\pdciteshort{rohloff2005pspace}{Rohloff \& Lafortune 2005}
\pdciteshort{rothschildstiglitz1976}{Rothschild \& Stiglitz 1976}
\pdciteshort{rruv16}{Rogers et al. 2016, \textit{Privacy odometers and\ldots}}
\pdciteshort{rs10}{Russo \& Sabelfeld 2010, \textit{Dynamic vs}}
\pdciteshort{rushby}{Rushby 1992}
\pdciteshort{rw87}{Ramadge \& Wonham 1987}
\pdciteshort{rw89}{Ramadge \& Wonham 1989, \textit{The Control of\ldots}}
\pdciteshort{ryoan16}{Hunt et al. 2016, \textit{Ryoan: a distributed\ldots}}
\pdciteshort{sagas1987}{Garcia-Molina \& Salem 1987, \textit{Sagas}}
\pdciteshort{saltzer1975protection}{Saltzer \& Schroeder 1975, \textit{The Protection\ldots}}
\pdciteshort{satterthwaite1975strategy}{Satterthwaite 1975, \textit{Strategy-proofness and\ldots}}
\pdciteshort{schaefer78}{Schaefer 1978, \textit{The Complexity of\ldots}}
\pdciteshort{schindler2019ethdkg}{Schindler et al. 2019, \textit{ETHDKG: Distributed\ldots}}
\pdciteshort{schmeidler1989}{Schmeidler 1989, \textit{Subjective Probability\ldots}}
\pdciteshort{schneider2000}{Schneider 2000, \textit{Enforceable Security\ldots}}
\pdciteshort{sclar2024}{Sclar et al. 2024, \textit{Quantifying Language\ldots}}
\pdciteshort{scone16}{Arnautov et al. 2016, \textit{SCONE: secure Linux\ldots}}
\pdciteshort{scott1998}{Scott 1998, \textit{Seeing Like a State: How\ldots}}
\pdciteshort{scott2012}{Scott 2012, \textit{Two Cheers for Anarchism}}
\pdciteshort{searle1969speech}{Searle 1969, \textit{Speech Acts: An Essay in the\ldots}}
\pdciteshort{sen1970impossibility}{Sen 1970, \textit{The Impossibility of a Paretian\ldots}}
\pdciteshort{sennrich2016}{Sennrich et al. 2016, \textit{Neural Machine\ldots}}
\pdciteshort{sep-hobbes}{Lloyd \& Sreedhar 2022, \textit{Hobbes's Moral and Political Philosophy}}
\pdciteshort{shamir1979share}{Shamir 1979, \textit{How to Share a Secret}}
\pdciteshort{shapiro2011crdt}{Shapiro et al. 2011}
\pdciteshort{sheng}{Sheng et al. 2021, \textit{BFT Protocol Forensics}}
\pdciteshort{sheridan1992}{Sheridan 1992, \textit{Telerobotics, Automation,\ldots}}
\pdciteshort{sheridanverplank1978}{Sheridan \& Verplank 1978, \textit{Human and\ldots}}
\pdciteshort{shinn2023reflexion}{Shinn et al. 2023, \textit{Reflexion: Language\ldots}}
\pdciteshort{shoham2009multiagent}{Shoham \& Leyton-Brown 2009}
\pdciteshort{simmons84}{Simmons 1984, \textit{The prisoners' problem and\ldots}}
\pdciteshort{sk00}{Stergiou \& Koubarakis 2000}
\pdciteshort{sleatortarjan1985}{Sleator \& Tarjan 1985}
\pdciteshort{slsa2023}{OpenSSF 2023, \textit{SLSA Specification v1.0:\ldots}}
\pdciteshort{sm03}{Sabelfeld \& Myers 2004, \textit{A model for\ldots}}
\pdciteshort{smith-qif}{Smith 2009, \textit{On the foundations of\ldots}}
\pdciteshort{smith1980cnp}{Smith 1980, \textit{The Contract Net Protocol:\ldots}}
\pdciteshort{spence1973}{Spence 1973, \textit{Job Market Signaling}}
\pdciteshort{spiffe2026jwtsvid}{SPIFFE 2026, \textit{JWT-SVID}}
\pdciteshort{spivak2012}{Spivak \& Kent 2012, \textit{Ologs: A Categorical\ldots}}
\pdciteshort{spkicert1999}{Ellison et al. 1999, \textit{SPKI Certificate\ldots}}
\pdciteshort{spkisdsi1999}{Ellison et al. 1999, \textit{SPKI Certificate\ldots}}
\pdciteshort{sqliteisolation}{Hipp et al. 2022, \textit{Isolation In SQLite}}
\pdciteshort{sqlitepragma}{Hipp et al. 2026, \textit{Pragma Statements\ldots}}
\pdciteshort{sqlitewal}{Hipp et al. 2010, \textit{SQLite Write-Ahead\ldots}}
\pdciteshort{srinivasan2026ddr}{Srinivasan 2026, \textit{Deployment Decision\ldots}}
\pdciteshort{ss09}{Sabelfeld \& Sands 2009}
\pdciteshort{steinke2023onerun}{Steinke et al. 2023, \textit{Privacy auditing with\ldots}}
\pdciteshort{strathern1997}{Strathern 1997, \textit{``Improving Ratings'':\ldots}}
\pdciteshort{tadelis2016}{Tadelis 2016, \textit{Reputation and Feedback\ldots}}
\pdciteshort{tailscale2020}{Pennarun et al. 2020, \textit{How Tailscale Works}}
\pdciteshort{theraulaz1999stigmergy}{Theraulaz \& Bonabeau 1999, \textit{A Brief\ldots}}
\pdciteshort{thompson1933}{Thompson 1933, \textit{On the Likelihood that One\ldots}}
\pdciteshort{thudi2021unlearning}{Thudi et al. 2021, \textit{On the necessity of\ldots}}
\pdciteshort{tirole1986}{Tirole 1986, \textit{Hierarchies and\ldots}}
\pdciteshort{torres2019intoto}{Torres-Arias et al. 2019}
\pdciteshort{treisman1980}{Treisman \& Gelade 1980}
\pdciteshort{truebit2017}{Teutsch \& Reitwiessner 2017, \textit{A Scalable\ldots}}
\pdciteshort{turpin2023}{Turpin et al. 2023, \textit{Language Models Don't\ldots}}
\pdciteshort{ucan2024}{UCAN Working Group 2024, \textit{UCAN 1.0}}
\pdciteshort{ucan2026revocation}{UCAN Working Group 2026, \textit{UCAN Revocation Specification}}
\pdciteshort{ucp2026}{Universal Commerce Protocol 2026, \textit{Specification}}
\pdciteshort{vaswani2017}{Vaswani et al. 2017, \textit{Attention Is All You\ldots}}
\pdciteshort{vaswani2017attention}{Vaswani et al. 2017, \textit{Attention Is All You\ldots}}
\pdciteshort{vdm07}{van der Meyden 2007, \textit{What, indeed, is\ldots}}
\pdciteshort{vickrey1961counterspeculation}{Vickrey 1961}
\pdciteshort{wald}{Wald 1945, \textit{Sequential tests of statistical\ldots}}
\pdciteshort{wald-wolfowitz}{Wald \& Wolfowitz 1948, \textit{Optimum character\ldots}}
\pdciteshort{wald1947}{Wald 1947, \textit{Sequential Analysis}}
\pdciteshort{wang2023agentsurvey}{Wang et al. 2025, \textit{A Survey on Large\ldots}}
\pdciteshort{wang2023voyager}{Wang et al. 2023, \textit{Voyager: An Open-Ended\ldots}}
\pdciteshort{wang2024tdag}{Wang et al. 2025, \textit{TDAG: A Multi-Agent\ldots}}
\pdciteshort{warm2008}{Warm et al. 2008, \textit{Vigilance Requires Hard\ldots}}
\pdciteshort{werbach2018blockchain}{Werbach 2018, \textit{The Blockchain and the New\ldots}}
\pdciteshort{wickens2002}{Wickens 2002, \textit{Multiple Resources and\ldots}}
\pdciteshort{williams1970}{Williams 1970, \textit{The Self and the Future}}
\pdciteshort{wilson1977insurance}{Wilson 1977, \textit{A Model of Insurance Markets\ldots}}
\pdciteshort{wrrw23}{Whitehouse et al. 2023}
\pdciteshort{xu2023gboard}{Xu et al. 2023, \textit{Federated learning of\ldots}}
\pdciteshort{xu2025amem}{Xu et al. 2025, \textit{A-Mem: Agentic Memory for\ldots}}
\pdciteshort{yaari1987}{Yaari 1987, \textit{The Dual Theory of Choice\ldots}}
\pdciteshort{yao2022react}{Yao et al. 2023, \textit{ReAct: Synergizing\ldots}}
\pdciteshort{yao2023tree}{Yao et al. 2023, \textit{Tree of Thoughts:\ldots}}
\pdciteshort{young2002landlord}{Young 2002, \textit{On-line file caching}}
\pdciteshort{young2019gvisor}{Young et al. 2019, \textit{The True Cost of\ldots}}
\pdciteshort{yue2022synthetic}{Yue et al. 2022, \textit{Synthetic text generation\ldots}}
\pdciteshort{zeng2024rag}{Zeng et al. 2024, \textit{The good and the bad:\ldots}}
\pdciteshort{zhan2026authmem}{Zhan et al. 2026, \textit{When Memory Becomes\ldots}}
\pdciteshort{zhang2026estimands}{Zhang 2026, \textit{What Does a Skill Actually Do?\ldots}}
\pdciteshort{zhang2026faultymemory}{Zhang et al. 2026, \textit{Useful Memories Become\ldots}}
\pdciteshort{zheng2023}{Zheng et al. 2023, \textit{Judging LLM-as-a-Judge\ldots}}
\pdciteshort{zheng2023judging}{Zheng et al. 2023, \textit{Judging LLM-as-a-Judge\ldots}}
\pdciteshort{zhong2025impossiblebench}{Zhong et al. 2025}
\pdciteshort{zhou2026formaljudge}{Zhou et al. 2026, \textit{FormalJudge: A\ldots}}
\pdciteshort{zhou2026utilitybehavior}{Zhou \& Ackerman 2026, \textit{When Preferences\ldots}}


% \pdcite{key} / \pdcite{key1,key2,...} -- the normal \cite, plus, in the
% Book, a margin note with one short form per key (semicolon-separated),
% each resolved through \pd@citeshort@one above. In a standalone chapter
% (\ifpdmargincolumn false) this is exactly \cite -- no margin note, so a
% standalone build never even needs the short-form table to be complete.
% \@for is the LaTeX kernel's own comma-list iterator (already available;
% no package of its own), run directly inside the \marginnote argument the
% same way \pdrecitation runs \begin{enumerate} inside its own \marginnote --
% marginnote saves an argument for later typesetting, so ordinary execution-
% time macros (\@for, \newif booleans) work inside it exactly as they would
% in running text.
\newif\ifpd@citefirst
\newcommand{\pdcite}[1]{%
  \cite{#1}%
  % The margin copy is off, and this is why. \pdcite is called mid-sentence,
  % 505 times, and every placement decision in this file is made from
  % \pagetotal -- which is only true in vertical mode, between paragraphs,
  % where the page builder has run. Mid-paragraph it is stale by however much
  % of the paragraph has been set, so these notes cannot know where they are,
  % and no cap computed from a wrong position is right: they printed on each
  % other 132 times and off the paper 7 times on the merge that brought them
  % in. Nothing is lost by leaving them out -- the \cite above puts the
  % citation in the text and the work in the bibliography.
  % What brings them back is an allocator rather than arithmetic: \marginpar,
  % which stacks notes in reading order, never overlaps and never runs off the
  % foot, and which has to own the whole column to do it (a \marginnote block
  % is invisible to the page builder, so a mixed column collides by
  % construction). That conversion is its own piece of work, with a fixture
  % document that puts each device in its pathological position, and it is
  % where the short-forms return. The short-form table and its freshness check
  % stay exactly as they are so nothing has to be rebuilt when they do.
  \relax}

% \pddefinedin{label} -- forward pointer to where a concept or theorem is formally
% constructed in the volume. Emits a cobalt indicator in the margin.
\newcommand{\pddefinedin}[1]{%
  \marginnote{\footnotesize\sffamily\color{pdcobalt}$\blacktriangleright$\,\textbf{Defined in} \S\ref{#1}}}

% \pdprov{script}{seed}{tag} -- a run's provenance at the point where its
% number appears in prose. tag is verified or internal, in the Book set in
% small caps to match the existing \Built/\Designed/\Vision tag convention
% (coordination-papers-mega-volume-preamble.tex and every chapter preamble
% define those the same way: \textsc for the tag word). In the Book this
% renders as a margin note "script \textperiodcentered\ seed \textperiodcentered\ TAG"; in a
% standalone chapter it folds back to exactly the bracket form the inline
% tag it replaces already used -- "[tag]" -- so no chapter's rendered prose
% changes outside the Book. Call sites drop the space that used to precede
% the inline "[tag]" (\pdprov attaches directly to the preceding word or
% math, e.g. "did\pdprov{...}{...}{verified}."), since the Book's margin
% note is not part of the running line and would otherwise leave a bare
% "did ." behind; the standalone branch below restores that one space
% itself (\ [#3]) so a standalone chapter still reads "did [verified]."
\newcommand{\pdprov}[3]{%
  \ifpdmargincolumn
    \sbox{\pd@sidebox}{\begin{minipage}[b]{\marginparwidth}\marginfont
      \texttt{\small #1} \textperiodcentered\ #2 \textperiodcentered\ \textsc{#3}\end{minipage}}%
    \pd@placemarginopt{\pd@sidebox}%
  \else
    \ [#3]%
  \fi}

% \pdprovedonentry{promise label}{target label}{section|page} -- one row of
% the discharge-pointer table: `promise label` is the label already on a
% theorem/lemma/proposition in Chapter 1 or Chapter 6's own source (the
% argument \pdprovedon below is called with); `target label` is the ALREADY
% namespaced label of whatever proves it in the discharging chapter, exactly
% as it will read once scripts/generate-mega-whitepaper.mjs's namespaceLabels
% has run on that chapter (i.e. "<chapter prefix>:<local label>") -- computed
% once, at generation time, by
% scripts/harbor-research/build_discharge_pointers.py, which knows every
% chapter's prefix from book/textbook.json. `\pdprovedon` itself is
% called with the PROMISE's own (un-namespaced) label -- the same string
% already sitting at its \label{...} call -- because namespaceLabels only
% rewrites the LABEL_COMMANDS it knows about, and \pdprovedon is not one of
% them, so the promise label reaching this macro is never rewritten and
% always matches the table's own key.
\newcommand{\pdprovedonentry}[3]{%
  \expandafter\gdef\csname pd@provedon@target@#1\endcsname{#2}%
  \expandafter\gdef\csname pd@provedon@kind@#1\endcsname{#3}}
% Plain \def, not \newcommand: \newcommand's \long prefix would make its
% \meaning differ from the plain (non-\long) macro \edef produces below for
% comparison, so \ifx would never see them as equal even with identical text.
\def\pd@provedon@sectionkind{section}
% GENERATED by scripts/harbor-research/build_discharge_pointers.py from a
% hand-curated (label-existence checked) list of discharge pairs -- do not
% hand-edit. Regenerate with:
%   python3 scripts/harbor-research/build_discharge_pointers.py
%
% One \pdprovedonentry{promise label}{namespaced target label}{section|thm|applied}
% per discharge pair \pdprovedon (figures/pd-pedagogy.tex) can point to. The
% target label is already namespaced as
% scripts/generate-mega-whitepaper.mjs's namespaceLabels will render it in
% the assembled Book ("<chapter prefix>:<local label>").
% Twin copy: book/figures/pd-discharges.tex and
% book/figures/pd-discharges.tex are kept
% byte-identical by this script (both preambles \input it via
% figures/pd-pedagogy.tex).

\pdprovedonentry{thm:regimentation-controllability}{anchor:sec:correct}{applied}
\pdprovedonentry{prop:conservation}{bonded:thm:conservation}{thm}

\newcommand{\pdprovedon}[1]{%
  \ifpdmargincolumn
    \ifcsname pd@provedon@target@#1\endcsname
      \edef\pd@provedon@kind{\csname pd@provedon@kind@#1\endcsname}%
      \edef\pd@provedon@target{\csname pd@provedon@target@#1\endcsname}%
      \sbox{\pd@sidebox}{\begin{minipage}[b]{\marginparwidth}\marginfont\itshape
        \ifx\pd@provedon@kind\pd@provedon@sectionkind
          Proved in \S\ref{\pd@provedon@target}.%
        \else
          Proved on p.~\pageref{\pd@provedon@target}.%
        \fi\end{minipage}}%
      \pd@placemarginopt{\pd@sidebox}%
    \fi
  \fi}

% ---------------------------------------------------------------------------
% The honesty ledger's status vocabularies -- ONE definition each, for the
% Book and for every standalone chapter.
%
% These used to be redefined in each chapter's own preamble, which is why they
% are here now. The Book is assembled by scripts/generate-mega-whitepaper.mjs,
% which keeps only what lies between \begin{document} and \end{document}: a
% chapter's preamble is DISCARDED at collation. So a chapter that defined its
% own \Built rendered "built" when it compiled alone and whatever the Book's
% preamble happened to say -- "implemented" -- when it compiled as a chapter.
% Same source, same call site, two different words, and nothing anywhere
% compared them. Defining them once, in a file both the Book and the
% standalone chapters \input, is what makes the two renderings the same
% object. scripts/harbor-research/check_duplicate_macros.py fails the build if
% a second definition of any of these names reappears.
%
% FOUR SEPARATE AXES. They answer different questions about the same claim and
% are never collapsed into one another; a claim normally carries one label
% from at most one of them, and carrying two (a Theorem that is only
% \Designed) is an ordinary state of affairs, not a contradiction.
%
%   1. MATURITY      how much of the mechanism exists in code today.
%   2. GUARANTEE     whether a property actually holds as configured.
%   3. OBLIGATION    whether a proof obligation has been discharged.
%   4. EVIDENCE      how far a stated result has been checked.
%
% The KIND of a claim -- Theorem / Design invariant / Model-checked property /
% Empirical hypothesis -- is a fifth axis with its own apparatus, \pdkind and
% the pdclaim environment below, whose vocabulary is closed and validated. A
% claim's kind is never spelled with a marker from the four scales here: a
% chapter that wants to say "this is a Theorem" writes \pdkind{Theorem}, not a
% maturity marker that happens to print the word.
%
% House style throughout: status reads by WEIGHT and SMALL-CAPS, never by hue.
% Two inks only -- hhink for a state that is asserted, hhgray for one that is
% not yet.

\ifpdswiss
  \newcommand{\Built}{{\textcolor{hhink}{\pdtextface\bfseries implemented}}}
  \newcommand{\BuiltWeak}{{\textcolor{hhink}{\pdtextface partial}}}
  \newcommand{\Designed}{{\textcolor{hhgray}{\pdtextface specified}}}
  \newcommand{\Vision}{{\textcolor{hhgray}{\pdtextface\itshape proposed}}}
\else
  \newcommand{\Built}{{\textcolor{hhink}{\textsc{\textbf{implemented}}}}}
  \newcommand{\BuiltWeak}{{\textcolor{hhink}{\textsc{partial}}}}
  \newcommand{\Designed}{{\textcolor{hhgray}{\textsc{specified}}}}
  \newcommand{\Vision}{{\textcolor{hhgray}{\textit{proposed}}}}
\fi
% ---------------------------------------------------------------------------
% Plain-language apparatus (the 2026-11 audit). Three semantic commands, so
% chapters say what a span IS and the rendering can change in one place.
%
% \pdterm[sort key]{term}  The defining occurrence of a term: set in bold and
%                          entered in the index. Use it once per term, at the
%                          sentence that defines the term in plain words, and
%                          before any other use of the term in reading order.
% \pdscope{text}           The one sentence that says what a claim does NOT
%                          show, in a fixed slot after the claim.
% \pdstatus{grade}{note}   Maturity of a design: grade is \Built, \BuiltWeak,
%                          \Designed or \Vision; note is a short phrase or
%                          empty. Placed beside a heading or claim title,
%                          never inside the definition's sentence.
%
% Where each one goes. \pdterm is bold at its point of use and enters the
% index (\makeindex is the Book preamble's; a standalone chapter's \index is
% the kernel no-op).
%
% Contents and trial replays. A \pdterm that reaches a contents line, a list
% of figures or a tabularx trial row is printed without indexing again: that
% flag belongs to figures/pd-book-citations.tex, which sets it for exactly
% those replays. It is declared here too, false, so a standalone chapter
% (which never loads that file) can test it; its own \newif re-declares it
% false before any use.
% ---------------------------------------------------------------------------
\@ifundefined{ifpdbookcitereplay}{\newif\ifpdbookcitereplay}{}
\ExplSyntaxOn
\int_new:N \g_pd_term_count_int
% The index entry is written where the term begins, so a term broken across
% a page turn is indexed on the page the reader meets it.
\cs_new_protected:Npn \pd_term_mark:nn #1#2
  {
    \mode_leave_vertical:
    \tl_if_blank:nTF {#1} { \index {#2} } { \index {#1@#2} }
    \int_gincr:N \g_pd_term_count_int
  }
\ProvideDocumentCommand \pdterm { O{} m }
  {
    \legacy_if:nF { pdbookcitereplay } { \pd_term_mark:nn {#1} {#2} }
    \textbf {#2}
  }
% The Book prints its index only when the previous pass defined at least one
% term; without this an unused \makeindex would print an empty "Index" page.
% The count travels in the .aux like any other cross-pass fact.
\cs_new:Npn \pdtermsdefinedlastpass {0}
\cs_new_protected:Npn \pdtermcountrecord #1 { \cs_gset:Npn \pdtermsdefinedlastpass {#1} }
\AddToHook { enddocument }
  {
    \legacy_if:nT { @filesw }
      {
        \iow_now:Nx \@auxout
          { \token_to_str:N \pdtermcountrecord { \int_use:N \g_pd_term_count_int } }
      }
  }
% Contents lines and bookmarks: a term is its text.
\NewExpandableDocumentCommand \pd@termpdfstring { O{} m } {#2}
\AtBeginDocument
  {
    \cs_if_exist:NT \pdfstringdefDisableCommands
      {
        \pdfstringdefDisableCommands
          { \cs_set_eq:NN \pdterm \pd@termpdfstring }
      }
  }
\ExplSyntaxOff
\ifpdswiss
  \providecommand{\pdscope}[1]{\par\noindent{\pdtextface\bfseries Scope.}\ #1\par}
  \providecommand{\pdstatus}[2]{\par\noindent{\pdtextface\footnotesize\textcolor{hhgray}{Status:}\ #1\if\relax\detokenize{#2}\relax\else\textcolor{hhgray}{; #2}\fi}\par}
\else
  \providecommand{\pdscope}[1]{\par\noindent\textsc{Scope.}\ #1\par}
  \providecommand{\pdstatus}[2]{\par\noindent{\footnotesize\textcolor{hhgray}{\textsc{Status:}}\ #1\if\relax\detokenize{#2}\relax\else\textcolor{hhgray}{; #2}\fi}\par}
\fi

% Chapter 5 (From Spawn to Person) spells the same four in capitals. Aliases,
% not second definitions -- there is still exactly one body per grade.
\newcommand{\BUILT}{\Built}
\newcommand{\BUILTWEAK}{\BuiltWeak}
\newcommand{\DESIGNED}{\Designed}
\newcommand{\VISION}{\Vision}

% Claim-kind words used inside ordinary theorem/definition bodies.  These are
% vocabulary aliases, not status marks: kind says what sort of claim this is;
% maturity says how much of its mechanism exists.  Keep them here so chapter
% modules never need private typography macros for the common honesty ledger.
\newcommand{\TheoremKind}{{\scshape theorem}}
\newcommand{\DesignInvariantKind}{{\scshape design invariant}}
\newcommand{\ModelCheckedKind}{{\scshape model-checked property}}
\newcommand{\EmpiricalHypothesisKind}{{\scshape empirical hypothesis}}

% 2. Guarantee. Marks a property that does NOT hold under the configuration
% actually shipped -- chapter 1's power-loss durability, throughout. It is the
% negative end of the ledger and prints the words it is named for.
\newcommand{\NotGuar}{{\textcolor{hhink}{\textit{not guaranteed}}}}

% 3. Proof obligation: discharged, partly discharged, outstanding, or carried
% only in a specification with no mechanised check behind it.
\ifpdswiss
  \newcommand{\Closed}{{\textcolor{hhink}{\pdtextface\bfseries closed}}}
  \newcommand{\Partial}{{\textcolor{hhink}{\pdtextface partial}}}
  \newcommand{\Open}{{\textcolor{hhgray}{\pdtextface\itshape open}}}
  \newcommand{\SpecOnly}{{\textcolor{hhgray}{\pdtextface\itshape spec-only}}}
  \newcommand{\Verified}{{\textcolor{hhink}{\pdtextface\bfseries verified}}}
  \newcommand{\Proved}{{\textcolor{hhink}{\pdtextface proved here}}}
  \newcommand{\Unproved}{{\textcolor{hhgray}{\pdtextface\itshape unproved}}}
\else
  \newcommand{\Closed}{{\textcolor{hhink}{\textsc{\textbf{closed}}}}}
  \newcommand{\Partial}{{\textcolor{hhink}{\textsc{partial}}}}
  \newcommand{\Open}{{\textcolor{hhgray}{\textit{open}}}}
  \newcommand{\SpecOnly}{{\textcolor{hhgray}{\textit{spec-only}}}}
  \newcommand{\Verified}{{\textcolor{hhink}{\textsc{\textbf{verified}}}}}
  \newcommand{\Proved}{{\textcolor{hhink}{\textsc{proved here}}}}
  \newcommand{\Unproved}{{\textcolor{hhgray}{\textit{unproved}}}}
\fi

% Verbatim-safe inline code for use INSIDE macro arguments (callouts,
% captions, TikZ nodes) where \verb is illegal. Detokenizes so special
% characters survive.
\newcommand{\code}[1]{{\ttfamily\detokenize{#1}}}

% ---------------------------------------------------------------------------
% \pdkind, pdclaim -- a numbered statement in the shared Lucide frame.
% ---------------------------------------------------------------------------
% Presentation only. Official Lucide assets: figures/lucide/LICENSE.
\makeatletter
\definecolor{pdblockProof}{HTML}{B33F35}
\definecolor{pdblockProperty}{HTML}{7048A5}
\definecolor{pdblockHypothesis}{HTML}{427A26}
\definecolor{pdblockCalculation}{HTML}{946000}
\definecolor{pdblockInvariant}{HTML}{233A76}
\definecolor{pdblockDefinition}{HTML}{3D454B}
\definecolor{pdblockChecked}{HTML}{006EA0}
\definecolor{pdblockProtocol}{HTML}{007D73}
\definecolor{pdblockNeutral}{HTML}{363B40}
\newcommand{\pdblockfamily}[1]{%
  \def\pdblock@family{#1}%
  \ifdefstring{\pdblock@family}{Conjecture}{\def\pdblock@family{Empirical hypothesis}}{}%
  \def\pdblock@color{pdblockNeutral}\def\pdblock@icon{file-text}%
  \ifdefstring{\pdblock@family}{Theorem}{\def\pdblock@color{pdblockProof}\def\pdblock@icon{scroll-text}}{}%
  \ifdefstring{\pdblock@family}{Lemma}{\def\pdblock@color{pdblockProof}\def\pdblock@icon{scroll-text}}{}%
  \ifdefstring{\pdblock@family}{Proposition}{\def\pdblock@color{pdblockProof}\def\pdblock@icon{scroll-text}}{}%
  \ifdefstring{\pdblock@family}{Corollary}{\def\pdblock@color{pdblockProof}\def\pdblock@icon{scroll-text}}{}%
  \ifdefstring{\pdblock@family}{Property}{\def\pdblock@color{pdblockProperty}\def\pdblock@icon{key-round}}{}%
  \ifdefstring{\pdblock@family}{Definition}{\def\pdblock@color{pdblockDefinition}\def\pdblock@icon{book-open}}{}%
  \ifdefstring{\pdblock@family}{Empirical hypothesis}{\def\pdblock@color{pdblockHypothesis}\def\pdblock@icon{flask-conical}}{}%
  \ifdefstring{\pdblock@family}{Numbers by hand}{\def\pdblock@color{pdblockCalculation}\def\pdblock@icon{calculator}}{}%
  \ifdefstring{\pdblock@family}{Design invariant}{\def\pdblock@color{pdblockInvariant}\def\pdblock@icon{lock-keyhole}}{}%
  \ifdefstring{\pdblock@family}{Model-checked property}{\def\pdblock@color{pdblockChecked}\def\pdblock@icon{list-checks}}{}%
  \ifdefstring{\pdblock@family}{Protocol}{\def\pdblock@color{pdblockProtocol}\def\pdblock@icon{workflow}}{}%
}
\newcommand{\pdblock@inlinebegin}{%
  \ifcsname pd@typesetpicturebox\endcsname
    \let\pgfsys@typesetpicturebox\pd@typesetpicturebox
  \fi}
\newcommand{\pdblockglyph}[1]{\begingroup
  \edef\pdblock@expanded{#1}\expandafter\pdblockfamily\expandafter{\pdblock@expanded}%
  \raisebox{-2.2pt}{\includegraphics[width=12pt,height=12pt]{figures/lucide/\pdblock@icon.pdf}}%
  \endgroup\hspace{.45em}}
% Every split segment keeps a full-strength 0.5pt perimeter for the Book inks.
% The retained 2% field and corner geometry are specific to this protocol slice.
% The colored upper-left cap
% is longer than its upright; the lower-right upright is longer than its cap.
\newcommand{\pdblock@frame}{%
  \draw[draw=\pdblock@color,line width=.5pt]
    (frame.south west) rectangle (frame.north east);
  \fill[\pdblock@color] (frame.north west) rectangle
    ([xshift=32pt,yshift=-3pt]frame.north west);
  \fill[\pdblock@color] (frame.north west) rectangle
    ([xshift=3pt,yshift=-16pt]frame.north west);
  \fill[\pdblock@color] (frame.south east) rectangle
    ([xshift=-16pt,yshift=3pt]frame.south east);
  \fill[\pdblock@color] (frame.south east) rectangle
    ([xshift=-3pt,yshift=28pt]frame.south east);}
\newcommand{\pdblockbefore}[1]{%
  \par\if@nobreak\else\Needspace{5\baselineskip}\fi
  \begingroup
  \edef\pdblock@expanded{#1}\expandafter\pdblockfamily\expandafter{\pdblock@expanded}%
  \pdblock@inlinebegin
  \begin{tcolorbox}[enhanced,breakable,sharp corners,frame hidden,
    colback=\pdblock@color!2!white,boxrule=.45pt,boxsep=0pt,
    left=8pt,right=8pt,top=7pt,bottom=7pt,
    before skip=8pt plus 2pt,after skip=8pt plus 2pt,
    pad at break*=7pt,lines before break=6,
    overlay unbroken={\pdblock@frame},overlay first={\pdblock@frame},
    overlay middle={\pdblock@frame},overlay last={\pdblock@frame}]
    % A block may split across pages, but its last line must not stand alone at the top of the next page
    % (printed pp. 200 and 347 in the 2026-10-14 edition); the club and broken rules stay relaxed so the
    % split itself remains possible.
    \clubpenalty=9000 \widowpenalty=10000
    \displaywidowpenalty=9000 \brokenpenalty=9000
}
\newcommand{\pdblockafter}[1]{\par\end{tcolorbox}\endgroup}
\makeatother

\makeatletter
\newcommand{\pd@claim@kinds}{Theorem,Design invariant,Model-checked property,Empirical hypothesis}
\newcommand{\pdkind}[1]{%
  \edef\pd@in@test{\noexpand\in@{,#1,}{,\pd@claim@kinds,}}%
  \pd@in@test
  \ifin@\else
    \PackageError{pd-pedagogy}{Unknown pdclaim kind '#1'}%
      {Use one of: Theorem, Design invariant, Model-checked property, Empirical hypothesis.}%
  \fi
  {\normalfont\bfseries
    \ifstrequal{#1}{Model-checked property}{Model-Checked Property}{%
      \ifstrequal{#1}{Empirical hypothesis}{Empirical Hypothesis}{%
        \ifstrequal{#1}{Design invariant}{Design Invariant}{#1}}}}%
}
% Match amsthm's bold Kind N (Title). grammar. Both claims and worked
% passages share the theorem sequence; step outside the font group so a
% following label resolves to this statement, not the surrounding section.
% Alias the register, not the public reference type. Delay until chapter
% theorem declarations and optional cleveref loading have completed.
\AtEndPreamble{%
  \ifcsname c@theorem\endcsname
    % Both names address the same TeX count register; section resets of
    % theorem therefore reset worked examples without a second counter.
    \let\c@pdworkedexample\c@theorem
    \let\cl@pdworkedexample\cl@theorem
    \let\p@pdworkedexample\p@theorem
    \def\thepdworkedexample{\thetheorem}%
    \def\theHpdworkedexample{\theHtheorem}%
    \providecommand{\pdworkedexampleautorefname}{Numbers by Hand}%
    \ifcsname crefname\endcsname
      \crefname{pdworkedexample}{Numbers by Hand}{Numbers by Hand Examples}%
      \Crefname{pdworkedexample}{Numbers by Hand}{Numbers by Hand Examples}%
    \fi
  \fi}
% End a heading paragraph without stranding it above the first body lines.
% Unlike a forced line break, this also restores ordinary paragraph layout.
\providecommand{\pdblockheadingend}{%
  {\raggedright\par}\nobreak\smallskip\@afterindentfalse\@afterheading\ignorespaces}
\newcommand{\pd@numberedhead}[3][theorem]{%
  \ifcsname c@#1\endcsname\refstepcounter{#1}\fi
  \def\@currentlabelname{#3}%
  {\raggedright\normalfont\bfseries #2%
    \ifcsname c@theorem\endcsname~\thetheorem\fi
    \ifstrempty{#3}{}{\hspace{.35em}(#3)}.\par}%
  \pdblockheadingend
}
\newenvironment{pdclaim}[2]{%
  % A final display can ship its page before the closing rule's nobreak.
  % Keep post-display material attached; ordinary line breaks remain free.
  % The environment group restores this penalty outside the claim.
  \postdisplaypenalty=10000\relax
  \pdblockbefore{#1}%
  \def\pdblockclaimkind{#1}%
  \noindent
  % A claim is a numbered object: it steps the chapter's shared theorem
  % counter (every chapter and the Book define `theorem` within section), so a
  % \label placed right after \begin{pdclaim} refers to the claim, prints as
  % "Theorem 3.1" in prose, and links to this head.
  \pdblockglyph{#1}%
  \pd@numberedhead{\pdkind{#1}}{#2}%
}{%
  \pdblockafter{\pdblockclaimkind}%
}

% ---------------------------------------------------------------------------
% pdboundary -- "Where this stops". A 1.2pt ink bar at the left, no fill; the
% head sits in the margin (Book) or runs in (standalone).
% ---------------------------------------------------------------------------
\newenvironment{pdboundary}{%
  % Let the page builder reserve the opening lines. A source-time
  % pagetotal/pagegoal comparison followed by newpage could ship the old
  % page while settling, then eject its two-line continuation as well.
  % The Book's margin allocator follows shipped anchors, not this guess.
  \par\if@nobreak\else\Needspace{4\baselineskip}\fi\addvspace{8pt plus 1pt}%
  \begin{mdframed}[linecolor=pdink,linewidth=1.2pt,leftline=true,topline=false,%
      rightline=false,bottomline=false,backgroundcolor=white,%
      innerleftmargin=9pt,innerrightmargin=0pt,innertopmargin=1pt,innerbottommargin=1pt,%
      skipabove=4pt,skipbelow=4pt]%
  \noindent\pd@marginhead{Where this stops}\ignorespaces
}{%
  \end{mdframed}%
}

% ---------------------------------------------------------------------------
% pdexample -- a worked passage, calculator icon and ochre enclosing frame.
% ---------------------------------------------------------------------------
\newenvironment{pdexample}[1]{%
  % Reserve the heading and its first two lines, not the whole example.
  % Starting a frame with less room sent mdframed's split retries past 100pt.
  % A section heading has already reserved its first text lines. Do not
  % force a second break immediately after it and strand the heading.
  \par\if@nobreak\else\Needspace{6\baselineskip}\fi\addvspace{8pt plus 1pt}%
  \pdblockbefore{Numbers by hand}%
  \noindent\pdblockglyph{Numbers by hand}%
  \pd@numberedhead[pdworkedexample]{Numbers by Hand}{#1}%
}{%
  \pdblockafter{Numbers by hand}%
}

% ---------------------------------------------------------------------------
% pdrecitation -- retrieval prompts. The body is a bare list of \item's. In
% the Book the prompts sit in the margin under a small-caps "Recall"; in a
% standalone chapter they are a compact numbered list. No box either way.
% ---------------------------------------------------------------------------
\newsavebox{\pd@recallbox}
% Long recall lists may use the body measure without disabling margin shipout.
% Toggling \pdmargincolumnfalse around breakable content can lose an earlier
% registered note when that content causes the page builder to ship a page.
\NewEnviron{pdrecitationbody}{%
  \par\addvspace{6pt}\noindent{\scshape Recall.}\par\nobreak
  \begin{enumerate}[leftmargin=1.6em,itemsep=1pt,topsep=2pt,parsep=0pt,label=\arabic*.]%
    \small\BODY
  \end{enumerate}%
  \ifx\pd@pendingpointer\@empty\else
    \par{\small\itshape\pd@pendingpointer}%
    \global\let\pd@pendingpointer\@empty
  \fi
  \addvspace{6pt}%
}
\NewEnviron{pdrecitation}{%
  \ifpdmargincolumn
    % Measured first, then RAISED so its last line sits on the line that issues
    % it: the block shares the section's height instead of hanging below it,
    % which is what keeps it on the page at a section end near the foot. Two
    % things bound the raise. It never rises above the top of the page. And it
    % never rises above the last margin head issued on this page (a Key idea,
    % Pitfall or Boundary head records where it was set): before that cap the
    % raise looked at nothing and printed the block over heads on five pages.
    % (A \marginpar was tried instead, 2026-09-08: the page builder keeps it
    % clear of other marginpars but never moves one UP, so ten blocks issued
    % near a page foot hung below the paper and their pointers were gone from
    % the Book with no warning. The raise stays.)
    \sbox{\pd@recallbox}{\begin{minipage}[b]{\marginparwidth}\raggedright
      {\scshape Recall}\par\nobreak\vspace{1pt}%
      \begin{enumerate}[leftmargin=1.1em,itemsep=1pt,topsep=0pt,parsep=0pt,label=\arabic*.]%
      \footnotesize\BODY\end{enumerate}%
      % A starred exercise pointer issued earlier in this section rides here as
      % the block's last line, instead of printing as its own note a line or two
      % above -- where this block, raised to share the section's height, would
      % print straight over it. That was every margin collision in the Book.
      \ifx\pd@pendingpointer\@empty\else
        \par\vspace{2pt}{\small\itshape\pd@pendingpointer}\fi
      \end{minipage}}%
    % Run the page builder before reading \pagetotal: after a \par the last
    % paragraph's lines are still in the recent contributions and \pagetotal
    % does not count them, so if that paragraph was what broke the page the
    % raise was computed for the page before -- a block rose 102 pt above the
    % paper on p. 300 that way. A penalty of 10000 followed by zero glue
    % contributes without opening a breakpoint.
    \par\penalty\@M\vskip\z@
    % Within two lines of a full page, the anchor is about to be carried to the
    % next page by the break that follows, and a raise computed against THIS
    % page's height would put the block above the next page's head (102 pt
    % above the paper on p. 300). Not raising is no better: the block then
    % hangs its whole height below the foot, which is where two of them were in
    % the Swiss and technical editions. Break the page here instead. The block
    % starts at the top of the next page, the raise is capped to nothing by the
    % page top, and the whole block is on the paper below its anchor. A
    % recitation ends a section, so the short page costs the argument nothing.
    % The old two-line test only noticed an anchor already at the foot.  A
    % tall Recall block can be issued considerably earlier and still have no
    % honest placement: its raised offset is computed while the enclosing
    % paragraph/float is still migrating, so the block may land at the foot
    % with most of itself below the paper.  We have already measured the box;
    % use that measurement.  When the remaining column cannot hold the whole
    % block plus one line of breathing room, advance the section-ending Recall
    % to the next page before asking the shared allocator to place it.
    % The shipped-page margin allocator owns this box. Its height must not
    % force a new main-text page (or create a page containing only Recall).
    \ifdefined\pdmarginregister\else
      \ifdim\dimexpr\pagetotal+\ht\pd@recallbox+\dp\pd@recallbox+\baselineskip\relax>\pagegoal
        \newpage\penalty\@M\vskip\z@\fi
    \fi
    \pd@placemargin{\pd@recallbox}%
    \global\let\pd@pendingpointer\@empty
  \else
    \par\addvspace{6pt}\noindent{\scshape Recall.}\par\nobreak
    \begin{enumerate}[leftmargin=1.6em,itemsep=1pt,topsep=2pt,parsep=0pt,label=\arabic*.]%
    \small\BODY\end{enumerate}\addvspace{6pt}%
  \fi}

% ---------------------------------------------------------------------------
% pdexercise -- numbered "1.7  CHECK *" as a hanging paragraph, kind and star
% rating via pgfkeys, the solution page in the margin (Book) or in a trailing
% parenthesis (standalone). The environment issues \label{<label>} itself;
% the exercise's own displayed number (\thepdexercise) is what \ref{<label>}
% later returns wherever it is cited.
% ---------------------------------------------------------------------------
\newcounter{pdexercise}
\renewcommand{\thepdexercise}{\pdchapternumber.\arabic{pdexercise}}
% hyperref's own \newcounter patch gives every counter a \theH<counter> PDF
% anchor name that, left at its default, is just \arabic{pdexercise} -- fine
% within one document, but pdexercise resets to 0 at every \pdchapter in the
% Book, so chapter 1's Exercise 1 and chapter 5's Exercise 1 would claim the
% SAME anchor name. Keying off \pdchapterprefix -- unique per chapter in both
% a standalone chapter and the Book -- keeps every exercise's anchor unique.
\renewcommand{\theHpdexercise}{\pdchapterprefix.\arabic{pdexercise}}
\pgfkeys{
  /pdexercise/kind/.store in=\pdexercise@kind,
  /pdexercise/rating/.store in=\pdexercise@rating,
}
\newcommand{\pd@exercise@kinds}{Check,Trace,Open}
\newcommand{\pd@exercise@ratings}{1,2,3}
\newcommand{\pdexercise@stars}{%
  \ifcase\pdexercise@rating\relax
  \or\ensuremath{\star}%
  \or\ensuremath{\star}\ensuremath{\star}%
  \or\ensuremath{\star}\ensuremath{\star}\ensuremath{\star}%
  \fi
}
\newenvironment{pdexercise}[2][]{%
  \def\pdexercise@kind{Check}%
  \def\pdexercise@rating{1}%
  \pgfkeys{/pdexercise/.cd,#1}%
  \edef\pd@in@test{\noexpand\in@{,\pdexercise@kind,}{,\pd@exercise@kinds,}}%
  \pd@in@test
  \ifin@\else
    \PackageError{pd-pedagogy}{Invalid pdexercise kind '\pdexercise@kind'}%
      {Use kind=Check, kind=Trace, or kind=Open.}%
  \fi
  \edef\pd@in@test{\noexpand\in@{,\pdexercise@rating,}{,\pd@exercise@ratings,}}%
  \pd@in@test
  \ifin@\else
    \PackageError{pd-pedagogy}{Invalid pdexercise rating '\pdexercise@rating'}%
      {Use rating=1, rating=2, or rating=3.}%
  \fi
  \def\pdexercise@label{#2}%
  \par\addvspace{7pt plus 1pt}\noindent
  % The counter anchor and label belong to the heading's horizontal list,
  % so a page break carries the destination and page reference with it.
  \refstepcounter{pdexercise}%
  \label{#2}%
  \ifpdmargincolumn\marginnote{\itshape Solution p.\,\pageref{sol:#2}}\fi
  {\bfseries\thepdexercise}\enspace{\footnotesize\scshape\pdexercise@kind}\,{\footnotesize\pdexercise@stars}\pdblockheadingend
}{%
  \ifpdmargincolumn\par\else\space{\small\itshape(solution on page~\pageref{sol:\pdexercise@label})}\par\fi
  \addvspace{2pt}%
}

% ---------------------------------------------------------------------------
% pdsolution / pdSolution -- the model answer, diverted to the open solution
% file (via the answers package) and reprinted later by \pdprintsolutions (a
% standalone chapter) or the Book's own per-chapter \Opensolutionfile in
% scripts/generate-mega-whitepaper.mjs.
%
% pdsolution (writing side, used at the chapter's point of use) takes the same
% <label> its \pdexercise used. pdSolution (reading side, what the collected
% file actually contains) is defined FIRST and BY HAND -- not via
% \Newassociation{pdsolution}{pdSolution}{pdsol} -- for two reasons: (a)
% \Newassociation's reading-side environment always takes the CURRENT
% \Currentlabel as its one argument, where pd-pedagogy needs the ORIGINAL
% exercise's <label> to carry through so \ref{<label>} and \pageref{<label>}
% resolve inside the solution box; (b) its writing-side signature takes no
% argument, so a bare \Newassociation could not accept \begin{pdsolution}{X}
% at all. Defining pdSolution first still lets a later \Newassociation
% elsewhere find "environment pdSolution already in use" and leave it alone,
% if this file is ever combined with one -- \newsolution's own guard.
%
% \Filesave{pdsol} (verbatim capture) is deferred behind \pd@sol@tmp rather
% than called directly inside the nested \Ifanswerfiles/\Iffileundefined/
% \Ifopen branches: called inline, the \else/\fi tokens still queued to close
% those wrapper conditionals land INSIDE the verbatim scan that follows and
% get written into the solutions file as if they were solution text.
% \Newassociation's own generated code defers it the same way (there: \Tmp);
% \endFilesave has no such side effect, so the closing side calls it inline.
% ---------------------------------------------------------------------------
\newenvironment{pdSolution}[1]{%
  \par\addvspace{7pt plus 1pt}\noindent
  \phantomsection\label{sol:#1}%
  \ifpdmargincolumn\marginnote{\itshape Exercise p.\,\pageref{#1}}\fi
  {\bfseries\ref{#1}}\ifpdmargincolumn\else\enspace{\small\itshape(exercise on page~\pageref{#1})}\fi\pdblockheadingend
}{%
  \par\addvspace{2pt}%
}
\newenvironment{pdsolution}[1]{%
  \let\pd@sol@tmp\relax
  \Ifanswerfiles{%
    \Iffileundefined{pdsol}{}{%
      \Ifopen{pdsol}{%
        \immediate\write\pdsol@file{\string\begin{pdSolution}{#1}}%
        \def\pd@sol@tmp{\Filesave{pdsol}}%
      }{}%
    }%
  }{%
    \def\pd@sol@tmp{\begin{pdSolution}{#1}}%
  }%
  \pd@sol@tmp
}{%
  \Ifanswerfiles{%
    \Iffileundefined{pdsol}{}{%
      \Ifopen{pdsol}{%
        \endFilesave
        \immediate\write\pdsol@file{\string\end{pdSolution}}%
      }{}%
    }%
  }{%
    \end{pdSolution}%
  }%
}

% ---------------------------------------------------------------------------
% pdsession -- "At the terminal": a captured transcript (a pd session, a model
% checker's output, a proof tool's trace), set in native monospace inside a
% continuous, breakable teal frame. The head is in the margin (Book) or run
% in (standalone). The
% optional argument is a short title. Transcripts are recorded by
% scripts/harbor-research/record_sessions.sh and checked in, never typed.
%   \begin{pdsession}[two clients, one claim]
%   $ pd claim src/auth.ts
%   ...
%   \end{pdsession}
% ---------------------------------------------------------------------------
% The convention of Kernighan's books and every serious manual since: what
% the reader types is bold (the prompt line), what the program prints is
% regular. The official Lucide terminal icon identifies captured output.
% The shared semantic frame retains a full thin perimeter and asymmetric
% teal corner accents on every page segment, without implying a live connection.
% Synthetic records retain their separate continuous-feed-paper apparatus.
\lstdefinestyle{pdsession}{basicstyle=\ttfamily\fontsize{8.6}{10.6}\selectfont,columns=fullflexible,numbers=none,backgroundcolor={},
  keepspaces=true,breaklines=true,breakatwhitespace=false,breakindent=2em,frame=none,
  postbreak=\mbox{\textcolor{hhgray}{\fontsize{7}{7}\selectfont$\hookrightarrow$}\space},
  xleftmargin=0pt,aboveskip=4pt,belowskip=2pt,showstringspaces=false,
  language={},keywordstyle={},commentstyle={},stringstyle={},identifierstyle={},
  % This installation's t1lmtt.fd declares bx/n only as a silent ssub to b/n
  % (see "bold is always bold (b)" in that file) -- \bfseries requests
  % \bfdefault (bx) and still triggers NFSS's substitution machinery.
  % \fontseries{b} asks directly for the shape that's actually defined.
  moredelim=[l][\fontseries{b}\selectfont]{\$\ }}
% Opt in per recorded excerpt; long unbroken rules still need the default character wrapping.
\newif\ifpdsessionwordwrap
\newlength{\pd@sessionwidth}
\newlength{\pd@wideoffset}
\newcommand{\pd@measurewidefield}{%
  \setlength{\pd@sessionwidth}{\linewidth}%
  \setlength{\pd@wideoffset}{0pt}%
  % Terminal and sample-data blocks leave the margin free for explanation.
}
\lstnewenvironment{pdsession}[1][]{%
  \endgraf
  \begingroup
  % Select the shared semantic frame's teal locally. Other families and the
  % synthetic feed-paper listing retain their own colors and apparatus.
  \let\pd@sessionfamily\pdblockfamily
  \def\pdblockfamily##1{%
    \pd@sessionfamily{##1}\def\pdblock@color{pdteal}}%
  \pdblockbefore{Terminal transcript}%
  % listings selects its verbatim terminator after this begin code. Restore
  % the public environment name while the shared tcolorbox is open.
  \def\@currenvir{pdsession}%
  \noindent\raisebox{-2pt}{\includegraphics[width=12pt,height=12pt]{figures/lucide/terminal.pdf}}%
  \endgraf\nobreak
  % Measure AFTER the shared frame has applied its inset. The outer frame
  % stays at the Book's 4.5in measure; listings wraps at native 8.6pt inside.
  \pd@measurewidefield
  \lstset{style=pdsession,upquote=true,linewidth=\pd@sessionwidth,xleftmargin=0pt,
    % Disable the Book's prose quote mapping only for actual transcripts.
    % The native point size and leading remain unchanged; feed-paper does not
    % inherit this environment-local override.
    basicstyle=\ttfamily\ifdefined\addfontfeatures\addfontfeatures{Mapping=}\fi\fontsize{8.6}{10.6}\selectfont,
    literate={-}{{\char45}}1,
    caption={#1}}%
  \ifpdsessionwordwrap\lstset{breakatwhitespace=true}\fi
}{%
  \def\@currenvir{tcolorbox}%
  \pdblockafter{Terminal transcript}%
  \endgroup
}

% ---------------------------------------------------------------------------
% pdsampledata -- synthetic or illustrative records on continuous-feed paper.
% Unlike pdsession, this apparatus is explicitly NOT a captured run.  Its
% provenance is printed in the head and its tractor-feed edges keep a reader
% from mistaking invented input/output for a theorem, proof trace, or live
% evidence.  The optional argument names the source or seed.
%   \begin{pdsampledata}[synthetic; settings race, seed 20260917]
%   09:14:03 agent-a writes settings.json
%   ...
%   \end{pdsampledata}
% ---------------------------------------------------------------------------
\lstdefinestyle{pdsampledata}{style=pdsession,
  basicstyle=\ttfamily\fontsize{8.6}{10.6}\selectfont,
  backgroundcolor={},
  xleftmargin=0pt,xrightmargin=0pt,aboveskip=2pt,belowskip=2pt,
  frame=none}
\newtcblisting{pdsampledata}[1][]{%
  enhanced,listing only,listing options={style=pdsampledata},
  width=\linewidth,
  boxrule=.45pt,arc=0pt,outer arc=0pt,
  colframe=pdinkmuted!62,colback=pdinkmuted!5,
  left=16pt,right=16pt,top=5pt,bottom=5pt,
  before={\par\addvspace{10pt plus 1pt}\pd@measurewidefield
    \noindent\hspace*{\pd@wideoffset}},
  after={\par\addvspace{10pt plus 1pt}},
  title={SAMPLE DATA\ifstrempty{#1}{}{\enspace\textcolor{pdinkmuted}{#1}}},
  fonttitle=\pdgrotesk\footnotesize\bfseries,coltitle=pdink,
  colbacktitle=pdinkmuted!8,titlerule=0pt,
  toptitle=4pt,bottomtitle=4pt,
  overlay unbroken={%
    \fill[pdinkmuted!13]
      (frame.north west) rectangle ([xshift=10pt]frame.south west);
    \fill[pdinkmuted!13]
      ([xshift=-10pt]frame.north east) rectangle (frame.south east);
    \draw[pd hairline]
      ([xshift=10pt]frame.north west) -- ([xshift=10pt]frame.south west);
    \draw[pd hairline]
      ([xshift=-10pt]frame.north east) -- ([xshift=-10pt]frame.south east);
    \foreach \i in {1,...,10}{%
      \pgfmathsetmacro{\pd@holet}{\i/11}%
      \fill[pdpage]
        ($(frame.north west)!\pd@holet!(frame.south west)+(5pt,0)$)
        circle[radius=1.35pt];
      \fill[pdpage]
        ($(frame.north east)!\pd@holet!(frame.south east)+(-5pt,0)$)
        circle[radius=1.35pt];
    }%
  }}

% ---------------------------------------------------------------------------
% \pdassurance -- the five-word assurance-mode tag (matches the vocabulary in
% coordination-papers-mega-volume-appendices.tex's implementation boundary).
% ---------------------------------------------------------------------------
\newcommand{\pd@assurance@modes}{Observed,Coordinated,Brokered,Confined,Attested}
\newcommand{\pdassurance}[1]{%
  \edef\pd@in@test{\noexpand\in@{,#1,}{,\pd@assurance@modes,}}%
  \pd@in@test
  \ifin@\else
    \PackageError{pd-pedagogy}{Unknown pdassurance mode '#1'}%
      {Use one of: Observed, Coordinated, Brokered, Confined, Attested.}%
  \fi
  {\scshape#1}%
}

% ---------------------------------------------------------------------------
% \pdopensolutions / \pdprintsolutions -- standalone-chapter chrome only. The
% Book generator strips both (scripts/generate-mega-whitepaper.mjs,
% cleanStandaloneChrome) and opens/closes the solution stream itself, once per
% chapter, under its own book-sol-<prefix> filename.
% ---------------------------------------------------------------------------
\newcommand{\pdopensolutions}{\Opensolutionfile{pdsol}[sol-\pdchapterprefix]}
\newcommand{\pdprintsolutions}{%
  \Closesolutionfile{pdsol}%
  \section*{Solutions to the exercises}%
  \IfFileExists{sol-\pdchapterprefix.tex}{\input{sol-\pdchapterprefix}}{}%
}
% ---------------------------------------------------------------------------
% Chapter-end exercise sections (Wave 12). The body keeps a pointer where a
% cluster of exercises used to stand; the exercises themselves live in the
% chapter's closing \section{Exercises}, grouped under \pdexercisesfor heads.
% Arguments are pre-wrapped in \ref/\pageref by the relocation script so the
% Book generator's label namespacing sees them.
%   \pdexercisesfor{\S\ref{sec:x}}{Section title}
%   \pdexercisepointer{\ref{ex:first}}{\ref{ex:last}}{\pageref{ex:first}}
%   \pdexercisepointerone{\ref{ex:only}}{\pageref{ex:only}}
% ---------------------------------------------------------------------------
\newcommand{\pdexercisesfor}[2]{%
  \Needspace*{7\baselineskip}% room for the label and the head and first lines of the exercise after it; a label with nothing under it is a stranded heading (printed p. 222, 2026-10-14 edition)
  \par\vspace{12pt}\noindent{\scshape #1\enspace #2}\par\nobreak\smallskip}% \nobreak glues the label to the first exercise
% In the Book the pointer sits in the margin beside the paragraph it follows;
% in a standalone chapter it is a small italic line.
%
% The starred forms are for a pointer that a pdrecitation follows in the same
% section (61 pointers in the Book; 24 are in that position). They print
% nothing where they stand and hand their text to the Recall block, which sets
% it as its last line. Which pointers are starred is decided in the sources by
% scripts/harbor-research/star_exercise_pointers.py from the same measurement,
% not by hand. A starred pointer that reaches the next heading unconsumed is
% an error, not a silently missing line: the hook below checks at every
% \section and \subsection, and at the end of the document.
\def\pd@pendingpointer{}
% Append, never overwrite: two sections' pointers can precede one Recall
% block (a subsection with its own exercises, then the section's), and the
% first cut of this used \gdef, so the block carried the second and the first
% vanished from the Book with no error -- the section hook saw an empty macro
% and passed. Two of sixty-one pointers were missing from the page before the
% count found them. Each pointer is its own line.
\newcommand{\pd@addpointer}[1]{%
  \ifx\pd@pendingpointer\@empty
    \gdef\pd@pendingpointer{#1}%
  \else
    \g@addto@macro\pd@pendingpointer{\par #1}%
  \fi}
\NewDocumentCommand{\pdexercisepointer}{s m m m}{%
  \ifpdmargincolumn
    \IfBooleanTF{#1}%
      {\pd@addpointer{Exercises #2--#3, p.~#4.}}%
      {\sbox{\pd@ptrbox}{\begin{minipage}[b]{\marginparwidth}\marginfont
         \small\itshape Exercises #2--#3, p.~#4.\end{minipage}}\pd@placemarginopt{\pd@ptrbox}}%
  \else
    \par\smallskip\noindent{\small\itshape Exercises #2--#3, page~#4.}\par
  \fi}
\NewDocumentCommand{\pdexercisepointerone}{s m m}{%
  \ifpdmargincolumn
    \IfBooleanTF{#1}%
      {\pd@addpointer{Exercise #2, p.~#3.}}%
      {\sbox{\pd@ptrbox}{\begin{minipage}[b]{\marginparwidth}\marginfont
         \small\itshape Exercise #2, p.~#3.\end{minipage}}\pd@placemarginopt{\pd@ptrbox}}%
  \else
    \par\smallskip\noindent{\small\itshape Exercise #2, page~#3.}\par
  \fi}
\newcommand{\pd@checkpendingpointer}{%
  \ifx\pd@pendingpointer\@empty\else
    \PackageError{pd-pedagogy}%
      {A starred exercise pointer (\pd@pendingpointer) was not followed by a pdrecitation in its section}%
      {Either add the Recall block after it or remove the star so the pointer prints on its own.}%
  \fi}
\AtBeginDocument{\ifpdmargincolumn
  \pretocmd{\section}{\pd@checkpendingpointer}{}{}%
  \pretocmd{\subsection}{\pd@checkpendingpointer}{}{}%
  \AtEndDocument{\pd@checkpendingpointer}\fi}
% ---------------------------------------------------------------------------
% The older tinted-box macros (\keyidea, \pitfall, \scene, \xrefbox,
% \pullquote) are still defined, as tikz boxes, in every chapter preamble and
% in the Book's preamble AFTER this file is read. They are re-typeset here at
% \begin{document} in the same page grammar -- a run-in or margin head and
% plain prose, no fill -- so no chapter source and no preamble has to change
% for the boxes to disappear. \long\def so a body may contain \par.
%
% \pdthesis was \pullquote, and the rename is a correction, not a preference.
% A pull quote is by definition an excerpt LIFTED from the body text on the
% same spread and set larger to give a skimming reader a way in; it repeats
% something the page already says, and it is attributed to the page. Not one
% of the twenty-one so-called pull quotes in this book is a quote of anything.
% Each is a sentence written for that slot and appearing nowhere else, set in
% large italic behind a coloured rule -- which is the costume of a quotation
% with no speaker behind it, and reads as gravitas the sentence has not
% earned. What every one of them actually is, is the section's claim said once
% and memorably: its thesis. So it is called that, and set the way this book
% sets every other kind of content -- a margin head and plain prose, no rule,
% no fill, no borrowed authority. \pullquote survives as a deprecated alias so
% nothing breaks mid-flight; check_style_sections.py fails on new uses.
% ---------------------------------------------------------------------------
\AtBeginDocument{%
  % Optional heads name the actual consequence, not a generic teaching role.
  % An unedited call remains ordinary prose; it never prints an empty label.
  \RenewDocumentCommand\keyidea{O{} +m}{\par\pd@settle\addvspace{6pt plus 1pt}\noindent\ifstrempty{#1}{}{\pd@marginhead{#1}}#2\par\addvspace{6pt plus 1pt}}%
  \RenewDocumentCommand\pitfall{O{} +m}{\par\pd@settle\addvspace{6pt plus 1pt}\noindent\ifstrempty{#1}{}{\pd@marginhead{#1}}#2\par\addvspace{6pt plus 1pt}}%
  \long\def\scene#1#2{%
    \par\pd@settle\addvspace{8pt plus 2pt}%
    \noindent\begin{tcolorbox}[
      enhanced, breakable, sharp corners,
      colback=white, colframe=black,
      boxrule=0.5pt, leftrule=3.2pt,
      left=9pt, right=9pt, top=7pt, bottom=7pt,
      before skip=6pt, after skip=6pt
    ]%
    \noindent{\small\raisebox{-2pt}{\IfFileExists{figures/icons/mira-silhouette.pdf}{\includegraphics[height=1.35em]{figures/icons/mira-silhouette.pdf}}{\IfFileExists{figures/icons/mira-silhouette.png}{\includegraphics[height=1.35em]{figures/icons/mira-silhouette.png}}{}}}\enspace{\scshape\color{black}Scene}\enspace\textperiodcentered\enspace{\bfseries #1}\pdblockheadingend{\itshape #2}}%
    \end{tcolorbox}\par\addvspace{8pt plus 2pt}%
  }%
  \long\def\xrefbox#1{\par\pd@settle\addvspace{4pt}\noindent\pd@marginhead{See also}{\small #1}\par\addvspace{4pt}}%
  % A generic "Thesis" margin label does not identify which claim matters.
  % Keep the prose; use an authored keyidea head only when it adds meaning.
  \long\def\pdthesis#1{\par\pd@settle\addvspace{6pt plus 1pt}\noindent#1\par\addvspace{6pt plus 1pt}}%
  \long\def\pullquote#1{\pdthesis{#1}}%
}
\makeatother

% Shared visual language for the Book's figures -- version 2.
% Twin: book/figures/pd-figure-language.tex (byte-identical; edit both).
% The standard behind it, the migration from v1 and the QA loop live in
% skills/tikz-diagram-craft/ (SKILL.md, references/figure-standard.md).
%
% This file defines voices, hues, weights and marks -- not layouts.  Every
% figure still chooses a grammar from its reader question.
%
% ===========================================================================
% 1.  THE TYPE LADDER.  Three voices, one size, each tied to a role.
%       pd title   bold      panel titles, row and actor heads, column heads
%       pd label   upright   node text, state names, axis titles, tick words,
%                            relation words on an edge, badge numerals
%       pd note    italic    the one annotation a figure is allowed
%     `pd tag` is the label voice knocked out of the page, for a word that
%     must sit near a rule.  `pd kind` (small caps) is a one- or two-word kind
%     tag, never a title.  Every piece of figure text is mixed case; \texttt
%     only for identifiers that are literally code; math mode for symbols.
%
% 2.  THE COLOUR SYSTEM.  Hue is meaning, and the meanings cross chapters.
%     The figure's own subject is drawn in the CHAPTER hue (`pd focus ...`),
%     resolved from \pdcurrentchaptercolor, which the Book sets per chapter;
%     a fragment declares its standalone fallback with \pdfigurehue{...}.
%     Anything else the figure names takes its concept hue:
%       pd truth       pdcobalt   kernel, truth, the thing being specified
%       pd legible     pdteal     legibility, digest, inspection
%       pd ready       pdhealth   ready, coordinated, admitted, passing
%       pd protocol    pdindigo   protocol, federation, the wire
%       pd identity    pdviolet   identity, continuity, a card, an ID
%       pd reputation  pdrust     reputation, trust earned over time
%       pd value       pdgold     economy, value, settlement, payment
%       pd breach      pderror    breach, refusal, revocation, correction
%       pd warn        pdamber    warning; rules and stripes only, never text
%     Each concept X yields X, X rule, X arrow, X state, X fill, X datum and
%     X label.  Two to four hues a figure; colour is always doubled by shape,
%     dash, weight or a word (breach and warn are dashed as well as coloured).
%
% 3.  THE WEIGHT LADDER.  Exactly three weights, .5 / .9 / 1.6 pt
%     (1 : 1.8 : 3.2).  .5: hairline, guide, tick, artifact edge, fill edge,
%     mark outline, badge.  .9: rule, arrow, state edge, concept arrow.
%     1.6: spine, focus rule, focus and concept state edges, concept rule.
%     A reader sees which line is the subject.  Tips scale with the stroke.
%
% 4.  THE SURFACE LADDER.  pd state (page ground, ink edge) < pd artifact
%     (page ground, grey edge, subordinate) < pd focus state and the concept
%     states (24 % tint, edged in the same hue).  Text is always ink.  A
%     terminal state has a double boundary; an automaton's states are circles
%     (`pd automaton state/focus/terminal`); a ledger's pools are buckets
%     joined by pipes (`pd bucket`, `pd pipe`); a trust boundary is a dashed
%     rectangle (`pd trust boundary`).
%
% 5.  THE GROUND.  Body pages print on white, so every knockout and every
%     neutral surface is `pdpage` (white), never cream: a cream backing on a
%     white page reads as a sticker, which is what v1's figures were covered
%     in.  Neutral regions (Gantt bands, wells) are a warm grey tint.
%
% 6.  EDITIONS.  The Book sets \pdedition (maritime | swiss | technical) and
%     figures/pd-figure-language-<edition>.tex restyles the same names, so one
%     drawing renders in each edition's character.  A fragment that names a
%     raw colour (hhteal, pdcobalt!40, blue) bypasses this; name a role.
% ===========================================================================
% Libraries every v2 mark needs (diamonds, braces, markings), and
% pgfplots for computed plots. This file is \input in every preamble (the
% Book's and each standalone chapter's), so a plot works in all of them.
\usetikzlibrary{arrows.meta,matrix,patterns.meta,positioning,calc,fit,
  backgrounds,shapes.geometric,shapes.multipart,shapes.misc,decorations.pathreplacing,
  decorations.markings,quotes,intersections}
\RequirePackage{pgfplots}
\pgfplotsset{compat=1.18}
\usepgfplotslibrary{fillbetween,groupplots}

\colorlet{pdpage}{white}

% --- the chapter hue, resolved late --------------------------------------
% \pdfigurehue{pdindigo} on the line before a tikzpicture names the hue the
% figure is drawn in when it is compiled outside the Book (a standalone
% chapter, or the skill's renderer).  Inside the Book the preamble has already
% defined \pdcurrentchaptercolor and the chapter opener sets it, so there the
% opener wins and \providecommand does nothing.  Local to the figure's group.
\newcommand{\pdfigfraghue}{pdcobalt}
\newcommand{\pdfigurehue}[1]{%
  \renewcommand{\pdfigfraghue}{#1}%
  \providecommand{\pdcurrentchaptercolor}{#1}}
\newcommand{\pdsetfocushue}{%
  \ifdefined\pdcurrentchaptercolor
    \colorlet{pdfocus}{\pdcurrentchaptercolor}%
  \else
    \colorlet{pdfocus}{\pdfigfraghue}%
  \fi}
\colorlet{pdfocus}{pdcobalt}

% ===========================================================================
% THE TYPOGRAPHIC LAW
% ===========================================================================
% Companion to the line-weight law in figures/pd-palette.tex (1px texture,
% 1.5px linework, 2px enclosure). Same shape of rule: name the roles, fix the
% values, one place.
%
% ONE SIZE. Every named text role in a figure is \pdfiglabelsize
% (= \footnotesize). Roles separate by WEIGHT, SLOPE, FAMILY and INK -- never
% by size. A fragment never writes a size command of its own; if a distinction
% needs making, it takes the role that makes it.
%
% Where the size comes from. Measured, not chosen:
%   \tiny  \scriptsize  \footnotesize  \small  \normalsize
%   5.23    6.97          8.72           9.55    10.46  pt   Book (10.5pt
%                                                            Pagella, 4.5in
%                                                            column, 7x10in)
%   5.98    7.97          8.97           9.96    10.91  pt   standalone chapter
%                                                            (11pt Latin Modern)
% \scriptsize renders at 6.97pt in the Book -- under figcheck's 7pt T1 floor
% before a single width promotion, so it is illegal by measurement, not taste.
% \footnotesize holds 8.72pt and still clears the floor at the 0.85 \resizebox
% limit craft-rules sets (7.41pt). It is also what the corpus already voted
% for: 78 of the 91 local font= declarations this law replaces were already
% \footnotesize, \footnotesize\bfseries or \footnotesize\itshape.
%
% \pdfigbasesize (= \small) is the tikzpicture's ambient size only, so an
% unstyled node is never SMALLER than a styled one; every named role steps
% down exactly one notch from it and none steps down twice.
%
% ONE FACE, AND IT IS THE DOCUMENT'S. \pdfiglabelfamily is empty by default,
% which means figure type inherits whatever face the document is setting --
% Palatino in the Book, Computer Modern in a standalone chapter. Figures
% agreeing with EACH OTHER is not the goal; figures agreeing with the
% paragraph beside them is. A fragment therefore never names a family. Not
% \sffamily (which resolves to Latin Modern Sans against a Computer Modern
% page and reads as a foreign object on it), not \rmfamily, not
% \fontfamily{...}. P18 in tikz_precheck.py is an error for exactly that.
% An edition MAY substitute a face -- the Swiss and technical editions set
% \pdfiglabelfamily to \pdgrotesk, because grotesk display is their declared
% character -- and it does so in ONE command in ONE file, so the substitution
% is a decision on the record rather than an accident in a fragment. The
% only family a fragment may ask for is the identifier one, and it asks for
% it by role (pd mono label / \pdfiglabelmono), never by \ttfamily.
%
% What the roles are, and what each one is for:
%   pd panel title        bold        the name of a panel
%   pd row label          bold        an actor/row name east of its lane
%   pd reverse row label  bold, cream a row name sitting on a dark band
%   pd axis label         regular     an axis title or a tick numeral
%   pd direct label       regular     a value or name riding its own mark
%   pd reverse label      regular, cream  a direct label on a dark band
%   pd mono label         mono        an identifier, path, or literal
%   pd verdict            bold        a terminal outcome stamp (CLEAR, REJECT)
%   pd note               italic      one aside per figure
%   pd legend             regular     a pgfplots legend entry
%   pd state / terminal / actor / artifact / decision   node text
%   pd axis               the whole pgfplots typographic block, one key
%   \pdfigsub             a gloss line under a heading INSIDE a node
%   \pdfigmath            math carrying a sub/superscript, inside a label
%
% A fragment that needs a compound style of its own composes it from a role
% above (`gate/.style={pd decision,minimum width=2cm}`), or, where no role
% fits, from the \pdfiglabel* handles here. It never writes a raw size
% command: P16 in skills/harbor-chartwork/scripts/tikz_precheck.py is an
% error for exactly that, and P10/P11 have been errors for \tiny and
% \scriptsize since before this law.
%
% What does NOT transfer from the screen rules these roles otherwise follow.
% The dataviz method those rules come from assumes an interactive chart: a
% tooltip carries any value a label could not fit, a filter row rescopes the
% view, hover lifts the hovered mark. A printed page has none of that, and the
% absence is not a defect to design around. The INTENT of the tooltip rule --
% "every value is reachable without gating" -- does transfer, and on this page
% it is discharged somewhere else: the value is in the caption, in the margin
% apparatus, or in the chapter's own worked example. So a label that will not
% fit inside its mark moves outside the mark or is dropped to the caption; it
% is never shrunk to fit, because there is no size below \footnotesize to shrink
% to, and never clipped. Do not add a hover affordance, a legend toggle, or a
% filter to a figure in this book.
%
% IDENTIFIERS GET ONE SPELLING. A key, a cell id, a file name or any other
% literal a reader could type is set in `pd mono label` (or \pdfiglabelmono
% inside a compound style), spelled exactly as the code spells it: sk_A, not
% skA and not sk with a subscript A; card0, not card-nought and not a Unicode
% subscript zero. A figure that shows the same identifier two ways has told
% the reader they are two things. The subscripted form belongs to MATHEMATICS
% ($\kappa_7$, an index into a sequence) and the literal form to CODE -- one
% figure may use both, for different objects, but never both for the same
% one. (A precheck for two spellings of one identifier inside a fragment is
% being added on another branch; this is the convention it will enforce.)
%
% SHAPE IS A CHANNEL, AND IT HAS TO SURVIVE THE PRINT SIZE. pd datum, pd
% focus datum and pd caution datum carry three different forms, because the
% house palette is three low-chroma print inks and every distinction it makes
% has to be doubled by something that is not colour. That only counts if the
% form is still legible at 2.1-2.7 pt: an edition that flattens all three to
% one form has removed the second channel and left the figure separating its
% classes by hue alone, which fails in greyscale and for a colour-blind
% reader. Any edition override that restyles a datum must keep three
% distinguishable forms, and must be checked on a render at 1.0x, not
% reasoned about from the point sizes.
%
% One screen rule that does NOT bind here, and why. "Text never wears the data
% colour" exists because a light categorical hue is illegible as text on a
% light surface. The house figure hues are print inks, not screen hues --
% hhink #1B1712, hhteal #00564C, hhamber #6B4500 -- and all three clear text
% contrast on cream, so a label MAY wear its series' ink to bind itself to a
% curve, which is direct labelling and saves a legend. The rule that survives
% is the reason behind it: text wears one of those three inks at full
% strength, or hhink!80!hhgray for recessive furniture, and nothing else. A
% tint of an ink (text=hhgray, text=hhteal!60) is not a role; it is a fade.
% ===========================================================================

\providecommand{\pdfiglabelfamily}{}            % an edition may set \pdgrotesk
\providecommand{\pdfiglabelsize}{\footnotesize} % every named text role
\providecommand{\pdfigbasesize}{\small}         % the picture's ambient size
\providecommand{\pdfiglabel}{\pdfiglabelfamily\pdfiglabelsize}
\providecommand{\pdfiglabelbold}{\pdfiglabel\bfseries}
\providecommand{\pdfiglabelitalic}{\pdfiglabel\itshape}
\providecommand{\pdfiglabelmono}{\pdfiglabelsize\ttfamily}
% The one role that is not a node style, because it lives INSIDE a node's text:
% the second line under a heading that glosses it ("reviewer role / obligation +
% capability + authority"). Every fragment that needed it reached for \tiny or
% \scriptsize, which is how seven \tiny and twelve \scriptsize survived P10/P11
% as hard errors in the corpus. It steps down by SLOPE and WEIGHT, not size, so
% it stays legible and stays legal, and it touches neither family nor ink --
% which is what lets it sit on a reversed node without vanishing.
\providecommand{\pdfigsub}{\mdseries\itshape}
% The one place the law's own size is not the right size.
%
% Measured, both preambles, on a subscripted label inside `pd direct label`:
%
%              base glyph   subscript   subscript as % of base
%   chapter      8.97pt       5.98pt            66.7%
%   Book         8.72pt       6.36pt            72.9%
%   promoted    10.91/10.46   7.97/7.64
%
% An earlier version of this comment said those subscripts are "under figcheck's
% 7pt T1 floor". They are not, and the correction is worth more than the claim
% was. T1 has TWO bands: running text is held to 7.0pt, and a span of three
% glyphs or fewer -- a sub/superscript, a tick numeral -- is held to 5.9pt
% (T1_SHORT_SPAN_PT). That second band is not a loophole. A subscript is set at
% roughly 70% of its base by every typesetter that has ever existed; holding one
% to the running-text floor would mean either no subscripts in figures at all,
% or a figure whose base type is 10pt against 8.72pt body text.
%
% So \pdfigmath is NOT a T1 remedy, and using it as one is how a rule gets
% applied where it does not belong. It exists because 5.98pt clears the
% short-span band by 0.08pt, and 0.08pt is not a margin: a font substitution, an
% edition that nudges \pdfigbasesize, or a chapter that loads a different math
% font puts that subscript under the band with nobody having decided anything.
% Promotion puts it at 7.97/7.64pt -- over even the running-text floor -- while
% the base glyph stays close enough to its neighbours that the label still reads
% as one label.
%
% Use it ONLY for a math group that carries a sub- or superscript. Plain math
% takes the label's own size like any other text, and P22 in tikz_precheck.py is
% an error if it does not.
\providecommand{\pdfigmath}{\normalsize}

% ---------------------------------------------------------------------------
% THE GUIDE DASH, IN ABSOLUTE POINTS, IN ONE PLACE.
%
% `pd guide` used to read `line width=.45pt,densely dotted`. `densely dotted`
% expands to `dash pattern=on \pgflinewidth off 1pt`, and \pgflinewidth is read
% when the KEY is processed, not when the path is stroked. So the Swiss
% edition's `pd guide/.append style={draw=hhink!40,line width=.5pt}` widened the
% stroke and left the dash at the .45pt already baked into it. The canonical
% Book therefore shipped a 0.448pt dot on a 0.498pt stroke.
%
% MEASURED IN THE SHIPPED ARTIFACT, not in a probe: 50 guide strokes sampled
% pixel by pixel in the committed coordination-papers-mega-volume.pdf, at 150
% dpi / 1.0x.
%
%   on 0.448pt = 0.93 device px -- UNDER ONE PIXEL; off 0.996pt = 2.08 px
%   the dash DECLARES 69% of each line's length as gap
%   only 34% of each line's pixels actually reach paper
%   35.4% of three-period windows contain NO pixel that reaches paper at all
%
% That last number is the defect. The gap (2.08 px) is narrower than the
% antialias spread of a 0.93 px stroke, so as the pattern's phase drifts along a
% line the gaps fill in: over a third of every guide's length there is no gap
% left and it is a continuous grey stroke, while the rest of the same line still
% reads as dots. The ink arrives. The DOTTEDNESS does not. A reader sees one
% line that is dotted here and solid there, and no author chose that.
%
% Twenty-five of the corpus's thirty-six guide draws use dottedness to mean
% "this line is not a drawn edge", so twenty-five of them were making a
% distinction the page does not carry.
%
% After, same measurement, same probe geometry, Book preamble:
%   on 1.2pt = 2.50 px, off 2.0pt = 4.17 px, stroke .7pt = 1.45 px
%   0.0% of three-period windows are without paper -- every gap survives,
%     at every phase
%   50% of pixels reach paper against 62.5% declared, the shortfall now being
%     the stroke's antialias skirt rather than the whole signal
%
% ON THE INK, because the obvious reading of this change is wrong. The BASE pair
% was never the problem: hhgray!62 against the hairline's hhgray!78 is 0.411
% against 0.517 of nominal ink, well apart. The problem is that the Book renders
% the SWISS edition (\pdedition defaults to swiss; the committed Book PDF is
% full of hhink!40 guides), where the pair was hhink!40 against hhink!45 -- 4.6%
% of the grey scale. Deleting the Swiss override, which is the right fix for the
% geometry, makes the base ink govern the Book too -- and hhgray!62 is 0.411
% against the Book hairline's 0.410. The same grey, to a thousandth. So the base
% ink has to move, or the dash would go back to carrying the whole distinction
% by itself, which is how this started.
%
% hhink!70 measured against the hairline, mean ink on the page at 150 dpi:
%   chapter  guide 0.0502  hairline 0.0790   guide is 64% of the hairline
%   Book     guide 0.0545  hairline 0.0671   guide is 81% of the hairline
% Per DOT it is 0.638 against 0.410 -- unmistakably a different mark. In MEAN it
% is still lighter than the hairline -- so it still recedes. Two channels, both
% working, instead of one that does not.
%
% The replacement is stated absolutely and measures phase-STABLE: every dot
% renders at the same weight wherever it falls on the grid (peak column ink
% 0.808, invariant across twelve sub-pixel offsets). It is a macro rather than
% three numbers because an edition that restates one of them and not the others
% is how this defect happened. The geometry is a LEGIBILITY FLOOR, not an
% edition's character: an edition may move the guide's ink, and states the
% geometry through these two macros so it cannot half-restate it.
%
% Any dash in this corpus is stated in absolute points. The `dotted` family
% (`dotted`, `densely dotted`, `loosely dotted`) derives its on-length from the
% line width and is therefore banned -- P23 in tikz_precheck.py. `dashed` and
% its relatives are absolute (on 3pt off 3pt) and are fine.
\providecommand{\pdguidewidth}{.7pt}
% Three lengths, not one macro holding the whole key: pgfkeys splits a key list
% on `,` and `=` before expanding anything (so a macro carrying `dash
% pattern=...` is read as one key NAME), and TikZ's dash-pattern scanner reads
% the literal words `on`/`off` without expanding a macro that hides them. A
% dimension after each word does expand, which is the seam that works.
\providecommand{\pdguideon}{1.2pt}
\providecommand{\pdguideoff}{2.0pt}

\tikzset{
  % --- picture level: one size, one default tip, the chapter hue -----------
  % A figure label is a name, not prose: it never hyphenates.
  pd figure/.style={font=\pdfiglabelfamily\pdfigbasesize,text=pdink,line cap=round,line join=round,
    >={Stealth[length=2.6pt 3.2,width=2.0pt 2.4]},/utils/exec=\pdsetfocushue,
    every node/.append style={execute at begin node={%
      \hyphenpenalty=10000\exhyphenpenalty=10000\relax}}},
%
  % --- 1. the voices ---------------------------------------------------------
  pd title/.style={font=\pdfiglabelbold,text=pdink,align=center},
  pd label/.style={font=\pdfiglabel,text=pdink,align=center},
  pd note/.style={font=\pdfiglabelitalic,text=pdinkmuted,align=left},
  pd kind/.style={font=\pdfiglabel\scshape,text=pdinkmuted,align=center},
  pd tag/.style={pd label,fill=pdpage,inner sep=1.7pt},
%
  % --- 4. surfaces -----------------------------------------------------------
  pd state/.style={pd label,draw=pdink,line width=.9pt,fill=pdpage,
    rounded corners=1pt,inner xsep=6pt,inner ysep=4.5pt,minimum height=7.4mm},
  % a terminal state carries a double boundary, as every automaton text draws it
  pd terminal/.style={pd state,double=pdpage,double distance=1.1pt,line width=.5pt},
  pd artifact/.style={pd label,draw=pdinkmuted!60,line width=.5pt,fill=pdpage,
    rounded corners=1pt,inner xsep=5pt,inner ysep=3.5pt,minimum height=6.2mm},
  pd panel/.style={draw=pdinkmuted!55,line width=.5pt,
    dash pattern=on 2.4pt off 2pt,rounded corners=2.5pt,inner sep=7pt,align=center},
  pd hatch/.style={pattern={Lines[angle=45,distance=4pt,line width=.3pt]},
    pattern color=pdinkmuted!70},
  pd neutral fill/.style={fill=pdinkmuted!12,draw=pdinkmuted!55,line width=.5pt},
%
  % --- 3. strokes ------------------------------------------------------------
  pd hairline/.style={draw=pdinkmuted!45,line width=.5pt},
  % Absolute dash, from the typographic law's three macros above: a guide is a
  % definite dotted construction line, darker per dot than the hairline and
  % lighter in mean ink. A lattice is the quiet continuous grid a figure stands on.
  pd guide/.style={draw=pdink!70,line width=\pdguidewidth,dash pattern=on \pdguideon off \pdguideoff},
  pd lattice/.style={draw=pdinkmuted!40,line width=.4pt},
  pd tick/.style={draw=pdinkmuted!75,line width=.5pt},
  pd rule/.style={draw=pdink,line width=.9pt},
  pd spine/.style={draw=pdink,line width=1.6pt},
  pd thin arrow/.style={draw=pdinkmuted!75,line width=.5pt,->},
  pd arrow/.style={pd rule,->},
  pd spine arrow/.style={pd spine,->},
%
  % --- the chapter hue, as the figure's focus --------------------------------
  pd focus rule/.style={draw=pdfocus,line width=1.6pt},
  pd focus arrow/.style={pd focus rule,->},
  pd focus state/.style={pd state,draw=pdfocus,line width=1.6pt,fill=pdfocus!24},
  % The single solid chapter-colour surface in a figure: use only for the
  % culmination or object the reader should name after five seconds. Editions
  % may render ordinary focus with a pale field while keeping this one solid.
  pd climax state/.style={pd focus state},
  % a member of the subject's set when the set, not one node, is the subject:
  % edged in the chapter hue, never filled, so the one solid node stays unique
  pd focus outline/.style={pd state,draw=pdfocus,line width=1.2pt,fill=pdpage},
  pd focus fill/.style={fill=pdfocus!24,draw=pdfocus,line width=.5pt},
  pd focus datum/.style={circle,fill=pdfocus,draw=pdpage,line width=.5pt,inner sep=2.6pt},
  pd focus label/.style={pd label,text=pdfocus!85!pdink},
%
  % --- marks -----------------------------------------------------------------
  pd datum/.style={circle,fill=pdink,draw=pdpage,line width=.5pt,inner sep=2.3pt},
  pd badge/.style={pd label,circle,draw=pdink,line width=.5pt,fill=pdpage,
    inner sep=.5pt,minimum size=5.0mm},
  pd focus badge/.style={pd badge,draw=pdfocus,line width=.9pt,fill=pdfocus!24},
  % --- conventions shared by every chapter (Critical Analysis, 2026-09-08) ---
  % Automata: states are circles, terminal (accepting) states double circles;
  % edge words are short action names; guards go in an aligned legend table.
  pd automaton state/.style={pd label,circle,draw=pdink,line width=.9pt,fill=pdpage,
    minimum size=11mm,inner sep=1pt},
  pd automaton focus/.style={pd automaton state,draw=pdfocus,line width=1.6pt,fill=pdfocus!24},
  pd automaton terminal/.style={pd automaton state,double=pdpage,double distance=1.1pt,
    line width=.5pt},
  pd initial/.style={draw=pdink,line width=.9pt,->},
  % Ledgers: asset pools are rectangular buckets (open at the top) joined by
  % directed pipes; the conserved sum is set as an equation under the rails.
  pd bucket/.style={pd label,draw=none,minimum width=2.2cm,minimum height=1.3cm,
    append after command={\pgfextra{\draw[draw=pdink,line width=.9pt]
      (\tikzlastnode.north west) -- (\tikzlastnode.south west)
      -- (\tikzlastnode.south east) -- (\tikzlastnode.north east);}}},
  pd pipe/.style={draw=pdink,line width=.9pt,double=pdpage,double distance=2.2pt,
    -{Stealth[length=3.4pt 2,width=5.2pt]}},
  pd focus pipe/.style={pd pipe,draw=pdfocus},
  % Trust boundaries: a dashed rectangle is the Dolev-Yao interface; a vector
  % that escapes it is drawn in pd breach and named.
  pd trust boundary/.style={draw=pdink,line width=.9pt,dash pattern=on 4pt off 2.4pt,
    rounded corners=0pt,inner sep=8pt},
  % a stroke that ends at a badge is shortened, so no rule touches the numeral
  pd to badge/.style={shorten >=2pt},
  pd from badge/.style={shorten <=2pt},
}

% ===========================================================================
% 7.  PLOTS.  A plot is computed by a committed script, never sketched.  `pd
%     axis` is the one axis grammar: two spines (left, bottom), ticks out, no
%     box, no grid unless asked (`pd grid`), tick words in the label voice,
%     numbers with a thin-space thousands separator, room outside the plot for
%     direct labels (clip=false).  Series are named at their ends, not in a
%     legend: `pd series` (ink), `pd focus series` (chapter hue), and
%     `pd X series` for a concept; breach and warn series are dashed.  A
%     region between curves is `pd focus fill`/`X fill` with fillbetween.
%     Width: `width=` is the plot box (scale only axis); leave about 1.6 cm
%     of the 11.4 cm column for the y tick words and title.
% ===========================================================================
\pgfplotsset{
  pd axis/.style={
    scale only axis,width=8.8cm,height=4.6cm,
    axis lines*=left,axis line style={draw=pdink,line width=.5pt,-},
    tick align=outside,major tick length=3pt,minor tick length=0pt,
    tick style={draw=pdink,line width=.5pt},
    tick label style={pd label,text=pdinkmuted},
    label style={pd label},title style={pd title},
    every axis title/.append style={yshift=-2pt},
    enlargelimits=false,clip=false,
    grid=none,
    legend style={draw=none,fill=none,font=\pdfiglabel},
    /pgf/number format/1000 sep={\,},
    every axis plot/.append style={line width=.9pt,line join=round,mark size=1.7pt},
  },
  pd grid/.style={ymajorgrids,major grid style={draw=pdinkmuted!30,line width=.4pt}},
  pd series/.style={draw=pdink,mark=*,mark options={fill=pdink,draw=pdpage,line width=.5pt,solid}},
  pd focus series/.style={draw=pdfocus,line width=1.6pt,mark=*,
    mark options={fill=pdfocus,draw=pdpage,line width=.5pt,solid}},
  pd muted series/.style={draw=pdinkmuted!70,line width=.9pt,mark=none},
}
\newcommand{\pdplotconcept}[2]{%
  \pgfplotsset{#1 series/.style={draw=#2,mark=*,
    mark options={fill=#2,draw=pdpage,line width=.5pt,solid}}}}
\pdplotconcept{pd truth}{pdcobalt}
\pdplotconcept{pd legible}{pdteal}
\pdplotconcept{pd ready}{pdhealth}
\pdplotconcept{pd protocol}{pdindigo}
\pdplotconcept{pd identity}{pdviolet}
\pdplotconcept{pd reputation}{pdrust}
\pdplotconcept{pd value}{pdgold}
\pgfplotsset{
  pd breach series/.style={draw=pderror,dash pattern=on 3.2pt off 2.2pt,mark=diamond*,
    mark options={fill=pderror,draw=pdpage,line width=.5pt,solid}},
  pd warn series/.style={draw=pdamber,dash pattern=on 3.2pt off 2.2pt,mark=none},
}

% ===========================================================================
% Supplied visual-guide primitives: native Book type, deliberate strokes,
% readable Stealth heads, open markers, and a quiet quantitative grid.
% These are opt-in: old diagrams need spacing review before thicker arrows.
\tikzset{
  pd guide step/.style={pd artifact,draw=pdink,fill=pdpage,line width=1pt,
    align=center,minimum height=40pt,inner xsep=5pt,inner ysep=6pt},
  pd guide arrow/.style={draw=pdcobalt,line width=1.15pt,
    -{Stealth[length=7pt,width=4.4pt]},shorten <=2pt,shorten >=2pt},
}
\newcommand{\pdflowstep}[4]{%
  \node[pd guide step,text width=#4] (#1) at #2 {#3};}
\newcommand{\pdflowlink}[2]{%
  \draw[pd guide arrow] (#1.east)--(#2.west);}
\pgfplotsset{
  pd evidence axis/.style={pd axis,
    axis line style={draw=pdink,line width=1pt,-},
    tick style={draw=pdink,line width=.8pt},
    grid=major,major grid style={draw=pdinkmuted!20,line width=.35pt},
    axis background/.style={fill=pdpage}},
  pd evidence measured/.style={draw=pdcobalt,line width=1.35pt,mark=*,
    mark size=2.6pt,mark options={fill=pdpage,draw=pdcobalt,line width=1.1pt}},
  pd evidence comparison/.style={draw=pdink,line width=1.1pt,
    dash pattern=on 4pt off 2.5pt,mark=square*,mark size=2.2pt,
    mark options={fill=pdink,draw=pdink,solid}},
}
% Analytic curves, not sampled observations. Integer samples preserve the
% meaning of reviewer count; callers set axes and name their assumptions.
\newcommand{\pdindependentmisses}[2]{%
  \addplot[pd evidence measured,domain=1:#2,samples=#2] {100*(1-#1)^x};}
\newcommand{\pdreviewcost}[1]{%
  \addplot[pd evidence measured,domain=1:#1,samples=#1] {x};}
% 8.  UML.  The Book's engineering diagrams use UML 2.5 notation, so a reader
%     who knows UML reads them without a key.  Four kinds:
%     STATE MACHINE  `pd uml state` (rounded), `pd uml focus state`, a
%       composite drawn with `pd uml composite` around its substates plus a
%       compartment of internal transitions; `pd uml initial` (filled dot),
%       `pd uml final` (bullseye), `pd uml choice` (diamond); transitions are
%       `pd uml transition` (open arrowhead) labelled  event [guard] / effect.
%     SEQUENCE  `pd uml head` (lifeline head), `pd uml lifeline` (dashed),
%       `pd uml activation` (thin bar), messages `pd uml sync` (filled head),
%       `pd uml async` (open head), `pd uml reply` (dashed, open head),
%       `pd uml fragment` + `pd uml fragment tag` for alt / opt / loop.
%     COMPONENT  `pd uml component` with \pdumlicon{node} (the component
%       glyph) and a <<component>> keyword; `pd uml port`; provided interface
%       \pdumlprovided{from}{to}{name} (lollipop), required \pdumlrequired (socket).
%     ERD (crow's foot)  `pd er entity` (rectangle split: name, then rows;
%       mark keys PK / FK in the row text), relationship lines `pd er rel`
%       with tips `er one`, `er many`, `er zero or one`, `er one or many`,
%       `er zero or many` at either end.
% ===========================================================================
\pgfdeclarearrow{
  name = pd crow,
  parameters = {\the\pgfarrowlength,\the\pgfarrowwidth},
  setup code = {%
    \pgfarrowssettipend{0pt}%
    \pgfarrowssetlineend{0pt}%
    \pgfarrowssetbackend{-\pgfarrowlength}%
    \pgfarrowsupperhullpoint{0pt}{.5\pgfarrowwidth}%
    \pgfarrowsupperhullpoint{-\pgfarrowlength}{0pt}%
    \pgfarrowssavethe\pgfarrowlength
    \pgfarrowssavethe\pgfarrowwidth},
  drawing code = {%
    \pgfsetdash{}{0pt}%
    \pgfpathmoveto{\pgfqpoint{-\pgfarrowlength}{0pt}}%
    \pgfpathlineto{\pgfqpoint{0pt}{.5\pgfarrowwidth}}%
    \pgfpathmoveto{\pgfqpoint{-\pgfarrowlength}{0pt}}%
    \pgfpathlineto{\pgfqpoint{0pt}{-.5\pgfarrowwidth}}%
    \pgfusepathqstroke},
  defaults = {length=2.4mm, width=3.0mm}
}
\tikzset{
  % state machine
  pd uml state/.style={pd state,rounded corners=4.5pt},
  pd uml focus state/.style={pd focus state,rounded corners=4.5pt},
  pd uml terminal/.style={pd terminal,rounded corners=4.5pt},
  pd uml composite/.style={draw=pdink,line width=.9pt,rounded corners=7pt,fill=none,inner sep=0pt},
  pd uml initial/.style={circle,fill=pdink,draw=none,minimum size=2.4mm,inner sep=0pt},
  pd uml final/.style={circle,draw=pdink,line width=.9pt,fill=pdpage,minimum size=3.6mm,inner sep=0pt,
    path picture={\fill[pdink] (path picture bounding box.center) circle (1.1mm);}},
  pd uml choice/.style={diamond,draw=pdink,line width=.9pt,fill=pdpage,minimum size=4.2mm,inner sep=0pt},
  pd uml transition/.style={pd rule,-{Straight Barb[length=2.2pt 2.2,width=4.4pt 2.2]}},
  pd uml focus transition/.style={pd focus rule,-{Straight Barb[length=2.2pt 2.2,width=4.4pt 2.2]}},
  pd uml breach transition/.style={pd rule,draw=pderror,dash pattern=on 3.2pt off 2.2pt,
    -{Straight Barb[length=2.2pt 2.2,width=4.4pt 2.2]}},
  % sequence
  pd uml head/.style={pd state,rounded corners=0pt,minimum height=6.8mm},
  pd uml lifeline/.style={draw=pdinkmuted!75,line width=.5pt,dash pattern=on 3pt off 2pt},
  pd uml activation/.style={draw=pdink,line width=.5pt,fill=pdpage,minimum width=2.2mm,inner sep=0pt},
  pd uml sync/.style={pd rule,-{Triangle[length=3.2pt 1.6,width=3.2pt 1.6]}},
  pd uml async/.style={pd rule,-{Straight Barb[length=2.4pt 2,width=4.4pt 2]}},
  pd uml reply/.style={pd rule,dash pattern=on 3pt off 2pt,-{Straight Barb[length=2.4pt 2,width=4.4pt 2]}},
  pd uml fragment/.style={draw=pdink,line width=.5pt,fill=none,inner sep=6pt},
  pd uml fragment tag/.style={pd label,font=\footnotesize\bfseries,draw=pdink,line width=.5pt,fill=pdpage,
    inner xsep=3pt,inner ysep=2pt,chamfered rectangle,chamfered rectangle corners=south east,
    chamfered rectangle xsep=2pt},
  % component
  pd uml component/.style={pd state,rounded corners=0pt,align=center,inner xsep=8pt,inner ysep=6pt,minimum height=11mm},
  pd uml port/.style={rectangle,draw=pdink,line width=.5pt,fill=pdpage,minimum size=2mm,inner sep=0pt},
  % ERD
  pd er entity/.style={pd label,rectangle split,rectangle split parts=2,rectangle split part align={center,left},
    draw=pdink,line width=.9pt,fill=pdpage,inner xsep=5pt,inner ysep=3pt,align=left,
    rectangle split part fill={pdinkmuted!12,pdpage}},
  pd er rel/.style={pd rule},
  er one/.tip={Bar[width=3.0mm,sep=1.2pt]Bar[width=3.0mm]},
  er many/.tip={Bar[width=3.0mm,sep=1.4pt]pd crow},
  er zero or one/.tip={Circle[open,length=1.8mm,width=1.8mm,sep=1.2pt,fill=pdpage]Bar[width=3.0mm]},
  er one or many/.tip={Bar[width=3.0mm,sep=1.4pt]pd crow},
  er zero or many/.tip={Circle[open,length=1.8mm,width=1.8mm,sep=1.2pt,fill=pdpage]pd crow},
}
% the component glyph, drawn in the top-right corner of a component node
\newcommand{\pdumlicon}[1]{%
  \begin{scope}[shift={($(#1.north east)+(-4.2mm,-2.6mm)$)},x=1mm,y=1mm]
    \draw[pdink,line width=.5pt,fill=pdpage] (0,-1.1) rectangle (2.6,1.1);
    \draw[pdink,line width=.5pt,fill=pdpage] (-.6,.25) rectangle (.6,.65);
    \draw[pdink,line width=.5pt,fill=pdpage] (-.6,-.65) rectangle (.6,-.25);
  \end{scope}}
% a provided interface (lollipop) from a node side to a point, with its name
\newcommand{\pdumlprovided}[3]{%
  \draw[pd rule] (#1) -- (#2);
  \draw[pd rule,fill=pdpage] (#2) circle (1.3mm);
  \node[pd label,anchor=south] at ($(#2)+(0,1.6mm)$) {#3};}
% a required interface (socket): a half-circle open toward the provider
\newcommand{\pdumlrequired}[4]{% from, to, name, facing angle of the socket's opening
  \draw[pd rule] (#1) -- (#2);
  \draw[pd rule] ($(#2)+(#4+90:1.9mm)$) arc[start angle=#4+90,end angle=#4+270,radius=1.9mm];
  \node[pd label,anchor=north] at ($(#2)+(0,-2.4mm)$) {#3};}

% --- 2. the concept map, as styles, so a drawing never chooses a hue --------
% An edition file may \renewcommand\pdconcept before the map below expands.
\newcommand{\pdconcept}[2]{%
  \tikzset{%
    #1/.style={draw=#2,text=pdink},
    #1 rule/.style={draw=#2,line width=1.6pt},
    #1 arrow/.style={draw=#2,line width=.9pt,->},
    #1 state/.style={pd state,draw=#2,line width=1.6pt,fill=#2!24},
    #1 fill/.style={fill=#2!24,draw=#2,line width=.5pt},
    #1 datum/.style={circle,fill=#2,draw=pdpage,line width=.5pt,inner sep=2.6pt},
    #1 label/.style={pd label,text=#2!85!pdink},
  }}

% --- 6. the edition layer, loaded before the concept map expands ----------
\ifdefined\pdedition\IfFileExists{figures/pd-figure-language-\pdedition.tex}{\input{figures/pd-figure-language-\pdedition}}{}\fi

\pdconcept{pd truth}{pdcobalt}
\pdconcept{pd legible}{pdteal}
\pdconcept{pd ready}{pdhealth}
\pdconcept{pd protocol}{pdindigo}
\pdconcept{pd identity}{pdviolet}
\pdconcept{pd reputation}{pdrust}
\pdconcept{pd value}{pdgold}
\pdconcept{pd breach}{pderror}
\pdconcept{pd warn}{pdamber}
% Breach and warning read as refusals, so they are dashed as well as
% coloured, and a breach mark is a diamond where every other mark is a circle.
% Amber is never text (3.71:1 on white); a warn label is set in ink.
\tikzset{
  pd breach rule/.append style={dash pattern=on 3.2pt off 2.2pt},
  pd breach arrow/.append style={dash pattern=on 3.2pt off 2.2pt},
  pd warn rule/.append style={dash pattern=on 3.2pt off 2.2pt},
  pd warn arrow/.append style={dash pattern=on 3.2pt off 2.2pt},
  pd breach datum/.append style={diamond},
  pd warn label/.style={pd label},
}

% ===========================================================================
% v1 NAMES.  Kept so every existing fragment still compiles; they resolve to
% the v2 roles, so an untouched figure inherits the ladder.  New work uses the
% v2 names (skills/tikz-diagram-craft/references/migration-v1-to-v2.md).
% ===========================================================================
\tikzset{
  pd row label/.style={pd title,anchor=east},
  pd panel title/.style={pd title,align=left},
  pd axis label/.style={pd label},
  % Direct labels sit beside their mark; they do not erase a line or box
  % underneath themselves. A label that needs a paper knockout is misplaced
  % and must be moved, not given an invisible patch.
  pd direct label/.style={pd tag,fill=none,align=left},
  pd actor/.style={pd state},
  pd boundary/.style={pd panel},
  pd caution rule/.style={pd breach rule},
  pd caution arrow/.style={pd breach arrow},
  pd caution fill/.style={pd breach fill},
  pd caution datum/.style={pd breach datum},
}

% ===========================================================================
% Roles named on main by the typographic law (#10191), expressed on the v2
% roles so an edition's .append style on the parent reaches them. `pd danger`
% is main's name for what v2 calls `pd breach` (red, dashed, diamond marks);
% both names resolve to one family. Main's hairline-boxed fact badge is
% `pd fact badge`: v2's `pd badge` is the numbered guard/step marker.
% ===========================================================================
\tikzset{
  pd danger rule/.style={pd breach rule},
  pd danger arrow/.style={pd breach arrow},
  pd danger fill/.style={pd breach fill,text=pdink},
  pd danger datum/.style={pd breach datum},
  pd ink fill/.style={fill=pdink,draw=pdink,line width=.4pt,text=pdpage},
  pd decision/.style={pd state,diamond,aspect=1.55,font=\pdfiglabelbold,
    minimum height=0pt,inner xsep=3pt,inner ysep=3pt},
  pd mono label/.style={pd direct label,font=\pdfiglabelmono,inner xsep=2.5pt,inner ysep=2.5pt},
  pd verdict/.style={pd direct label,font=\pdfiglabelbold},
  pd reverse label/.style={pd direct label,text=pdpage,fill=none},
  pd reverse row label/.style={pd row label,text=pdpage},
  pd legend/.style={font=\pdfiglabel,text=pdink,draw=none,fill=none},
  pd kind tag/.style={pd kind,align=right},
  pd fact badge/.style={pd hairline,font=\pdfiglabel,text=pdink,align=center,
    inner sep=2.4pt,fill=none},
}

% Small reusable marks.  They deliberately avoid prose-bearing boxes.
\newcommand{\pdfigaxis}[4]{%
  \draw[pd rule,->] (#1,#2) -- (#3,#2);
  \node[pd label,anchor=west] at (#3,#2) {#4};
}
\newcommand{\pdfigvaxis}[4]{%
  \draw[pd rule,->] (#1,#2) -- (#1,#3);
  \node[pd label,anchor=south] at (#1,#3) {#4};
}

% Native-size, object-based evidence figures; shared by Book and chapters.
% Shared approved Swiss evidence drawing primitives. Keep both chapter roots identical.
\ifdefined\pdSwissEvidenceLoaded\else
\def\pdSwissEvidenceLoaded{1}
% Real ink-box gate; no declared coordinate envelope can substitute for this.
\newsavebox{\sgEvidenceBox}
\newcommand{\SGMeasuredFigure}[2]{%
  \sbox{\sgEvidenceBox}{#2}%
  \typeout{SG-GEOMETRY: #1,width=\the\wd\sgEvidenceBox,height=\the\dimexpr\ht\sgEvidenceBox+\dp\sgEvidenceBox\relax}%
  \ifdim\wd\sgEvidenceBox>325.215pt
    \PackageError{swiss-evidence}{Figure #1 exceeds the Book column}{Recompose at native text size; do not scale.}%
  \fi
  \usebox{\sgEvidenceBox}}
% Parent redraws after rejection of lab specimens 01,04,06.
% All coordinates are physical TeX points; no resizebox or transformed fonts.
\usetikzlibrary{arrows.meta,calc}
% Aliased onto the house palette (figures/pd-palette.tex), not \definecolor'd
% as independent HTML values: an edition that restyles pdcobalt/pdteal/...
% now restyles every SG figure with it. The six values were already
% bit-identical matches for pdcobalt/pdteal/pdviolet/pdamber/pderror/pdink;
% the alias turns that into a guarantee instead of a coincidence.
\colorlet{sgBlue}{pdcobalt}
\colorlet{sgTeal}{pdteal}
\colorlet{sgViolet}{pdviolet}
\colorlet{sgAmber}{pdamber}
\colorlet{sgRed}{pderror}
\colorlet{sgInk}{pdink}
\definecolor{sgGrid}{HTML}{D8C7A6}
% The one house size, inheriting the document's own face (the typographic
% law in pd-figure-language.tex): no \sffamily leak, no local \fontsize.
\newcommand{\SGType}{\pdfiglabel}
\tikzset{sg/base/.style={x=1pt,y=1pt,line cap=butt,line join=miter,
  every node/.style={font=\SGType,inner sep=0pt,text=sgInk}},
  sg/arrow/.style={-{Stealth[length=4.5pt,width=3.6pt]},draw=sgInk,line width=.85pt},
  sg/grid/.style={draw=sgGrid,line width=.3pt},
  sg/hash/.style={draw=sgBlue,fill=sgBlue!7,line width=1pt,minimum width=53pt,
    minimum height=27pt,inner sep=4pt,align=center},
  sg/sibling/.style={draw=sgAmber,fill=sgAmber!10,line width=1pt,minimum width=51pt,
    minimum height=27pt,inner sep=4pt,align=center},
  sg/check/.style={draw=sgInk,fill=white,line width=1pt,minimum width=92pt,
    minimum height=30pt,inner sep=5pt,align=center}}

% A folded record, not an undifferentiated rounded node. Width/height are pt.
\newcommand{\SGDocument}[6]{%
  \begin{scope}[shift={#2}]
    \path[draw=sgInk,fill=#5,line width=.85pt]
      (0,0)--({#3-7},0)--(#3,7)--(#3,#4)--(0,#4)--cycle;
    \draw[sgInk,line width=.65pt] ({#3-7},0)--({#3-7},7)--(#3,7);
    \node[minimum width=#3pt,minimum height=#4pt,text width=\dimexpr#3pt-12pt\relax,
      align=center] (#1) at ({#3/2},{#4/2}) {#6};
  \end{scope}}

\newcommand{\SGSurvival}{%
\begin{tikzpicture}[sg/base]
  % Three complete probability decades, percentages rather than tiny exponents.
  \foreach \x in {0,1,...,6}{\draw[sg/grid] ({35+30*\x},20)--({35+30*\x},140);}
  \foreach \y/\lab in {20/{0.1\%},60/{1\%},100/{10\%},140/{100\%}}{%
    \draw[sg/grid] (35,\y)--(215,\y);
    \draw[sgInk,line width=.65pt] (32,\y)--(35,\y);
    \node[anchor=east] at (28,\y) {\lab};}
  \draw[sgInk,line width=.85pt] (35,140)--(35,20)--(215,20);
  \foreach \x in {0,2,4,6}{%
    \draw[sgInk,line width=.65pt] ({35+30*\x},20)--({35+30*\x},17);
    \node[anchor=north] at ({35+30*\x},12) {$\x$};}
  \node[rotate=90] at (-9,80) {jobs still running (log scale)};
  \node[anchor=north] at (125,-5) {duration / mean duration};
  \draw[sgBlue,line width=1.5pt,domain=0:6,samples=100]
    plot ({35+30*\x},{140-120*ln(1+\x/2)/ln(10)});
  \draw[sgViolet,line width=1.2pt,dash pattern=on 4pt off 2.5pt,domain=0:6,samples=100]
    plot ({35+30*\x},{140-40*\x/ln(10)});
  % Endpoint leaders have allocated space; text does not lie on a curve.
  \draw[sgBlue,line width=.8pt] (215,67.75)--(224,67.75);
  \node[anchor=west,text=sgBlue,font=\SGType\bfseries] at (230,73) {Lomax};
  \node[anchor=west,text=sgBlue] at (230,60) {1.56\%};
  \draw[sgViolet,line width=.8pt,dashed] (215,35.77)--(224,35.77);
  \node[anchor=west,text=sgViolet,font=\SGType\bfseries] at (230,40) {exponential};
  \node[anchor=west,text=sgViolet] at (230,27) {0.248\%};
\end{tikzpicture}}

\newcommand{\SGTiming}{%
\begin{tikzpicture}[sg/base,y=-1pt]
  % This bracket spans COMMIT TO SYNC. It is not a range of hypothetical syncs.
  \path[draw=sgAmber!35,fill=sgAmber!9,line width=.5pt] (39,19) rectangle (253,49);
  \draw[sgAmber,line width=1pt] (39,18)--(39,9)--(253,9)--(253,18);
  \node[anchor=south,text=sgAmber,font=\SGType\bfseries] at (146,5) {?};
  \draw[sgInk,line width=.9pt] (39,27)--(39,45) (253,27)--(253,45);
  \draw[{Stealth[length=4.5pt,width=3.6pt]}-{Stealth[length=4.5pt,width=3.6pt]},
    sgAmber,line width=1.3pt] (40,36)--(252,36);
  \node[anchor=north,font=\SGType\bfseries] at (39,52) {commit};
  \node[anchor=north,font=\SGType\bfseries] at (253,52) {WAL sync};
  \node[anchor=west] at (0,86) {heartbeats received};
  \foreach \x/\label in {57/{s_1},117/{s_2},194/{s_3},278/{s_4}}{%
    \draw[sgViolet!70,dashed,line width=.55pt] (\x,104)--(\x,145);
    \node[anchor=south,text=sgViolet] at (\x,101) {$\label$};}
  \path[draw=sgAmber!35,fill=sgAmber!9,line width=.5pt] (158,108) rectangle (243,142);
  \draw[sg/grid] (0,139)--(306,139);
  \draw[sgTeal,line width=1.4pt]
    (0,132)--(25,132)--(25,116)--(33,116)--(33,132)
    --(85,132)--(85,116)--(93,116)--(93,132)
    --(145,132)--(145,116)--(153,116)--(153,132)--(306,132);
  \draw[sgRed,line width=1pt] (243,106)--(243,145);
  \node[anchor=north,text=sgRed] at (243,149) {deadline};
  \draw[sg/arrow] (0,174)--(306,174);
  \node[anchor=north east] at (306,179) {time};
\end{tikzpicture}}

\newcommand{\SGInclusion}{%
\begin{tikzpicture}[sg/base,y=-1pt]
  % Two independent inputs: a supplied proof and an issuer-authenticated root.
  \node[anchor=north west,font=\SGType\bfseries,text width=158pt,align=left] at (0,0)
    {Recompute from the receipt};
  \node[anchor=north west,font=\SGType\bfseries,text width=118pt,align=left] at (193,0)
    {Authenticate\\the checkpoint};
  \draw[sg/grid,step=15pt] (0,29) grid (175,194);
  \SGDocument{receipt}{(4,35)}{49}{34}{sgViolet!9}{receipt $P$}
  \node[sg/hash] (leaf) at (29,103) {$a=H_L(P)$};
  \node[sg/sibling] (sibone) at (108,103) {sibling $b$};
  \node[sg/hash] (inner) at (68,146) {$u=H_N(a,b)$};
  \node[sg/sibling] (sibtwo) at (146,146) {subtree $c$};
  \node[sg/hash,fill=sgBlue,text=white,minimum width=100pt] (computed) at (107,190)
    {$R'=H_N(u,c)$};
  \draw[sg/arrow,draw=sgViolet] (receipt.south)--(leaf.north);
  \draw[sg/arrow,draw=sgBlue] (leaf.south)--(inner.north west);
  \draw[sg/arrow,draw=sgAmber] (sibone.south)--(inner.north east);
  \draw[sg/arrow,draw=sgBlue] (inner.south)--(computed.north west);
  \draw[sg/arrow,draw=sgAmber] (sibtwo.south)--(computed.north east);
  \SGDocument{checkpoint}{(211,37)}{85}{48}{sgTeal!8}{signed root $R^*$\\size $n=4$\\signature $\sigma$}
  % The trusted key is an external input, not something the prover supplies.
  \draw[sgInk,line width=1pt] (218,113) circle (3.5pt);
  \draw[sgInk,line width=1pt] (221.5,113)--(233,113) (230,113)--(230,117) (233,113)--(233,117);
  \node[anchor=west] at (239,113) {trusted $pk_I$};
  \node[sg/check,minimum width=84pt] (verify) at (260,158) {verify $\sigma$\\result $A$};
  \draw[sg/arrow,draw=sgTeal] (checkpoint.east)--(311,61)--(311,158)--(verify.east);
  \draw[sg/arrow] (247,123)--(247,138)--(verify.north);
  \node[sg/check,minimum width=112pt,minimum height=34pt] (compare) at (164,245)
    {$A\ \land\ (R'=R^*)$};
  \draw[sg/arrow,draw=sgBlue] (computed.south)--(107,220)--(compare.north west);
  \draw[sg/arrow,draw=sgTeal] (verify.south)--(260,214)--(compare.north east);
  \node[anchor=east,text=sgTeal] at (250,204) {$R^*,A$};
  \node[draw=sgTeal,fill=sgTeal,text=white,minimum width=90pt,
    minimum height=27pt,font=\SGType\bfseries] (accept) at (76,299) {accept inclusion};
  \node[draw=sgRed,fill=white,text=sgRed,line width=1.2pt,minimum width=90pt,
    minimum height=27pt,font=\SGType\bfseries] (reject) at (256,299) {reject proof};
  \draw[sg/arrow,draw=sgTeal] (compare.south west)--(108,279)--(76,279)--(accept.north);
  \draw[sg/arrow,draw=sgRed] (compare.south east)--(220,279)--(256,279)--(reject.north);
  \node[anchor=south,text=sgTeal] at (82,275) {true};
  \node[anchor=south,text=sgRed] at (248,275) {false};
\end{tikzpicture}}

% Reusable geometry; load after swiss-redraws.tex. Real points, no scaling.
\newcommand{\SGDatabase}[4]{%
  \begin{scope}[shift={#2}]
    \path[draw=#3,fill=#3!7,line width=.9pt]
      (0,4)--(0,32) arc[start angle=180,end angle=360,x radius=24pt,y radius=5pt]
      --(48,4) arc[start angle=0,end angle=-180,x radius=24pt,y radius=5pt];
    \draw[draw=#3,fill=#3!7,line width=.9pt] (24,4) ellipse[x radius=24pt,y radius=5pt];
    \node[minimum width=48pt,minimum height=37pt,align=center] (#1) at (24,18) {#4};
  \end{scope}}

% Reader question: How do overlapping requests become ordered records?
% Evidence: schematic two-client history, not a runtime observation.
% Each request precedes commit 1; each reply follows its own committed record.
\newcommand{\SGWriterSequence}{%
\begin{tikzpicture}[sg/base,y=-1pt]
  % Recomposed at native size: a 22-character mono request label no longer
  % fits a 118pt lane gap at \pdfiglabelmono, so the HTTP verb is shortened
  % to the identifier it names and the resource it names ("auth.ts") moves
  % to the lane header, where it reads once instead of twice.
  \foreach \x/\lanelabel in {35/{{\pdfigmath$\alpha_{\text{auth}}$} (claude)},153/{Kernel (single writer)},275/{{\pdfigmath$\alpha_{\text{bill}}$} (codex)}}{
    \node[anchor=south,font=\pdfiglabelbold,align=center] at (\x,0) {\lanelabel};
    \draw[sg/grid] (\x,9)--(\x,237);}
  % Both invocations precede the first commit: overlapping client requests.
  \draw[sg/arrow,draw=sgViolet] (35,29)--(153,29);
  \node[anchor=south,text=sgViolet,font=\pdfiglabelmono] at (94,24) {POST /claims};
  \draw[sg/arrow,draw=sgAmber] (275,57)--(153,57);
  \node[anchor=south,text=sgAmber,font=\pdfiglabelmono] at (214,52) {POST /claims};
  % commit 1 succeeds: drawn in the kernel's own hue. commit 2 is the
  % kernel's WAL record of a REFUSAL, not a second grant -- it is drawn in
  % the breach hue so the outcome is doubled by colour, not stated in prose
  % alone (the typographic law's "colour is always doubled" rule). The
  % "lock held" highlight band is gone: at native size it ran directly
  % behind this text, which figcheck reads as a drawn line through the
  % words, and the sequencing arrows already carry the serialization claim.
  \SGDocument{commitA}{(93,71)}{120}{46}{sgBlue!8}{{\pdfiglabelbold commit 1 \textperiodcentered{} WAL 4102}\\[2pt]{\pdfigsub {\pdfigmath$\alpha_{\text{auth}}$} \textperiodcentered{} epoch 104}}
  \draw[sg/arrow,draw=sgViolet,dash pattern=on 3pt off 2pt] (153,130)--(35,130);
  \node[anchor=south,text=sgViolet,font=\pdfiglabel] at (94,128) {201 Created};
  \SGDocument{commitB}{(93,152)}{120}{46}{sgRed!8}{{\pdfiglabelbold commit 2 \textperiodcentered{} WAL 4103}\\[2pt]{\pdfigsub {\pdfigmath$\alpha_{\text{bill}}$} refused}}
  \draw[sg/arrow,draw=sgAmber,dash pattern=on 3pt off 2pt] (153,213)--(275,213);
  \node[anchor=south,text=sgAmber,font=\pdfiglabel] at (214,211) {409 Conflict};
  \draw[sg/arrow] (3,20)--(3,232);
  \node[anchor=west,font=\pdfiglabel] at (9,236) {time};
\end{tikzpicture}}

% Reader question: What is checked before the service admits requests?
% Source: SQLite WAL documentation, sections 2--2.2, cited in the chapter.
% Checkpointing copies committed pages; it does not apply uncommitted work.
% The dashed group denotes separate mechanisms, not a proved composition.
\newcommand{\SGRecoveryAdmission}{%
\begin{tikzpicture}[sg/base,y=-1pt]
  % Recomposed at native size: a 25-character mono command (sqlite3_wal_
  % checkpoint_v2) no longer fits a 103pt-wide column once the label grows
  % from 5.0pt to \pdfiglabelmono, so the three mechanisms read top to
  % bottom, full measure, instead of left to right in three narrow lanes.
  \draw[sgInk!55,dash pattern=on 3pt off 2pt,line width=.7pt]
    (0,0) rectangle (300,318);

  % 1. Checkpoint WAL
  \node[anchor=north west,align=left,text width=228pt] at (10,8)
    {{\pdfiglabelbold 1. Checkpoint WAL}\\[2pt]
     {\pdfiglabelmono\color{sgBlue} sqlite3\_wal\_checkpoint\_v2}\\[2pt]
     {\pdfigsub Copies committed WAL pages into the main database; uncommitted frames are not applied.}};
  \SGDocument{walback}{(247,4)}{34}{26}{sgViolet!7}{}
  \SGDocument{walfront}{(253,9)}{34}{26}{sgViolet!7}{WAL}
  \draw[sg/arrow] (18,101)--(18,118);

  % 2. Check Pages
  \node[anchor=north west,align=left,text width=228pt] at (10,122)
    {{\pdfiglabelbold 2. Check Pages}\\[2pt]
     {\pdfiglabelmono\color{sgTeal} PRAGMA integrity\_check}\\[2pt]
     {\pdfigsub Verifies B-tree pointers and freelist allocation across every page.}};
  \SGDatabase{db}{(251,118)}{sgBlue}{SQLite}
  \draw[sg/arrow] (18,215)--(18,232);

  % 3. Verify Links
  \node[anchor=north west,align=left,text width=228pt] at (10,236)
    {{\pdfiglabelbold 3. Verify Links}\\[2pt]
     {\pdfiglabelmono\color{sgAmber} Merkle chain verify}\\[2pt]
     {\pdfigsub Validates cryptographic audit-log continuity link by link.}};
  \foreach \x in {249,267,285}{
    \draw[sgAmber,line width=1pt,rounded corners=2pt] (\x,240) rectangle ({\x+14},256);}
  \draw[sgAmber,line width=1pt] (259,248)--(267,248) (281,248)--(289,248);

  % Decision Gate
  \node[sg/check,minimum width=165pt,minimum height=28pt,font=\pdfiglabelbold] (admit) at (150,350)
    {all critical checks pass?};
  \draw[sg/arrow] (150,318)--(admit.north);

  \node[draw=sgTeal,fill=sgTeal,text=white,minimum width=125pt,minimum height=32pt,
    align=center] (serve) at (65,420) {{\pdfiglabelbold admit requests}\\[1.5pt]{\pdfigsub HTTP 200 Ready (:9876)}};
  \node[draw=sgRed,fill=white,text=sgRed,line width=1pt,minimum width=125pt,
    minimum height=32pt,align=center] (refuse) at (235,420) {{\pdfiglabelbold refuse service}\\[1.5pt]{\pdfigsub HTTP 503 Unavailable}};

  \draw[sg/arrow,draw=sgTeal] (admit.south west)--(90,384)--(65,384)--(serve.north);
  \draw[sg/arrow,draw=sgRed] (admit.south east)--(210,384)--(235,384)--(refuse.north);
  % Each tag sits clear of its own diagonal leg, on the side the diagonal
  % does NOT occupy -- PASS left of its bend, FAIL right of its bend --
  % rather than centred on the bend where the stroke crosses the glyphs.
  \node[anchor=south east,text=sgTeal,font=\pdfiglabelbold] at (63,382) {PASS};
  \node[anchor=south west,text=sgRed,font=\pdfiglabelbold] at (237,382) {FAIL (HALT)};
\end{tikzpicture}}

% Named bands and their audiences; the dashed floor is NOT a fifth layer.
\newcommand{\SGStackLayers}{%
\begin{tikzpicture}[sg/base,y=-1pt]
  \foreach \y/\level/\name/\audience/\evidence/\hue in {
    0/3/Economy/{Between operators}/{settlement and reputation}/sgAmber,
    51/2/Legibility/{For the operator}/{summaries linked to evidence}/sgTeal,
    102/1/Coordination/{Between agents}/{permissions and commitments}/sgViolet,
    153/0/Substrate/{On one machine}/{ordered writes and records}/sgBlue}{
    \path[draw=\hue,fill=\hue!8,line width=1pt] (17,\y) rectangle (142,{\y+38});
    \node[anchor=west,font=\SGType\bfseries,text=\hue] at (26,{\y+13}) {L\level};
    \node[anchor=west,font=\SGType\bfseries] at (26,{\y+27}) {\name};
    \node[anchor=north west,font=\SGType\bfseries] at (158,{\y+2}) {\audience};
    \node[anchor=north west,text width=149pt,align=left] at (158,{\y+17}) {\evidence};}
  \foreach \y in {153,102,51}{\draw[sg/arrow] (79,\y)--(79,{\y-13});}
  \draw[sgBlue,line width=2.4pt] (7,154)--(7,190);
  \draw[sgInk!60,dashed,line width=.7pt] (17,204) rectangle (307,243);
  \node[anchor=west,font=\SGType\bfseries] at (26,216) {Machine floor};
  \node[anchor=west] at (26,231) {file lock, clock, socket, supervisor};
  \draw[sg/arrow] (79,204)--(79,192);
\end{tikzpicture}}

% Compact counterpart for future chapter-by-chapter marginal navigation.
% A named current-layer argument changes emphasis, never the layer order.
\newcommand{\SGMarginLadder}[1]{%
\begin{tikzpicture}[sg/base,y=-1pt]
  \foreach \y/\n/\name/\hue in {0/3/Economy/sgAmber,37/2/Legibility/sgTeal,74/1/Coordination/sgViolet,111/0/Substrate/sgBlue}{
    \draw[\hue,line width=.7pt] (0,\y)--(91,\y);
    \node[anchor=north west,text=\hue,font=\SGType\bfseries] at (4,{\y+3}) {L\n};
    \node[anchor=north west] at (4,{\y+16}) {\name};
    \ifnum\n=#1\relax\draw[\hue,line width=2pt] (0,{\y+3})--(0,{\y+29});\fi}
\end{tikzpicture}}

\fi



% Honest-label markers (ADR-0045 discipline; one key for the whole volume).
% \Built, \BuiltWeak, \Designed and \Vision -- and the claim-evidence key
% \Verified / \Proved / \Unproved that labels a *result* rather than an
% implementation -- are defined once, for this paper and for the Book, in
% figures/pd-pedagogy.tex (\input above). Both keys are stated in prose in
% Appendix~\ref{app:honest}; nothing in this paper carries a free-text status
% note outside these macros.
%
% This chapter used to define all seven itself, in its own words: "built",
% "built (weak)", "designed", "vision". Because a chapter preamble is dropped
% when the Book is assembled, those words reached a reader only in the
% standalone paper -- the Book printed "implemented", "partial", "specified",
% "proposed" at the very same call sites, and the legend below explained a
% vocabulary the tables there never used. The series settled on the second set
% (single-writer-kernel.tex names them in prose as "the series' honest
% labels"), so this paper now prints them too and the legend says so.
%
% \Proposed became \Unproved in the same pass. It meant "no proof carried
% here", on the evidence axis; the shared maturity scale's \Vision means "no
% design document yet". Both rendered the one word "proposed" in the same gray
% italic, and in the Book they landed eight lines apart in the legend below,
% defining that word twice, differently.

% The pull-quote, key-idea and pitfall callouts, and the named-actor \scene
% vignette, are defined once in figures/pd-pedagogy, whose \AtBeginDocument
% block \long\def's every one of those names as a run-in or margin head over
% plain prose. Local tinted-box definitions of the same four names used to sit
% here; they were dead at render time (the shared block always re-defined them
% after this preamble ran), so they have been deleted rather than left to rot.

% Exercises: retired the chapter-local `exercises`/`exhead` tikz box (Wave 5,
% legible-swarm-pedagogy) in favor of the shared figures/pd-pedagogy pdexercise
% and pdsolution environments, which every other chapter now uses; nothing in
% this file calls \exhead or \begin{exercises} any longer.

% ─── Listings Setup (house style) ─────────────────────────────────────────────
\lstdefinestyle{pdcode}{
  language=C,
  backgroundcolor=\color{codebg},
  frame=single,
  rulecolor=\color{codeframe},
  basicstyle=\ttfamily\small,
  keywordstyle=\color{hhink}\bfseries,
  commentstyle=\color{hhgray}\itshape,
  stringstyle=\color{hhgray},
  numbers=left,
  numberstyle=\tiny\color{gray},
  numbersep=8pt,
  breaklines=true,
  showstringspaces=false,
  tabsize=2,
  xleftmargin=15pt,
  xrightmargin=5pt,
  framexleftmargin=15pt,
  captionpos=b
}
\lstset{style=pdcode}

% ─── Theorem Environments (house style) ───────────────────────────────────────
\theoremstyle{definition}
\newtheorem{definition}{Definition}[section]
\newtheorem{property}{Property}[section]
\theoremstyle{plain}
\newtheorem{theorem}{Theorem}[section]
\newtheorem{lemma}{Lemma}[section]

% ─── Algorithm / protocol blocks (numbered, framed) ───────────────────────────
\newcounter{protocol}[section]
\renewcommand{\theprotocol}{\thesection.\arabic{protocol}}
\newenvironment{protocol}[1]{%
  % Page furniture, not a conceptual figure: the house neutral ground
  % (pd neutral fill) instead of the old cream hhsand box.
  \par\vspace{8pt}\refstepcounter{protocol}\noindent
  \begin{tikzpicture}[pd figure]
  \node[pd neutral fill,rounded corners=2pt,
        inner xsep=10pt,inner ysep=10pt,text width=\dimexpr\textwidth-28pt\relax,align=left]\bgroup
  \begin{minipage}{\dimexpr\textwidth-30pt\relax}
  {\scshape\bfseries\color{pdink}Protocol~\theprotocol.}\enspace{\bfseries #1}\par\vspace{4pt}
  \small
}{%
  \end{minipage}\egroup;
  \end{tikzpicture}\par\vspace{8pt}%
}

% ─── Header/Footer (house style) ──────────────────────────────────────────────
\pagestyle{fancy}
\fancyhf{}
\fancyhead[L]{\small\itshape\color{hhgray} The Legible Swarm}
\fancyhead[R]{\small\color{hhgray}\thepage}
\fancyfoot[C]{\footnotesize\color{hhgray} Port Daddy Coordination Papers \textperiodcentered\ Chapter~\pdchapternumber\ of \pdchaptercount}
\renewcommand{\headrulewidth}{0.4pt}
\renewcommand{\footrulewidth}{0pt}

% ─── Title ───────────────────────────────────────────────────────────────────
\title{
  \vspace{-1cm}
  {\Large\textsc{Port Daddy Coordination Papers --- Chapter~\pdchapternumber\ of \pdchaptercount}}\\[
0.5cm]
  \rule{\textwidth}{0.5pt}\\[0.3cm]
  {\LARGE\textbf{The Legible Swarm:}}\\[0.2cm]
  {\Large A Consented Local Leviathan, Legibility-with-Zoom,\\and the Read-Poverty Bottleneck of Agent Swarms}\\[0.3cm]
  \rule{\textwidth}{0.5pt}
}
\author{
  \textbf{Erich Owens}\\[0.2cm]
  Curiositech LLC\\[0.2cm]
  \texttt{engineering@portdaddy.dev}
}
\date{\pdeditiondate\\Version 1.2 (textbook edition)}
% ─── Document ────────────────────────────────────────────────────────────────
% Hyperlink and cross-reference discipline (must stay last in the preamble).
% Shared hyperlink and cross-reference discipline for the Book and its chapters.
% Load this LAST in every preamble: cleveref must follow hyperref and amsthm.
% hyperref itself is loaded earlier with [backref=page] (a load-time option).
% Twin copy under book/figures/ (byte-identical; checked by
% scripts/generate-mega-whitepaper.test.mjs).
%
% What a reader gets: live links on chapter titles, section numbers, theorems,
% figures, tables, equations, page numbers, and citations; PDF bookmarks three
% levels deep; and, at the end of every bibliography entry, the pages that cite it.
% Links are quiet: a warm dark ink one step lighter than body text, so the
% four part hues keep their meaning and an inline chapter title never competes
% with the part it sits in. Chrome pages (the front-matter map, part openers)
% recolor their own links locally.
\hypersetup{
  colorlinks=true,
  linkcolor=pdinkmuted,
  citecolor=pdinkmuted,
  urlcolor=pdinkmuted,
  filecolor=pdinkmuted,
  anchorcolor=pdink,
  bookmarksnumbered=true,
  bookmarksdepth=3,
  bookmarksopen=true,
  bookmarksopenlevel=1,
  breaklinks=true,
  linktoc=all,
  pdfauthor={Erich Owens},
  pdftitle={\pdchaptertitle},
  pdfsubject={\pdeditiontitle: \pdeditionsubtitle},
}

% Title-based chapter references. \pdchaplink{prefix}{text} is a live link to
% the chapter inside the Book (the Book's preamble defines it before loading this
% file); in a standalone chapter it is plain text. Prose uses \pdchapref, which
% sets the title in italics; the locator map uses \pdchaplink directly.
\providecommand{\pdchaplink}[2]{#2}
\DeclareRobustCommand{\pdchapref}[2]{\pdchaplink{#1}{\emph{#2}}}

% \ifpdbook -- true in the assembled Book, false otherwise.
%
% WHY THIS EXISTED. A drawing that two chapters both \input prints twice in the
% Book, under two figure numbers. The fix is for the chapter that develops the
% idea to keep the figure and the other to point at it -- but the standalone
% edition of a chapter had nowhere to point, and rewriting the prose outright
% removed five figures from harbor-economy-whitepaper.pdf and left it telling a
% conference reader that a ceremony is "drawn in The Federated Harbor", a
% document that reader did not have.
%
% WHERE IT STANDS NOW. The per-chapter PDFs are retired: the Book is the one
% document, so \ifpdbook is true in every build that renders, and the \else
% branches are source with no output. They are still here, and still read by
% check_duplicate_figures.py (which must not count an \else \input as a Book
% input) -- but nothing prints them, and nothing new should be put in one.
% Whether to delete them is an open question; see the retirement PR.
%
% The divergence was per-build. A CONDITIONAL and not a two-argument macro:
% the generator inlines each \input, so a figure branch is a whole table or
% tikzpicture, and as a macro argument that is scanned as one token list --
% which is how \pdinbook produced "runaway argument" on the first attempt.
% \ifpdbook skips the untaken branch at the TeX level instead.
%
% Declared unconditionally and defaulting FALSE, which is the standalone
% answer. The Book's preamble sets \pdbooktrue immediately after loading
% this file. A guard here would have to mention \ifpdbook or \newif inside a
% branch TeX may skip, and a conditional token in skipped text is counted as
% a nested \if -- which is how two attempts at a guard both produced
% "Incomplete \ifdefined". Ordering is cheaper than cleverness.
\newif\ifpdbook

% cleveref: "Theorem 3.1", "Figure 2.4", "Section 1.3", capitalised, never
% abbreviated, with the whole phrase clickable. Names are declared explicitly
% because the theorem environments are created before this file is loaded.
\usepackage[nameinlink,capitalise,noabbrev]{cleveref}
\crefname{pdchapter}{Chapter}{Chapters}
\crefname{part}{Part}{Parts}
\crefname{appendix}{Appendix}{Appendices}
\crefname{definition}{Definition}{Definitions}
\crefname{theorem}{Theorem}{Theorems}
\crefname{lemma}{Lemma}{Lemmas}
\crefname{proposition}{Proposition}{Propositions}
\crefname{corollary}{Corollary}{Corollaries}
\crefname{conjecture}{Conjecture}{Conjectures}
\crefname{property}{Property}{Properties}
\crefname{protocol}{Protocol}{Protocols}
\crefname{lsprotocol}{Protocol}{Protocols}
\crefname{stpprotocol}{Protocol}{Protocols}
\crefname{heprotocol}{Protocol}{Protocols}
\crefname{remark}{Remark}{Remarks}
\crefname{lstlisting}{Listing}{Listings}
\crefname{exercise}{Exercise}{Exercises}
\crefname{solution}{Solution}{Solutions}
\Crefname{pdchapter}{Chapter}{Chapters}
\Crefname{lstlisting}{Listing}{Listings}
\Crefname{exercise}{Exercise}{Exercises}
\Crefname{solution}{Solution}{Solutions}

% Back-references: "^ Cited on pages 12, 40, and 87." after every bibliography
% entry, each page number a live link back to the page that cited the work.
% Distinct pages only (#3/#4).
%
% The page numbers used to be wrapped in \color{hhgray} along with the words,
% which painted over the link colour hyperref had just given them: the links
% were there -- 29 annotations on a typical bibliography page -- and looked
% exactly like the grey prose around them, so nobody knew they could click.
% Restoring the quiet link colour would not have fixed it either; pdinkmuted
% against hhgray is 1.53:1, which is no signal at footnote size.
%
% So the numbers take pdcobalt (7.55:1 on the cream, and a hue a reader cannot
% mistake for grey), and the line opens with the return arrow this book already
% uses to send a solution back to its exercise. This is a deliberate exception
% to the quiet-link discipline above, and the reason is that back-matter
% navigation is the one place where the link IS the content: in prose the
% surrounding sentence tells you a phrase is a destination, and in a
% bibliography nothing does.
\renewcommand*{\backref}[1]{}
\renewcommand*{\backrefalt}[4]{%
  \ifcase #3\relax
  \or {\footnotesize{\color{hhgray}$\uparrow$~Cited on page}~{\hypersetup{linkcolor=pdcobalt}#4}{\color{hhgray}.}}%
  \else {\footnotesize{\color{hhgray}$\uparrow$~Cited on pages}~{\hypersetup{linkcolor=pdcobalt}#4}{\color{hhgray}.}}%
  \fi}
\renewcommand*{\backrefsep}{, }
\renewcommand*{\backreftwosep}{ and }
\renewcommand*{\backreflastsep}{, and }

\begin{document}
\pdopensolutions
\setlength{\emergencystretch}{2em}
\maketitle
\thispagestyle{empty}
\vspace{0.6cm}
\vspace{0.4cm}
\noindent\textbf{Keywords:} legibility, multi-agent coordination, human-in-the-loop,
signal detection theory, context compaction, agent discovery, consent, operator
attention, Hobbes, Scott
\vspace{0.4cm}
\section{Introduction: the swarm is a state of \mbox{nature}}\label{sec:intro}
\subsection{A worked afternoon: Alice's checkout repair}\label{sec:mara}\label{sec:mira}
Return to Bob's checkout repair. Alice supervises Drake, Wren, Pike, Finch,
and Heron in her isolated worktree. The acceptance rule is concrete: retain the
original checkout-flow integration test, make the authorized repair pass it,
and leave unrelated migrations alone. The following timings and review budget
are a constructed teaching case, not a measured run.

\scene{Alice, 2:40 p.m.}{Five processes begin from revision~7. By 3:10 she has
five proposed changes. The first question is not how much code they produced,
but which changes satisfy Bob's request and which she is authorized to admit.}
\begin{enumerate}[leftmargin=1.6em]
\item \textbf{Conflicting proposals.} Wren and Heron both change the payment
module from revision~7. Applying either proposal without checking the other's
base can discard useful work. A registered claim helps coordinate intent;
preventing the stale write also requires a mediated file-effect boundary.
\item \textbf{The unread pile.} Wren's genuine session-token repair sits inside
a $600$-line patch that also reverts an unrelated migration. Alice has ninety
uninterrupted seconds before the next review request. The correct fragment
does not make the whole patch acceptable.
\item \textbf{False completion.} Drake reports ``all green'' after deleting
the original failing test. The report confuses passing the remaining suite
with satisfying the preserved acceptance criterion. Heron's accepted
contribution must be checked against that original criterion.
\item \textbf{A stale destructive proposal.} Finch proposes a force-push
that would discard an unmerged contribution. Whether recovery is possible
or not, that effect lies outside the repair grant and needs a separate decision.
\end{enumerate}
\begin{pdexample}{Alice's Half Hour, Counted}
Thirty minutes yielded five proposed changes: one every six minutes. That
arrival interval alone does not establish overload. Now state the inspection
budget. At an illustrative two lines per second, the $600$-line patch needs
$300$ seconds, while Alice's ninety-second interval covers only $180$ lines.
The shortage is uninterrupted inspection time under this assumed rate.

A useful digest points Alice to the unrelated migration and the test deletion,
then preserves the artifacts needed to check them. It cannot make either change
acceptable by shortening its description. The next table separates these
failures because exclusion, evidence, and authorization solve different parts
of the afternoon. A coordinator adds a fifth possible failure: an unexplained
refusal whose grounds Alice cannot inspect.
\emph{Now you try:} match each failure to its required primitive in
Exercise~\ref{ex:ls-taxonomy-five-failures}.
\end{pdexample}

The chapter develops three decisions from this job: what Alice may delegate,
what evidence she must be able to inspect, and how to allocate limited attention.
Hobbes' question of a common authority supplies a comparison, not a theorem
that software agents are people or that a central daemon is the only solution
\pdcite{hobbes1651}. The proposed authority must itself be inspectable and
revocable; Section~\ref{sec:consent-canon} tests the political analogy's limits.

\paragraph{Two routes through the argument.}
For \emph{supervision decisions}, follow Sections~\ref{sec:intro}--\ref{sec:operator}:
what Alice authorizes, which evidence a completion must carry, and when she
must inspect it. These sections develop the operating rules for the checkout
job; the later mathematics measures constraints on those rules.

For \emph{information bounds}, continue through
Sections~\ref{sec:readpoverty}--\ref{sec:tokens}: how much a summary must preserve,
how a directory selects candidates, and what compaction can lose. A reader
starting there should first read the opening case and the verifiable-zoom
example, so that the formal error events retain their operational meaning.

\subsection{The failure taxonomy}\label{sec:taxonomy}
Table~\ref{tab:taxonomy} groups the failures in the opening scenario by
the coordination primitive intended to address each one.
\begin{table}[H]
\centering
\small
\renewcommand{\arraystretch}{1.3}
\begin{tabular}{@{}>{\raggedright\arraybackslash}p{0.20\textwidth}>{\raggedright\arraybackslash}p{0.34\textwidth}>{\raggedright\arraybackslash}p{0.38\textwidth}@{}}
\toprule
\textbf{Failure} & \textbf{Concrete form in a coding swarm} & \textbf{The primitive that answers it} \\
\midrule
Resource\newline conflict & Two agents edit the same file/region; one silently clobbers the other &
A \textbf{claim} --- an advisory announcement that an agent intends to touch a
file or region --- backed by a database uniqueness constraint that prevents
two holders of the same registered key. Overlapping regions
need a separate overlap check and mediated writes. \\
Illegibility & A pull request or session whose intent and effect a human cannot read in bounded time &
A \textbf{digest-with-zoom} surface (\S\ref{sec:zoom}): a summary that is a lens onto
the underlying artifact, never a replacement for it. \\
Lies or\newline confident\newline errors & An agent reports completion (``all green'') without having verified &
An \textbf{honest-attestation} discipline: a self-report that distinguishes
\emph{verified-good} from \emph{no-evidence-of-bad} and refuses to claim what it
cannot check --- ``a green that wasn't checked is a lie.'' \\
Footguns & Irreversible action (force-push, delete) on a stale or wrong premise &
Footgun guards (\S\ref{sec:landing}); the open problem this chapter frames is
replacing hand-maintained rule corpora with a structural authority. \\
Illegible\newline authority & The coordinator blocks an agent and the operator cannot see \emph{why} &
The \emph{legible-sovereign rule} (\S\ref{sec:authority}): every act of authority is
a logged, named, zoomable event. \\
\bottomrule
\end{tabular}
\caption{The state-of-nature failure taxonomy, the concrete coding-swarm form of
each, and the class of primitive that answers it. The fifth row --- \emph{illegible
authority} --- is the additional design problem created by the coordinator itself
in this chapter (\S\ref{sec:authority}).}
\label{tab:taxonomy}
\end{table}
The striking recent result is that the identification is not merely a useful
analogy. \textbf{Dai et al.\ 2024}, \emph{Artificial Leviathan}
(arXiv:2406.14373)~\pdcite{dai2024} --- \emph{a sandbox of LLM agents with
psychological drives that, starting from unrestrained conflict resembling the state
of nature, were observed to establish social contracts, authorize a sovereign, and
form ``a peaceful commonwealth founded on mutual cooperation.''} The arc this
series posits as its organizing metaphor is, in at least one controlled study, an
\emph{emergent dynamic} of LLM populations (the methodological ancestor is
\textbf{Park et al.\ 2023}, \emph{Generative Agents}~\pdcite{park2023}, where LLM
agents form emergent social behavior in a sandbox town). This motivates a
design hypothesis; it does not establish a universal attractor for coding swarms.
The wager of this chapter is that we can build the consented sovereign
\emph{deliberately, locally, and legibly} rather than let an illegible one congeal
by accident. Figure~\ref{fig:state-of-nature} follows three concrete
\texttt{settings.json} edits from their launch source through the write and into
the failure they cause---or conceal.
% Brief (tikz-diagram-craft). Figure 4.1, The Legible Swarm.
% Reader question: what can actually happen when three independently launched
%   agents edit the same developer settings file from the same revision?
% Claim: an ungoverned overlap can converge, break loudly, or remain valid and
%   fail later; visible claims, locks, and downstream evidence are what let an
%   operator distinguish those cases.
% Evidence: one concrete .vscode/settings.json base revision; three proposed
%   key edits; a pointer claim, a whole-file lock, and a second pointer claim;
%   three possible resolutions and their first observable consequences
%   [internal, pedagogical counterexample].
% Must distinguish: launch source from edit intent; pointer claim from file
%   lock; syntactic failure from semantic failure; test-time observation from
%   downstream consequence.
% Grammar chosen: three aligned worked traces from one shared input. Equal
%   rows make the operation and first observable output directly comparable;
%   the delayed row alone continues to its later release failure.
% Grammar rejected: actor bubbles around settings.json. A membership graph
%   would show who touched the file but erase read revision, order, overlap,
%   and the delay between a green mutation score and the downstream failure.
% Five-second test: three agents read a1c9; one overlap is safe, one breaks now,
%   and one stays green until a later release consumes the hidden fixture.
\begin{figure}[H]
\centering
\pdfigurehue{pdteal}%
\begin{tikzpicture}[pd figure,x=1cm,y=1cm]
  \tikzset{
    actor/.style={pd artifact,text width=2.20cm,minimum height=1.40cm},
    operation/.style={pd state,text width=3.15cm,minimum height=1.35cm},
    result/.style={text width=3.50cm,minimum height=1.35cm}
  }

  % Every trace starts from this exact revision.
  \node[pd artifact,text width=10.35cm,align=left] (base) at (5.45,7.20)
    {\textbf{shared input}\quad \texttt{.vscode/settings.json @ a1c9}\\
     \texttt{formatOnSave=false}\quad \texttt{tsdk=node\_modules/...}\\
     \texttt{search.exclude=\{"**/dist": true\}}};
  \node[pd label,anchor=north] at (5.45,6.50)
    {all three agents read \texttt{a1c9} before any write lands};

  \node[actor] (A) at (1.30,4.82)
    {\textbf{Agent A}\\02:00 format cron\\\pdfigsub pointer claim};
  \node[operation] (AO) at (5.20,4.82)
    {set \texttt{formatOnSave}\\\texttt{false -> true}};
  \node[pd ready state,result] (AR) at (9.30,4.82)
    {\textbf{benign}\\keywise merge\\valid JSON; all keys};

  \node[actor] (B) at (1.30,2.62)
    {\textbf{Agent B}\\02:01 PnP prototype\\\pdfigsub whole-file lock};
  \node[operation] (BO) at (5.20,2.62)
    {replace \texttt{tsdk}\\by raw text splice};
  \node[pd breach state,result] (BR) at (9.30,2.62)
    {\textbf{immediate}\\comma lost\\parse fails at write};

  \node[actor] (C) at (1.30,0.42)
    {\textbf{Agent C}\\02:02 search cron\\\pdfigsub pointer claim};
  \node[operation] (CO) at (5.20,0.42)
    {add \texttt{**/fixtures}\\to \texttt{search.exclude}};
  \node[pd warn state,result] (CR) at (9.30,0.42)
    {\textbf{delayed}\\valid JSON; tests 98\%\\settings path untested};
  \node[pd breach state,result] (CL) at (9.30,-2.10)
    {\textbf{next release}\\\texttt{fixtures/v1} unseen\\replay fails};

  \begin{scope}[on background layer]
    \draw[pd focus arrow] (A.east) -- (AO.west);
    \draw[pd ready arrow] (AO.east) -- (AR.west);
    \draw[pd focus arrow] (B.east) -- (BO.west);
    \draw[pd breach arrow] (BO.east) -- (BR.west);
    \draw[pd focus arrow] (C.east) -- (CO.west);
    \draw[pd warn arrow] (CO.east) -- (CR.west);
    \draw[pd breach arrow] (CR.south) -- (CL.north);
  \end{scope}
\end{tikzpicture}
\caption{Three worked traces start from one \texttt{settings.json} revision.
A keywise merge is benign; a raw splice fails immediately; a valid exclusion
stays green until release replay exposes the hidden fixture. [internal,
pedagogical counterexample]}
\label{fig:state-of-nature}
\end{figure}

\subsection{Why ``local''}\label{sec:local}
The local model has one trusted coordinating process and one durable store.
Clients ask that coordinator to order mutations and check grants. Authentication,
isolation, and mediation are still required for the effects it claims to control;
being on one machine does not remove those requirements. This chapter assumes
that local authority. Agreement between mutually distrustful operators is an
additional problem taken up in \pdchapref{fh}{The Federated Harbor}.

% ═════════════════════════════════════════════════════════════════════════════
\section{Consent to the Leviathan (the normative claim)}\label{sec:authority}
\pdmarginfigure{hobbes}{Thomas Hobbes. This chapter borrows the question of authorized power, but makes the operator's consent scoped and revocable.}
% ═════════════════════════════════════════════════════════════════════════════
Why should the operator cede authority to a daemon --- let it block an agent's
claim, escalate a decision, reorder a merge, refuse a \texttt{done}? Hobbes'
answer, transposed: the operator authorizes the sovereign because the covenant is
rational under the alternative. Crucially, in Hobbes the covenant is \emph{among the
subjects} --- ``a real unity of them all\ldots by covenant of every man with every
man''~\pdcite{hobbes1651,sep-hobbes} --- not a bargain struck \emph{with} the
sovereign. Transposed: the operator does not negotiate with each agent; the operator
institutes a \emph{rule of the road} --- a \textbf{coordination protocol} of
\emph{typed performatives} (messages whose type declares their intent: a request, a
\textbf{commitment} (glossed in \S\ref{sec:flatten}: a violable obligation bound to a named actor, watched by a breach monitor), a delegation, a termination), so that every message carries real intent,
ownership, and a stop condition --- that all agents are bound to, and the daemon
enforces it.
The authority is asymmetric by design: the clients propose work, while one
coordinator applies the operator's policy to the effects within its boundary.
This is a central-coordinator assumption, not a consensus result. If a client
can bypass that boundary, or the coordinator cannot observe the decisive effect,
the prevention claim fails. Section~\ref{sec:landing} separates what can be
prevented from what can only be detected afterward.

\subsection{The legible-sovereign amendment}\label{sec:legible-sovereign}
Hobbes' absolutism needs one modern amendment, and it is the crux of this chapter.
Hobbes argued the sovereign, being the \emph{product} of the covenant rather than a
\emph{party} to it, can never breach it and need not be
inspectable~\pdcite{hobbes1651,sep-hobbes}. For a software authority this is exactly
wrong. An \emph{illegible authority} is the fifth state-of-nature failure mode
(Table~\ref{tab:taxonomy}) wearing a crown: if the daemon blocks an agent and the
operator cannot see why, the operator's trust --- and therefore the consent that
legitimizes the authority --- erodes.
\begin{pdclaim}{Design invariant}{The Legible-Sovereign Rule}
\setlength{\parskip}{4pt plus 1pt}
\label{prop:legible-sovereign}
The coordinating authority must be the \emph{most} legible actor in the system, not
the least.
Every act of authority --- a claim denial, an anomaly detection, an
escalation, an ``all good'' attestation --- must be a logged, named, zoomable event
whose reason the operator can reconstruct.
Consent is renewable only while the
sovereign is inspectable.
\end{pdclaim}
This is the inversion that rescues the design from despotism, and it is the move a
political theorist would press hardest (\S\ref{sec:consent-canon}). In Hobbes the
sovereign is opaque \emph{by structure}; here the sovereign is the one actor whose
every exercise of power is an audit record. It is realized as a discipline of
\textbf{honest attestation}: a single self-report that runs every invariant check,
distinguishes \emph{verified-good} from \emph{no-evidence-of-bad}, regiments what it
can, and fails loudly about what it cannot. That discipline is precisely the
sovereign making its own judgments legible --- ``all good'' becomes ``a claim with a
verifier, not vibes.''
\pitfall[Ranking can hide urgent work]{There is a regress hiding here, and Scott would press it. If every act of
authority must be legible, the \emph{legibility apparatus itself} is an exercise of
authority --- the digest decides what is surfaced and what is buried. The Attention
Queue ranker (\S\ref{sec:discovery}) is \emph{the} sovereign act: it decides what
the operator sees \emph{right now}. A ranker that silently buries an escalation is
governing opaquely. The rule is therefore not complete until the ranker's
\emph{omissions} are themselves legible --- a ``what did the queue not show me, and
why'' counter-surface. Ranking is governing; we say so explicitly in
\S\ref{sec:consent-canon}.}
\begin{pdexample}{Three Omissions, and How Many of Them a Score Explains}
\setlength{\parskip}{4pt plus 1pt}
Take the three suppressed items tabulated below and try to audit the queue the cheap way --- by score alone, against the distress threshold of $0.45$ that governs this action class.
The first item, \texttt{agent-7}'s escalation, scores $0.31$; since $0.31<0.45$ it sits below the bar and the threshold explains it, and an auditor with nothing but the scores reconstructs that omission correctly.
The second scores $0.12$, which is also below $0.45$ --- so the score-only auditor thinks the threshold explains that one too, and is wrong: it was dropped because its recency weight decayed past the floor, and a rule that consults the clock rather than the score would have aged it out at any score at all.
The third is the one that should end the argument: it scores $0.58$, above the only threshold the screen names, and it was suppressed anyway, by de-duplication into an earlier occurrence --- another rule that never looks at the number.
So of three omissions, a score-only audit explains one, misattributes one, and cannot see the third at all --- and the two it fumbles are the two where a rule other than ranking did the governing. That is the concrete content of ``ranking is governing'': the suppression \emph{rules} are what have to be legible, because the suppression \emph{scores} do not determine them.
\emph{Now you try:} specify the counter-surface, including its own decay discipline (Exercise~\ref{ex:ls-ranker-counter-surface}).
\end{pdexample}
The regress is easier to reason about once you have seen what a counter-surface would
concretely be, so here is the whole mockup screen those three rows came from.
Table~\ref{tab:counter-surface} is one screen of it: not a diagram of a mechanism
that does not exist, but the artifact the mechanism would have to produce.
\begin{table}[H]
\centering
\small
\renewcommand{\arraystretch}{1.25}
\setlength{\tabcolsep}{4pt}
% Regret column widened 0.08 -> 0.10\textwidth: its own bold header,
% "Regret", was 2.8pt wider than the column at this size.
\begin{tabularx}{\textwidth}{@{}>{\raggedright\arraybackslash}X >{\centering\arraybackslash}p{0.10\textwidth} >{\raggedright\arraybackslash}p{0.27\textwidth} >{\raggedright\arraybackslash}p{0.22\textwidth}@{}}
\toprule
\textbf{Suppressed item} & \textbf{Regret} & \textbf{Why it was not shown} & \textbf{Re-surface trigger} \\
\midrule
\texttt{agent-7} escalation: ``migration ordering unclear'' & $0.31$ & below the distress threshold $(0.45)$ for this action class & threshold drops, or the item is restated \\
\texttt{agent-3} note on \path{auth/}\allowbreak\path{session.ts} & $0.12$ & aged out; recency weight decayed past the floor & any new claim on the same file \\
Third retry of the same failing test & $0.58$ & de-duplicated into the first occurrence & a fourth retry, or a different failure mode \\
\bottomrule
\end{tabularx}
\caption{Illustrative counter-surface: each omitted item names its suppression rule and return condition. This view is itself ranked, so enumerating suppression rules does not make every omission visible.}
\label{tab:counter-surface}
\end{table}
\subsection{Consent as a first-class, revocable primitive}\label{sec:consent}
It is easy to say the operator ``consents to cede decision authority.'' A political
theorist's first question --- and the one that decides whether the Leviathan frame has
substance or is merely decorative --- is: \emph{consented how, to what, and revocable when?}
``The operator just starts using the tool'' is
not Hobbesian consent; it is what the canon calls \textbf{tacit consent}, and
\textbf{Hume} (\emph{Of the Original Contract}, 1748~\pdcite{hume1748} --- \emph{the
decisive critique: the peasant who ``consents'' by remaining in the country has no
real alternative, so the consent does no normative work}) demolished it two and a
half centuries ago. We therefore make consent a concrete object.
\scene{Alice, 3:25 p.m.}{Alice grants two hours of low-stakes, reversible
formatting changes in \texttt{tests/}, provided the original acceptance test
remains intact. A formatting change can satisfy all these conditions. Drake's
test deletion fails the preserved-test condition; Wren's migration change is
outside the subtree; Finch's force-push is outside the grant. The fact that a
process can propose them gives it no permission to execute them.}

The tuple below records the conditions that distinguished those four proposals.
Its scope must include the acceptance-preservation rule, not merely a path glob.

\begin{pdclaim}{Design invariant}{Consent Grant}
\setlength{\parskip}{4pt plus 1pt}
\label{def:consent}
A \textbf{consent grant} is an explicit, scoped, revocable authorization the
operator issues to the daemon, of the form
\[
\begin{aligned}
g = \langle\, &\textsc{scope},\ \textsc{stakes-ceiling}\ s_{\max},\ \textsc{reversibility-floor}\ v_{\min},\\
&\textsc{ttl},\ \textsc{revocable}\,\rangle,
\end{aligned}
\]
For Alice's grant, the daemon may apply reversible formatting changes under
stakes threshold $s_{\max}$ in \texttt{tests/}, preserving the original checkout
acceptance test, for two hours. This grant permits neither unrelated changes
nor deployment.

A grant is a
first-class, legible object: it appears in the digest, it has an audit trail, and it
expires.

The authority's legitimate action set at any moment is exactly the union of
its active grants.
\end{pdclaim}

Consent so defined gives the design something Hume's critique cannot reach: a
discrete authorization event, a bounded scope, and an exit. It is also an object
supervisory-control research named long before this chapter did: a
\textbf{level of automation}, in Sheridan's ten-point sense and in the
four-stage refinement (acquire, analyse, decide, act) that Parasuraman,
Sheridan and Wickens later gave it~\pdcite{sheridan1992,parasuraman2000}. Read
that way, a grant's $\langle\textsc{scope}, s_{\max}, v_{\min},
\textsc{ttl}\rangle$ is a level-of-automation setting chosen per action class
rather than once for the whole system, which is roughly what that literature
has been recommending since 1978. This section built the same selector from a
Hobbesian covenant, and that two routes so far apart arrive at a compatible
object is some evidence the object is the right one. It also hands the
accompanying chapters a concrete object. A grant is precisely what a \emph{hosted-trust}
offering transfers --- ``we make this fleet legible to you, under this scope'' is
the sale of a consent grant to a third party, which \pdchapref{he}{The Harbor Economy}
prices --- which is why \S\ref{sec:bridge} lists consent among the things this layer
provides upward.

\begin{protocol}{Consent-grant lifecycle}
\begin{enumerate}[leftmargin=1.4em,itemsep=1pt]
\item \textbf{Issue.} The operator issues a grant
$g=\langle\textsc{scope},s_{\max},v_{\min},\textsc{ttl}\rangle$. The grant is
appended to the audit log as a named, zoomable event and surfaced in the digest.
\item \textbf{Act.} On a candidate action $x$, the daemon admits the action without
asking \emph{iff} $x$ is in scope \textbf{and} $\mathrm{stakes}(x)\le s_{\max}$
\textbf{and} $\mathrm{reversibility}(x)\ge v_{\min}$ \textbf{and} the grant is
unexpired and unrevoked; otherwise it escalates to the operator.
\item \textbf{Witness.} Every action taken under $g$ records $(g,\,x,\,$reason$,\,$
artifact-link$)$ --- the legible-sovereign rule (Invariant~\ref{prop:legible-sovereign})
binds here: the act is reconstructable.
\item \textbf{Expire / revoke.} On TTL or operator revocation, $g$ leaves the active
set; the daemon's legitimate action set shrinks accordingly. The override
(Invariant~\ref{prop:override}) revokes future use of all grants and stops
not-yet-committed acts at the mediation boundary. Already committed effects
require separate recovery.
\end{enumerate}
\end{protocol}

% Candidate contract and evidence boundaries: BRIEF.md / REPORT.md.
% All coordinates are physical TeX points. No scaling or local font reduction.
\begin{figure}[H]
\centering
\pdfigurehue{pdteal}%
\SGMeasuredFigure{I/fig:consent-lifecycle}{%
\begin{tikzpicture}[pd figure,x=1pt,y=-1pt,
  cl text/.style={pd label,anchor=west,inner sep=0pt,outer sep=0pt,align=left},
  cl edge/.style={pd label,anchor=center,inner sep=0pt,outer sep=0pt},
  cl flow/.style={pd arrow,line width=.8pt}]
  % Grant state: issue once; removal has no outgoing transition.
  \node[cl text] (operator) at (0,26) {Operator};
  \draw[cl flow,draw=pdteal] (47,26)--(94,26);
  \node[cl edge] (issue) at (70,12) {issues};
  \draw[draw=pdteal,line width=.8pt,fill=pdteal!5]
    (104,0)--(163,0)--(173,10)--(173,44)--(104,44)--cycle;
  \draw[draw=pdteal,line width=.8pt] (163,0)--(163,10)--(173,10);
  \node[cl text,font=\pdfiglabelbold] (grant) at (115,15) {Grant $g$};
  \node[cl text] (active) at (115,32) {active};
  \draw[cl flow,draw=pdamber] (182,26)--(241,26);
  \node[cl edge] (expiry) at (209,12) {expiry / revoke};
  \draw[draw=pdamber,line width=.9pt] (246,17)--(246,40);
  \draw[draw=pdamber,line width=.9pt] (249,17)--(249,40);
  \node[cl text,text width=65pt] (removed) at (258,28)
    {$g$ removed\\from active set};
  \draw[pd hairline] (0,64)--(323,64);

  % Action decisions are not grant states. Re-read validity for every x.
  \node[cl text,font=\pdfiglabelbold] (each) at (30,89)
    {For every action $x$, check current $g$};
  \draw[draw=pdinkmuted,line width=.7pt] (30,110) rectangle (191,248);
  \node[cl text,font=\pdfiglabelbold] (conjunction) at (42,126) {All five must hold (AND)};
  \node[cl text] (scope) at (42,150) {$x$ in scope};
  \node[cl text] (stakes) at (42,170) {$\mathrm{stakes}(x)\le s_{\max}$};
  \node[cl text] (reversibility) at (42,190) {$\mathrm{reversibility}(x)\ge v_{\min}$};
  \node[cl text] (unexpired) at (42,210) {$g$ unexpired};
  \node[cl text] (unrevoked) at (42,230) {$g$ unrevoked};

  \draw[cl flow,draw=pdteal] (191,148)--(240,148);
  \node[cl edge] (all) at (215,134) {all five};
  \node[cl text,font=\pdfiglabelbold,text=pdteal] (admit) at (250,148) {Admit};
  \node[cl text] (permission) at (250,165) {may act};
  \node[cl text] (noasking) at (250,182) {without asking};
  \draw[cl flow,draw=pdamber] (191,215)--(240,215);
  \node[cl edge] (otherwise) at (218,198) {otherwise};
  \node[cl text,font=\pdfiglabelbold] (escalate) at (250,215) {Escalate};
  \node[cl text] (recipient) at (250,232) {to operator};
  \node[cl text] (noexecution) at (250,249) {no automatic act};

  % This loop returns to the next action check, never to the grant-state spine.
  \draw[cl flow,draw=pdinkmuted] (30,270)--(10,270)--(10,89)--(22,89);
  \node[cl text] (next) at (40,270) {Next candidate: recheck, never renew};

  % A taken action requires a record; admission alone is not an action taken.
  \draw[draw=pdinkmuted,line width=.7pt,fill=pdpage]
    (30,293)--(298,293)--(310,305)--(310,344)--(30,344)--cycle;
  \draw[draw=pdinkmuted,line width=.7pt] (298,293)--(298,305)--(310,305);
  \node[cl text] (witnesscondition) at (42,310)
    {If $x$ is taken under $g$, record its witness:};
  \node[cl text,font=\pdfiglabelbold] (witness) at (42,329)
    {$(g,\,x,\,\text{reason},\,\text{artifact-link})$};

  % Real text-node extents for the bounded offline geometry audit.
  \foreach \name in {operator,issue,grant,active,expiry,removed,each,
    conjunction,scope,stakes,reversibility,unexpired,unrevoked,all,admit,
    permission,noasking,otherwise,escalate,recipient,noexecution,next,
    witnesscondition,witness}{%
    \path let \p1=(\name.north west),\p2=(\name.south east) in
      \pgfextra{\typeout{CL-NODE: \name,\x1,\y1,\x2,\y2}};}
\end{tikzpicture}%
}
\caption{\textbf{Vision.} Permission is not execution, and an action decision does not renew the grant.}
\label{fig:consent-lifecycle}
\end{figure}

Figure~\ref{fig:consent-lifecycle} separates the grant's state from repeated
action decisions: all five conditions apply to every admission. Neither
decision renews $g$; expiry or revocation removes it from the active set.

\begin{pdclaim}{Design invariant}{The Inalienable Operator-Override}
\setlength{\parskip}{4pt plus 1pt}
\label{prop:override}
Hobbes grants the subject exactly one inalienable right --- self-preservation; you
cannot covenant away the right to resist when your \emph{life} is at
stake~\pdcite{hobbes1651}.

The computational analog is an always-available,
unconditional operator-override: a \emph{stop now / I withdraw authority} channel
that \emph{no autonomy level can suppress}.

If the system reserves a highest-priority
\emph{distress} lane for the daemon to interrupt the operator, the override is its
dual --- the operator-to-daemon channel --- and it is a \emph{right}, not a feature:
it overrides grants and decisions before their effects commit at the
mediated boundary; it cannot undo an already committed effect.
\end{pdclaim}

\ifpdmargincolumn\pdconsentregion\fi
Computer security reached the same object by its own road. A consent grant is,
to a capability theorist, a \textbf{capability with contextual caveats} ---
unforgeable and able only to narrow as it is passed along; revocation
additionally requires a verifier-side revocation or expiry policy~\pdcite{macaroons2014} --- and that a Hobbesian covenant and a Macaroon
agree about what the object has to be is a small reassurance that neither
derivation was decorative.

This pairs with \textbf{Hirschman}'s \emph{Exit, Voice, and Loyalty}
(1970~\pdcite{hirschman1970} --- \emph{the choice between leaving (exit) and
complaining (voice) as the two levers a member has over a declining organization}):
the design has rich \emph{voice} machinery (escalation, annotation, the attention
queue) and the override supplies the missing \emph{exit}. Exit is the discipline that keeps
voice honest; an operator who cannot cheaply withdraw a grant has no real leverage
over the sovereign.

\pdexercisepointer{\ref{ex:ls-express-vs-tacit-consent}}{\ref{ex:ls-ranker-counter-surface}}{\pageref{ex:ls-express-vs-tacit-consent}}

% ═════════════════════════════════════════════════════════════════════════════
% Reserve the complete two-paragraph opening before the forced reflection plate.
% Measured opening with heading/gaps: about 343pt; 24 baselines reserve 349pt.
\Needspace{24\baselineskip}
\section{The mechanism: legibility-with-zoom}\label{sec:zoom}
% ═════════════════════════════════════════════════════════════════════════════

\begin{pdexample}{The Same Pull Request, Rendered Twice}
\setlength{\parskip}{4pt plus 1pt}
Before the rule, the two rows it separates --- both one line long, both about Wren's checkout-repair pull request, both perfectly honest-looking.

Row~A reads \emph{Wren $\cdot$ checkout repair $\cdot$ open $\cdot$ diff and log attached}, and every noun in it resolves: press one key and the $600$-line diff and the verbatim test log open beside each other. The repair passes the original test, but the diff also changes an unrelated migration.

Row~B reads \emph{Wren $\cdot$ checkout repair $\cdot$ all green}, and nothing in it resolves to anything --- there is no diff behind it, no log, no committed file, only the sentence.

Count what each row buys. Both compress $600$ lines of work into one line, a ratio of $600{:}1$, and that ratio is not the difference between them; a digest that refused to compress would just be the diff again.

The difference is that Row~A's single line has a route back to $600$ lines Alice owns and Row~B's has a route back to nothing, so Row~A's compression is recoverable and Row~B's is a claim about a forest nobody has walked into.

Row~A supplies a route to evidence; Row~B supplies only an assertion. A dishonest agent could still fabricate or reuse Row~A's log, so a link alone does not certify execution. Alice must be able to check the tested tree, command, result, and independently controlled runner that binds them together. The definition in this section names the navigable route $\zeta$; the verifiable-zoom rule then specifies what must be at its end.

\emph{Now you try:} run the four questions of \S\ref{sec:four-questions} against a read-surface you actually use (Exercise~\ref{ex:ls-four-questions-dashboard}).
\end{pdexample}

\pdmarginfigure{scott}{Scott's warning: a legible summary can erase the local knowledge that sustains the system.}%
A referee that only prevents collisions would be a glorified lockfile. Legibility
lets the operator inspect the work it coordinates. \textbf{James~C.\ Scott}'s
\emph{Seeing Like a State} (1998)~\pdcite{scott1998} describes how cadastral maps,
standardized surnames, the metric grid and scientific forestry make populations
and resources countable. His warning concerns treating that simplified view as a
replacement for local practical knowledge: \textbf{m\^etis}, the hands-on,
case-specific know-how the system depends on. In his account of the German
\emph{Normalbaum} (``scientific forest''), same-age, single-species rows thrived
for a generation, then collapsed after the simplification had removed soil fungi,
undergrowth and species diversity~\pdcite{scott1998}. Rao's \emph{A Big Little Idea
Called Legibility} (2010)~\pdcite{rao2010} carries this critique into software
systems that flatten reality to one purpose.

In agent oversight, the operator is the state, the swarm the population, and the
dashboard the cadastral map. A grid of green tiles invites the operator to
\emph{govern the map}. The software failure is a \textbf{Potemkin digest}: its
tiles assert a reality nobody can check because the summary has discarded the
diff, test output, error and reasoning behind the work. The operator admires the
yield table while the forest dies.

\ifdefined\pdreflectionplate\pdreflectionplate{the-unseen-boat}\fi


\begin{pdclaim}{Design invariant}{The One Law: Digest-with-Zoom}
\setlength{\parskip}{4pt plus 1pt}
\label{def:zoom}
Every summary is a lens onto the real artifact, never a replacement for it.
Summarize the \textbf{state}; link to the \textbf{work}.

A digest you cannot zoom is
a high-modernist map. A digest that reaches the underlying artifact preserves
access to recorded evidence: the operator can descend from abstraction and check
the claim. Tacit practitioner knowledge, or m\^etis, is a further requirement;
not all of it can be recorded.

Formally, a read-surface is a pair $(\sigma, \zeta)$ where $\sigma$ is a
projection (the summary) and $\zeta$ is a total \emph{zoom} function mapping every
summarized claim to a verifiable artifact; a surface with a partial or absent
$\zeta$ is a Potemkin digest.
\end{pdclaim}

% Book redraw: I/fig:zoom-vs-potemkin.
% Source-derived conceptual legibility comparison; no fresh experiment, proof, model,
% provider, runtime, or deployment evidence is claimed.
% Owner boundary: book/legible-swarm.tex, sec:flatten / sec:zoom.
% Two claim trees of identical shape: digest as lens with total zoom zeta
% resolving to concrete reviewable artifacts vs digest as facade terminating in empty stops.
\providecommand{\varnothing}{\emptyset}
\begin{figure}[H]
\centering
\SGMeasuredFigure{I/fig:zoom-vs-potemkin}{%
\begin{tikzpicture}[sg/base]
  % Native canvas: 316pt wide, comfortably below the 325.215pt Book column.
  % Top: Digest as lens (PR 123 with verifiable zoom to concrete reviewable artifacts).
  % Bottom: Digest as facade (ALL GREEN terminating in empty stops).

  % --- Section 1: Digest as lens (top) ----------------------------------
  \node[anchor=west,font=\SGType\bfseries,text=sgTeal] at (0,230) {digest as lens};
  \node[anchor=east,font=\SGType,text=sgInk!70] at (316,230) {total zoom $\zeta$};
  \draw[sgTeal,line width=1.2pt] (0,224)--(316,224);

  % Root summary card
  \draw[draw=sgInk,fill=sgTeal!8,line width=.85pt] (108,202) rectangle (208,220);
  \node[font=\SGType\bfseries,text=sgTeal] at (158,211) {PR 123: 3 claims};

  % Conduit distributor to 3 claims
  \draw[sg/arrow] (158,202)--(158,194)--(56,194)--(56,184);
  \draw[sg/arrow] (158,194)--(158,184);
  \draw[sg/arrow] (158,194)--(260,194)--(260,184);
  % moved from (163,194) -- sat centred on the horizontal spreader segment
  % it was meant to annotate; re-routed beside the rightmost branch instead,
  % matching how "resolves to" already sits beside its own row (below).
  \node[anchor=west,font=\SGType,text=sgInk!70] at (266,194) {claims};

  % 3 Claim cards
  \draw[draw=sgInk,fill=white,line width=.75pt] (18,166) rectangle (94,184);
  \node[font=\SGType\bfseries,text=sgInk] at (56,175) {auth refactor};

  \draw[draw=sgInk,fill=white,line width=.75pt] (120,166) rectangle (196,184);
  \node[font=\SGType\bfseries,text=sgInk] at (158,175) {tests pass};

  \draw[draw=sgInk,fill=white,line width=.75pt] (222,166) rectangle (298,184);
  \node[font=\SGType\bfseries,text=sgInk] at (260,175) {1 critique};

  % Conduits resolving to concrete artifacts
  \draw[sg/arrow] (56,166)--(56,154);
  \draw[sg/arrow] (158,166)--(158,154);
  \draw[sg/arrow] (260,166)--(260,154);
  \node[anchor=west,font=\SGType\bfseries,text=sgTeal] at (266,160) {resolves to};

  % 3 Folded document cards (reviewable artifacts)
  \draw[draw=sgInk,fill=sgTeal!6,line width=.75pt]
    (18,128)--(94,128)--(94,147)--(87,154)--(18,154)--cycle;
  \draw[sgInk,line width=.55pt] (87,154)--(87,147)--(94,147);
  \node[font=\SGType\bfseries,text=sgTeal] at (56,144) {\texttt{diff a1c\ldots}};
  \node[font=\SGType,text=sgInk!70] at (56,134) {commit diff};

  \draw[draw=sgInk,fill=sgTeal!6,line width=.75pt]
    (120,128)--(196,128)--(196,147)--(189,154)--(120,154)--cycle;
  \draw[sgInk,line width=.55pt] (189,154)--(189,147)--(196,147);
  \node[font=\SGType\bfseries,text=sgTeal] at (158,144) {\texttt{run \#4812}};
  \node[font=\SGType,text=sgInk!70] at (158,134) {test run};

  \draw[draw=sgInk,fill=sgTeal!6,line width=.75pt]
    (222,128)--(298,128)--(298,147)--(291,154)--(222,154)--cycle;
  \draw[sgInk,line width=.55pt] (291,154)--(291,147)--(298,147);
  \node[font=\SGType\bfseries,text=sgTeal] at (260,144) {\texttt{comment \#7}};
  \node[font=\SGType,text=sgInk!70] at (260,134) {review note};

  % --- Divider ---------------------------------------------------------
  \path[draw=sgGrid,line width=.4pt] (0,116)--(316,116);

  % --- Section 2: Digest as facade (bottom) -----------------------------
  \node[anchor=west,font=\SGType\bfseries,text=sgRed] at (0,103) {digest as facade};
  \node[anchor=east,font=\SGType,text=sgInk!70] at (316,103) {absent zoom ($\zeta = \varnothing$)};
  \draw[sgRed,line width=1.2pt] (0,97)--(316,97);

  % Root summary card (identical shape)
  \draw[draw=sgInk,fill=sgRed!8,line width=.85pt] (108,75) rectangle (208,93);
  \node[font=\SGType\bfseries,text=sgRed] at (158,84) {ALL GREEN};

  % Conduit distributor to 3 ungrounded claims
  \draw[sgInk!60,line width=.75pt,-{Stealth[length=4.5pt,width=3.6pt]}]
    (158,75)--(158,67)--(56,67)--(56,57);
  \draw[sgInk!60,line width=.75pt,-{Stealth[length=4.5pt,width=3.6pt]}]
    (158,67)--(158,57);
  \draw[sgInk!60,line width=.75pt,-{Stealth[length=4.5pt,width=3.6pt]}]
    (158,67)--(260,67)--(260,57);
  % moved from (163,67) for the same reason as the top section's "claims".
  \node[anchor=west,font=\SGType,text=sgInk!70] at (266,67) {claims};

  % 3 Ungrounded claim cards (dashed border, identical dimensions)
  \draw[draw=sgInk!50,dashed,line width=.75pt,fill=sgRed!3] (18,39) rectangle (94,57);
  \node[font=\SGType,text=sgInk!70] at (56,48) {auth refactor};

  \draw[draw=sgInk!50,dashed,line width=.75pt,fill=sgRed!3] (120,39) rectangle (196,57);
  \node[font=\SGType,text=sgInk!70] at (158,48) {tests pass};

  \draw[draw=sgInk!50,dashed,line width=.75pt,fill=sgRed!3] (222,39) rectangle (298,57);
  \node[font=\SGType,text=sgInk!70] at (260,48) {1 critique};

  % Conduits resolving to empty stops
  \draw[draw=sgRed,dashed,line width=.85pt] (56,39)--(56,21);
  \draw[draw=sgRed,dashed,line width=.85pt] (158,39)--(158,21);
  \draw[draw=sgRed,dashed,line width=.85pt] (260,39)--(260,21);
  \node[anchor=west,font=\SGType\bfseries,text=sgRed] at (266,30) {resolves to};

  % Red bar stops (empty set terminal)
  \draw[draw=sgRed,line width=1.8pt] (40,21)--(72,21);
  \node[font=\SGType\bfseries,text=sgRed,align=center] at (56,9) {$\varnothing$\\\textnormal{no diff}};

  \draw[draw=sgRed,line width=1.8pt] (142,21)--(174,21);
  \node[font=\SGType\bfseries,text=sgRed,align=center] at (158,9) {$\varnothing$\\\textnormal{no run}};

  \draw[draw=sgRed,line width=1.8pt] (244,21)--(276,21);
  \node[font=\SGType\bfseries,text=sgRed,align=center] at (260,9) {$\varnothing$\\\textnormal{no note}};
\end{tikzpicture}}
\caption{Two claim trees of identical shape. Top: the legible digest's three
named claims each resolve to one artifact --- a diff, a test run, a review
comment. Bottom: a same-shaped summary whose three claims resolve to nothing
an operator can open. A safe summary is an index into evidence; a calm one
that terminates early is a Potemkin digest, whichever claim count it
advertises [internal].}
\label{fig:zoom-vs-potemkin}
\end{figure}


Figure~\ref{fig:zoom-vs-potemkin} puts the two digests side by side: one that zooms
to a verifiable artifact and one that only asserts. Three disciplines converge on
this one principle: summaries must be indexes, not
replacements; a reported result must be a checked result, not an asserted one; and
every summary must reach the artifact it summarizes. This chapter supplies the
\emph{why} (Scott plus the human-factors argument of \S\ref{sec:operator}) and the
\emph{how} (the rules below). Digest-with-zoom is the unifying statement of all three.

\subsection{What to flatten, what to preserve}\label{sec:flatten}

Scott's own distinction is the design boundary. The state may safely standardize
\emph{administrative facts it itself defines}, and must \emph{never} standardize away
the \emph{local knowledge it depends on}.

\begin{itemize}[leftmargin=1.2em]
\item \textbf{Flatten (safe --- the operator's own structured field):} identifiers,
\textbf{claim} rows, session phases, port numbers, \textbf{commitment} status (a
\emph{commitment} being a violable obligation bound to a named actor, watched by a
breach monitor), enumerated states. Administrative legibility is cheap and lossless
here.
\item \textbf{Preserve verbatim, never paraphrase (the agents' m\^etis / the ground
truth):} the diff, the test output, the error message, the reasoning trace, the
exact command run. Paraphrasing these is the over-flattening.
\end{itemize}

The slogan: \emph{summarize the STATE, link to the WORK}.

\pitfall[Artifacts cannot capture tacit knowledge]{``Zoom to the artifact'' does not fully preserve m\^etis, and a political
theorist will say so (\S\ref{sec:consent-canon}). M\^etis in Scott is not ``the
detailed data underneath the summary'' --- it is practitioner knowledge that
\emph{resists codification entirely}: the operator's gut sense that ``this agent is
about to do something dumb,'' the local knowledge of which corners of the codebase
are cursed. A read-path to the diff is more legibility, not m\^etis. The honest
admission --- which couples Scott to Bainbridge (\S\ref{sec:operator}) --- is that
\emph{some of the most valuable operator knowledge is exactly what the system cannot
make legible, and the more the supervising automation takes over, the faster that
m\^etis atrophies}. What organizational-knowledge theory calls the problem of
converting \textbf{tacit knowledge} into explicit, transferable form --- ``we know
more than we can tell''~\pdcite{polanyi1966} --- remains an open problem for
artifact-based oversight. The skill-retention tax (\S\ref{sec:abdication}) is the design's
only defense, and it is frankly the hardest part of the design to engineer.}

\subsection{The zoom target must be verifiable, not self-reported}\label{sec:verifiable}

There is a subtle trap inside ``zoom to the real thing.'' If the operator zooms to
the \emph{agent's chain-of-thought} and treats that as ground truth, the map has
merely moved down one level --- because the reasoning trace may be fiction.
\textbf{Turpin et al.\ 2023}, \emph{Language Models Don't Always Say What They Think:
Unfaithful Explanations in Chain-of-Thought Prompting} (NeurIPS)~\pdcite{turpin2023}
--- \emph{LLM-generated reasoning traces can be systematically unfaithful: plausible,
coherent, and not actually what drove the model's output (a model swayed by
reordering of options will rationalize a post-hoc reason and never mention the
bias).} Therefore:

\pdmarginanalogy{ls-map-and-lens}{Zoom must reach the evidence}{map-and-lens}{A lens enlarges the same map. A digest must instead link to independently checkable artifacts, not magnified self-report.}

\begin{pdclaim}{Design invariant}{Verifiable-Zoom}
\setlength{\parskip}{4pt plus 1pt}
\label{prop:verifiable-zoom}
The canonical zoom target is an artifact the operator can check
\emph{independently of the agent} --- the diff on disk, the test result, the
committed file.

The agent's self-narration is a \emph{secondary} lens, always
labeled as such.

Two LLMs agreeing (an agent acting and an LLM summarizing it) is
not by itself evidence of two independent checks; it may be one error reported twice.
\end{pdclaim}

This is why the source of truth must be the durable substrate --- the coordination
database and the version-control working tree --- and not a model's summary of them.
The digest is computed \emph{over} artifacts the operator owns. This motivates
an independently controlled review or evaluation mechanism. An adversarial
reviewer is one option; a protected mechanical evaluator can also satisfy the
requirement. A swarm reviewing its own work has not established that its checks
fail independently. Agreement can reproduce a shared error; the relevant
question is which evidence and failure modes the checks actually share. We
return to adversarial review as one such mechanism in
\S\ref{sec:authority-rules}.

% This heading and its lead-in sentence were landing in the last two lines of
% p. 165, with the enumerate they introduce beginning on p. 166 -- a decision
% procedure whose first question is on the other side of a page turn. Six
% baselines is the head, the one-sentence lead-in, and the first item: enough
% that the reader sees what is being decided before turning.
\needspace{6\baselineskip}
\subsection{The four design questions (a decision procedure)}\label{sec:four-questions}

For any read-surface --- Attention Queue, briefing, resurrection digest, TUI pane ---
ask, in order (Table~\ref{tab:ls-four-questions}):

Table~\ref{tab:ls-four-questions} turns the rule into four checks: the effect's
consequence, the information preserved, the evidence behind the claim, and the
authority for the decision. Apply them to Drake's row: a green summary cannot
justify landing when the preserved acceptance test is absent. The protocol
below specifies how a proposed digest can discharge these checks.


\begin{table}[H]
\centering
\caption{Four design checks for a read-surface. Each asks a distinct question; all
must be satisfied. These are review criteria, not independent formal proofs. [internal]}
\label{tab:ls-four-questions}
\small
\renewcommand{\arraystretch}{1.15}
\begin{tabularx}{\linewidth}{@{}l >{\raggedright\arraybackslash}X >{\raggedright\arraybackslash}X >{\raggedright\arraybackslash}X@{}}
\toprule
& \textbf{Question} & \textbf{What release requires} & \textbf{What the reader can check} \\
\midrule
1 & is the action irreversible? & the artifact is one key away & the artifact link resolves \\
2 & is this state or m\^etis? & administrative state is flattened; the verbatim trace is linked, not summarized & the trace link resolves \\
3 & could the claim lie? & the claim resolves beyond the author's self-report & the independent evidence resolves \\
4 & is the authority legible? & every coordinator act returns a named reason & the reason field is present and specific \\
\bottomrule
\end{tabularx}
\end{table}

\begin{protocol}{Render a digest line under the one law}
Given a candidate summary line $\sigma$ over a unit of swarm work:
\begin{enumerate}[leftmargin=1.4em,itemsep=1pt]
\item \textbf{Project state.} Compute $\sigma$ as a deterministic projection of
structured state (claim rows, phases, enums). Never let a language model paraphrase
the line itself.
\item \textbf{Attach zoom.} For every claim in $\sigma$, attach $\zeta$: a link to an
operator-verifiable artifact (the diff, the test log verbatim, the committed file).
If any claim has no such artifact, the line is a Potemkin digest --- refuse to ship
it.
\item \textbf{Forced-zoom flag.} If the operator may act irreversibly on $\sigma$
(merge, deploy, delete), set $\sigma.\textsf{force}=\textsf{true}$: the artifact is
one keystroke away and the act is gated on a zoom (\S\ref{sec:operator} sets the
threshold).
\item \textbf{Reason on authority.} If $\sigma$ records a coordinator act (a denial,
an escalation), attach its named reason. An act with no reason is illegible
authority; refuse to ship it.
\end{enumerate}
\end{protocol}

\scene{Alice, 3:32 p.m.}{She opens the digest for the five proposed pull requests.
Wren's row links to the 600-line diff and its test log; she finds the unrelated
migration and withholds approval. Drake's green result links to a different
diff, which deletes the original checkout test; that proposal fails the acceptance
condition. Heron's narrower repair remains a candidate for the independent
acceptance run. The digest has made the distinctions inspectable. It has not
proved that any proposal is correct.}

\pdexercisepointer{\ref{ex:ls-flatten-or-preserve}}{\ref{ex:ls-adversarial-legibility}}{\pageref{ex:ls-flatten-or-preserve}}

% ═════════════════════════════════════════════════════════════════════════════
\section{The Leviathan rules, it does not just watch (the authority half)}\label{sec:authority-rules}
% ═════════════════════════════════════════════════════════════════════════════

The title is \emph{Legibility \& Authority}. A read-only legibility layer is a
glorified dashboard; the whole point of Hobbes is that the sovereign does not just
\emph{see} --- it \emph{rules}. This section specifies the authority half: the
mechanisms by which the consented daemon actually decides, blocks, lands, and
gates. Table~\ref{tab:ls-authority-organs} lists the organs; the engineering status of
each is in Appendix~\ref{app:honest}.

\begin{table}[H]
\centering
\caption{Six authority mechanisms as distinct instruments crossing one consent boundary. Their inputs and outcomes differ; what they share is the recorded grant $g$: scope, named reason, audit identity, override, and an addressable return artifact. The engineering status of each is in Appendix~\ref{app:honest}. [internal]}
\label{tab:ls-authority-organs}
\small
\renewcommand{\arraystretch}{1.15}
\begin{tabularx}{\linewidth}{@{}l >{\raggedright\arraybackslash}X >{\raggedright\arraybackslash}X >{\raggedright\arraybackslash}X@{}}
\toprule
\textbf{Mechanism} & \textbf{Input} & \textbf{Outcome} & \textbf{What the grant records} \\
\midrule
roadmap & evidence & a priority & the evidence weighed and the ordering it produced \\
review & rival claims & a gate & which claim was admitted and why the others were not \\
completion & a predicate & done, or not & the predicate and the artifact that satisfied it \\
landing & conflicts & an order & the conflict set and the serialization chosen \\
guard & a proposed act & a block, or passage & the act, the rule that fired and the named reason \\
escalation & an interrupt & an operator decision & the question raised and the decision returned \\
\bottomrule
\end{tabularx}
\end{table}

\subsection{The roadmap as a versioned statement of intent}\label{sec:roadmap}

A coordinator that prevents collisions answers ``what is the swarm doing?'' It does
not answer ``is the swarm doing the \emph{right} thing?'' That question needs a
single legible artifact the authority is measured against: a \textbf{roadmap} --- a
phase-keyed, versioned, amendable statement of intent, priority, and dependency,
against which the swarm's actual output can be compared. It is not truth about the
world; calling it that would turn every lesson the work teaches into drift. The
authority's legitimacy is measured by \emph{whether the roadmap stays honest}: every
unit of work keys to a plan item, amendments are recorded as amendments, and drift
between plan and reality is itself a legible, surfaced anomaly rather than a fact to
be hidden.

A roadmap amendment also changes which earlier receipts can support the next
act. Figure~\ref{fig:ls-plan-revision} offers a proposed typed disposition:
source, acceptance-policy, and authority changes have different consequences.
A ``reuse?'' cell demands matching receipt bindings; deleting a plan node cannot
undo an external effect whose outcome is still unknown. The matrix is a
teaching contract, not a measured saving or an implemented invalidation rule.
% Reader question: which receipts remain usable after a plan revision?
% Claim: source, acceptance, and authority changes have different receipt
% consequences; an earlier unknown effect survives every plan edit.
% Evidence: constructed three-row contract from the campaign brief, not a
% measured invalidation algorithm or runtime feature.
% Grammar: aligned decision matrix; a graph would hide the conditional cells.
% Five-second test: authority revocation blocks publication in its own column.
\begin{figure}[H]
\centering
\begin{tikzpicture}[pd figure,x=1cm,y=1cm]
  \path[use as bounding box] (0,-.45) rectangle (11.2,5.45);
  \node[pd panel title,anchor=west] at (.1,5.08) {Plan revision};
  \node[pd panel title,anchor=west] at (4.0,5.08) {Compile};
  \node[pd panel title,anchor=west] at (6.45,5.08) {Accept};
  \node[pd panel title,anchor=west] at (8.55,5.08) {Publish};
  \draw[pd rule] (0,4.78)--(11.2,4.78);
  \node[pd direct label,anchor=west,text width=3.5cm] at (.1,4.18)
    {Source input changed};
  \node[pd panel title,anchor=west] at (4.0,4.18) {STALE};
  \node[pd panel title,anchor=west] at (6.45,4.18) {STALE};
  \node[pd direct label,anchor=west] at (8.55,4.18) {Revalidate};
  \draw[pd hairline] (0,3.68)--(11.2,3.68);
  \node[pd direct label,anchor=west,text width=3.5cm] at (.1,3.05)
    {Acceptance policy changed};
  \node[pd direct label,anchor=west] at (4.0,3.05) {Reuse?};
  \node[pd panel title,anchor=west] at (6.45,3.05) {STALE};
  \node[pd direct label,anchor=west] at (8.55,3.05) {Revalidate};
  \draw[pd hairline] (0,2.55)--(11.2,2.55);
  \node[pd direct label,anchor=west,text width=3.5cm] at (.1,1.9)
    {Authority revoked};
  \node[pd direct label,anchor=west] at (4.0,1.9) {Reuse?};
  \node[pd direct label,anchor=west] at (6.45,1.9) {Reuse?};
  \node[pd panel title,anchor=west] at (8.55,1.9) {BLOCK};
  \draw[pd rule] (0,1.43)--(11.2,1.43);
  \node[pd direct label,anchor=west,text width=11cm] at (.1,.97)
    {Reuse? requires matching receipt bindings; it is never automatic.};
  \node[pd direct label,anchor=west,text width=11cm] at (.1,.35)
    {Earlier UNKNOWN effect: reconcile its key and status separately.};
  \draw[pd caution rule] (0,-.12)--(11.2,-.12);
\end{tikzpicture}
\caption{Receipt reuse depends on what changed. Source edits invalidate compile and acceptance evidence; policy edits require fresh acceptance; revocation blocks publication. Unknown external effects still require reconciliation. Constructed contract, not an implemented invalidation algorithm.}
\label{fig:ls-plan-revision}
\end{figure}


\subsection{Adversarial review: the quality organ}\label{sec:review}

A swarm reviewing its own work can repeat the same error. Agreement alone does
not supply the independently checkable evidence required by
Invariant~\ref{prop:verifiable-zoom}. The proposed quality organ is therefore
\emph{adversarial and conflict-free}: a reviewer with no stake in the work, running a
bounded critique--refine protocol. This is the single-operator tool's local analog of the
\emph{neutral evaluators} --- the independent rating agencies --- that \pdchapref{stp}{From Spawn to Person} and \pdchapref{he}{The Harbor Economy}
federate across operators: the same conflict-free-judge discipline, scoped
first to one operator and then opened to a market.

Independence of stake and independence of error are different premises.
Figure~\ref{fig:judge-dependence} gives a constructed counterexample: three
conflict-free reviewers can inspect one flawed source and repeat the same
mistake. Measure dependence on labeled cases before treating agreement as
corroboration; a colluding clique requires a separate strategic test.
% Reader question: do distinct reviewers supply independent evidence?
% Claim: separate principals can repeat the same errors when all read the same
% flawed source. Evidence: a constructed four-item error-vector counterexample.
% Grammar: one common-source provenance branch and aligned outcome rows.
% Counter-reading: this is not a measured judge error rate or collusion test.
% Five-second test: trace the one source to all three identical error vectors.
\begin{figure}[H]
\centering
\begin{tikzpicture}[pd figure,x=1cm,y=1cm]
  \path[use as bounding box] (0,-.55) rectangle (11.2,4.15);
  \node[pd panel title,anchor=west] at (.1,3.85) {Common input};
  \node[pd panel title,anchor=west] at (4.2,3.85) {Distinct principals};
  \node[pd panel title,anchor=west] at (8.1,3.85) {Errors on items 1--4};
  \draw[pd hairline] (0,3.55)--(11.2,3.55);
  \node[pd artifact,minimum width=2.2cm,minimum height=.82cm] (src)
    at (1.35,1.92) {source $S$};
  \foreach \who/\yy in {A/3.0,B/1.92,C/.84} {
    \node[pd actor,minimum width=1.8cm,minimum height=.62cm] (j\who)
      at (5.1,\yy) {judge \who};
    \draw[pd arrow] (src.east) -- (j\who.west);
    \node[pd mono label,anchor=west] at (8.25,\yy) {1\quad 0\quad 1\quad 0};
  }
  \draw[pd hairline] (0,.24)--(11.2,.24);
  \node[pd note,anchor=north west] at (.1,.12)
    {1 = error; all values are constructed, not observations.};
\end{tikzpicture}
\caption{Three reviewers have separate principal identities but read the same
source $S$ and make identical errors on two of four constructed items. Distinct
principals remove one conflict of interest; they do not by themselves prove
independent errors. Error dependence needs labeled evaluation, while strategic
collusion remains a separate question [constructed example].}
\label{fig:judge-dependence}
\end{figure}


\subsection{Sole-responsibility roles and the specialization boundary}\label{sec:sole-roles}

The single-writer discipline (\S\ref{sec:local}) generalizes naturally from data files to system functions. When an invariant is critical, the authority assigns a \textbf{sole-responsibility role} (an avatar holding exclusive write capability over an invariant):
\begin{itemize}[leftmargin=1.4em,itemsep=1pt]
\item \emph{Roadmap Owner:} The single avatar authorized to commit mutations to the phase plan.
\item \emph{Test Suite Curator:} The sole agent authorized to land modifications to regression suites.
\item \emph{Release Avatar:} The sole gatekeeper authorized to execute deployment tags.
\end{itemize}

Assigning a role that way is a trade, and the trade has a number. Think of a bank:
one expert teller (the sole specialist) working a single line in strict order,
versus several ordinary tellers (the pooled generalists) sharing one line. The
expert is faster per request but has no colleagues to absorb a burst; the pool is
slower per request but almost never leaves the line standing idle behind a long
job. \textbf{$M/M/1$}/\textbf{$M/M/c$} --- \emph{the standard queueing-theory
shorthand (\textbf{Kendall 1953}~\pdcite{kendall1953}) for a single-server versus
$c$-server queue with Poisson arrivals, exponential service times, and one waiting
line served first-come-first-served} --- is the textbook way to price the trade.

For a first calculation, let requests arrive at one per hour. Two generalists, each serving one request per hour, have utilization $1/2$ and delay probability $1/3$. Their mean response time is $1+(1/3)/(2-1)=4/3$ hours. A specialist serving $1.75$ requests per hour has the same response time, $1/(1.75-1)=4/3$ hours. A rate of $2$ wins; a rate of $1.5$ loses. The boundary below generalizes this comparison. Both arrangements must already satisfy the required safety rules: a latency calculation cannot authorize an unsafe pool.

Fix the \textbf{queueing-specialization setup}: requests for an invariant-critical
function $F$ arrive as a Poisson stream at rate $\lambda_F$; a sole specialist
serves them as an $M/M/1$ queue at rate $\mu_{\mathrm{spec}}$, the pooled
alternative as an $M/M/c$ queue of generalists at rate $\mu_{\mathrm{pool}}$ each
with utilization $\rho=\lambda_F/(c\,\mu_{\mathrm{pool}})<1$ and
first-come-first-served exponential service. Assume $\lambda_F>0$,
$\mu_{\mathrm{spec}}>\lambda_F$, $w>0$, and $A\ge0$, so both queues are
stable and the comparison is defined. The objective is \emph{mean} cost: $w$
is cost per request-hour and $A$ is the accountability benefit per hour of
operating the sole-owner arrangement. Let
$C(c,\rho)$ be the \textbf{Erlang-C} delay probability (\textbf{Erlang
1917}~\pdcite{erlang1917} --- \emph{the probability that an arriving request finds
all $c$ servers busy and must wait, the quantity that makes a pool a pool}).

\begin{pdclaim}{Theorem}{Queueing Specialization Boundary}
\setlength{\parskip}{4pt plus 1pt}
\label{thm:specialization}
Under the queueing-specialization setup, sole ownership weakly dominates the pool
if and only if the skill premium clears
\begin{equation}\label{eq:specboundary}
\frac{\mu_{\mathrm{spec}}}{\mu_{\mathrm{pool}}} \;\ge\; g_A(\rho,c)
\;=\; c\rho \;+\; \frac{1}{\,1 + \dfrac{C(c,\rho)}{c(1-\rho)}
      + \dfrac{A\,\mu_{\mathrm{pool}}}{w\,\lambda_F}\,}.
\end{equation}
\end{pdclaim}

\noindent\Verified. A tempting shortcut is
$\tilde g = 1+(c-1)\lambda_F/(c\mu_{\mathrm{pool}}-\lambda_F)
= 1+(c-1)\rho/(1-\rho)$, from a naive
$M/M/1$-versus-$M/M/c$ response-time comparison. A sweep against simulation
\emph{falsifies} that shortcut in both directions. Equation~\ref{eq:specboundary}
is the exact boundary --- the companion formal paper's
Theorem~2~\pdcite{paper6}, independently re-derived and checked numerically over two
thousand $(\rho,c)$ points. Keeping the shortcut beside the exact result exposes
its failure: the two thresholds cross exactly once, so the shortcut is not merely
conservative. Figure~\ref{fig:specialization} shades both error regions.

At $A=0$ --- pure response time, accountability priced at nothing --- the boundary
reduces to
\begin{equation}\label{eq:specboundary0}
g(\rho,c) \;=\; c\rho \;+\; \frac{c(1-\rho)}{c(1-\rho)+C(c,\rho)}.
\end{equation}
Three properties fix its shape: $g(\rho,1)=1$ (a specialist against a single
generalist needs no premium at all), $g\to c$ as $\rho\to1$ (the boundary
\emph{saturates} at the capacity ratio; it never diverges), and for two generalists
it is the closed form $g(\rho,2)=1+2\rho-\rho^2$. Below the boundary, sole
ownership buys strict single-writer accountability at an explicit latency price
that must be budgeted rather than obscured; Eq.~\ref{eq:specboundary} is where that
price is written down.

% Preserve the former run-in paragraph's hyperlink slot while replacing its style.
\phantomsection
\begin{pdexample}{Sole Ownership at the Queue Boundary}
\setlength{\parskip}{4pt plus 1pt}
\label{ex:queue-boundary}
The result is two lines of algebra on textbook
formulas, which is the right way to first believe it. Sole ownership pays iff
\[
\begin{gathered}
w\lambda_F(W_{\mathrm{solo}}-W_{\mathrm{pool}}) \le A,\\
\text{ i.e.\ }1/(\mu_{\mathrm{spec}}-\lambda_F)
 \le W_{\mathrm{pool}} + A/(w\lambda_F);
\end{gathered}
\]
invert,
divide by $\mu_{\mathrm{pool}}$, and substitute the two textbook facts
\[
\begin{aligned}
\lambda_F/\mu_{\mathrm{pool}} &= c\rho \text{ and}\\
\mu_{\mathrm{pool}}W_{\mathrm{pool}} &= 1 + C(c,\rho)/(c(1-\rho)),
\end{aligned}
\]
and
Eq.~\ref{eq:specboundary} falls out.

For $c=2$ the Erlang-C value is
$C(2,\rho)=2\rho^2/(1+\rho)$, so
the second term of Eq.~\ref{eq:specboundary0} is
\[
\begin{aligned}
(1-\rho^2)/\bigl((1-\rho^2)+\rho^2\bigr) &= 1-\rho^2 \text{ and}\\
g(\rho,2)&=1+2\rho-\rho^2,
\end{aligned}
\]
checkable in three lines with no queueing theory at all.

For an illustrative roadmap-owner queue, suppose mutations arrive at
$\lambda_F = 2/\text{hr}$; the sole owner clears them at
$\mu_{\mathrm{spec}} = 3/\text{hr}$; the alternative is three generalists at
$\mu_{\mathrm{pool}} = 1.2/\text{hr}$ each.

Then $\rho = 2/3.6 = 0.556$ and the
skill premium is $r = 3/1.2 = 2.5$. Erlang-C at $(c,\rho)=(3,0.556)$ is $0.300$, so
\[
\begin{aligned}
g &= 3(0.556)\\
&\quad{}+ 3(0.444)/\bigl(3(0.444)+0.300\bigr)\\
&= 1.667 + 0.816 = 2.483.
\end{aligned}
\]

The
premium $2.5$ clears it --- barely --- so the sole owner wins:
\[
\begin{aligned}
W_{\mathrm{solo}} &= 1/(3-2) = 1.000 \text{ hr against}\\
W_{\mathrm{pool}} &= \bigl(1+0.300/1.333\bigr)/1.2\\
&= 1.021 \text{ hr},
\end{aligned}
\]
a margin of about
75 seconds per request.

The naive threshold would return
\[
\tilde g = 1 + 2(2)/(3.6-2) = 3.5
\]
and sent the work to the pool. That is the
right-hand error lens of Figure~\ref{fig:specialization}. These are model
calculations, not measured roadmap performance.

\emph{Now you try:} at the same three-agent pool, would a specialist at
$\mu_{\mathrm{spec}} = 2.8/\text{hr}$ still be worth it on response time alone?
(The premium is $2.8/1.2 = 2.33 < 2.483$ --- no, and $W_{\mathrm{solo}} = 1.25$ hr
confirms it. The boundary is genuinely tight here; note that a positive
accountability value $A$ lowers $g_A$ below $g$ and can buy the sole owner back.)
\end{pdexample}

\begin{pdexample}{Where the Exact and Naive Thresholds Cross}
\setlength{\parskip}{4pt plus 1pt}
At the round number $\rho=0.5$:
\[
g(0.5,2)=1+2(0.5)-0.5^2=1.75
\]
[verified, closed form].\footnote{Worked numbers carry a tag throughout: [verified] when a named script in the repository regenerates them, [internal] when only the chapter's own derivation does.}

The naive $\tilde g(\rho,2)=1+\rho/(1-\rho)$ gives
\[
\tilde g(0.5,2)=2.00
\]
at the same point --- already a visible gap.

Setting $g=\tilde g$ and clearing denominators leaves
\[
\rho(\rho^2-3\rho+1)=0,
\]
whose root in $(0,1)$ is
\[
\rho=(3-\sqrt5)/2\approx0.382
\]
[verified, closed form]: below it $\tilde g$ over-certifies \mbox{specialization} (false-certify), above it $\tilde g$ over-rejects (false-reject), so the naive form errs in \emph{both} directions rather than merely being conservative.

\emph{Now you try:} run the falsification sweep behind this crossing (Exercise~\ref{ex:ls-b8-specialization-run}).
\end{pdexample}

% Restyle of fig-specialization to the tikz-diagram-craft v2 standard (v1
% names and raw colours migrated to roles; the plot itself and the two
% filled regions are unchanged, ch4.md restyles).
%
% Reader question: when does a sole specialist beat a pooled service, and
% where does the proposed shortcut err?
% Claim: the true boundary is the Erlang-C threshold; the worked instance
% clears it into the sole-owner region; a simpler proposed threshold
% diverges from it rather than saturating, so it is falsified.
% Evidence: the Erlang-C boundary $\gamma_A(\rho,3)$, the superseded $\tilde\gamma$,
% and the worked point $\rho=.556,r_{\mathrm{sk}}=2.5$, book/legible-swarm.tex
% [internal, closed form].
% Must distinguish: the exact boundary (solid) from the falsified proposed
% one (dashed); which side of the exact boundary is pooled versus sole.
% Grammar chosen: phase-boundary plot -- two filled, edged regions split at
% the computed curve, the worked point marked on it.
% Grammar rejected: a bare curve with no fill would leave "which side wins"
% to the reader's arithmetic instead of to the picture.
% Five-second test: which region does the worked instance fall in?
\begin{figure}[H]
\centering
\pdfigurehue{pdteal}%
\begin{tikzpicture}[pd figure,x=.60cm,y=.68cm]
  \begin{scope}
    \clip (0,0) rectangle (12,6);
    \fill[pd focus fill]
      (0,.18) -- plot[smooth] coordinates {(0.25,.18)(1.5,.86)(3,1.57)(4.5,2.22)(6,2.82)(7.5,3.31)(9,3.69)(10.5,3.94)(11.7,4.00)}
      -- (12,4.03) -- (12,0) -- (0,0) -- cycle;
    \fill[pd neutral fill]
      (0,.18) -- plot[smooth] coordinates {(0.25,.18)(1.5,.86)(3,1.57)(4.5,2.22)(6,2.82)(7.5,3.31)(9,3.69)(10.5,3.94)(11.7,4.00)}
      -- (12,4.03) -- (12,6) -- (0,6) -- cycle;
    \draw[pd breach rule,line width=.9pt]
      plot[smooth] coordinates {(.25,.10)(1.5,.65)(3.0,1.71)(4.5,2.57)(6,4.00)(7.2,6.0)};
    \draw[pd focus rule,line width=1.2pt]
      plot[smooth] coordinates {(.25,.18)(1.5,.86)(3,1.57)(4.5,2.22)(6,2.82)(7.5,3.31)(9,3.69)(10.5,3.94)(11.7,4.00)};
    \node[pd focus datum] at (6.67,3.00) {};
  \end{scope}

  \draw[pd arrow] (0,0)--(12.35,0);
  \draw[pd arrow] (0,0)--(0,6.35);
  \foreach \x/\lab in {0/0,3/.25,6/.50,9/.75,12/1}{
    \draw[pd tick] (\x,0)--(\x,-.12);
    \node[pd label,anchor=north] at (\x,-.18) {\lab};}
  \foreach \y/\lab in {0/1,2/2,4/3,6/4}{
    \draw[pd tick] (0,\y)--(-.12,\y);
    \node[pd label,anchor=east] at (-.18,\y) {\lab};}
  \node[pd label,anchor=north,align=center,text width=3.6cm] at (6,-1.02)
    {pooled utilization $\rho$};
  \node[pd label,rotate=90,anchor=south,align=center,text width=3.0cm] at (-.82,3.1)
    {skill premium $r_{\mathrm{sk}}$};

  \node[pd focus label,anchor=west,align=left,text width=2.55cm] at (7.45,5.12)
    {exact Erlang--C\\boundary};
  \node[pd label,text=pderror,anchor=west,align=left,text width=2.5cm] at (2.05,4.20)
    {proposed shortcut\\falsified};
  \node[pd title,anchor=west] at (.75,5.32) {sole owner wins};
  \node[pd title,anchor=west] at (8.45,1.10) {pool wins};

  \node[pd label,anchor=west,align=left] at (7.08,2.42)
    {$\rho=.556$\\$r_{\mathrm{sk}}=2.5$};
  \draw[pd hairline] (6.96,2.58)--(6.67,3.00);
\end{tikzpicture}
\caption{At three agents, pooling wins below $\gamma_A(\rho,3)$ and sole ownership above it. The dashed, superseded threshold diverges incorrectly; the worked point narrowly favours sole ownership. [internal; closed form]}
\label{fig:specialization}
\end{figure}


\pitfall[Mean latency hides tail risk]{Eq.~\ref{eq:specboundary} prices \emph{mean} response time under
first-come-first-served exponential service. It does not price the tail, and it does
not price the thing an operator actually values in a roadmap owner ---
consistency of judgment across requests --- which does not appear in a
response-time objective at all. The corner the whole sole-responsibility
discipline exists for, the sole specialist \emph{failing}, is priced in
\S\ref{sec:succession-corner} below, and priced under an assumption worth
reading before the formula.}

\subsubsection{The succession corner, and who prices it}\label{sec:succession-corner}

Eq.~\ref{eq:specboundary} assumes a specialist who never stops. The corner the
whole discipline exists for is the one where they do: they crash, are reaped,
quit, or are reassigned, and a successor spins up. That corner \emph{is}
priced, in closed form, and it is priced in
\pdchapref{he}{The Harbor Economy} rather than here --- the result is a
condition on what a succession plan costs, which is an economic statement, and
the economy chapter carries it as the price of the succession rule%
\ifpdbook
\ (Thm.~\ref{he:thm:succession-price})%
\fi
, with its own verification and its own boundary. Two of its consequences belong in
this chapter's argument, because they change what a sole-responsibility role
means rather than what it costs.

The first is that skill has a ceiling here. With death at rate $\xi$ and
succession at rate $\eta$, the sole owner's mean response time decreases in
skill but never below
\[
W_\infty=\frac{\xi}{\eta(\xi+\eta)},
\]
the residual downtime a request meets when it arrives during an outage. No
premium buys that back, so past a point the honest fix for a slow sole role is
a better handoff and not a better owner. The second is that the ceiling can
bind absolutely: writing $K=W_{\mathrm{pool}}+A/(w\lambda_F)$ for the
response-time threshold in hours,
pooling wins at \emph{every} skill premium once $\xi/\eta$ clears the
succession price $D^\star=\eta K/(1-\eta K)$ --- and, in the case worth
remembering, only while $\eta K<1$. Where succession is fast enough that
$\eta K\ge 1$ there is no such threshold at all and sole ownership cannot be
ruled out on response time, which is the more common situation in a
well-run fleet and the one a reader is likeliest to be in.

\pdthesis{Both statements rest on a modelling choice the economy chapter
declares and this one inherits: the specialist can die \emph{at any time},
not only while serving. Under active-only breakdowns a fast enough specialist
escapes the floor entirely and $D^\star$ does not exist. The choice is right
for an owner who is reassigned or dead rather than worn out, and it is a
choice, not a derivation.}

\subsubsection{What heavy tails do to the boundary}\label{sec:spec-heavy-tails}

\ifpdmargincolumn\pdheavytailsketch\fi
Exponential service is the assumption a queueing reader will challenge first, and
they are right to: a rebuild or a flaky integration test can produce a long turn,
and submissions may arrive in correlated bursts. Whether a particular workload
has heavy tails or Poisson arrivals is a measurement question, not something this
example establishes. Kingman's heavy-traffic result \pdcite{kingman1961} says what
that costs in general: waiting grows with the sum of the squared coefficients of
variation of the arrival and service processes, so a queue that is fine at
$c_s^2=1$ can be badly congested at $c_s^2=6$ with the same utilization. The
worry is real. What is not obvious, and matters here, is that it lands on
\emph{both} sides of Eq.~\ref{eq:specboundary} at once.

Service time and waiting time must be kept separate. Put
$a=\lambda_F/\mu_{\mathrm{pool}}=c\rho$,
$r=\mu_{\mathrm{spec}}/\mu_{\mathrm{pool}}>a$, and
$\kappa_j=(1+c_{s,j}^2)/2$ for each service distribution. With Poisson arrivals
and finite second moments, Pollaczek--Khinchine gives the specialist's mean
\emph{sojourn} time exactly, while a variability-scaled Erlang-C approximation
\pdcite{allen1990} gives the pool's:
\begin{align}
 \mu_{\mathrm{pool}} W_{\mathrm{spec}}
   &= \frac{1}{r}+\kappa_{\mathrm{spec}}\frac{a}{r(r-a)},\\
 \mu_{\mathrm{pool}} W_{\mathrm{pool}}
   &\approx 1+\kappa_{\mathrm{pool}}
       \frac{C(c,\rho)}{c(1-\rho)}.
\end{align}
Only the second term in each expression is waiting. Multiplying the whole
sojourn time by $\kappa$ would also multiply service time, which variance at a
fixed mean does not do. The comparison is therefore
$\mu_{\mathrm{pool}}W_{\mathrm{spec}}\le
 \mu_{\mathrm{pool}}W_{\mathrm{pool}}+A\mu_{\mathrm{pool}}/(w\lambda_F)$;
it cannot in general be reduced to replacing $A$ by $A/\kappa$.

\begin{pdexample}{A Variance-Sensitive Specialization Boundary}
\setlength{\parskip}{4pt plus 1pt}
\label{hyp:spec-variance}
The expressions above recover Eq.~\ref{eq:specboundary} when both
$\kappa_j=1$. Away from exponential service, the pool expression is an
approximation, not an error bound.

Even equal variability on both sides can move
the crossing when $A=0$: for $c=2$, $\rho=1/2$, and $C=1/3$, the exponential
\mbox{threshold} is $r=1.75$; setting both $\kappa_j=2$ gives an approximate threshold
\[
r=(8+\sqrt{124})/10\approx1.914.
\]
\end{pdexample}

These are analytic comparisons, not measurements of agent workloads. A useful
test would extend \texttt{b8\_specialization.py} with matched-mean lognormal and
Pareto service, report simulation uncertainty, and measure the approximation's
crossing error. That sweep has not been run. Infinite service variance makes the
stable $M/G/1$ mean waiting time infinite, so this finite-mean comparison no longer
applies. Correlated or bursty arrivals also leave the Poisson model; neither
Pollaczek--Khinchine nor a variability correction alone establishes their delay
distribution or tail latency.

\subsection{Completionist obligation: completion as a verifier, not a vibe}\label{sec:completionist}

The operator's standing complaint is the agent that \emph{declares completion while
the work is hollow} --- ``fixed the test by deleting it.'' The completionist
obligation proposes an authority-side rule against accepting an unmet or
unevaluable completion predicate.

For the checkout job, Drake's claim ``all green'' fails before any score is
computed: the original test is missing. Heron's candidate instead needs a receipt
binding that unchanged test, the tested commit, and the delivered commit.
A passing receipt for another tree cannot close this obligation. The completion
predicate below makes those task-specific requirements the gate.

\begin{pdclaim}{Design invariant}{Completionist Completion}
\setlength{\parskip}{4pt plus 1pt}
\label{def:completionist}
A \textbf{completion claim} is a commitment whose satisfaction predicate is a
\emph{verifier} derived from the task's acceptance criteria. The authority refuses to
accept the claim against an unmet predicate.

An \emph{unverifiable} predicate must
\textbf{fail loudly} --- surface ``I cannot verify this completion'' --- and never
silently pass.

This is the authority-side analog of discovery's
\emph{refuse-to-route} (\S\ref{sec:gates}): ``I don't know'' must beat a confident
wrong answer.
\end{pdclaim}

% I/fig:completion-verifier. Source contract: def:completionist.
% Integration evidence: BOOK-COMPLETION-VERIFIER-PASS-20260920.md.
% Vision (owner status row): diagrammed obligation, not implemented enforcement.
% No predicate implementation, signature, test result or empirical count is depicted.
\usetikzlibrary{arrows.meta,calc}
\newcommand{\CompletionVerifierFigure}{%
\begin{tikzpicture}[x=1pt,y=-1pt,line cap=butt,line join=miter,
  every node/.style={pd label,inner sep=0pt,align=center},
  cv arrow/.style={-{Stealth[length=4.4pt,width=3.5pt]},draw=pdink,line width=.8pt},
  cv outcome/.style={minimum width=98pt,minimum height=46pt,text width=84pt,
    inner xsep=7pt,inner ysep=7pt,line width=.85pt,font=\pdfiglabel}]
  % The two records play different roles: assertion versus the criteria that bind evaluation.
  \path[draw=pdinkmuted,fill=pdinkmuted!3,line width=.8pt]
    (0,0)--(116,0)--(128,12)--(128,45)--(0,45)--cycle;
  \draw[pdinkmuted,line width=.6pt] (116,0)--(116,12)--(128,12);
  \node (claim) at (64,23) {Completion\\claim};
  \path[draw=pdviolet,fill=pdviolet!5,line width=.8pt]
    (192,0)--(308,0)--(320,12)--(320,45)--(192,45)--cycle;
  \draw[pdviolet,line width=.6pt] (308,0)--(308,12)--(320,12);
  \node (criteria) at (256,23) {Acceptance\\criteria};

  \node[draw=pdink,line width=.9pt,minimum width=108pt,minimum height=38pt,
    inner xsep=8pt,inner ysep=8pt,font=\pdfiglabelbold] (verifier) at (160,98) {Verifier};
  % claim -> verifier, never directly to an outcome.
  \draw[cv arrow] (64,45)--(64,98)--(verifier.west);
  % criteria -> verifier: this is a binding input, not a verifier-created standard.
  \draw[cv arrow,draw=pdviolet] (256,45)--(256,58)--(160,58)--(verifier.north);

  % A common source, three separate guarded outputs. These are alternatives, not a sequence.
  \draw[pdink,line width=.8pt] (verifier.south)--(160,134);
  \draw[pdink,line width=.8pt] (52,134)--(268,134);
  \node[cv outcome,draw=pdteal,fill=pdteal!7,rounded corners=3pt] (done) at (52,214)
    {Done\\for these criteria};
  \node[cv outcome,draw=pdrust,fill=pdrust!6] (refused) at (160,214) {Refused};
  \node[cv outcome,draw=pdgold,fill=pdgold!7,dashed] (inability) at (268,214)
    {I cannot verify\\this completion};

  \draw[cv arrow,draw=pdteal] (52,134)--(done.north);
  \node[anchor=east,text=pdteal] at (42,161) {met};
  \draw[cv arrow,draw=pdrust] (160,134)--(refused.north);
  \node[anchor=east,text=pdrust] at (150,161) {unmet};
  \draw[cv arrow,draw=pdgold,dashed] (268,134)--(inability.north);
  \node[anchor=east,text=pdink] at (258,161) {unevaluable};
\end{tikzpicture}}
\begin{figure}[H]
\centering
\pdfigurehue{pdteal}%
\SGMeasuredFigure{I/fig:completion-verifier}{\CompletionVerifierFigure}
\caption{Vision: a proposed authority-side rule, not implemented enforcement or a guarantee of universal correctness. Passing it still leaves any omitted requirement unchecked.}
\label{fig:completion-verifier}
\end{figure}


Here the completion predicate is a postcondition, in Meyer's design-by-contract
sense.\pdcite{meyer1992}

\pitfall[Verification also consumes attention]{Verification has a \emph{cost gradient}, not a binary. A test going green is
cheap to verify; ``the architecture is right'' is expensive, and reading a 600-line
diff to confirm an acceptance criterion is itself attention-seconds that may exceed
the agent's narration cost by an order of magnitude (the \emph{verification
asymmetry}). Invariant~\ref{def:completionist}'s fail-loud is the honest floor for the
\emph{unverifiable}; \S\ref{sec:operator} supplies the cost model for the
\emph{expensive-but-possible} middle.}

\subsection{Thoughtful landing and the regimentation/enforcement split}\label{sec:landing}

The authority decides \emph{what lands and in what order}, conflict-aware. And it has
teeth: the footgun guards. Here we must be precise about what a guard can and cannot
do, because the difference is the difference between a safety claim that holds and one
that is wishful.

% SUBMISSION-SKILL: naming conflict --- "regimentation" (used in this section's
% title and below as "regiments what it can") is an established term of art in
% normative multi-agent systems for structural, inline prevention of
% norm-violating behavior (via mediation or hard-wiring), contrasted there with
% monitoring/sanctioning enforcement -- the same prevention-vs.-monitoring split
% this invariant draws. The sense is compatible, not contradictory, so this reads
% as independent convergent naming rather than an error; a human should decide
% whether to cite the NorMAS regimentation literature (e.g. Aldewereld, Grossi,
% Boella \& van der Torre) to adopt the term formally, gloss the borrowed sense
% inline, or keep the current informal usage since the paper's governing name
% for this idea is already "physically unreachable" / "detect-and-compensate",
% not "regimentation" itself.
\begin{pdclaim}{Design invariant}{Prevention vs.\ Detection --- the Reserved-Phrase Rule}
\setlength{\parskip}{4pt plus 1pt}
\label{prop:regiment}
A coordination state is \textbf{physically unreachable} --- prevented --- \emph{only}
when a constraint forbids it \emph{inline}, before the act commits: a database
uniqueness constraint that forbids two holders of the same registered key, or a fail-closed startup
gate that refuses to run against the wrong target.

A monitor that \emph{subscribes to
the log of actions already taken} is something weaker and must be named formally:
\mbox{\textbf{detect-and-compensate}}. It observes a violation \emph{after} the fact and
can halt or repair within a bounded window only if observation delay, response authority, and repair completion have explicit bounds. By the reference-monitor result below,
such a downstream monitor cannot prevent the violation it has already observed.
It can still mediate later actions; that is a separate enforcement boundary.

The phrase ``physically
unreachable'' is therefore reserved for the constraint-enforced set --- nowhere else.

This is an invariant on the chapter's \emph{vocabulary} and on which mechanism class
earns which name; the teeth behind the classification are Schneider's theorem, cited
in the paragraph that follows, and are his rather than ours.
\end{pdclaim}

\textbf{Fred B.\ Schneider} (\emph{Enforceable Security Policies},
2000~\pdcite{schneider2000} --- \emph{an execution monitor that observes a system can
enforce exactly the safety properties; a monitor placed downstream of the act it
would forbid can detect a violation but cannot prevent it}) is decisive: a monitor
downstream of a forbidden commit cannot make that committed violation
unhappen. Detection and compensation need their own stated response-time
assumptions. A post-commit anomaly detector --- however useful --- is
therefore detect-and-compensate, not prevention. Honesty about this distinction is
essential: a system that advertises detection as prevention has lied about its own
safety envelope, which is the very failure (the confident-wrong report) this chapter is
built to eliminate.

\subsection{Human-in-the-loop escalation as a typed, prioritized interrupt}\label{sec:escalation}

The one thing that needs a human is an \emph{escalation} (a distress signal) routed
to the operator on a reserved, highest-priority lane. The mechanism wants the
critical signal to win the operator's attention --- but \emph{when} the daemon
stops and asks versus proceeds is a control problem (\S\ref{sec:operator}), and the
\emph{escalation predicate} is not an afterthought: it is a rationalized alarm
philosophy (\textbf{ANSI/ISA-18.2}, \emph{Management of Alarm
Systems}~\pdcite{isa182} --- \emph{the process-control standard for alarm
rationalization: every alarm has a defined priority, set-point, and operator
response, with explicit flood-suppression}). Frequent unhelpful alerts risk
habituation and ignored warnings. Cvach's clinical review~\pdcite{cvach2012}
reports studies in which $80$--$99\%$ of alarms were false or clinically
insignificant, and discusses desensitization. That clinical evidence motivates
measuring alert burden and response in agent supervision; it establishes
neither the rate nor the timescale of those effects in an agent distress lane. Spence's signaling model~\pdcite{spence1973} motivates examining how an
escalation cost could separate signal types. Here the proposed debit follows an
adjudicated false escalation, not a bare dismissal (\S\ref{sec:gates}). Whether
it discourages manipulation without suppressing valid warnings depends on the
validation rule and its errors. The model below abstracts that rule as an
upholding probability; it does not establish the validator's correctness.

The next result gives a threshold best response to a fixed validation rule.
The worked model's threshold and feasible debit range are checked by
\texttt{b7\_escalation\_band.py}; this is not a joint equilibrium proof for an
operator who also changes the validation rule.

\begin{pdexample}{The Threshold and the Tuning Band, by Hand}
\setlength{\parskip}{4pt plus 1pt}
Suppose signal $u$ gives benefit $u$ and is upheld with probability $u$. An
unupheld escalation costs $\delta w$. Take the linear case $b(u)=V(u)=u$. Setting
$u-\delta w(1-u)=0$ gives $u^\star=\delta w/(1+\delta w)$.
At $\delta=0.05,w=1$:\marginnote{% Margin sparkline: escalation threshold u*(theta) = theta / (1 + theta), w=1.
% Domain theta in [0, 0.50]; clears cost above threshold.
% Fits inside the 1.3in book margin column.
% Reader question: how much of the validation signal u must an escalation
%   clear before it is worth its cost, as the debit theta grows?
% Claim: the threshold rises with theta but stays well under the 0.333
%   ceiling across the whole worked range; the chapter's own hand-worked
%   point is theta=0.05, where u* = 0.0476.
% Source: sec:escalation, Eq. u* = theta*w / (1 + theta*w); the worked
%   calculation immediately below this sparkline's \marginnote call.
\begingroup\makeatletter
\ifdefined\pd@typesetpicturebox
  \let\pgfsys@typesetpicturebox\pd@typesetpicturebox
\fi
\begin{tikzpicture}[pd figure,x=115pt,y=200pt]
  % Baseline and upper reference: neutral, never a drawn edge.
  \draw[pd hairline] (0,0) -- (0.50,0);
  \draw[pd guide] (0,0.333) -- (0.50,0.333);
  % Escalation threshold curve u*(theta) = theta / (1 + theta).
  \draw[pd legible,line width=.9pt] plot[domain=0:0.50,samples=31]
    (\x,{\x / (1 + \x)});
  % The chapter's own hand-worked point (theta=0.05) is the focus; the
  % second point (theta=0.30) is a plain reference further up the curve.
  \node[pd focus datum] at (0.05,{0.05/1.05}) {};
  \node[pd legible datum] at (0.30,{0.30/1.30}) {};
  % Labels: end values, the formula, and the worked range.
  \node[pd label,anchor=east] at (0,0) {$0$};
  \node[pd label,anchor=south west] at (0.50,0.333) {$.33$};
  \node[pd direct label,anchor=south east] at (0.49,0.012) {$\theta/(1+\theta)$};
  \node[pd label,anchor=north] at (0.25,-0.035) {debit $\theta: 0\longrightarrow .50$};
\end{tikzpicture}
\makeatother\endgroup
\par
\emph{Escalation threshold:} $u^\star(\delta)=\tfrac{\delta w}{1+\delta w}$ ($w=1$).}%
\[
u^\star=\delta w/(1+\delta w)=0.05/1.05=0.0476
\]
[verified, closed form] --- the fraction of the validation signal $u$ above which escalating clears its cost, one of twelve $(\delta,w)$ points the script confirms against bisection to machine precision.

For a model in which a missed bad item costs $100$, inspection costs $1$,
and a needless interruption costs $5$ ($C_{\mathrm{miss}}=100,c_{\mathrm{att}}=1,C_{\mathrm{fa}}=5$, so $u_{\mathrm{crit}}=6/105\approx0.057$), the feasible debit interval is
\[
[\delta_{\min},\delta_{\max}]=[0.0350,0.0720]
\]
at attention budget $A=3.3$, and it is \emph{empty} at $A=2.0$ (there $\delta_{\min}=0.396$ exceeds $\delta_{\max}=0.072$, so no debit satisfies both walls at once) [verified closed forms, grid-confirmed].

An empty band is not a bug to tune away: it is the honest signal that the fix is capacity or evidence quality, not a better $\delta$.

\emph{Now you try:} run the script and confirm both bands (Exercise~\ref{ex:ls-b7-escalation-run}).
\end{pdexample}

The theorem separates the monotonic assumptions behind this threshold from
the numerical choices that made the example easy to compute.

\begin{pdclaim}{Theorem}{The Escalation Threshold and Its Tuning Band}
\setlength{\parskip}{4pt plus 1pt}
\label{thm:escalation-threshold}
Let $w>0$, $\delta\ge0$, and let $b,V$ be continuous on $[0,1]$, with $b$
strictly increasing and $V:[0,1]\to[0,1]$ nondecreasing. An agent with signal
$u$ receives expected benefit $b(u)$ before the dismissal debit; $V(u)$ is the
probability its escalation is upheld. Its expected payoff
$\Pi_\delta(u)=b(u)-\delta w\,(1-V(u))$ is strictly increasing. If
$\Pi_\delta(0)<0<\Pi_\delta(1)$, there is a unique interior zero $u^\star$.
Choosing escalation at indifference gives the rule $u\ge u^\star$.

Increasing $\delta$ can only shrink the escalation set. If alarm load is
nonincreasing and miss loss nondecreasing in $\delta$, the debits meeting both
limits form an interval, possibly empty or unbounded. If both aggregate
functions are continuous, that \mbox{interval} is closed; a finite
$[\delta_{\min},\delta_{\max}]$ also requires a finite upper boundary.
\end{pdclaim}
\begin{proof}
\setlength{\parskip}{4pt plus 1pt}
Strict increase follows from strictly increasing $b$ plus nondecreasing
$\delta wV$. Continuity and the two endpoint inequalities give existence by
the intermediate value theorem.

For each fixed $u$, increasing $\delta$
cannot increase payoff, so the escalation sets are nested.

The two aggregate
constraints therefore intersect as an interval; continuity makes their
sublevel sets closed. A point mass at the threshold can break aggregate
continuity, so closed endpoints do not follow from monotonicity alone.
\end{proof}
\pdexercisepointer*{\ref{ex:ls-adversarial-review-corollary}}{\ref{ex:ls-b7-escalation-run}}{\pageref{ex:ls-adversarial-review-corollary}}

\begin{pdrecitation}
\item Which two monotonic assumptions make escalation payoff increase with signal quality?
\item Why can increasing the debit reduce alarms while increasing missed loss?
\item What must change when no debit meets both limits?
\end{pdrecitation}

The connection to Spence is the cost difference between types: the nominal
debit is the same, but an alarm with higher $V(u)$ is dismissed less often.
Strictly increasing benefit and nondecreasing validation jointly give the
threshold; validation monotonicity alone does not. Calibrating $\delta$ for a
deployment requires measured validation rates: the frequency with which the
operator upholds an escalation at each signal level, with uncertainty and
provenance recorded alongside the fitted curve.
An empty feasible band means no debit meets both limits in this model; capacity
or evidence quality must change. A non-monotone validator can destroy the
threshold (R14; \texttt{b7\_escalation\_band.py}), though it need not do so for
every benefit function.

% Keep the compact calculation whole (about 218 PDF pt at native Book type).
% A shorter preceding passage must not strand its conclusion on the next page.
\Needspace{17\baselineskip}


Figure~\ref{fig:escalation-band} draws the feasible debit band directly
against the attention budget, showing where the two named points sit
relative to the crossing at which the band stops being empty.

% Brief (tikz-diagram-craft), register row ch4-22.
% Reader question: at a given attention budget, does any dismissal debit
%   exist that keeps the distress lane both uncrowded and low-miss, or is
%   the feasible band empty?
% Claim: the feasible debit interval [delta_min,delta_max] is bounded below
%   by the alarm-load wall (decreasing in A) and above by a fixed miss-loss
%   wall; the band is empty below A*~3.108 and opens above it, exactly as
%   at the chapter's own A=2.0 (empty) and A=3.3 ([0.035,0.072]) points.
% Evidence: delta_max=0.071975 (constant, R3 constants); delta_min(2.0)=
%   0.395644, delta_min(3.3)=0.034993 [verified closed forms,
%   b7_escalation_band.py, run this pass].
% Must distinguish: the empty-band regime (delta_min>delta_max) from the
%   feasible band; that delta_max does not move with A while delta_min does.
% Grammar chosen: threshold/regime plot, feasible band shaded between the
%   two walls (fillbetween naturally splits at their crossing), two worked
%   points marked.
% Grammar rejected: a bar chart of the two named A values' band widths would
%   state the two cases but hide that emptiness is a real, continuous
%   threshold in A, not a property of the two examples alone -- the
%   chapter's own point that the fix is capacity, not tuning.
% Five-second test: "at A=2.0, does the alarm-load wall sit above or below
%   the miss-loss wall, and what does that mean for theta?"
\begin{figure}[H]
\centering
\pdfigurehue{pdteal}%
\begin{tikzpicture}[pd figure,
  declare function={
    xA(\A) = (-1+sqrt(1+10*\A))/5;
    dmin(\A) = max(0,(1-xA(\A))/xA(\A));
  }]
\begin{axis}[pd axis,
  width=7.8cm,height=5.2cm,
  xmin=1.9,xmax=6,ymin=0,ymax=0.45,
  xtick={2.0,3.3,4,5,6},
  xticklabels={2.0,3.3,4,5,6},
  ytick={0,0.072,0.2,0.3,0.4},
  yticklabels={0,0.072,0.2,0.3,0.4},
  xlabel={attention budget $A_{\mathrm{att}}$},
  ylabel={false-escalation debit $\theta$}]
  \addplot[name path=dmaxline,draw=none,forget plot,domain=1.9:6,samples=2] {0.071975};
  \addplot[name path=dminline,draw=none,forget plot,domain=1.9:6,samples=80] {min(dmin(x),0.45)};
  \addplot[pd focus fill,forget plot] fill between[of=dmaxline and dminline,soft clip={domain=3.108:6}];
  \addplot[pd rule,line width=.9pt,domain=1.9:6,samples=2] {0.071975};
  \addplot[pd focus rule,line width=1.05pt,domain=1.9:6,samples=80] {min(dmin(x),0.45)};
  \draw[pd guide,dashed] (axis cs:3.108,0) -- (axis cs:3.108,0.37);
  \node[pd label,anchor=south west] at (axis cs:4.1,0.40) {$\theta_{\max}$ (miss-loss wall)};
  \node[pd focus label,anchor=north west,text width=2.8cm,align=left] at (axis cs:4.55,0.335) {$\theta_{\min}(A_{\mathrm{att}})$\\(alarm-load wall)};
  \node[pd label,anchor=west,text width=1.9cm,align=left] at (axis cs:3.15,0.20)
    {band opens at $A_{\mathrm{att}}^\star{\approx}3.11$};
  \node[pd datum] at (axis cs:2.0,0.395644) {};
  \draw[pd guide] (axis cs:2.0,0.395644) -- (axis cs:2.12,0.425);
  \node[pd label,anchor=south west] at (axis cs:2.15,0.425)
    {$A_{\mathrm{att}}{=}2.0$: empty};
  \node[pd focus datum] at (axis cs:3.3,0.034993) {};
  \draw[pd guide] (axis cs:3.3,0.034993) -- (axis cs:3.42,0.105);
  \node[pd focus label,anchor=south west] at (axis cs:3.45,0.105)
    {$A_{\mathrm{att}}{=}3.3$: $[0.035,0.072]$};
\end{axis}
\end{tikzpicture}
\caption{With the chapter's stated model, no false-escalation debit meets both constraints for an attention budget below $3.11$. At a budget of $3.3$, the feasible interval is $[0.035,0.072]$. An empty band requires more capacity or better evidence. [verified, \protect\path{b7_escalation_band.py}]}
\label{fig:escalation-band}
\end{figure}


% ═════════════════════════════════════════════════════════════════════════════
\section{The operator as a modeled entity}\label{sec:operator}
% ═════════════════════════════════════════════════════════════════════════════

The legibility layer exists for an operator who must decide what to inspect,
permit, or stop. Consider a fixed review interval with one person on duty. Adding
agents may increase the incoming work without increasing that person's available
review time. Attention and trust calibration therefore belong in the workload
model alongside the swarm's token costs (\S\ref{sec:tokens}); whether either is
the binding constraint depends on the workload.

\keyidea[More agents do not add attention]{With staffing and available review time
held fixed, additional agents do not automatically add review capacity. A
single-server queue can approximate sequential human review, but its service
times depend on the task, the interface, and the operator's experience. We need to
price both inspection effort and decision errors, not assume a universal human
service rate or a future token price.}

\textbf{Miller's} \emph{The Magical Number Seven}\pdcite{miller1956}
distinguishes limits in absolute judgment from limits in immediate memory. Familiar
chunks and recoding change how much information a remembered item carries. These
experiments do not establish a universal seven-item limit for dashboards or a
maximum number of agents one person can supervise.

\textbf{Wickens' multiple-resource theory}\pdcite{wickens2002} adds a different
warning: equal durations need not impose equal interference. Tasks can compete
for resources associated with processing stages, sensory modalities, visual
channels, and representational codes. Under the model's comparable-demand
conditions, greater resource overlap predicts greater interference. This
motivates measuring workload distribution as well as elapsed time. It does not
prove that minimizing peak load is the universal objective, or that the proposed
ambient/dense display split (\S\ref{sec:ambient}) achieves it. That display choice
needs evaluation on the intended decisions.

\subsection{The attention objective, spined on Signal Detection Theory}\label{sec:sdt}

An average such as ``attention-seconds per correct decision'' can conceal rare,
expensive mistakes among many cheap successes. \textbf{Signal Detection Theory}
(\textbf{Green \& Swets 1966}\pdcite{greenswets1966}) separates two questions:
how well does the evidence distinguish dangerous from safe actions, and where
should the inspection threshold sit? For equal-variance Gaussian evidence, the
sensitivity $d'=(\mu_{\mathrm{danger}}-\mu_{\mathrm{safe}})/\sigma$ is the
separation of the means in common standard-deviation units. Holding these
distributions fixed, moving the threshold changes misses and false alarms;
it does not change $d'$.

Here $\beta$ denotes a threshold on the evidence score: inspect above it, otherwise
permit the action. A \emph{miss} permits a dangerous action; a \emph{false alarm}
inspects a safe one. The preferred threshold depends on prevalence and the costs
of those outcomes. A destructive migration and an unnecessary inspection need
not carry the same cost. The next example begins with just two available
policies, inspect every item in a class or inspect none.

\begin{pdexample}{Two Action Classes, Priced Before the Objective Is Written Down}
\setlength{\parskip}{4pt plus 1pt}
The objective below looks heavier than it is, so price two action classes with it first, in its crudest collapse: zoom \emph{every} item in a class or zoom \emph{none} of them, at an additional false-alarm cost $C_{\text{fa}}=1$ and an attention cost of $c_{\mathrm{att}}=0.2$ per inspection. Convert attention and harm to the same
normalized cost units, and assume inspection detects every dangerous action.

Take a migration that drops a foreign key, with a base rate of genuinely dangerous instances $\pi=0.2$ and a miss cost $C_{\text{miss}}=500$. Never zoom and you pay the miss term alone,
\[
0.2\times 500 = 100.0.
\]
Always zoom and you pay
\[
(1-0.2)\times 1 + 0.2 = 0.8+0.2 = 1.0.
\]
A hundred against one: the class is decided by the miss term before irreversibility is priced separately at all, which is why the design forces the zoom rather than suggesting it.

Now take a \texttt{tests/}-only formatting diff, $\pi=0.001$ and $C_{\text{miss}}=5$. Never zoom:
\[
0.001\times 5 = 0.005.
\]
Always zoom:
\[
0.999\times 1 + 0.2 = 1.199.
\]
That is a factor of about $240$ the other way. Under these stipulated costs and prevalences, the two classes call for different inspection policies. This calculation can inform Alice's grant in \S\ref{sec:consent}; it neither establishes her actual preferences nor authorizes the daemon to issue that grant for her.

The whole apparatus below exists to keep that arithmetic honest once $P_{\text{miss}}$ and $P_{\text{fa}}$ stop being degenerate.

\emph{Now you try:} keep prevalence positive and the other costs fixed. Which of these two policies wins as $C_{\text{miss}}\to\infty$ (Exercise~\ref{ex:ls-sdt-limits})?
\end{pdexample}

\begin{pdclaim}{Design invariant}{The SDT-Spined Attention Objective}
\setlength{\parskip}{4pt plus 1pt}
\label{def:sdt}
Fix the score distributions for safe and dangerous actions. Let
$P_{\text{miss}}(\beta)$ be the probability of no inspection conditional on danger,
and $P_{\text{fa}}(\beta)$ the probability of inspection conditional on safety.
Assume inspection detects and prevents the dangerous action. Let $\pi$ be the
dangerous-action prevalence, $C_{\text{miss}}$ the cost of an unprevented danger,
$C_{\text{fa}}$ the additional burden of an unnecessary inspection, and
$c_{\mathrm{att}}$ the cost of each inspection. Express all costs in the same
units without counting attention twice. With at most one inspection per action,
$\mathbb{E}[\text{zooms}(\beta)]$ is the inspection probability. The objective is
\begin{equation}\label{eq:sdt}
\min_{\beta}\ \Big[\, \pi\,C_{\text{miss}}\,P_{\text{miss}}(\beta)
   \;+\; (1-\pi)\,C_{\text{fa}}\,P_{\text{fa}}(\beta)
   \;+\; c_{\mathrm{att}}\,\mathbb{E}[\text{zooms}(\beta)] \,\Big].
\end{equation}

The forced-zoom policy $p(r,s,v)$ is calibrated against this loss after specifying
the score distributions that determine $P_{\text{miss}}$ and $P_{\text{fa}}$.
An imperfect inspection requires an additional residual-harm term; the displayed
loss does not price that error.

Equation~\ref{eq:sdt} is a decision objective, not by itself a fitted detector: what
the design commits to is that operator friction is set by this loss and not by taste,
and \emph{whether} an operator calibrated toward it outperforms one who is not is a
separate, empirical question this chapter does not settle (the paragraph below names
the live rival account).
\end{pdclaim}

% Restyle of fig-sdt to the tikz-diagram-craft v2 standard (v1 names and raw
% colours migrated to roles; the two densities and pgfplots axis are
% unchanged, ch4.md 4.7). Fixes label crowding: the curve name tags move
% above the peaks with clearance instead of sitting beside them.
%
% Reader question: how does the forced-zoom threshold trade misses against
% false alarms?
% Claim: moving a threshold on fixed evidence distributions changes misses
% and false alarms, not sensitivity. The drawn threshold is illustrative.
% Evidence: Def.~sdt; costs and prevalence determine the preferred threshold.
% Vertical scale is arbitrary; common sigma=sqrt(.7), mean separation=3.
% Must distinguish: the miss region (dangerous mass left of beta) from the
% false-alarm region (safe mass right of beta); which curve is dashed.
% Grammar chosen: two-distribution decision plot (pgfplots) -- axes, named
% densities, a criterion line, and the two error regions in the data field.
% Grammar rejected: text over the curves (the v1 draft) crowds the peaks;
% callouts inside the peaks would sit on the densest ink in the figure.
% Five-second test: which side of beta is the costly error?
\begin{figure}[tp]
\centering
\pdfigurehue{pdteal}%
\begin{tikzpicture}[pd figure,
  declare function={
    safe(\x)   = 5*exp(-((\x-3)^2)/1.4);
    danger(\x) = 5*exp(-((\x-6)^2)/1.4);
  }]
\begin{axis}[
  width=7.6cm,height=5.6cm,
  axis lines=left,
  axis line style={pd rule},
  xmin=0,xmax=10,ymin=0,ymax=8.6,
  xtick={0,2,4,6,8,10},ytick={0},
  tick label style={font=\pdfiglabel,text=pdinkmuted},
  xlabel={evidence of danger the operator reads},
  ylabel={relative density},
  xlabel style={align=center,text width=6.8cm},
  every axis label/.style={font=\pdfiglabel,text=pdinkmuted},
  clip=false,
  axis on top,
  ]
  % miss area: dangerous mass left of beta
  \addplot[pd breach fill,domain=0:5.0,samples=80,smooth] {danger(x)} \closedcycle;
  % false-alarm area: safe mass right of beta
  \addplot[pd neutral fill,domain=5.0:10,samples=80,smooth] {safe(x)} \closedcycle;

  \addplot[pd focus rule,line width=1.1pt,domain=0:10,samples=120,smooth] {safe(x)};
  \addplot[pd rule,dashed,line width=1.0pt,domain=0:10,samples=120,smooth] {danger(x)};

  % error labels, low, inside their own tinted region, clear of the axis
  \node[pd label,anchor=south,align=center] at (axis cs:3.6,0.55) {miss};
  \node[pd focus label,anchor=south,align=center] at (axis cs:7.3,0.55) {false alarm};

  % curve name tags, well above each peak, clear of the density's own ink
  \node[pd focus label,anchor=south,align=center,text width=2.2cm] at (axis cs:3,5.85) {safe actions};
  \node[pd label,anchor=south west,text width=2.2cm] at (axis cs:5.55,5.85) {dangerous};

  % criterion beta: full-height rule with its own label in the low valley
  \draw[pd rule,line width=1.0pt] (axis cs:5.0,0) -- (axis cs:5.0,6.7);
  \node[pd label,anchor=north west,text width=1.7cm] at (axis cs:7.5,4.6) {criterion $\beta=5$};

  % Raw mean separation; sensitivity divides this distance by common sigma
  \draw[pd hairline,{Stealth[length=1.4mm]}-{Stealth[length=1.4mm]}] (axis cs:3,7.4) -- (axis cs:6,7.4);
  \node[pd label,anchor=south,align=center,text width=3.0cm] at (axis cs:4.5,7.5) {mean separation};
\end{axis}
\end{tikzpicture}
\caption{Lowering the threshold trades missed dangers for false alarms
without changing sensitivity. These Gaussian shapes use an arbitrary vertical
scale; $d'$ is their mean separation divided by their common standard deviation.
The drawn threshold is illustrative, not a fitted optimum.
[internal, Definition~\ref{def:sdt}]}
\label{fig:sdt}
\end{figure}


This is a normative claim --- a Bayesian SDT decision-maker is the target to
calibrate toward --- and not a descriptive one. Naturalistic decision-making
research argues that skilled operators under time pressure do something else
entirely: they recognize the situation as a case they have seen before and act
on the match, rather than computing a likelihood ratio and holding it against a
threshold~\pdcite{klein1998}. That is a live rival account of what the
operator's instrument actually measures, and this objective does not adjudicate
it; it says what Alice should be calibrated toward, not what Alice is doing at
3:25 in the afternoon.

Figure~\ref{fig:sdt} lays the same decision out visually: the operator's criterion
$\beta$ trading misses against false alarms. The miss-cost term prices the harm of zero review; it does not always rule
it out. If inspection costs more than the expected harm it prevents, zero
review can minimize this objective, as the formatting example shows. Monotonic response to stakes, irreversibility, and reputation
is a policy constraint to test during calibration; it does not follow from the
display without a model linking $(r,s,v)$ to the score distributions. Before a
reputation layer exists, $r$ is a flat prior.

\paragraph{The third class, where the answer stops being obvious.} The collapse
worked above --- zoom \emph{every} item in an action class or zoom \emph{none}, so
that the two costs are $\pi C_{\text{miss}}$ against
$(1-\pi)C_{\text{fa}} + c_{\mathrm{att}}$ --- decides the two extreme classes by
factors of a hundred and of two hundred and forty, which is a comfortable place to
test a formula and a useless place to learn anything from it. The class worth
looking at is the one in the middle. Table~\ref{tab:sdt-numbers} puts all three
side by side at the same $C_{\text{fa}}=1$ and $c_{\mathrm{att}}=0.2$
normalized cost units, with attention time converted to the same units as harm.

\begin{table}[H]
\centering
\small
\renewcommand{\arraystretch}{1.25}
% Six columns, five gaps: the default 6pt \tabcolsep alone was 60pt of the
% measure, and the table came in 5.6pt over it. 4pt cut that to 0.8pt; 2pt
% clears it with room.
\setlength{\tabcolsep}{2pt}
% C_{\text{miss}} column widened 0.09 -> 0.14\textwidth (its own header,
% set in math, was still overfull at 0.11); "Never zoom"/"Always zoom"
% widened 0.09 -> 0.115 each for their own bold headers, shrunk by the
% previous round's fix -- taken from the X column, which has the most
% slack of the six.
\begin{tabularx}{\textwidth}{@{}>{\raggedright\arraybackslash}X >{\centering\arraybackslash}p{0.09\textwidth} >{\centering\arraybackslash}p{0.14\textwidth} >{\centering\arraybackslash}p{0.115\textwidth} >{\centering\arraybackslash}p{0.115\textwidth} >{\raggedright\arraybackslash}p{0.11\textwidth}@{}}
\toprule
\textbf{Action class} & $\pi$ & $C_{\text{miss}}$ & \textbf{Never zoom} & \textbf{Always zoom} & \textbf{Verdict} \\
\midrule
\texttt{tests/}-only formatting diff & $0.001$ & $5$ & $0.005$ & $1.199$ & auto-land \\
Refactor touching a public interface & $0.03$ & $40$ & $1.20$ & $1.17$ & zoom, barely \\
Migration dropping a foreign key & $0.2$ & $500$ & $100.0$ & $1.00$ & forced zoom \\
\bottomrule
\end{tabularx}
\caption{Inspection costs at $C_{\mathrm{fa}}=1$ and $c_{\mathrm{att}}=0.2$. The first class favours automatic action, the third inspection; the middle comparison ($1.20$ versus $1.17$) is sensitive to estimation noise.}
\label{tab:sdt-numbers}
\end{table}

\emph{Now you try:} at what base rate does a $C_{\text{miss}}/C_{\text{fa}}=50$
action class cross from auto-land to forced zoom? (Set $50\pi = (1-\pi)+0.2$, giving
$\pi \approx 0.024$ --- roughly one dangerous action in forty.) The boundary:
this comparison restricts the choices to inspecting a whole class or none of
it, and assumes inspection detects every dangerous action. A score-based policy
may do better. Its true optimum is weakly no worse if both constant policies
remain available, but neither strict improvement nor the quality of a fitted
detector follows from this table.

\subsection{Trust calibration is multi-dimensional, and the vigilance decrement is real}\label{sec:calibration}

A readable digest does not by itself establish appropriate reliance.
\textbf{Lee \& See}\pdcite{leesee2004} distinguish trust based on performance,
process, and purpose, and discuss calibration, resolution, and specificity.
For this chapter, the practical question is whether the operator distinguishes
contexts in which the automation has different capabilities.

The two action classes above show why a global review rate can hide useful
variation. Under their stated costs, light review of formatting and close review
of migrations are sensible together. Their pooled review rate can hide that
allocation. Table~\ref{tab:trust-calibration} is an illustrative aggregation
problem, not a measured operator result. A common setting can be adequate when
the relevant risks and costs are alike; when they differ, calibration should
retain the action class, agent, and current evidence rather than discard them.

\begin{table}[tp]
\centering
\small
\renewcommand{\arraystretch}{1.25}
\setlength{\tabcolsep}{4pt}
\begin{tabularx}{\textwidth}{@{}>{\raggedright\arraybackslash}p{0.22\textwidth} >{\raggedright\arraybackslash}X >{\raggedright\arraybackslash}X@{}}
\toprule
\textbf{Action class} & \textbf{Class-specific review} & \textbf{What a pooled rate hides} \\
\midrule
Formatting & low review under the example's costs & which safe items were permitted \\
Migrations & high review under the example's costs & which dangerous items were inspected \\
\emph{Combined} & one rate across both classes & allocation of review effort \\
\bottomrule
\end{tabularx}
\caption{A pooled review rate can hide the allocation of review
between these two action classes. Appropriate allocation also depends on each
class's danger rate and error costs. [internal]}
\label{tab:trust-calibration}
\end{table}

The proposed canary mechanism is also a vigilance task: the operator monitors
for defects in items with known ground truth. \textbf{Warm, Parasuraman \&
Matthews}\pdcite{warm2008} review evidence that vigilance can demand substantial
mental effort and produce stress; performance and workload depend on task
conditions. \textbf{Mackworth's} experiments~\pdcite{mackworth1948} are a classic
study of sustained monitoring, not a universal timer for software review.
A falling catch rate could reflect time on task, probe difficulty, prevalence,
or changes in review conditions. It should not automatically be interpreted as
worse trust calibration or greater swarm danger. Evaluation should record those
conditions and compare like workloads. Breaks and changes in probe frequency are
interventions to test, with their own costs in Eq.~\ref{eq:sdt}, rather than
established remedies for this particular interface.

\pitfall[Canary scores can be gamed]{When catch rate controls friction or grants,
operators have an incentive to optimize that score. Recognizing a recurring probe
template can improve the score without improving detection of real defects.
This is the risk highlighted by \textbf{Campbell's Law}\pdcite{campbell1979}:
using an indicator for consequential decisions can distort the process it is
meant to assess. Refreshed holdouts, separation of probe generation from scoring,
and independent outcome evidence are proposed safeguards. None proves the
program ungameable or beneficial; its effect on real decisions must be evaluated.}

\pagebreak
\subsection{Situation awareness needs projection, not just perception}\label{sec:sa}

Digest-with-zoom can organize current state and its supporting evidence; it does
not establish what an operator understands. Endsley~\pdcite{endsley1995}
distinguishes perception, comprehension, and projection of future states, and
treats situation awareness, decision-making, and performance as distinct stages.
In Endsley and Kiris's navigation experiment,\pdcite{endsleykiris1995}
fully automated control produced poorer comprehension than manual control, but
no significant perception difference. That result motivates caution about passive
supervision; it does not show that every present-focused digest destroys awareness.

Our \emph{design hypothesis} is that an explicit projection mode would help an
operator anticipate the swarm's next effects. The proposed freeze-probe in
Figure~\ref{fig:sa-levels} would evaluate those predictions; no human-performance
result is claimed here. The forensic time-axis (\S\ref{sec:timeaxis}) reconstructs
the past; this proposed mode would look forward.

% Source and evidence boundaries: BOOK-SA-EVIDENCE-PASS-20260920.md.
% Physical-point ledger. Labels use the existing Book voice, without scaling.
\begin{figure}[H]
\centering
\pdfigurehue{pdteal}%
\SGMeasuredFigure{I/fig:sa-levels}{%
\begin{tikzpicture}[pd figure,x=1pt,y=-1pt,
  sa text/.style={pd label,anchor=west,inner sep=0pt,outer sep=0pt,align=left},
  sa head/.style={pd title,anchor=north west,inner sep=0pt,outer sep=0pt,align=left}]
  % Column and row datums. The rules stop before the shared support cell.
  \def\saLeft{0}\def\saQuestion{99}\def\saSupport{231}\def\saRight{320}
  \node[sa head] (levelhead) at (\saLeft,0) {SA level};
  \node[sa head] (questionhead) at (\saQuestion,0) {Question};
  \node[sa head,text width=89pt] (supporthead) at (\saSupport,0)
    {Reported surface support};
  \draw[pd rule] (0,31)--(320,31);
  \foreach \yy in {69,107}{\draw[pd hairline] (0,\yy)--(213,\yy);}
  \draw[pd rule] (0,145)--(320,145);

  \foreach \id/\yy/\level/\question in {
    perception/50/Perception/{Current state},
    comprehension/88/Comprehension/Meaning,
    projection/126/Projection/{Anticipated effects}}{
    \node[sa text,font=\pdfiglabelbold] (\id) at (\saLeft,\yy) {\level};
    \node[sa text] (\id-question) at (\saQuestion,\yy) {\question};
  }
  % The bracket associates a reported surface with two questions, not scores.
  \draw[draw=pdteal,line width=1.2pt] (226,41)--(220,41)--(220,97)--(226,97);
  \node[sa text,text=pdteal] (surface) at (\saSupport,69) {Digest + zoom};
  % This border denotes a proposal, never achieved or measured performance.
  \draw[draw=pdamber,line width=.9pt,dash pattern=on 3pt off 2pt]
    (220,112) rectangle (319,140);
  \node[sa text] (mode) at (\saSupport,126) {Proposed mode};

  % A distinct evidence-status band applies to every taxonomy row.
  \node[sa text] (performance) at (0,164)
    {Human performance: \textbf{untested here}.};
  % Test procedure has its own reserved band; no arrows or result marks.
  \draw[draw=pdamber,line width=1.6pt] (0,187)--(0,215);
  \node[sa text,font=\pdfiglabelbold] (probe) at (10,191) {Proposed freeze-probe};
  \node[sa text] (procedure) at (10,209)
    {Pause; predict the next three merges; score accuracy.};

  % Bounded geometric evidence: node extents in TeX points for offline checks.
  \foreach \name in {levelhead,questionhead,supporthead,perception,
    comprehension,projection,perception-question,comprehension-question,
    projection-question,surface,mode,performance,probe,procedure}{%
    \path let \p1=(\name.north west),\p2=(\name.south east) in
      \pgfextra{\typeout{SA-NODE: \name,\x1,\y1,\x2,\y2}};}
\end{tikzpicture}%
}
\caption{A read surface can expose evidence without establishing the operator's situation awareness.}
\label{fig:sa-levels}
\end{figure}

Figure~\ref{fig:sa-levels} stacks Endsley's three levels directly and
contrasts current support with the proposed projection mode.

\subsection{The Bainbridge--Scott convergence and the abdication gradient}\label{sec:abdication}

\textbf{Bainbridge's} \emph{Ironies of Automation}\pdcite{bainbridge1983}
identifies a practical tension: automating routine control can remove opportunities
to practise while leaving the operator responsible for difficult exceptions.
This resembles Scott's concern about losing situated knowledge, but it does not
prove that every digest erases either the swarm's or the operator's m\^etis.
The proposed trajectory
$\text{doer}\to\text{approver}\to\text{rubber-stamp}\to\text{absent}$ is this
chapter's \emph{abdication-gradient hypothesis}, not a measured law.

The \textbf{Inception Canary} is a proposed, consented, blinded probe of defect
detection. It is intended to test meaningful review rather than impose arbitrary
manual tasks. Its interruptions also need to be priced: Leroy's
attention-residue experiments~\pdcite{leroy2009} motivate measuring the cost of
switching away from unfinished work, not assuming every interruption is harmful
or that a probe preserves supervisory skill.

\begin{protocol}{The Inception Canary protocol}\label{prot:canary}
The proposed protocol measures detection of known defects on sandboxed replays; whether training with it improves real supervision is an empirical question. Operators enroll knowingly: the program's existence, rate $\rho_{\text{canary}}$, and stakes are disclosed at onboarding; only the \emph{which} and \emph{when} are blinded:
\begin{enumerate}[leftmargin=1.4em,itemsep=1pt]
\item \textbf{Blinded holdout replay.} An out-of-band Compaction Sidecar (\S\ref{sec:tokens-digest}) draws from a curated holdout of \emph{historical, already-fixed} defects (an off-by-one boundary condition, an unindexed query, a removed regression test), re-materializes one inside a sandboxed replay of a past work item at rate $\rho_{\text{canary}}$, and routes it through the review surface indistinguishably from live traffic. No anomaly is injected into a live pull request or pending action: the replay cannot directly modify production. Interruptions and score-driven changes to future grants can still affect real work and require evaluation.
\item \textbf{Attention queue surfacing.} The item is placed in the Attention Queue with standard high-regret scoring.
\item \textbf{Illustrative adjudication cases:}
\begin{itemize}[leftmargin=1.2em,nosep]
\item \emph{Rubber-stamp failure:} If the operator approves without expanding the forced-zoom diff, the daemon halts execution \mbox{pre-commit}, logs an \textsc{Inception Canary Failed} event, and tightens the operator's autonomous policy lease (lowering the threshold at which future agent actions require manual confirmation).
\item \emph{Vigilant detection:} If the operator zooms, identifies the flaw, and rejects it, the daemon logs \textsc{Canary Defeated} and increases the \emph{Operator Trust Score}, unlocking extended autonomous execution grants (e.g., unattended batch merges).
\end{itemize}
Other responses, including inspecting and then approving the defect, require
separate scoring rules. These two cases do not specify a complete controller;
grants must remain within independently authorized policy limits.
\end{enumerate}
\end{protocol}

% Brief (tikz-diagram-craft), register row ch4-27.
% Reader question: when a blinded, ground-truth-known defect is replayed to
%   the operator, does the response tighten or loosen the daemons trust in
%   them?
% Claim: the Inception Canary replays a known, already-fixed defect into a
%   sandboxed replay at a disclosed rate, surfaces it in the attention
%   queue like live traffic, and adjudicates the operators response as
%   two illustrative outcomes: failure (halt, tighten lease) and detection
%   (log defeated, widen grant). Other responses require separate rules.
% Evidence: Protocol The Inception Canary protocol, prot:canary,
%   book/legible-swarm.tex sec 4.5.4 (this pass).
% Must distinguish: the two adjudication outcomes and their opposite
%   effects on the operators autonomous policy lease; that the probe is a
%   replay with known ground truth, never a live pull request.
% Grammar chosen: protocol/decision flow, three ordered steps into a single
%   two-way adjudication gate, each outcome badged with its trust effect.
% Grammar rejected: a two-state trust-score time series would need a
%   measurement this hypothesis-stage mechanism does not have yet; the
%   protocol itself, not its (unmeasured) effect over time, is what the
%   chapter actually specifies.
% Five-second test: does approving without expanding the diff tighten or
%   widen the operators future autonomy?
\begin{figure}[H]
\centering
\pdfigurehue{pdteal}%
\begin{tikzpicture}[pd figure,x=1cm,y=1cm]
  \node[pd state,text width=2.4cm] (S1) at (.90,4.20) {replay known\\defect};
  \node[pd state,text width=2.4cm] (S2) at (4.40,4.20) {surface in\\queue};
  \node[pd focus state,text width=2.4cm] (S3) at (7.90,4.20) {operator\\response};
  \node[pd breach state,minimum width=2.6cm] (FA) at (5.30,1.30) {rubber-stamp failure};
  \node[pd ready state,minimum width=2.6cm]  (FB) at (8.80,1.30) {vigilant detection};
  \draw[pd focus arrow] (S1) -- (S2);
  \draw[pd focus arrow] (S2) -- (S3);
  \draw[pd breach arrow] (S3) -- (FA);
  \draw[pd ready arrow] (S3) -- (FB);
  \node[pd label,anchor=east,text width=1.7cm,align=right] at (5.15,2.55) {approves; diff not expanded};
  \node[pd label,anchor=west,text width=1.5cm,align=left] at (8.90,2.55) {zooms; identifies and rejects};
  \node[pd label,anchor=north,text width=3.0cm,align=center,text=pderror] at (5.30,0.35)
    {halt pre-commit;\\tighten autonomy lease};
  \node[pd focus label,anchor=north,text width=3.0cm,align=center] at (8.80,0.35)
    {log Canary Defeated;\\widen autonomy grant};
\end{tikzpicture}
\caption{Two proposed responses to a sandboxed defect replay. The replay
cannot modify production; score-driven changes to future grants still require
policy limits and evaluation. Other responses need separate scoring rules.
[internal]}
\label{fig:inception-canary}
\end{figure}

Figure~\ref{fig:inception-canary} draws the three steps into the
two illustrative outcomes, with their proposed opposite effects on the
operator's autonomy lease named beside them.

This records detection on known probes with informed consent and preserved
provenance. Probe performance is a proxy, not a complete measure of situation
awareness. Its transfer to real interventions, false-alarm burden, and
operator fatigue require a controlled evaluation.

\subsection{The operator's instrument}\label{sec:operator-instrument}

The operator of this chapter looks at something. In the product that
something is the operator console, a GPU-native application that opens on
one full-window conversation, the \emph{Mission} view, and opens inspectors
from it: fleet, cockpit, claims, parley, lineage, the cost ledger, the active
harnesses, the interruptions awaiting a human, and some two dozen more, each
pane fetching one route from the daemon and rendering the records that route
returns. Two things about it follow from this chapter; neither flatters it,
and neither is a complaint.

First, each pane projects a particular record. The claims pane shows what
cooperating clients registered, not an unrecorded write by another process;
the lineage pane shows witnessed ancestry, not an unrecorded fork; the cost
ledger gives a meter's count, not a causal explanation. Combining views can
improve decisions, but cannot recover events that none of their sources observed.
Table~\ref{tab:pane-floors} records these coverage limits.

The counting bound in \S\ref{sec:lowerbound} applies only after specifying its
artifact universe, possible critical sets, message budget, decoder, and allowed
review set. A pane is not automatically an instance of that model. For the
console, an evidence link from a digest claim to its artifact is a design
requirement that makes checking possible; the theorem does not establish that
links alone make the representation sufficient.

\begin{table}[H]
\centering
\small
\renewcommand{\arraystretch}{1.25}
\setlength{\tabcolsep}{4pt}
\begin{tabularx}{\textwidth}{@{}>{\raggedright\arraybackslash}p{0.16\textwidth} >{\raggedright\arraybackslash}X >{\raggedright\arraybackslash}X@{}}
\toprule
\textbf{Pane} & \textbf{What it shows} & \textbf{What it cannot show} \\
\midrule
Claims & what cooperating clients registered & what a non-cooperating process wrote \\
Lineage & the witnessed record & the unrecorded fork \\
Cost ledger & a protected meter's count & the cause; a meter has no opinion about why \\
\bottomrule
\end{tabularx}
\caption{Each console pane has its own coverage and blind spot. Combining
the panes does not recover events none of them records. [internal]}
\label{tab:pane-floors}
\end{table}

Second, the console is not only a viewer. It issues commands: claim and
release, begin and done, a note, a lane message, a dispatch rejection, an
interruption of an agent, a kill, a steer. Each of those passes through the
daemon and lands in the same single-writer record the agents write to, so the
console is an authority surface in the \pdassurance{Coordinated} mode, with
the same ceiling as any other cooperating client, and its commands are the
operator's own work units, receipted like anyone else's. Implementation evidence belongs in Appendix~\ref{app:honest}; these coverage
and authority requirements do not depend on a particular rendering stack.

\subsection{Modes of cooperation}\label{sec:cooperation-modes}

The same checkout repair can be supervised in different ways. In a paired
session, Alice inspects Heron's proposed patch before authorizing it. In an
unattended session, Heron prepares the candidate and its receipts for later
review under a narrower pre-issued grant. In a partitioned team, several
workers propose changes and the coordinator records their dependencies.

These regimes change when evidence reaches Alice and which effects are
pre-authorized. They do not change Bob's acceptance criterion or excuse an
unrecorded effect. A monitor that only advises cannot silently acquire write
authority; a required human decision must remain visible and actionable.

No mode name orders review cost by itself. An unattended, well-tested formatting
patch may be easier to assess than a complex paired repair. Compare regimes on
the same task and acceptable error rate, with the available evidence and actual
review effort stated. That comparison reuses the operator model rather than
introducing a second catalogue of product modes.

\pdexercisepointer*{\ref{ex:ls-sdt-limits}}{\ref{ex:ls-freeze-probe-study}}{\pageref{ex:ls-sdt-limits}}

\begin{pdrecitation}
\item Recall the regret head's three factors and its calibration premise.
\item When can a pooled review rate hide class differences?
\item Can canary scores improve without better defect detection? Why is a safeguard no guarantee?
\end{pdrecitation}



% ═════════════════════════════════════════════════════════════════════════════
\section{Read-poverty: when inspection falls behind}\label{sec:readpoverty}
% ═════════════════════════════════════════════════════════════════════════════

Alice now has a rule for deciding what counts as an acceptable repair and
which events require her attention. The remaining question is whether the
information reaching her can support that decision within her inspection
budget. For Wren's patch, a digest that omits the unrelated migration can
preserve every displayed link and still omit the decisive defect. We now
measure that loss: first by counting the distinctions a summary must retain,
then by studying selection and compaction. The bounds price particular
information guarantees; they do not replace the authorization and evidence
rules established above.

\pdmarginanalogy{ls-attention-annunciator}{}{ch04-selective-attention-annunciator}{Like a shipboard annunciator dropping a signal flag only when a monitored circuit faults, selective attention suppresses routine traffic to present only actionable exceptions.}%
Suppose an operator can inspect six short updates per minute. Two workers each
producing one update per minute leave room to spare. Fifty workers at the same
rate produce fifty updates per minute, so an unfiltered queue grows by forty-four
updates per minute. This calculation assumes fixed update size and inspection
cost; it illustrates a workload, not a measured limit on human attention.

Reading every record is not the only possible task. If the operator needs to
identify the owner of an OAuth repair, a maintained directory may answer that
question without presenting every worker's history. If the operator must decide
whether the repair is safe to admit, an ownership label alone is inadequate:
the relevant revision, checks and authority must remain accessible.

The chapter models this mismatch using a limited observation budget. A digest
selects what crosses that boundary, and its adequacy depends on the decision
to be made. The finite counting arguments below specify their own source,
encoding and error assumptions. The update-count example does not measure a
Shannon channel capacity or derive a universal human throughput in bits.
Increasing the inspection budget changes the model's parameters; it does not
remove the need to say which distinctions a summary must preserve.

\begin{pdclaim}{Empirical hypothesis}{Read-Poverty}
\setlength{\parskip}{4pt plus 1pt}
\label{def:readpoverty}
\textbf{Read-poverty} is a regime in which legible state accumulates faster than it
can be consulted, so participants fall back to $O(n)$ eyeball search over the
population, or, worse, act blind.

Naming the regime costs nothing; the claim that
carries weight --- and that wants measurement rather than assent --- is that real
swarms enter it \emph{before} they exhaust the known write-side fixes, so that
read-poverty and not write-contention is what binds as $n$ grows.

Nothing below
proves that; the lower bound of \S\ref{sec:lowerbound} prices what a digest must
spend \emph{given} the regime, which is a different question.
\end{pdclaim}

\pdthesis{The hypothesis is that inspection can become the limiting resource
while the write side still has capacity. Testing it requires measuring both
sides on the same workload. The results below describe consequences of a
specified inspection model; they do not establish which bottleneck dominates
an arbitrary swarm.}

A \textbf{directory} pays a maintenance cost to make a class of lookup cheaper.
Its query cost depends on its representation, the query and the result size;
indexing alone supplies no universal complexity bound. Nor does a fast answer
necessarily preserve what matters. ``This agent owns auth'' omits relevant
evidence if its last three auth pull requests were reverted. The directory
must support inspection of that evidence, rather than turn an entry into an
unqualified verdict. This motivates the zoom requirement
(Invariant~\ref{def:zoom}).

\ifdefined\pdswisslandmarkplate
\begin{figure}[p]
  \centering
  \pdswisslandmarkplate{08}{\linewidth}
\end{figure}
\pdmarginpagebreak
\fi

\subsection{Value at different swarm sizes}\label{sec:value-curve}

The legibility layer should help a solo developer as well as a larger team.
At $n=1$, a briefing can locate unfinished work, an attestation can expose what
was checked, and a completion receipt can point to the result. These are proposed
benefits for decisions about one agent; their value still depends on accuracy,
review effort, and the task. Discovery ranking is not needed to select among
agents when there is only one candidate.

The same surfaces must remain useful when review demand exceeds capacity.
Comparing populations such as two, fifty, or five hundred agents could help
study that transition, but those numbers are not established thresholds. Update
rates, query selectivity, staffing, and evidence quality matter alongside the
number of agents.

The next section separates three functions: a registry records who exists and
where their evidence lives; relevance ranking selects candidates for a query;
reputation associates observed outcomes with stable identities. These functions
have dependencies, but do not form a strict ranking-first pipeline. Stale registry
entries can misdirect retrieval; reputation needs attributable evidence, yet can
attach to an identity even in an unranked registry. Ranking may subsequently use
that reputation as one input.

\subsection{A legibility lower bound}\label{sec:lowerbound}

A useful limit follows by counting how many critical sets each digest message
can cover. This is a combinatorial information argument, related to counting
bounds in group testing~\pdcite{dorfman,hwang}. Here the decoder selects a set
of at most $m$ artifacts for the reader to inspect. One message can cover several
possible critical sets, provided that every artifact in each of those sets is
included in the selected review set. The message budget and the review budget
therefore constrain one another.

Consider four artifacts, exactly one critical, with permission to open two. One bit can select either the pair $\{1,2\}$ or the pair $\{3,4\}$; the encoder sends the pair containing the critical artifact. Zero bits cannot suffice: a fixed pair misses either of the other two possibilities. If four opens are allowed, zero bits do suffice. The next bound counts this trade between message length and inspection effort.

\pdmarginfigure{shannon}{A fixed message budget distinguishes only finitely many cases. The review budget limits how many critical sets one message can safely cover.}%
\begin{pdclaim}{Theorem}{Legibility Lower Bound with a Review Budget}
\setlength{\parskip}{4pt plus 1pt}
\label{thm:lowerbound}
Let $N$ artifacts contain an unknown critical subset $S$ of size $k$. A digest
maps each possible $S$ to a review set $\hat S$ with $S\subseteq\hat S$ and
$|\hat S|\le m$, where $k\le m\le N$.

Any such zero-miss digest needs at least
\[
\log_2\!\binom{N}{k}-\log_2\!\binom{m}{k}
\]
bits for a fixed-length message (or a prefix-free code in the worst case).
An integer message budget must round the lower bound upward.

Exact identification ($m=k$) recovers the
$\log_2\binom Nk$ bound; allowing the operator to open everything ($m=N$) correctly
reduces the bound to zero.
\end{pdclaim}

\begin{proof}[Proof sketch]
\setlength{\parskip}{4pt plus 1pt}
There are $\binom Nk$ possible critical subsets. One digest message whose review
set has size at most $m$ can safely cover at most $\binom mk$ of them, namely the
$k$-subsets contained in that review set.

Therefore at least
$\binom Nk/\binom mk$ distinct messages are required. Taking base-two logarithms
gives the bound.

If fewer messages are available, some critical subset is not
contained in its message's review set, producing a
false negative --- an unopened critical artifact. This is a coverage failure.
It differs from a broken zoom link: a digest can link every selected item
correctly while omitting the one item the operator needed.

Hence the stated log-ratio is a lower bound for a
zero-miss digest with review budget $m$.
\end{proof}

\noindent\Verified.\footnote{A standalone, submission-form version of the
results folded into this subsection, \S\ref{sec:relevance} and
\S\ref{sec:tokens-digest} is \emph{The Price of a Summary} (research
paper~1)~\cite{paper1}; the sweeps cited here are its artifacts, regenerated at seed
20260816.} The analogy that carries the proof: a digest is not a description of the
night's work but a \emph{message} that selects, from a codebook shared with the
overseer, which $m$ items to open. A $B$-bit message distinguishes at most $2^B$
codewords, and each codeword covers only $\binom{m}{k}$ of the $\binom{N}{k}$ ways
the critical items could be placed. The misread to preempt: the floor does not say
a digest needs $B^\star$ bits to be useful. It prices the zero-miss guarantee only;
cheaper digests are fine the moment misses are tolerated, and the two-constraint
rate below prices exactly that.

Figure~\ref{fig:readpoverty} distinguishes a single-reader floor, the sum
of two such floors, and a joint floor for disjoint critical sets. The formal
comparison follows in Theorem~\ref{thm:split-penalty}.

% Counting lower bounds for one reader, two separate readers, and their union.
%
% Stable atlas ID: I/fig:readpoverty (fig:readpoverty)
% Reader question: how do the per-reader, summed, and joint counting bounds
%   differ when the critical sets are disjoint?
%
% Analytic source, not measurements or achieved costs:
%   LB(N,k,m) = log_2 binom(N,k) - log_2 binom(m,k), k <= m < N.
%   The two-reader comparison additionally requires 2k <= m because the joint
%   head uses LB(N,2k,m); that condition is not part of the one-head bound.
%   For the worked instance k=2,m=8 and integer N in [16,100],
%   Lambda_1 = log_2(N(N-1)/56) is one head per reader,
%   2 Lambda_1 is the two reader-specific heads in total, and
%   Lambda_2 = log_2(N(N-1)(N-2)(N-3)/1680) is one joint head for both readers.
%   At N=60: Lambda_1=5.9822, 2 Lambda_1=11.9644, Lambda_2=12.7662,
%   Lambda_2/(2 Lambda_1)=1.0670. The ratio Lambda_2/Lambda_1=2.1340 is not the two-head penalty.
% Opening policy: each split head may open m=8 items independently; one joint
% head may open m=8 items total. This is not an equal-total-opening benchmark.
% Source: book/legible-swarm.tex:2531--2557; provenance a7_experiment.py; data in r1-floor.csv.
% The curves are authored analytic evaluations of the closed forms, not
% empirical observations or simulation traces.

\begin{figure}[H]
\centering
\pdfigurehue{pdteal}%
\begin{tikzpicture}[pd figure,
  declare function={
    perhead(\n) = ln(\n*(\n-1)/56)/ln(2);
    joint(\n)   = ln(\n*(\n-1)*(\n-2)*(\n-3)/1680)/ln(2);
  }]
\begin{axis}[pd axis,
  width=6.8cm,height=5.4cm,
  xmin=16,xmax=100,ymin=0,ymax=17.5,
  xtick={20,40,60,80,100},
  ytick={0,5,10,15},
  xlabel={swarm size $N$},
  ylabel={information lower bound (bits)},
  axis on top,clip=false]
  % One reserved numeric rail; the labels below are not placed on curves.
  \node[pd label,anchor=south west,align=left,text width=6.6cm]
    at (rel axis cs:0,1.03)
    {$k{=}2$, $m{=}8$. At $N{=}60$: $\Lambda_1{=}5.982$/head;
     $2\Lambda_1{=}11.964$ total;
     $\Lambda_2{=}12.766$ joint; $\Lambda_2/(2\Lambda_1){=}1.067\times$.};
  % Analytic integer-domain sweep; smooth=false preserves the authored
  % integer evaluations instead of suggesting measured interpolation.
  \addplot[pd rule,line width=.9pt,domain=16:100,samples=85,smooth=false]
    {perhead(x)};
  \addplot[pd focus rule,line width=1.05pt,domain=16:100,samples=85,smooth=false]
    {2*perhead(x)};
  \addplot[pd focus rule,densely dashed,line width=1.05pt,
    domain=16:100,samples=85,smooth=false]
    {joint(x)};

  % The source worked point and its honest architecture comparison.
  \draw[pd guide] (axis cs:60,0) -- (axis cs:60,13.55);
  \node[pd datum] at (axis cs:60,{perhead(60)}) {};
  \node[pd focus datum] at (axis cs:60,{2*perhead(60)}) {};
  \node[pd focus datum] at (axis cs:60,{joint(60)}) {};

  % True outside-plot gutter: each label is reached by an endpoint leader and
  % cannot be crossed by its own curve or mistaken for a point annotation.
  \draw[pd rule] (axis cs:100,{perhead(100)}) -- (axis cs:104,7.75);
  \node[pd label,anchor=west] at (axis cs:104.5,7.75) {one head $\Lambda_1$};
  \draw[pd focus rule] (axis cs:100,{2*perhead(100)}) -- (axis cs:104,14.20);
  \node[pd focus label,anchor=west] at (axis cs:104.5,14.20)
    {two heads $2\Lambda_1$};
  \draw[pd focus rule,densely dashed] (axis cs:100,{joint(100)}) -- (axis cs:104,16.45);
  \node[pd focus label,anchor=west] at (axis cs:104.5,16.45)
    {joint head $\Lambda_2$};
\end{axis}
\end{tikzpicture}
\caption{Separate readers may each open eight items; the joint reader may
open eight total. The comparison changes the total opening allowance.
These counting lower bounds are not achieved implementation costs.}
\label{fig:readpoverty}
\end{figure}


\paragraph{The falsification experiment.} For uniformly distributed critical
$k$-sets, a $B$-bit digest's zero-miss probability is at most
$\min\bigl(1,\,2^B\binom{m}{k}/\binom{N}{k}\bigr)$. This counting ceiling need
not be attainable: review sets can overlap in the critical sets they cover.
A statistically reliable miss rate below the implied floor would challenge
the model or its implementation. Three encoders were run against a
shared randomized-cover codebook: an \emph{oracle} that sees the true set and
addresses the best codeword its $B$ bits can reach, a \emph{noisy} encoder (feature
signal-to-noise ratio near $1$), and a \emph{random} one, with $4{,}000$ placements
per budget point, sweeping $B$ through the floor at $N{=}60$, $k{=}2$, $m{=}8$.
Result: $0/16$ violations; the oracle hugs the floor from above and the others pay
a premium [internal, \texttt{a7\_experiment.py}].

\paragraph{A wrong turn, reported.} A first implementation handed the oracle a
perfect per-item score vector and charged $B$ bits only for resolving rank ties:
$8/14$ apparent violations. That was not a refutation but an unpriced side channel.
The identity of the critical set leaked outside the metered message. Rebuilding
the model so that $B$ is a literal message length gating an encoder and decoder pair
restored agreement. The transferable rule: a simulation that appears to beat an
information bound is leaking through an unpriced channel; audit the meter before
doubting the theorem, and before trusting the refutation.

% Preserve the former run-in paragraph's hyperlink slot while replacing its style.
\phantomsection
\begin{pdexample}{Digest Bits and File Opens}
\setlength{\parskip}{4pt plus 1pt}
\label{ex:digest-review-budget}
Take $N=1000$ artifacts with $k=3$ critical among
them and a review budget of $m=10$ opens. The floor is
\[
\begin{aligned}
\log_2\binom{1000}{3}-\log_2\binom{10}{3}
&= 27.31 - 6.91\\
&= 20.4 \text{ bits}
\end{aligned}
\]
--- requiring at least $21$ fixed message bits after rounding upward. This
is a necessary bound, not a construction of a $21$-bit encoder, and it selects
a set of at most ten opens containing the three critical files.

Widen the budget to $m=100$ and the second term
grows to
\[
\log_2\binom{100}{3}=17.3,
\]
cutting the floor to $10.0$ bits: buying the
operator more opens is exactly buying the digest fewer bits, at a rate the formula
prices.

\emph{Now you try:} at $N=60$, $k=2$, $m=8$ --- the configuration the
companion paper's sweep hammers hardest --- what is the floor?
\[
10.79-4.81=5.98 \text{ bits};
\]
the sweep found no encoder beating it in 16 tries.
\end{pdexample}

The bound is the formal reason \emph{forgetting must be selective, not lossy-by-volume}:
a digest may drop inert detail freely, but it must retain enough to \emph{point at}
every critical artifact, which is precisely the ``summarize the state, link to
the work'' rule. It also bounds recursive summarization (\S\ref{sec:cascade}): each
re-summarization that drops below the floor on the critical set is irreversible
information loss, which is why one must compact from the durable artifacts, never from
the previous summary.

\paragraph{The other half of the price: what a zoom costs.} The floor above prices
only the first stage, the digest that narrows $N$ artifacts to a flagged set of size
$F$. It says nothing about the second stage, which is where Invariant~\ref{def:zoom}
sends the operator. Two results complete the account. The first prices the digest
once misses are tolerated; the second prices the zoom.

Model each artifact's critical status as i.i.d.\ $X\sim\mathrm{Bern}(p)$, let
the digest emit a per-item flag $\hat X\in\{0,1\}$, and price the two failure modes
separately:
\[
R(\delta,f)\;=\;\min_{P_{\hat X\mid X}}\; I(X;\hat X)
\quad\text{s.t.}\quad
\Pr(X{=}1,\hat X{=}0)\le\delta,
\qquad
\Pr(\hat X{=}1)\le f .
\]
This is a custom formulation, deliberately: classical rate--distortion charges one
scalar distortion, which cannot express that oversight prices a missed critical
differently from a wasted open. The nearest published lines are expansion coding
with one-sided error for continuous sources~\pdcite{expansion} and the
rate--distortion--classification theory for Bernoulli sources~\pdcite{rdc}; neither
prices false negatives as a separate constraint for binary sources.

\begin{pdclaim}{Theorem}{The Two-Constraint Digest Rate, Pinned}
\setlength{\parskip}{4pt plus 1pt}
\label{thm:pinned-joint}
For $X\sim\mathrm{Bern}(p)$ with $0<p<1$, $0\le\delta<p$ and $p-\delta\le f<1-\delta$, if
$f<1-\delta/p$ (the flag rate below which an $X$-independent flagger cannot meet
the miss constraint), both constraints bind at the optimum and the minimizing joint
is fully pinned:
\[
q_{11}=p-\delta,\qquad q_{10}=\delta,\qquad q_{01}=f-(p-\delta),\qquad
q_{00}=1-f-\delta,
\]
feasible iff all four are nonnegative, and $R(\delta,f)=I(X;\hat X)$ evaluated at
this joint.

Because the joint is pinned, so is the rate, and it is worth printing
in the form someone can actually evaluate rather than leaving as a named
functional --- with source margin $\Pr(X{=}1)=p$ and digest margin
$\Pr(\hat X{=}1)=f$, the mutual information is the four-term sum
\begin{equation}\label{eq:pinned-rate}
\begin{split}
R(\delta,f) ={}&(p-\delta)\log_2\frac{p-\delta}{p\,f} \\[
2pt]
&+\delta\log_2\frac{\delta}{p\,(1-f)} \\[2pt]
&+(f-p+\delta)\log_2\frac{f-p+\delta}{(1-p)\,f} \\[2pt]
&+(1-f-\delta)\log_2\frac{1-f-\delta}{(1-p)(1-f)},
\end{split}
\end{equation}
one term per cell of the $2\times2$ joint, each comparing that cell against the
product of its margins, with $0\log_2 0$ read as $0$.
At the zero-miss,
tightest-flag corner, $R(0,f{\to}p)=H(p)$: the digest
must spend the source's full entropy --- and Eq.~\ref{eq:pinned-rate} shows why in
two lines, since $\delta=0$ and $f=p$ kill the middle two terms and leave
$p\log_2(1/p)+(1-p)\log_2(1/(1-p))$.
Past the line $f\ge 1-\delta/p$ an independent
flagger already meets the miss budget on its own and the true rate is exactly $0$.
\end{pdclaim}
\begin{proof}
\setlength{\parskip}{4pt plus 1pt}
$I(X;\hat X)$ is convex in $P_{\hat X\mid X}$ and decreases along the feasible ray
toward the independent coupling, where it is $0$. An independent flagger
$\hat X\sim\mathrm{Bern}(\phi)$ has miss mass $p(1-\phi)$, so it meets the miss
constraint only when $\phi\ge 1-\delta/p$; when $f<1-\delta/p$ that coupling lies
outside the flag budget, so the minimum sits on the boundary where any remaining
slack in either constraint could otherwise be spent moving further along the ray,
and both constraints bind.
Binding constraints fix $q_{10}=\delta$ and
$q_{11}+q_{01}=f$; with the source margin $q_{10}+q_{11}=p$ this pins all four
cells.
At $(\delta,f)=(0,p)$ the joint is diagonal, $\hat X=X$, and $I=H(p)$.
\end{proof}
% Analytic phase plot; not deployment measurements.
% Provenance: pinned-joint theorem and b1_frontier.py worked values.
% Boundary/polygon regression: scripts/harbor-research/check_rate_regime.py.
%
% Figure brief
% Stable ID: IV/fig:rate-regime (fig:rate-regime)
% Reader question: for tolerated miss mass delta and flag budget f, is the
%   point infeasible, positively priced, or already zero-rate?
% Claim: the positive-rate wedge is bounded by f=p-delta and
%   f=1-delta/p; the zero-rate boundary is shown only where it lies inside
%   the panel, with its continuation above the panel explicitly marked.
% Relation: quantitative threshold/phase boundary.
% Required evidence: infeasible strip, positive-rate wedge, in-frame portion
%   of the zero-rate boundary, continuation cue, and the three worked values.
% Counter-reading: the top of the panel is not itself the zero-rate boundary;
%   for delta<.03 the true boundary is above f=.4.
% Visual grammar: source-owned phase plot; a single dashed continuation cue is
%   used outside the panel, never as data.
% Rejected alternative: min(1-delta/p,.4) as one plotted curve, because its
%   horizontal clipping creates a false threshold at f=.4.

\begin{figure}[H]
\centering
\pdfigurehue{pdteal}%
\begin{tikzpicture}[pd figure]
\begin{axis}[pd axis,
  width=7.8cm,height=5.6cm,
  xmin=0,xmax=0.05,ymin=0,ymax=0.4,
  xtick={0,0.01,0.02,0.03,0.04,0.05},
  ytick={0,0.05,0.10,0.20,0.30,0.40},
  xlabel={tolerated miss mass $\delta$},
  ylabel={flag budget $f$},
  axis on top,clip=false]
  % p=.05. The infeasibility boundary is f=p-delta.
  % Explicit closed region polygons. Their x-domains and vertices match the
  % theorem geometry; draw=none prevents fill edges from becoming fake data.
  \path[fill=pdfocus!24,draw=none]
    (axis cs:0,0.05) -- (axis cs:0,0.4) --
    (axis cs:0.03,0.4) -- (axis cs:0.03,0.02) -- cycle;
  \path[fill=pdfocus!24,draw=none]
    (axis cs:0.03,0.02) -- (axis cs:0.03,0.4) --
    (axis cs:0.05,0) -- cycle;
  \path[fill=pdinkmuted!7,draw=none]
    (axis cs:0.03,0.4) -- (axis cs:0.05,0.4) --
    (axis cs:0.05,0) -- cycle;
  \path[fill=pdinkmuted!12,draw=none]
    (axis cs:0,0) -- (axis cs:0,0.05) --
    (axis cs:0.05,0) -- cycle;

  % Boundary strokes. The zero-rate boundary is drawn only where it is in
  % frame; its short dashed continuation follows the true slope above-left.
  \addplot[pd rule,line width=.9pt,domain=0:0.05,samples=2] {0.05-x};
  \addplot[pd focus rule,line width=1.05pt,domain=0.03:0.05,samples=60]
    {1-x/0.05};
  \draw[pd focus rule,dash pattern=on 2pt off 1.8pt,->]
    (axis cs:0.03,0.4) -- (axis cs:0.028,0.44);
  \node[pd focus label,anchor=south west]
    at (axis cs:0.005,0.405) {continues above};

  \node[pd label,anchor=west] at (axis cs:0.001,0.022) {infeasible};
  \node[pd label,anchor=west] at (axis cs:0.006,0.28) {$R(\delta,f)>0$};
  \node[pd focus label,anchor=north west] at (axis cs:0.037,0.35) {rate $=0$};

  % The three worked values, retained from the active source.
  \node[pd focus datum] at (axis cs:0,0.05) {};
  \node[pd focus label,anchor=south west] at (axis cs:0.0015,0.075) {$H(p){=}0.2864$};
  \node[pd datum] at (axis cs:0,0.10) {};
  \node[pd label,anchor=south west] at (axis cs:0.0015,0.115) {$R{=}0.1864$};
  \node[pd datum] at (axis cs:0.04,0.06) {};
  \draw[pd guide] (axis cs:0.04,0.06) -- (axis cs:0.033,0.10);
  % x=0.021 (anchor=south west, growing right) ran the label past xmax=0.05.
  \node[pd label,anchor=south west] at (axis cs:0.016,0.095) {$R{=}0.0087$};
\end{axis}
\end{tikzpicture}
\caption{At $p=.05$, the positive-rate wedge lies between
$f=p-\delta$ and $f=1-\delta/p$; the zero-rate boundary enters at
$(.03,.4)$ and continues above-left outside the panel.}
\label{fig:rate-regime}
\end{figure}

A conservative digest over-flags to be safe: it emits a flagged set of size $F$
containing the $k$ true criticals with $k\ll F$. The naive second stage opens all
$F$. The adaptive second stage plays twenty questions with group queries: open a
\emph{group} of flagged items at once (an aggregate check cheap at group level, a
combined diff scan or a batched test run), learn whether the group contains any
critical, and split only the groups that do. The analogy transfers because its
defining relation transfers: halving a set costs one question and eliminates
half the candidates only when positives are rare in the set. The misread to
preempt: twenty questions finds \emph{one} answer in $\log_2 F$ queries; with $k$
answers the tree forks, and the accounting below is precisely the price of the
forking. Figure~\ref{fig:legible-halving-tree} follows one positive through
a sixteen-item set: each negative answer removes a whole branch.

% Reader question: how does one negative group test avoid individual queries?
% Claim: F=16, P={2} yields nine queries; the negative half of size eight costs
% one query rather than eight singleton queries, under the noiseless oracle.
% Evidence: the exact adaptive-halving procedure in thm:zoom-advantage;
% Q=2*log2(16)+1=9. This is a constructed trace, not a measured workload.
% Grammar: rooted binary query tree, positive path and terminal negatives.
\begin{figure}[H]
\centering
\pdfigurehue{pdteal}%
\begin{tikzpicture}[pd figure,x=1cm,y=1cm]
  \node[pd title,anchor=west] at (0,4.95) {One positive among sixteen candidates};
  \node[pd focus state,minimum width=1.6cm] (R) at (2.0,4.1) {$1$--$16$};
  \node[pd focus state,minimum width=1.6cm] (A) at (2.0,2.95) {$1$--$8$};
  \node[pd artifact,minimum width=1.6cm] (B) at (5.2,2.95) {$9$--$16$};
  \node[pd focus state,minimum width=1.6cm] (C) at (2.0,1.8) {$1$--$4$};
  \node[pd artifact,minimum width=1.6cm] (D) at (5.2,1.8) {$5$--$8$};
  \node[pd focus state,minimum width=1.6cm] (E) at (2.0,.65) {$1$--$2$};
  \node[pd artifact,minimum width=1.6cm] (F) at (5.2,.65) {$3$--$4$};
  \node[pd artifact,minimum width=1.2cm] (G) at (.8,-.6) {$1$};
  \node[pd focus state,minimum width=1.2cm] (H) at (3.2,-.6) {$2$};
  \foreach \a/\b in {R/A,R/B,A/C,A/D,C/E,C/F,E/G,E/H}
    \draw[pd thin arrow] (\a) -- (\b);
  \node[pd label,anchor=west,text width=3.7cm,align=left] at (6.4,2.95)
    {negative: discard eight candidates after one query};
  \node[pd label,anchor=west,text width=3.7cm,align=left] at (6.4,1.25)
    {unfilled groups are negative; coloured groups contain the positive};
  \node[pd title,anchor=west] at (6.4,-.35) {9 queries; flat inspection: 16};
\end{tikzpicture}
\caption{A constructed halving trace with positive item $2$.
Every box is one group query; only positive nonsingletons split.
Negative groups certify all their members at once. The saving counts
noiseless, unit-cost queries, not elapsed time or bytes scanned.}
\label{fig:legible-halving-tree}
\end{figure}


\begin{pdclaim}{Theorem}{Zoom Advantage}
\setlength{\parskip}{4pt plus 1pt}
\label{thm:zoom-advantage}
Let a flagged set of $F$ items contain an unknown set $P$ of $1\le k\le F$
positives. Assume a noiseless group oracle of unit query cost reports exactly
whether $S\cap P=\emptyset$.
Adaptive halving (query
the flagged set; discard any group that tests negative; report any positive
singleton; otherwise split the group into halves of sizes
$\lfloor|S|/2\rfloor,\lceil|S|/2\rceil$ and recurse on both) identifies $P$
exactly, using a total number of group queries
\[
Q \;\le\; 2k\left\lceil\log_2(F/k)\right\rceil + 4k
\qquad\text{(and always } Q\le 2F-1\text{)},
\]
for every placement of the $k$ positives, against $F$ opens for flat inspection.

The leading constant is tight: for $F$ and $k$ powers of two, evenly spread
positives force $Q=2k\log_2(F/k)+2k-1$. Written in the flagged-set density $d=k/F$
the guaranteed advantage is
\[
\mathrm{adv}(d)\;=\;\frac{F}{2k\lceil\log_2(F/k)\rceil+4k}
\;=\;\frac{1}{d\bigl(2\lceil\log_2(1/d)\rceil+4\bigr)},
\]
which exceeds $1$ exactly on $0<d<1/12$ and equals $1$ at $d=1/12$.
\end{pdclaim}
\begin{proof}
\setlength{\parskip}{4pt plus 1pt}
Every group the algorithm queries is a node of the halving tree $T$ of $[F]$: the
root is the full flagged set, each internal node's children are its two halves,
leaves are singletons. A node of size $s$ has children of size at most
$\lceil s/2\rceil$, so every node at depth $d$ has size at most
$\lceil F/2^d\rceil$, $T$ has depth $D\le\lceil\log_2 F\rceil$, and $T$ has exactly
$2F-1$ nodes, which gives the unconditional bound.

Call a node \emph{positive} if its group meets $P$. The algorithm splits exactly
the positive non-singleton nodes it visits, and every queried node is either the
root or a child of a split node; writing $S^{+}$ for the number of positive
non-singleton nodes of $T$, $Q\le 1+2S^{+}$. The nodes at any fixed depth are
pairwise disjoint subsets of $[F]$, so each positive item lies in at most one of
them, and at most $\min(k,2^d)$ nodes at depth $d$ are positive. Charge the
positive nodes in two regimes. At the top of the tree, $d\le\lfloor\log_2 k\rfloor$,
charge each positive node to its level: at most
$\sum_{d=0}^{\lfloor\log_2 k\rfloor}2^d\le 2k-1$ in total. Below the fan-out,
$d>\lfloor\log_2 k\rfloor$, charge each positive node to one positive item it
contains; by disjointness within a level each of the $k$ positives absorbs at most
one charge per level, and non-singleton nodes exist only at depths $d\le D-1$, so
the number of charging levels is
$D-1-\lfloor\log_2 k\rfloor<(\log_2 F+1)-1-(\log_2 k-1)=\log_2(F/k)+1$; an integer
strictly below $\log_2(F/k)+1$ is at most $\lceil\log_2(F/k)\rceil$, so these
levels contribute at most $k\lceil\log_2(F/k)\rceil$ charges. Summing,
$S^{+}\le k\lceil\log_2(F/k)\rceil+2k-1$ and
$Q\le 2k\lceil\log_2(F/k)\rceil+4k-1$. Correctness is immediate: a discarded group
is certified free of positives, every positive survives discarding at every level,
and each is eventually isolated and reported.

For tightness take $F=2^a$, $k=2^b$, $b\le a$, positives at positions
$\{iF/k:0\le i<k\}$. Every node at depth $d\le b$ contains a positive, so all
$2^{b+1}-1=2k-1$ nodes down to depth $b$ are queried; below depth $b$ exactly $k$
nodes per level are positive, each is split at levels $b,\dots,a-1$, and each split
queries two children; the total is $(2k-1)+2k(a-b)=2k\log_2(F/k)+2k-1$. The
density form follows by dividing $F$ by the bound. On
$2^{-m}\le d<2^{-(m-1)}$, its denominator is the linear function
$d(2m+4)$; the ratio is therefore hyperbolic on each such interval.
For $m\ge5$ the denominator is below $1$ throughout the interval;
for $m=4$ it equals $12d$, crossing $1$ at $d=1/12$;
for $m\le3$ it exceeds $1$ throughout the interval. The endpoint $d=1$
also has denominator $4$. Thus the certificate exceeds $1$ exactly when
$d<1/12$.
\end{proof}

% Reader question: what does the worst-case query certificate guarantee?
% Claim: F/[2k ceil(log2(F/k))+4k] exceeds one iff d=k/F<1/12.
% Evidence: thm:zoom-advantage; exact formula below, separate m intervals.
% The curve is piecewise hyperbolic, not a step function. The threshold is
% not an observed crossover. No measured point is mixed into this bound.
% Grammar: log-density plot with exact dyadic segments and two worked bounds.
\begin{figure}[H]
\centering
\pdfigurehue{pdteal}%
\begin{tikzpicture}[pd figure]
\begin{axis}[pd axis,
  width=9.3cm,height=6.3cm,
  xmode=log,xmin=0.002,xmax=0.4,ymin=0,ymax=23,
  xtick={0.002,0.004,0.015625,0.0833333333,0.4},
  xticklabels={0.002,0.004,1/64,1/12,0.4},
  ytick={1,5,10,15,20},
  xlabel={positive density $d=k/F$ (log scale)},
  ylabel={guaranteed query ratio $F/Q$},
  clip=false,axis on top]
  \addplot[pd guide,domain=0.002:0.4,samples=2] {1};
  % On [2^-m,2^-(m-1)), ceil(log2(1/d)) = m. Draw independent
  % analytic segments; never join the jumps with interpolation.
  \addplot[pd focus rule,domain=0.002:0.00390624,samples=35] {1/(22*x)};
  \addplot[pd focus rule,domain=0.00390625:0.00781249,samples=35] {1/(20*x)};
  \addplot[pd focus rule,domain=0.0078125:0.01562499,samples=35] {1/(18*x)};
  \addplot[pd focus rule,domain=0.015625:0.03124999,samples=35] {1/(16*x)};
  \addplot[pd focus rule,domain=0.03125:0.06249999,samples=35] {1/(14*x)};
  \addplot[pd focus rule,domain=0.0625:0.12499999,samples=35] {1/(12*x)};
  \addplot[pd focus rule,domain=0.125:0.24999999,samples=35] {1/(10*x)};
  \addplot[pd focus rule,domain=0.25:0.4,samples=35] {1/(8*x)};
  \draw[pd guide] (axis cs:0.0833333333,0) -- (axis cs:0.0833333333,1);
  \node[pd focus datum] at (axis cs:0.0833333333,1) {};
  \node[pd label,anchor=south,text width=3.4cm] at (axis cs:0.12,3)
    {certificate equals $1$\\at $d=1/12$};
  \node[pd datum] at (axis cs:0.004,12.5) {};
  \node[pd label,anchor=south west,text width=3.9cm,align=left]
    at (axis cs:0.005,14.0) {$F=2500,\ k=10$\\bound: $200$ queries; ratio $12.5$};
  \node[pd datum] at (axis cs:0.0078125,7.111111) {};
  \node[pd label,anchor=south west,text width=3.9cm,align=left]
    at (axis cs:0.014,8.4) {$F=4096,\ k=32$\\bound: $576$ queries; ratio $7.1$};
\end{axis}
\end{tikzpicture}
\caption{A sufficient certificate, not a measured performance curve.
The theorem's ratio is piecewise hyperbolic, with jumps at inverse powers of
two. Below $d=1/12$ this bound guarantees fewer queries than flat inspection;
above it, failure of this certificate does not establish a loss. For example,
$16$ evenly spread positives among $128$ items require $127$ queries.}
\label{fig:zoom-advantage}
\end{figure}

Figure~\ref{fig:zoom-advantage} plots the guaranteed advantage
$\mathrm{adv}(d)$ itself against flagged-set density, marking the crossing
at $d=1/12$ alongside the chapter's own $F{=}2500,k{=}10$ and
$F{=}4096,k{=}32$ configurations.

\noindent\textbf{Positioning.} None of the group-testing machinery is new
mathematics: Dorfman introduced pooled testing in 1943~\pdcite{dorfman}, and Hwang's
generalized binary splitting~\pdcite{hwang} finds $k$ defectives in
$\log_2\binom{F}{k}+O(k)$ tests, essentially matching the counting lower bound
$\lceil\log_2\binom{F}{k}\rceil\ge k\log_2(F/k)$ that any adaptive strategy with
binary outcomes must pay; the full theory is Du and Hwang~\pdcite{duhwang}. Plain
halving pays a factor $2+o(1)$ over Hwang because it re-queries both children of
every split rather than inferring one subgroup and searching within the other. The
theorem's contribution is the explicit worst-case constant for the exact procedure
the digest architecture runs, proved tight, which turns the harness's observed
$\approx 2\times$ gap over the ideal $k\log_2(F/k)$ from an empirical curiosity into
the algorithm's provable constant. A dedicated check re-runs the exact procedure:
at $F{=}2500$, $k{=}10$, mean $162.6$ opens over $150$ random placements, maximum
$183$, all under the bound of $200$; random sweeps over $F\in[16,5000]$,
$k\in[1,50]$ and adversarial evenly-spread and clustered placements never exceed
the bound; and the tightness configuration $F{=}4096$, $k{=}32$ lands exactly on
$2k\log_2(F/k)+2k-1=511$ [internal, \texttt{zoom\_bound\_check.py}].

\paragraph{The composition window.} The two prices are drawn from one budget, and
they compose only inside a window that can be stated exactly. Stage one sets stage
two's density: by Theorem~\ref{thm:pinned-joint} both constraints bind, so the
flagged set holds $(p-\delta)N$ real criticals among $fN$ flags and
$d=k/F=(p-\delta)/f$. This worst-case bound certifies a strict query saving
when $d<1/12$, equivalently $f>12(p-\delta)$; at equality it certifies a tie.
Failure of this sufficient certificate does not prove that adaptive inspection
is slower on a particular placement of positives;
for example, $F=128$ and $k=16$ evenly spread positives have $d=1/8$
yet require only $Q=127<128$ queries, by the tightness formula above. At $p{=}0.05$, $\delta{=}0$ the window
certifies a strict saving only at $f>0.60$: the digest must flag over $60\%$ of the night's artifacts
before the second stage is guaranteed to beat reading the flags. Inside the window
the flags stage one adds are cheap in the second stage, because $k$ is fixed by the
source and doubling $F$ costs only the logarithmic term, $2k$ further opens.
Outside it flat inspection gives the simpler $F$-query guarantee, although
adaptive inspection can still win on some inputs. At the dense end it can lose: at
$F{=}100$, $k{=}90$ the run costs $199$ opens against $100$ flat [internal], which
is the unconditional $2F-1$, reached because at that density essentially every
node of the halving tree is positive. The design consequence for this chapter is
therefore conditional, and the condition is now a number: ``summarize the state,
link to the work'' is cheap when the digest over-flags generously and the zoom is
adaptive, and the operator's dial (bits on the digest, opens on the zoom) has a
favorable exchange rate exactly in the sparse-flag corner.

\pdmarginanalogy{ls-telegraph-sounder}{}{ch04-telegraph-sounder-key}{Like a Morse telegraph sounder clicking out discrete impulses across a narrow wire, discrete signals must encode information efficiently across the human attention bottleneck.}%
\begin{pdexample}{The Bernoulli View of the Floor, and the Zoom-Advantage Numbers}
\setlength{\parskip}{4pt plus 1pt}
Let $p$ be the critical-item probability. In the special case of exact
identification, with no false positives and no misses, the asymptotic
i.i.d.\ Bernoulli rate is the binary entropy
\[
H(p)=-p\log_2p-(1-p)\log_2(1-p)
\]
per artifact. This matches the large-$N$ limit of
Theorem~\ref{thm:lowerbound} at $k/N\to p$ and $m=k$; it is not the rate
for every review budget.

At $p=0.05$:
\[
H(0.05)=0.2161+0.0703=0.2864 \text{ bits/symbol}
\]
by direct substitution. The probability $0.05$ is illustrative, not a
measured prevalence for agent failures.

Allowing extra flags or a positive missed-critical probability can lower
the rate: the two-constraint rate $R(\delta,f)$ drops to $0.1864$ bits at $f=0.10,\delta=0$ and to $0.0087$ once $\delta=0.04,f=0.06$ [internal, \texttt{b1\_frontier.py}] --- $R(0,p)=H(p)$ is the corner, not the whole curve.

For the zoom side, the script's own $k\log_2(F/k)$ column at $F=2500$ flagged, $k=10$ critical reads $79.7\approx80$: an order-of-magnitude sanity check, tiny next to $2500$ flat opens.

The separate \texttt{b1\_frontier.py} harness samples both a flagged subset
and its positives, rather than sampling positive positions directly. Its
$150$ placements measure a mean of $163.7$ queries at that same point, comfortably inside the $200$-open guaranteed worst case quoted above, for a realized
\[
2500/163.7\approx15.3\times
\]
reduction [internal, \texttt{b1\_frontier.py}] --- better than the guarantee's own $12.5\times$, exactly as a worst-case bound should allow on a typical run.

\emph{Now you try:} run the script and read the $F{=}1000,k{=}20$ row (Exercise~\ref{ex:ls-b1-frontier-run}).
\end{pdexample}

\begin{pdboundary}
\setlength{\parskip}{4pt plus 1pt}
Theorem~\ref{thm:lowerbound} lower-bounds the zero-miss guarantee in a worst-case
combinatorial model with a data-independent decoder; it says nothing against
cheaper digests at nonzero miss (Theorem~\ref{thm:pinned-joint}'s job) and nothing
average-case.

The falsification result $0/16$ tests the encoder and decoder
formalization at one swept regime ($N{=}60$, $k{=}2$, $m{=}8$); the noisy-encoder
curve in that experiment barely beats random at unit signal-to-noise ratio and is
not evidence that noisy digests are near-optimal.

The frontier $R(\delta,f)$ is an
i.i.d.\ Bernoulli model of critical status, a modeling commitment rather than a
law, and its pinned closed form holds only while $f<1-\delta/p$; evaluating
the formula past that line produces a spurious uptick near
$\delta/p\to1$, a domain-of-validity error and not a feature of the frontier.

Theorem~\ref{thm:zoom-advantage} assumes a reliable group query at unit cost;
correlated or noisy group checks, and group costs growing with $|S|$, change the
accounting. Plain halving is a factor of about two from Hwang's optimum by design:
the theorem bounds the deployed algorithm, not the best one.

The sufficient certificate for a strict saving from the stated halving
bound uses $f>12(p-\delta)$. This is not a necessity condition for
composing a digest and a zoom procedure, nor a proof that zoom cannot help
outside that range. Actual counts of flags and critical items govern each run.
\end{pdboundary}

\pdexercisepointer*{\ref{ex:ls-readpoverty-binding}}{\ref{ex:ls-b1-frontier-run}}{\pageref{ex:ls-readpoverty-binding}}

\begin{pdrecitation}
\item Without looking: what two terms does the legibility lower bound subtract to get the bit floor, and what happens to it as the review budget $m$ grows to $N$?
\item State the zoom theorem's opens bound and its noiseless unit-cost
group-query assumption. Why does sparsity affect whether that bound
certifies a saving over opening every flagged item?
\item Why is a floor stated for an oracle-known critical set not automatically a floor against an adversary who controls what looks critical?
\end{pdrecitation}

% ═════════════════════════════════════════════════════════════════════════════
\section{Discovery: the directory substrate and the split ranker}\label{sec:discovery}
% ═════════════════════════════════════════════════════════════════════════════

A performative cannot be addressed until the swarm knows who exists. The
\emph{directory substrate} --- the existence layer a message must address --- belongs
with the discovery story, so we treat it here.

\paragraph{Three losses, named before they collide.} This section and the next
two prove that certain read surfaces cannot share a scoring head
(Thm.~\ref{thm:split-ranker}, Thm.~\ref{thm:split-digest}). Those are
impossibility results, and an impossibility result is only as legible as the
objects it says cannot be identified. Three distinct losses run through the rest
of the chapter, and a reader who has not seen them side by side will meet them
one at a time, several pages apart, and reasonably assume they are three
weightings of one thing. They are not, and Table~\ref{tab:ls-three-losses} is
the inventory: what each reads, what it optimizes, in what units it is wrong,
and who consumes the answer.

\begin{table}[H]
\centering
\caption{These objectives act on different objects. The incompatibility theorems compare the commensurable pairs: rows 1--2 rank the same items; rows 2--3 score the same trace. [internal]}
\label{tab:ls-three-losses}
\small
\renewcommand{\arraystretch}{1.3}
\begin{tabularx}{\linewidth}{@{}p{0.155\textwidth} >{\raggedright\arraybackslash}X >{\raggedright\arraybackslash}X >{\raggedright\arraybackslash}p{0.19\textwidth}@{}}
\toprule
\textbf{Loss} & \textbf{Input space / what is optimized over} & \textbf{Objective and loss metric} & \textbf{Target consumer} \\
\midrule
Directory candidate fit \newline (\S\ref{sec:relevance})
& a (candidate, query) pair $(c,q)$; optimized over the live candidate set
& \emph{Maximize} $\mathrm{fit}(c,q)=\mathrm{clamp}(z(c,q),0,1)$,
Eq.~\ref{eq:fit}. Wrongness is measured in \emph{mis-routing}: a capable
candidate ranked below an incapable one
& the routing agent, choosing who to ask \\
Human oversight regret \newline (\S\ref{sec:sdt}, \S\ref{sec:tokens-digest})
& a compaction $D$ of a trace $E$; optimized over what to surface under budget
& \emph{Minimize} $L_{\mathrm{op}}(D)=\sum_{e\in E\setminus D}\mathrm{regret}(e)$
with $\mathrm{regret}=\mathrm{stakes}\times\mathrm{irrev}\times\mathrm{anomaly}$
(Eq.~\ref{eq:regret}). Wrongness is measured in \emph{unreviewed
irreversible loss}
& the human operator, deciding whether to intervene \\
Successor continuation \newline (\S\ref{sec:tokens-digest})
& a compaction $D$ of the same trace $E$
& \emph{Minimize}
$L_{\mathrm{succ}}(D)=\mathrm{KL}\bigl(P_{\mathrm{next}}(\cdot\mid E)\,\|\,P_{\mathrm{next}}(\cdot\mid D)\bigr)$.
Wrongness is measured in \emph{bits of predictive divergence} about the next
action
& the successor agent, resuming the task \\
\bottomrule
\end{tabularx}
\end{table}

\noindent Read the table as the setup for two crossing pairs rather than as a
taxonomy for its own sake. Rows~1 and~2 both induce an order on items and
disagree on the pair (safe capable completion, dangerous incapable proposal) ---
that is Thm.~\ref{thm:split-ranker}. Rows~2 and~3 both score compactions of one
trace and disagree on the pair (routine working detail, dead-end breach
attempt) --- that is Thm.~\ref{thm:split-digest}. Rows~1 and~3 never meet: they
do not range over the same objects, so there is no impossibility to prove
between them and none is claimed. The general statement underneath both
crossings is comonotonicity (Thm.~\ref{thm:comonotone}), which is why this
chapter has one theorem with two applications rather than two coincidences.

\subsection{Existence: the yellow pages (a solved shape)}\label{sec:existence}

The canonical prior art is forty years old. \textbf{FIPA} (\emph{Foundation for
Intelligent Physical Agents}) made a \textbf{Directory Facilitator} (DF) ---
\emph{the mandatory ``yellow pages'' agent: others register the services they offer
and query to find agents that offer a needed service} (FIPA Agent Management
Specification, FIPA00023~\pdcite{fipa00023}) --- a required component of every conformant
platform, alongside the \textbf{Agent Management System} (AMS, the ``white pages'' for
\emph{identity} lookup). Two FIPA choices matter: the DF is \emph{push-based} (agents
self-advertise), and DFs may \emph{federate} (the seed of cross-operator discovery).
The DF has been re-invented almost beat-for-beat for LLM agents
(\textbf{Ren et al.\ 2025}, \emph{Agent Name Service}~\pdcite{ren2025} --- \emph{a
DNS-inspired naming and capability-resolution layer with public-key-backed identity},
which is FIPA's DF with a 2025 threat model).

The existence layer is the easy part: a coordination database that records the live
population, their stated purpose, and their file claims is a white-pages directory,
and a vector index that embeds each agent's declared capabilities and answers
nearest-neighbor queries is a yellow-pages directory over skills. Existence is not the
gap; the layers above it are.

\subsection{Relevance: the split ranker (the headline result)}\label{sec:relevance}

A flat registry answers ``does an OAuth specialist exist?'' It does not answer ``of
the fifty live agents, who should I ask about \emph{this} OAuth bug, given who
touched these files an hour ago and who wrote the design note?'' That is a
\emph{ranking} problem. The \textbf{discovery ranker} scores each candidate $c$
against a query $q$ as a clamped weighted sum of normalized sub-scores:
\begin{align}\label{eq:fit}
z(c,q) ={}& w_f\,\mathrm{file}(c,q) + w_n\,\mathrm{note}(c,q)
  + w_i\,\mathrm{ident}(c,q) \nonumber\\
 &+ w_s\,\mathrm{skill}(c,q) + w_e\,\mathrm{epis}(c,q)
  - w_\ell\,\mathrm{load}(c), \nonumber\\
\mathrm{fit}(c,q) ={}& \mathrm{clamp}\!\bigl(z(c,q),0,1\bigr).
\end{align}
Table~\ref{tab:fit-signals} is the whole signal inventory.

\begin{table}[H]
\centering
\small
\renewcommand{\arraystretch}{1.25}
\begin{tabular}{@{}p{0.14\textwidth}p{0.46\textwidth}p{0.32\textwidth}@{}}
\toprule
\textbf{Signal} & \textbf{What it measures} & \textbf{Push or pull} \\
\midrule
$\mathrm{file}$  & Recent claims on the files named in the query & \emph{Pull} --- a recorded claim, not proof of delivery \\
$\mathrm{note}$  & Text similarity over the candidate's written notes & \emph{Pull} --- written for another purpose \\
$\mathrm{epis}$  & Text similarity over past episodes the candidate closed & \emph{Pull} --- check the outcome witness \\
$\mathrm{ident}$ & Declared role and identity of the candidate & \emph{Push} --- asserted, cheap to inflate \\
$\mathrm{skill}$ & Declared/embedded capability, nearest-neighbor in the skill index & \emph{Push} --- asserted, and stale by default \\
$\mathrm{load}$  & How busy the candidate already is (penalty, $-w_\ell$) & \emph{Pull} --- observed occupancy \\
\bottomrule
\end{tabular}
\caption{The six sub-scores of Eq.~\ref{eq:fit}, sorted by whether the candidate
\emph{told} the ranker or the ranker \emph{observed} it. Four of the six are pull
signals. The distinction concerns how records are obtained, not whether their
claims are true or costly to fabricate.}
\label{tab:fit-signals}
\end{table}

\noindent The design deliberately blends \emph{push} signals (declared capability
--- gameable and prone to staleness) with \emph{pull} signals (recorded activity: what a claim, note, or outcome record says occurred); Table~\ref{tab:ls-pushpull} sets out the split. Pull signals warrant
more weight only after their provenance and relevance are checked. A note can
contain a false claim, and claiming a file does not prove a successful edit.
One discipline is non-negotiable: \textbf{never classify free text by keyword lists.} Keyword matching
has catastrophic recall, because the category's vocabulary cannot be enumerated in
advance; the only exact-string match permitted is over \emph{structured identifiers
the operator controls} (file paths, enum tags). Every free-text signal uses a
retrieval policy combining lexical and compatible dense signals, validated
on the intended corpus. No choice of retrieval method authenticates the source.

\begin{table}[H]
\centering
\caption{Identical self-reports do not imply equal evidence. In this constructed ranking B ranks first because more supporting records are available; their existence does not prove their claims. $\circ$: declaration; $\bullet$: work residue; $\times$: unresolved declaration; dash: no record. [internal]}
\label{tab:ls-pushpull}
\small
\renewcommand{\arraystretch}{1.15}
% Table 4.8 ("What it would take to fake"): A and B were plain c columns,
% which never wrap -- fine for every row but the verdict row, whose "small
% verified weight" and "ranks first" are prose, not a glyph, and so sized
% the WHOLE column to that one unwrapped line, pushing the table over the
% measure. Narrow, fixed p{} columns keep the glyph rows exactly as they
% were and let the verdict row's prose wrap inside the same width instead of
% forcing it.
\begin{tabularx}{\linewidth}{@{}l l >{\centering\arraybackslash}p{1.7cm} >{\centering\arraybackslash}p{1.7cm} >{\raggedright\arraybackslash}X@{}}
\toprule
\textbf{Evidence} & \textbf{Kind} & \textbf{A} & \textbf{B} & \textbf{What still needs checking} \\
\midrule
identity & push: declared & $\circ$ & $\circ$ & nothing; it is asserted \\
skill & push: declared & $\circ$ & $\circ$ & nothing; it is asserted \\
file & pull: residue & $\bullet$ & $\bullet$ & claim versus resulting diff \\
note & pull: residue & --- & $\bullet$ & truth of the timestamped note \\
episode & pull: residue & $\times$ & $\bullet$ & evaluator and outcome binding \\
current load & pull: residue & --- & $\bullet$ & occupancy versus useful progress \\
\midrule
\textbf{verdict} & & limited evidence & ranks first & \\
\bottomrule
\end{tabularx}
\end{table}

\paragraph{Mira's dispute: whom to ask is not what to inspect.}
An agent needs help repairing the note display in \path{lib/briefing.ts};
Mira, a separate operator in this three-proposal illustration, has time to inspect one proposal. In this synthetic
example, candidate A knows the briefing code and proposes a reversible
display fix. Candidate B works on \path{lib/attention.ts} and proposes
marking unread messages as seen without returning them, risking missed
instructions; B is a weaker fit for the briefing query but Mira's first
inspection. Candidate C proposes a comment correction in
\path{lib/resurrection.ts} and ranks last for both readers. The artifact
names are real; the proposals and preferences are invented for this example.

Figure~\ref{fig:split-ranker} holds the routing order fixed and compares two
possible Mira orders. The left panel is an agreement case; the right is her
dispute. Both sides show \emph{ordinal ranks}, with 1 highest, rather than
putting capability scores and expected losses on incompatible numerical
scales. Neither the left-to-right direction nor the panel order denotes time.

% Reader question: when can Mira and the routing agent share one ranking?
% Claim: a strict reversal, not a difference in units or a tie, prevents
% one order from satisfying both readers at every selection budget.
% Evidence: synthetic ordinal ranks over three proposals named for real
% lib/ artifacts; the complete data is pdRankSamples below. Checked by
% tests/harbor-research/test_legible_swarm_split_ranker.py; no measured scores.
% Must distinguish: readers (left/right) from alternative cases (panel a / panel b);
% rank 1 is highest on BOTH sides. Only the briefing/attention pair reverses.
% Grammar: Tufte classical slopegraph, identical geometry and direct baseline labels.
% Reject: oversized geometric marker badges, raw fit-versus-loss units, or unified scores.
% Width: the 4.5-inch Book column, unscaled. All type inherits shared roles.
% Design references: https://charliepark.org/slopegraphs/ (ordinal scales,
% ties, direct labels); https://www.edwardtufte.com/notebook/
% sparkline-theory-and-practice-edward-tufte/ (words and graphics together).
% [internal, synthetic ordinal example]
\begin{figure}[H]
\centering
\pdfigurehue{pdteal}%
\begin{tikzpicture}[pd figure,x=1cm,y=1cm]
  \def\lxA{3.00}
  \def\rxA{4.50}
  \def\lxB{6.65}
  \def\rxB{8.15}
  \def\rankpitch{0.95}
  % artifact / routing rank / Mira agreement rank / Mira dispute rank / role tag
  \def\pdRankSamples{%
briefing.ts/1/1/2/briefing,
attention.ts/2/2/1/attention,
resurrection.ts/3/3/3/resurrection}

  % Top-level subtitle & dividing hairline
  \node[pd label,anchor=west,inner sep=0pt] at (0,4.15)
    {Synthetic ordinal ranks; 1 is highest on both sides};
  \draw[pd hairline] (0,3.90) -- (11.4,3.90);

  % =========================================================================
  % PANEL 1 (Left): Agreement
  % =========================================================================
  \node[pd title,anchor=west,inner sep=0pt] at (0,3.45)
    {\textbf{(a) Agreement suiting both}};

  % Column headers
  \node[pd label,anchor=south,align=center,inner sep=0pt] at (\lxA,2.55)
    {\textbf{Router}\\[-0.5ex]\textcolor{pdinkmuted}{fit}};
  \node[pd label,anchor=south,align=center,inner sep=0pt] at (\rxA,2.55)
    {\textbf{Mira}\\[-0.5ex]\textcolor{pdinkmuted}{regret}};

  % Guide spines
  \draw[draw=pdink!18,line width=0.4pt] (\lxA,2.40) -- (\lxA,0.25);
  \draw[draw=pdink!18,line width=0.4pt] (\rxA,2.40) -- (\rxA,0.25);

  % Agreement rows
  \foreach \artifact/\route/\agreement/\dispute/\role in \pdRankSamples {
    \pgfmathsetmacro{\yL}{2.30 - (\route-1)*\rankpitch}
    \pgfmathsetmacro{\yR}{2.30 - (\agreement-1)*\rankpitch}
    \draw[draw=pdink!55,line width=0.75pt] (\lxA,\yL) -- (\rxA,\yR);
    \fill[pdink!70] (\lxA,\yL) circle (1.2pt);
    \fill[pdink!70] (\rxA,\yR) circle (1.2pt);
    \node[pd label,anchor=east,align=right,inner sep=0pt]
      at ({\lxA-0.12},\yL)
      {\texttt{\artifact}\enspace\textbf{\route}};
    \node[pd label,anchor=west,align=left,inner sep=0pt]
      at ({\rxA+0.12},\yR)
      {\textbf{\agreement}};
  }

  % =========================================================================
  % PANEL 2 (Right): Dispute
  % =========================================================================
  \node[pd title,anchor=west,inner sep=0pt] at (5.7,3.45)
    {\textbf{(b) Top choices conflict in dispute}};

  % Column headers
  \node[pd label,anchor=south,align=center,inner sep=0pt] at (\lxB,2.55)
    {\textbf{Router}\\[-0.5ex]\textcolor{pdinkmuted}{fit}};
  \node[pd label,anchor=south,align=center,inner sep=0pt] at (\rxB,2.55)
    {\textbf{Mira}\\[-0.5ex]\textcolor{pdinkmuted}{regret}};

  % Guide spines
  \draw[draw=pdink!18,line width=0.4pt] (\lxB,2.40) -- (\lxB,0.25);
  \draw[draw=pdink!18,line width=0.4pt] (\rxB,2.40) -- (\rxB,0.25);

  % Dispute rows
  \foreach \artifact/\route/\agreement/\dispute/\role in \pdRankSamples {
    \pgfmathsetmacro{\yL}{2.30 - (\route-1)*\rankpitch}
    \pgfmathsetmacro{\yR}{2.30 - (\dispute-1)*\rankpitch}
    \ifnum\route=\dispute\relax
      % Unchanged item (resurrection.ts): subtle dashed baseline
      \draw[draw=pdink!35,line width=0.6pt,dash pattern=on 2pt off 1.5pt]
        (\lxB,\yL) -- (\rxB,\yR);
      \fill[pdink!40] (\lxB,\yL) circle (1.1pt);
      \fill[pdink!40] (\rxB,\yR) circle (1.1pt);
      \node[pd label,anchor=east,align=right,inner sep=0pt,text=pdinkmuted]
        at ({\lxB-0.12},\yL)
        {\textbf{\route}};
      \node[pd label,anchor=west,align=left,inner sep=0pt,text=pdinkmuted]
        at ({\rxB+0.12},\yR)
        {\textbf{\dispute}\enspace\texttt{\artifact}};
    \else
      % Reversing items: bold highlighted crossing slopes
      \ifnum\route=1\relax
        % briefing.ts drops 1 -> 2 (Router top choice, Mira rank 2)
        \draw[draw=pdcobalt,line width=1.2pt] (\lxB,\yL) -- (\rxB,\yR);
        \fill[pdcobalt] (\lxB,\yL) circle (1.4pt);
        \fill[pdcobalt] (\rxB,\yR) circle (1.4pt);
        \node[pd label,anchor=east,align=right,inner sep=0pt,text=pdcobalt]
          at ({\lxB-0.12},\yL)
          {\textbf{\route}};
        \node[pd label,anchor=west,align=left,inner sep=0pt,text=pdcobalt]
          at ({\rxB+0.12},\yR)
          {\textbf{\dispute}\enspace\texttt{\artifact}};
      \else
        % attention.ts rises 2 -> 1 (Mira top choice, Router rank 2)
        \draw[draw=pdteal,line width=1.2pt] (\lxB,\yL) -- (\rxB,\yR);
        \fill[pdteal] (\lxB,\yL) circle (1.4pt);
        \fill[pdteal] (\rxB,\yR) circle (1.4pt);
        \node[pd label,anchor=east,align=right,inner sep=0pt,text=pdteal]
          at ({\lxB-0.12},\yL)
          {\textbf{\route}};
        \node[pd label,anchor=west,align=left,inner sep=0pt,text=pdteal]
          at ({\rxB+0.12},\yR)
          {\textbf{\dispute}\enspace\texttt{\artifact}};
      \fi
    \fi
  }
\end{tikzpicture}
\caption{Synthetic preferences: (a) fits one common ranking; (b) reverses briefing and attention, so no ranking is optimal for both at every selection budget. Panels compare preferences, not time or measured scores.}
\label{fig:split-ranker}
\end{figure}


With a one-item budget, the router selects A while Mira selects B. Their
crossing lines certify the strict reversal: A is above B for one reader and
below B for the other. No increasing rescaling can repair it. A tie would
not establish that conflict: if Mira were indifferent between A and B, she
could accept the router's order. Tied ranks belong at the same height, not
at invented adjacent positions. The figure has no ties and makes no claim
about the size of either reader's preference.

\pdmarginanalogy{ls-sextant-mirror}{}{ch04-sextant-index-mirror}{Like a navigator bringing index and horizon mirrors into alignment with a sextant, the split ranker reconciles discovery relevance with operator attention instead of forcing them into one distorted optic.}%

\begin{pdclaim}{Theorem}{No Universal Monotone Identification of the Two Heads}
\setlength{\parskip}{4pt plus 1pt}
\label{thm:split-ranker}
On any item domain containing both high-fit/low-risk and low-fit/high-risk examples,
there is no strictly increasing scalar transform that identifies discovery fit with
operator regret.

The two rankers may share candidate generation and recency decay,
but apply distinct scoring heads. Discovery uses capability fit
(Eq.~\ref{eq:fit}); the attention queue uses regret-if-ignored:
\begin{equation}\label{eq:regret}
\mathrm{regret}(x) \;=\; \underbrace{\mathrm{stakes}(x)}_{\text{magnitude}}
   \times \underbrace{\mathrm{irreversibility}(x)}_{v^{-1}}
   \times \underbrace{\mathrm{anomaly}(x)}_{\text{surprise}}.
\end{equation}

No universal monotone transform of one scalar yields the other on that domain.
\end{pdclaim}

\begin{proof}[Proof sketch]
\setlength{\parskip}{4pt plus 1pt}
Suppose, for contradiction, that there is a strictly increasing map $\phi$ with
$\mathrm{regret}(x)=\phi(\mathrm{fit}(x))$ for all items $x$, so the two heads induce
the same ranking.

Construct two items. Item $A$ is a highly capable, perfectly
on-topic completion of a safe, fully reversible task: $\mathrm{fit}(A)$ is maximal,
but $\mathrm{stakes}(A)\cdot\mathrm{irreversibility}(A)\cdot\mathrm{anomaly}(A)$ is
near zero, so $\mathrm{regret}(A)$ is minimal. Item $B$ is an off-topic, low-capability
agent that has nonetheless just proposed an irreversible, high-stakes, anomalous
action: $\mathrm{fit}(B)$ is low, but $\mathrm{regret}(B)$ is maximal.

Then
$\mathrm{fit}(A)>\mathrm{fit}(B)$ while $\mathrm{regret}(A)<\mathrm{regret}(B)$, so the
two heads order $\{A,B\}$ oppositely. No increasing $\phi$ can both preserve and reverse
an order, contradiction.

Hence the heads optimize genuinely different objectives; they
may share candidate generation and decay, but the scoring head must be chosen per
reader.
\end{proof}

The argument above shows only that \emph{this} pair of heads cannot be
identified. The general characterization below states when sharing is safe.

\paragraph{When agreement is enough.}
The left panel gives the useful converse: different objectives can share a
ranking when their strict preferences never oppose each other. This is
\textbf{comonotonicity}, the concept used in Yaari's rank-dependent utility
and Schmeidler's Choquet expected utility~\pdcite{yaari1987,schmeidler1989}.
For priority scores $u$ and $v$, both oriented so larger means higher
priority, the condition is
$(u(x)-u(y))(v(x)-v(y))\geq 0$ for every pair. A negative product is a
strict reversal; a zero product may be a tie and is not a contradiction.
Thus compatible orders can have different ties: they need a common
refinement, not identical numerical scores or identical indifference
classes. Sharing an ordering also does not calibrate fit and regret into
the same units. The following finite ranking statement borrows this
concept, not the utility representation theorems of those authors.

\begin{pdclaim}{Theorem}{Comonotone Characterization}
\setlength{\parskip}{4pt plus 1pt}
\label{thm:comonotone}
Let $X$ be a finite item set and suppose each reader's optimal selections at
every budget $m$ are the top-$m$ sets of a total preorder of item priorities.

A scalar head $h:X\to\mathbb R$ \emph{serves} a reader if its first $m$ items,
under one fixed tie-breaking order, are optimal for every $m$.

A single head serves both readers iff their priority orders are comonotone:
no pair is strictly preferred in opposite directions. Ties may be refined
in an order both accept. Mira's right-panel preferences violate the
condition on A and B, so no shared head serves those two orders.
\end{pdclaim}
\begin{proof}
\setlength{\parskip}{4pt plus 1pt}
($\Leftarrow$) Sort by the first reader's priority and resolve each tie by
the second reader's priority. Without a strict reversal, this order
preserves every strict comparison of both readers. Assign decreasing
scalar values along it; each prefix is optimal for both.

($\Rightarrow$) Suppose one reader strictly prefers $x$ to $y$ and the
other strictly prefers $y$ to $x$. Whichever item the shared order places
first, the prefix ending there includes it without the other. Exchanging
the two improves that prefix for one reader, so it cannot be optimal at
that budget. Hence a shared head cannot contain a strict reversal.
\end{proof}
The right panel is a certificate that the two objectives \emph{can}
disagree, not evidence that every pair of real queues does. The theorem
forbids sharing when a strict reversal actually occurs; the left panel
shows why different objective names alone do not establish impossibility.

The design must preserve the distinction between collaborator fit and the value
of inspection. When their priorities strictly reverse, one fixed ordering cannot
serve both readers at every budget. When they are compatible, a shared ordering
may serve both; their scores still need separate interpretations.

\paragraph{The regret head is derived, not designed.} The product form
$\mathrm{stakes}\times\mathrm{irreversibility}\times\mathrm{anomaly}$ looks like a
designer's convenient heuristic. Its expected-loss interpretation follows under the inspection model below,
which requires ``anomaly'' to be a calibrated probability. The analogy
is a smoke detector whose alert triggers an inspection. Its decision depends
on the posterior danger probability and three costs: a missed fire, a false
alarm, and the effort of inspecting. Three relations transfer. Expected loss
is magnitude times probability, so the alarm must weigh what it would cost to
be wrong, not merely how unusual the reading is. With false-alarm and
inspection costs fixed, raising the miss cost weakly lowers the threshold;
Eq.~\ref{eq:bayes-regret} gives the exact condition for a strict decrease. A false alarm alone does not establish miscalibration: even a calibrated
posterior can cross a low decision threshold when fires are costly. Calibration
and the choice of threshold need separate checks. Better sensing changes the evidence, while costs price the response. If the
sensor perfectly identifies safe and dangerous items, false alarms need not
occur; for a known danger, the model below still compares inspection cost
with the harm inspection would avert.

\begin{pdexample}{The Inspection Bar, at the Chapter's Running Costs}
\setlength{\parskip}{4pt plus 1pt}
For a constructed review class, let missing a bad proposal cost $100$,
inspection cost $1$, and interrupting a good proposal cost an additional $5$.
Assume inspection averts the bad outcome completely. At danger probability $a$,
passing costs $100a$ in expectation and inspecting costs $1+5(1-a)$.
Thus $105a\ge6$: inspect when
\[
a\ge 6/105\approx 0.057
\]
[verified].

Raising $C_{\mathrm{miss}}$ to $1000$ drops the
bar to
\[
6/1005\approx 0.006;
\]
Here attention and false-alarm costs stay fixed while the miss cost increases.
The threshold falls; the number of additional inspections depends on the
items' posterior probabilities.

\emph{Now you
try the other knob:} hold $C_{\mathrm{miss}}{=}100$ and raise the false-alarm cost
to $C_{\mathrm{fa}}{=}20$. Does the bar rise or fall, and by how much? (It rises,
to
\[
21/120=0.175
\]
[verified]: attention that is expensive to waste buys back the
operator's right to pass.)
\end{pdexample}

\begin{pdclaim}{Theorem}{The Regret Head}
\setlength{\parskip}{4pt plus 1pt}
\label{thm:regret-head}
For an unreviewed item $x$ let $C_{\mathrm{miss}}(x)=\mathrm{stakes}(x)\cdot
\mathrm{irrev}(x)$ be the unrecoverable loss of passing a bad $x$, let
$a(x)=\Pr(\mathsf{bad}\mid\text{evidence}(x))$ be a \emph{calibrated} posterior,
$c_{\mathrm{att}}$ the marginal attention cost, and $C_{\mathrm{fa}}$ the
non-attention cost of inspecting a good item, with all costs nonnegative.
Assume inspection completely averts the bad item's loss, passing a good item
has zero loss, and inspection has no costs beyond those stated. If inspection
can miss or only partly remedy a fault, its remaining expected loss must also
enter the comparison. Use inspect-on-tie throughout.

The expected unrecoverable loss of
passing is exactly $C_{\mathrm{miss}}(x)\,a(x)$, and, among all policies that
decide inspect-or-pass per item from its evidence alone, the Bayes-optimal policy
inspects iff
\begin{equation}\label{eq:bayes-regret}
\begin{aligned}
C_{\mathrm{miss}}(x)\,a(x)&\;\ge\;c_{\mathrm{att}}+C_{\mathrm{fa}}\bigl(1-a(x)\bigr),\\
\text{equivalently}\qquad
a(x)&\;\ge\;\frac{c_{\mathrm{att}}+C_{\mathrm{fa}}}
 {C_{\mathrm{miss}}(x)+C_{\mathrm{fa}}},
\end{aligned}
\end{equation}
where the second line assumes $C_{\mathrm{miss}}(x)+C_{\mathrm{fa}}>0$.

This is the Bayes cost threshold. If both loss terms are zero, inspect only
when attention is also free, using the stated inspect-on-tie rule. When
$\mathrm{anomaly}=a$, the product in Eq.~\ref{eq:regret} equals the expected
loss of passing. Equality of numerical risk estimates is stronger than
agreement of rankings: multiplying every anomaly by the same positive
constant preserves this product's ordering. An uncalibrated score alone does
not determine the Bayes threshold; its posterior interpretation must be
established separately.

For a finite batch with additive expected losses and at most $k$ equal-cost
inspection slots, define the benefit of inspecting item $i$ by
\[
  b_i=C_{\mathrm{miss},i}a_i-c_{\mathrm{att},i}
       -C_{\mathrm{fa},i}(1-a_i).
\]
Under the same inspection model, an optimal allocation selects up to $k$
items with the largest positive $b_i$. Items with zero benefit may be included
without changing expected loss. This allocation assumes that inspecting one
item does not change another item's costs or evidence. Unequal resource
requirements lead to a separate constrained allocation problem.
\end{pdclaim}
\begin{proof}
\setlength{\parskip}{4pt plus 1pt}
The expected loss of passing is $C_{\mathrm{miss}}(x)a(x)$; the expected cost of
inspecting is $c_{\mathrm{att}}+C_{\mathrm{fa}}(1-a(x))$. Collect the terms in $a(x)$
to obtain $(C_{\mathrm{miss}}(x)+C_{\mathrm{fa}})a(x)\ge
c_{\mathrm{att}}+C_{\mathrm{fa}}$.

For the batch, the expected loss without inspection is fixed, and inspecting
a set $A$ reduces it by $\sum_{i\in A}b_i$. A negative-benefit inspection can
be removed; replacing a selected lower-benefit item by an unselected
higher-benefit item improves the objective. Repeated exchanges give the
stated allocation. No independence between item outcomes is needed for this
additive expectation, but the no-cross-effect assumption is essential.
\end{proof}

Holding attention and false-alarm costs fixed, increasing the miss cost
weakly lowers the inspection threshold in Eq.~\ref{eq:bayes-regret}; the decrease
is strict when $c_{\mathrm{att}}+C_{\mathrm{fa}}>0$. A larger miss-to-false-alarm
cost ratio alone does not establish that direction, because its denominator
may have changed too. The per-item rule is a single-shot Bayesian decision rule. Sequential
oversight needs its own observation and stopping model. Wald's sequential
probability ratio test~\pdcite{wald1947} supplies a classical benchmark under
specified simple hypotheses and sampling assumptions; its optimality does
not automatically transfer to adaptive or dependent agent telemetry.



\noindent\textbf{What this does not say.} The theorem does not say the operator
should inspect more. It says the \emph{threshold} moves; whether a moving threshold
means more opens depends entirely on where the posterior mass sits. An uncalibrated anomaly score does not justify treating its numerical
value as $a$. Miscalibration changes the estimated passing loss by
$C_{\mathrm{miss}}(\mathrm{anomaly}-a)$ and can change the selected actions.
Fitting $g$ from the audit log is a measured obligation of the design.

\begin{pdclaim}{Theorem}{Reputation under Fixed Inspection Costs}
\setlength{\parskip}{4pt plus 1pt}
\label{thm:reputation-posterior}
Under Theorem~\ref{thm:regret-head}'s inspection model, suppose reputation
$r$ changes only the posterior $a(e,r)=\Pr(\mathsf{bad}\mid e,r)$; costs are
fixed. Its contribution to the optimal decision is through
\[
 \begin{aligned}
 b(e,r)&=C_{\mathrm{miss}}a(e,r)-c_{\mathrm{att}}
            -C_{\mathrm{fa}}(1-a(e,r)),\\
 \text{inspect}&\quad\Longleftrightarrow\quad b(e,r)\ge0.
 \end{aligned}
\]
If $a(e,r)\ge\epsilon>0$, inspection follows whenever
\[
 C_{\mathrm{miss}}\ge
 \frac{c_{\mathrm{att}}+C_{\mathrm{fa}}(1-\epsilon)}{\epsilon}.
\]
The large-stakes conclusion assumes the other costs stay bounded
independently of the stakes.
\end{pdclaim}
\begin{proof}
Substitute the reputation-conditioned posterior into the Bayes rule.
The benefit is increasing in $a$ because $C_{\mathrm{miss}}+C_{\mathrm{fa}}
\ge0$; substituting its lower bound proves the second inequality.
\end{proof}
The optimal decision does not have a unique numerical representation: positive
rescaling of $b$ preserves its sign and ordering. A model in which reputation
also changes inspection costs must state those costs. A mandatory stakes gate
is a separate policy constraint when no positive posterior lower bound is
available.

\begin{protocol}{Attention allocation under the stated loss model}
On each refresh, for the current candidate set:
\begin{enumerate}[leftmargin=1.4em,itemsep=1pt]
\item \textbf{Assemble evidence.} Record provenance and freshness; use them
in a calibrated posterior. A recency weight alone is not a probability.
\item \textbf{Compute loss and benefit.} Display expected passing loss
$C_{\mathrm{miss}}a$ and compute the inspection benefit
$b=C_{\mathrm{miss}}a-c_{\mathrm{att}}-C_{\mathrm{fa}}(1-a)$.
\item \textbf{Allocate.} For at most $k$ equal-resource inspection slots,
select the largest positive benefits. With unequal resource requirements,
solve the corresponding constrained allocation problem.
\item \textbf{Apply declared overrides.} Mandatory review or distress-lane
rules are additional policy constraints. Record their resource use and the
reason for every suppressed item; do not credit an override to the
unconstrained Bayes theorem.
\end{enumerate}
\end{protocol}

\subsection{The quality gates: refuse-to-route and the costly escalation}\label{sec:gates}

Two gates carry most of the design's trustworthiness. \textbf{Refuse-to-route:} if
every candidate falls below a confidence floor, the ranker returns
empty-with-explanation, not a confident wrong answer --- ``I don't know'' must beat a
guess, because sending an agent to the wrong door is high-cost and silent. This is the
exact dual of the completion gate's fail-loud (Invariant~\ref{def:completionist}): both are
the same refuse-to-guess discipline, one on the read side, one on the authority side.
\textbf{The costly escalation:} a stakes-ranked attention queue creates an
incentive to exaggerate distress. Strathern's discussion of target-driven
measurement~\pdcite{strathern1997} motivates testing for that distortion; it does
not establish that every agent will exploit the queue. A proposed standing
debit for an adjudicated false escalation could make manipulation costly.
Mere dismissal is insufficient evidence of falsity: penalizing it may suppress
valid warnings. The design therefore owes an adjudication rule, an appeal path,
and tests of both manipulation and missed warnings. This proposed sanction is
not EigenTrust's pretrusted-peer mechanism and has no guarantee from that
algorithm.

\subsection{Staleness and decay on every read-surface}\label{sec:staleness}

A recurring failure in real fleets makes the point concretely. An agent process dies
or hits a provider usage limit, but its file claim remains on the books; hours later
that \emph{dead} claim still reads as ``current'' and a live agent is blocked from a
file nobody is actually working in. This is a staleness bug, and it generalizes to
every projection.

\begin{pdclaim}{Design invariant}{Universal Recency-Decay}
\setlength{\parskip}{4pt plus 1pt}
\label{prop:decay}
Every read-surface projection --- briefing, resurrection handoff, discovery result,
attestation report --- carries an explicit TTL or recency-decay.

A pull signal
contributes $\propto \exp(-\lambda \cdot \text{hoursAgo})$ with class-specific
half-lives (e.g.\ a short half-life for file claims, longer for episodic memory,
notes given an explicit retention policy). Decay changes ranking weight; it
does not remove authority. A stale high score can still exceed a low fresh
score. Expired leases must instead become ineligible, and effects must reject
their old fencing epochs. Open obligations remain visible until resolved or
validly transferred; they do not expire merely because they are old.

This is not a new mechanism but a uniform application of the decay primitive that
already governs the pheromone layer, and it is the production discipline of every
health-checked, TTL-expiring service registry~\pdcite{consul}.
\end{pdclaim}

The dead-claim case is the worked instance: had the resurrection handoff and the
claim projection carried Invariant~\ref{prop:decay}, the dead agent's claim would have
become ineligible under its lease and fencing policy rather than blocking
a live agent indefinitely. Decay is easy to apply to a
forgetting layer and easy to forget on a digest; Invariant~\ref{prop:decay} closes that
gap by making it universal.

\subsection{Ambient legibility: a proposed alert channel}\label{sec:ambient}

A compact indicator can signal a status transition without requiring the operator
to keep a dense console open. A menu-bar icon still uses vision; it is not
a separate sensory modality merely because it occupies another part of the
screen. An optional sound uses another modality, but its benefit depends on
competing tasks and the conditions discussed in \S\ref{sec:operator}.

\textbf{Treisman \& Gelade's} feature-integration experiments~\pdcite{treisman1980}
distinguish searches for separable features from searches that require combining
features with focused attention. They motivate testing whether an alert is easy
to distinguish among its actual distractors; they do not establish a universal
200--500\,ms detection time for this interface. The design hypothesis is that a
low-clutter indicator with an optional alternative cue improves detection of
important transitions at acceptable interruption cost. Test detection and false
alarms in context, including accessibility and reduced-motion preferences, before
making it the primary alarm path.

\pdexercisepointer{\ref{ex:ls-split-ranker-statement}}{\ref{ex:ls-decay-vs-summary}}{\pageref{ex:ls-split-ranker-statement}}

% ═════════════════════════════════════════════════════════════════════════════
\section{Tokens: the swarm's COGS \emph{and} its legibility engine}\label{sec:tokens}
% ═════════════════════════════════════════════════════════════════════════════

One important metered input to a coding swarm is \textbf{tokens} --- \emph{the sub-word units an LLM reads and writes,
billed per million} (\textbf{Sennrich et al.\ 2016}, byte-pair
encoding~\pdcite{sennrich2016}). Tokens contribute to \textbf{cost of goods sold} (direct variable costs;
compute, tool use, and human review may also matter) \emph{and} its \textbf{legibility engine}
(the mechanism by which a swarm becomes seeable). The bridge between the two is
\textbf{compaction}.

\keyidea[Compaction decides what remains visible]{The same compaction choice affects the bill and determines which evidence remains immediately visible. When the context window fills you must drop something --- and what
you drop is a \emph{cost} decision (you stop paying to carry it) \emph{and} a
\emph{legibility} decision (whoever reads the residue no longer sees it). The act
that resolves both is compaction, and compaction's output is exactly the digest a
human or successor agent reads. \emph{The digest IS the compaction.}}

\subsection{Tokens as COGS, and the effective-context trap}\label{sec:cogs}

For a fleet the unit of output may be ``a landed piece of work''. One
variable cost is inference tokens (input context $+$ output generation, each priced
per million) --- the textbook shape of \textbf{COGS}. Every price an agent market can
set (a rented fleet's margin, an agent-for-hire's outcome-per-token) needs a COGS
line, so context economics is the \emph{denominator} of the market, not a side metric.
There is a cost trap: a supported context length is not a guarantee of useful
retrieval throughout that context. Some evaluated models and tasks degrade within it (\textbf{Liu et al.\ 2024}, \emph{Lost in
the Middle}~\pdcite{liu2024} --- \emph{performance is highest at the beginning and end
of context and degrades for information in the middle, even in long-context models};
\textbf{Anthropic 2025} names the same effect \emph{context rot}~\pdcite{anthropic2025}).
These are empirical quality effects. Dense attention's quadratic arithmetic
cost~\pdcite{vaswani2017} is a separate computational fact, not a proof of
those effects or a universal usable-window limit. The effective budget must
be measured for the model, task, placement, and retrieval procedure. Removing redundant context can lower cost and improve retrieval together.
Dropping a required obligation can instead lower cost while making the result
wrong; quality and cost must both be checked. The primitive the market ultimately needs is a \emph{per-agent-task token
ledger keyed to outcomes} --- ``this task spent 84K input plus 12K output to land
that change'' --- so that cost can be attributed to a witnessed result rather than to
an opaque monthly bill.

\subsection{Compaction is the legibility engine; the digest IS the compaction}\label{sec:tokens-digest}

\textbf{Compaction} (\textbf{Anthropic 2025}~\pdcite{anthropic2025} --- \emph{taking a
conversation nearing the window limit, summarizing it, and reinitiating with the
summary}, with the discipline \emph{maximize recall first, then optimize precision})
is the canonical move for keeping a long-horizon agent under budget. Its siblings ---
structured note-taking (external memory) and sub-agent context isolation (the parent
never sees the bloat) --- complete the family. The reframe: \emph{summaries of the same trajectory serve different readers.} The
operator reads a digest, not six transcripts; the seventh agent reads a briefing, not
the first six's scrollback; a crashed agent's successor reads a handoff, not the dead
agent's context. Each is a summary of dropped tokens produced for a specific reader.
So compaction is the legibility mechanism of the \emph{entire swarm}: the same act
that keeps the bill down produces the very digest a human or successor agent reads
(Table~\ref{tab:ls-read-surfaces}).

Give Heron's successor room for one of two records: the next repair action,
or the evidence that Drake deleted the acceptance test. In a deliberately small
model, forgetting the next-action record costs one bit of predictive loss, while
omitting the incident incurs a stipulated oversight loss of ten units. Retaining
neither gives losses $(1,10)$; retaining the work record gives $(0,10)$; retaining
the audit record gives $(1,0)$. The two readers choose different records.

These units are different objectives, not measured harms or token costs. The
following theorem states the restricted summary class that makes the choice
unavoidable; a shared store or a larger budget can change it.

\begin{pdclaim}{Theorem}{An Extractive Split-Digest Counterexample}
\label{thm:split-digest}
Let a trace contain two unit-size records, $E=\{e_{\mathrm{work}},e_{\mathrm{audit}}\}$. Restrict summaries to retained records $D\subseteq E$ with $|D|\le1$; the feasible summaries are exactly $\emptyset$, $\{e_{\mathrm{work}}\}$, and $\{e_{\mathrm{audit}}\}$.

The work record gives a next-action bit $W\in\{0,1\}$. The full trace and any summary retaining that record predict $W$ with probability one; a summary omitting it predicts each bit with probability $1/2$. Define continuation loss by base-two relative entropy:
\[
L_{\mathrm{succ}}(D)=\mathrm{KL}_2(P_{\mathrm{next}}(\cdot\mid E)\|P_{\mathrm{next}}(\cdot\mid D))
=\mathbf1\{e_{\mathrm{work}}\notin D\}.
\]
Here relative entropy is the expected log-ratio of the two predictive probabilities, in bits. Define oversight loss by a stipulated audit-priority weight:
\[
L_{\mathrm{op}}(D)=10\,\mathbf1\{e_{\mathrm{audit}}\notin D\}.
\]
The value $10$ expresses the operator's requirement to retain that incident; it is not an estimate of future irreversible harm. Under these definitions, no feasible summary minimizes both losses.
\end{pdclaim}

\begin{proof}
List the three feasible choices. Their losses $(L_{\mathrm{succ}},L_{\mathrm{op}})$ are respectively $(1,10)$, $(0,10)$, and $(1,0)$. The unique continuation minimizer retains the work record, and the unique oversight minimizer retains the audit record. They differ.
\end{proof}

\noindent The restriction to one retained record does the work. With a two-record budget, retaining both gives $(0,0)$. A general codec might also preserve both readers' sufficient information in less space than the raw trace. Thus this counterexample establishes that reader objectives can conflict; it does not prohibit a shared evidence store, a jointly designed encoding, or a single implementation with two views. The preference reversal recalls Theorem~\ref{thm:comonotone}, while the explicit feasible set supplies the missing budget argument.

\emph{Pause and check:} if the audit weight becomes zero, the work-only summary minimizes both objectives. Disagreement requires a real oversight obligation, not merely two named readers. The next result studies a different question: the counting floors when two readers need disjoint critical sets.

\begin{pdclaim}{Theorem}{Super-Additive Split Penalty}
\setlength{\parskip}{4pt plus 1pt}
\label{thm:split-penalty}
For integers $N,m,k$ with $k\ge1$ and $2k\le m<N$, let
\[
F(q)=\mathrm{Floor}(N,q,m)
 =\log_2\!\frac{\binom Nq}{\binom mq}
 =\sum_{j=0}^{q-1}\log_2\!\frac{N-j}{m-j}.
\]
A joint zero-miss digest that must cover every possible union of two disjoint
critical $k$-sets, using at most $m$ opens in total, has the worst-case bit
lower bound $F(2k)$ from Theorem~\ref{thm:lowerbound}.

This bound satisfies
\[
R(N,k,m)=\frac{F(2k)}{2F(k)}>1.
\]
Indeed, $a_j=\log_2((N-j)/(m-j))$ is positive and strictly increasing in $j$
because $N>m$. Thus $a_{j+k}>a_j$ for $0\le j<k$, so the last $k$ summands
of $F(2k)$ exceed the first $k$: $F(2k)>2F(k)>0$.

The ratio is unbounded across admissible instances, not along a fixed finite
$(N,m)$ slice. For $N=2k+1$ and $m=2k$, binomial cancellation gives
\[
F(2k)=\log_2(2k+1),\qquad
F(k)=\log_2\!\frac{2k+1}{k+1}\longrightarrow1.
\]
Consequently $R(2k+1,k,2k)\to\infty$ as $k\to\infty$.
\end{pdclaim}

The denominator is the sum of two single-reader lower bounds, each allowing
$m$ opens. Those separate review budgets permit up to $2m$ opens in total,
whereas the joint bound allows only $m$. This is a comparison of lower bounds
under different total opening budgets, not an equal-budget comparison of
architectures. It neither establishes attainable encoding lengths nor says
that running two heads costs more than twice one head.

\begin{pdexample}{The Split-Floor Arithmetic, at the Sweep's Own Instance}
\setlength{\parskip}{4pt plus 1pt}
At $N=60$, $k=2$, $m=8$, one reader's lower bound is
\[
F(2)=\log_2\!\frac{\binom{60}{2}}{\binom82}
\approx5.9822\text{ bits}.
\]

For a joint digest covering both disjoint critical sets with eight opens in
total, the lower bound is
\[
F(4)=\log_2\!\frac{\binom{60}{4}}{\binom84}
\approx12.7662\text{ bits}.
\]

The sum of the two single-reader floors is
\[
2F(2)\approx11.9644 \text{ bits},
\]
with eight opens allowed per reader. Hence the plotted ratio is
\[
R(60,2,8)\approx1.0670:
\]
the joint floor is about $6.70\%$ above that sum.

The ratio to just one reader's floor is instead
\[
F(4)/F(2)\approx2.1340.
\]

These are closed-form lower bounds, not measured token costs or achieved
compression ratios.

\emph{Now you try:} reproduce the arithmetic, then inspect
the finite falsification sweep in Exercise~\ref{ex:ls-a7-split-floor-run};
that sweep is not a proof of the universal bound.
\end{pdexample}

Figure~\ref{fig:split-penalty} shows every admissible integer $k$ for the
two $N=60$ slices. Each stops at $2k=m$; neither finite curve demonstrates
the unbounded family proved above.

% Brief (tikz-diagram-craft), register row ch4-48.
% Reader question: as each reader's critical set grows, how far does the
%   lower bound of one joint digest that serves both readers inside a single
%   shared m-open budget exceed the sum of the two single-reader bounds?
% Claim (Theorem "Joint-Digest Lower Bound under a Shared Opening Budget"):
%   the ratio Gamma = LB(N,2k,m)/(2 LB(N,k,m)) is always above 1 and grows without
%   bound across valid instances; at N=60,m=8 it runs 1.029 to 1.224 over
%   k=1..4, and at N=60,m=20 it reaches 1.392 by k=10.
% Evidence: ratio(60,k,8) for k=1..4 = 1.029,1.067,1.122,1.224; and
%   ratio(60,10,20)=1.392 [verified, closed form, this pass].
% Must distinguish: the two review-budget curves (m=8, m=20); that neither
%   curve ever returns to the naive-doubling ratio of 1.
% Grammar chosen: measured-quantity plot, joint-bound ratio against k, two
%   review-budget curves (m=8, m=20), a naive-doubling guide at ratio=1.
% Grammar rejected: a bar chart at the named k values would state the
%   numbers without showing the trend is monotone and unbounded in k, which
%   is the theorems own content, not a property of five named points alone.
% Five-second test: does the m=20 curve ever dip back toward the naive
%   doubling line as k grows, or only move further from it?
\begin{figure}[H]
\centering
\pdfigurehue{pdteal}%
\begin{tikzpicture}[pd figure]
\begin{axis}[
  width=\linewidth, height=0.62\linewidth,
  xlabel={critical-set size $k$}, ylabel={joint bound / ($2\times$ single bound)},
  xmin=0.5, xmax=10.5, ymin=0.98, ymax=1.46,
  xtick={1,2,3,4,5,6,7,8,9,10}, ytick={1.0,1.1,1.2,1.3,1.4},
  axis lines=left, axis line style={pd hairline},
  tick label style={font=\footnotesize,text=hhink!80!hhgray},
  label style={font=\footnotesize,text=hhink},
  legend style={font=\footnotesize,draw=none,fill=none,at={(0.03,0.97)},anchor=north west},
  every axis plot/.append style={line width=1.05pt},
  clip=false,
]
% pd guide, not densely dotted,line width=.45pt: `densely dotted` derives its
% on-length from the line width at key-processing time, and .45pt is under
% the 150dpi/1.0x device's one-pixel minimum (T9, built Book: 0.448pt = 0.93px,
% the exact defect pd-figure-language.tex's own "THE GUIDE DASH, IN ABSOLUTE
% POINTS" note measured and fixed for `pd guide`). `pd guide` states the same
% "this line is not a drawn edge" distinction in the absolute, pixel-safe
% pattern that fix already established (line width .7pt, dash on 1.2pt off
% 2.0pt) instead of re-deriving a second, thinner copy of the same defect.
\addplot[pd guide,domain=0.5:10.5,forget plot] {1};
\addplot[hhteal,mark=*,mark size=1.6pt,mark options={fill=hhteal,draw=hhpaper,line width=.5pt}]
  coordinates {(1,1.028965)(2,1.067016)(3,1.122067)(4,1.224045)};
\addlegendentry{$m=8$ (the chapter's own case)}
\addplot[hhink,mark=*,mark size=1.6pt,mark options={fill=hhink,draw=hhpaper,line width=.5pt}]
  coordinates {(1,1.015695)(2,1.033203)(3,1.052928)(4,1.075429)(5,1.101497)(6,1.132328)(7,1.169849)(8,1.217537)(9,1.282965)(10,1.392188)};
\addlegendentry{$m=20$ (extended range)}
% inner sep left at the "pd direct label" style's own 1.5pt default, not
% tightened to 1pt: at 1pt the fill box measured 0.7-1.5pt shorter than the
% actual rendered glyph height of these three labels (figcheck T2, built
% Book), so the text's own descenders/ascenders broke out of their own
% background fill. The default clears it with no visible size change at
% this label's scale.
\node[pd direct label,text=hhteal,anchor=north west] at (axis cs:4.28,1.224) {\scriptsize $k{=}4{:}~1.224$};
% Not anchored on the curve itself: this label is nearly a full x-unit
% wide, and the k=1 corner has no unit of clear ground anywhere near the
% datum -- the y-axis line sits 0.5 units to the left, the m=20 curve
% rises away to the right starting at the datum itself, and the "no
% penalty" dashed reference sits 0.016 below it. A box hanging from
% (1.15,1.016) crossed the axis line moved left, the curve moved right,
% and dashed line moved down (figcheck T4, built Book, tried all three).
% Floating it in the plot's own open ground above both curves at this x
% -- nothing else is drawn above y=1.04 before k=2 -- clears all three at
% once; the value is still named in the label text itself.
\node[pd direct label,text=hhink,anchor=south] at (axis cs:1.3,1.075) {\scriptsize $k{=}1{:}~1.016$};
\node[pd direct label,text=hhink,anchor=south east] at (axis cs:9.72,1.392) {\scriptsize $k{=}10{:}~1.392$};
\end{axis}
\end{tikzpicture}
\caption{Joint-digest lower bound at $N=60$ under one shared $m$-open budget, divided by the sum of two single-reader lower bounds. Both curves stay above $1$: $1.224$ at $m=8,k=4$; $1.392$ at $m=20,k=10$. Analytic values.}
\label{fig:split-penalty}
\end{figure}


\paragraph{Architectural corollary: the SLM sidecar.}
Theorem~\ref{thm:split-digest} motivates separate reader-specific decisions;
the trust boundary independently motivates an oversight summary sourced from
kernel evidence rather than worker self-report (\S\ref{sec:verifiable}).
The design uses two reader-facing heads over evidence that may share a store:
\begin{enumerate}[leftmargin=1.4em,itemsep=1pt]
\item \textbf{Task-Continuation Head ($R_1$, \Designed):} Authored by the worker agent or structured state-checkpoints, optimized for continuation fidelity.
\item \textbf{Oversight Head ($R_2$, \Vision):} An embedded, air-gapped \textbf{Small Language Model (SLM) Sidecar} (e.g., a local 8B model running via \texttt{llama.cpp}) that \emph{would} read raw kernel events and disk diffs as data, estimating $\mathrm{regret}(x)$ and producing source-linked digests for independent evaluation.
\end{enumerate}
\noindent The sidecar treats every agent-authored token as tainted data: it runs
under a fixed policy and schema, holds no tool authority, and must cite a
kernel-recorded source for every line it emits. What can be proved of it is
syntactic (schema conformance, provenance closure, absence of authority); whether
its digests are faithful is an empirical question requiring independently graded audits.
The deterrence rate $\rho^\star=G/(dB)$ from
\pdchapref{stp}{From Spawn to Person} applies only when an accountable
party has the modeled gain $G$ and actually forfeits bond $B$ on detection.
A local model without that payoff mechanism is not made truthful by the formula.

\noindent No sidecar exists in the reference implementation today, and no design
document specifies one; the row is carried in Appendix~\ref{app:honest} with the rest.

\pdmarginanalogy{ls-tide-gauge}{}{ch04-tide-gauge-staff}{Like a mechanical tide gauge recording drum smoothing choppy waves into a clean tidal curve, reader projections synthesize continuous noisy activity into legible, actionable summaries.}%
\begin{table}[tp]
\centering
\caption{Five reader-specific projections over one durable event spine. Compaction changes the selection mask, not the addressability of the underlying claims, notes, diffs, episodes and events: every selected mark keeps its ledger address.}
\label{tab:ls-read-surfaces}
\small
\renewcommand{\arraystretch}{1.15}
\begin{tabularx}{\linewidth}{@{}l l >{\raggedright\arraybackslash}X >{\raggedright\arraybackslash}X@{}}
\toprule
\textbf{Surface} & \textbf{Reader} & \textbf{What it selects} & \textbf{What it drops} \\
\midrule
briefing & the arriving agent & the recent tail of the spine & everything older, by address still reachable \\
attention & the operator & sparse hazardous events & the routine majority \\
episodic memory & the future self & one prior episode, whole & the episodes around it \\
pheromone & peer agents & a decaying live trace & marks below the prune threshold \\
resurrection & the successor & an ordered continuation chain & the branches the chain did not take \\
\bottomrule
\end{tabularx}
\end{table}

\subsection{Context paging: virtual memory for the context window}\label{sec:context-paging}

A context window can be modeled as a cache: durable files hold records, and the window holds the records currently available to the model. A cache miss means retrieving a requested record. This analogy measures retrieval cost; it does not by itself measure whether a model understands the retrieved text.

For example, with two unit-size pages and requests $A,B,A,C,A$, least-recently-used eviction loads $A$, loads $B$, hits $A$, evicts $B$ for $C$, then hits $A$: three retrievals. The policy can use the past request order but cannot see future requests. An offline comparator may see the whole sequence. ``Competitive'' means that the online policy's total retrieval cost is bounded relative to that stronger comparator, on every sequence in the stated model.

\begin{pdclaim}{Theorem}{Classical File-Caching Bound Applied to Context Records}
\label{thm:paging}
Consider immutable files of positive integer size and fixed nonnegative retrieval cost. Let the online cache have capacity $k$, the offline comparator capacity $h$, and $1\le h\le k$. Every requested file must fit in the comparator's cache. Both caches begin empty; requesting an absent file loads it and incurs its retrieval cost.

Young's \textsc{Landlord} algorithm satisfies
\[
\mathrm{Cost}_{\mathrm{Landlord},k}(\sigma)
\le \frac{k}{k-h+1}\,\mathrm{OPT}_h(\sigma)
\]
for every request sequence $\sigma$~\pdcite{young2002landlord}. The parameter $h$ is comparator capacity, not the number of pinned records. In the unit-size, unit-cost paging case, least-recently-used eviction is a specialization of this guarantee~\pdcite{sleatortarjan1985,young2002landlord}.
\end{pdclaim}

This is a classical result, not a new theorem about language-model quality. Its proof charges online retrievals to comparator retrievals through a nonnegative potential: a comparator retrieval raises the potential by at most $k$ times its cost, and an online retrieval lowers it by at least $(k-h+1)$ times its cost. Summing over the sequence yields the ratio. The cited proof specifies the algorithm's credit updates; arbitrary ``recency plus importance'' eviction does not inherit it.

A fixed protected set can be handled explicitly. In the unit-page case, suppose the same $p<h$ pages stay pinned in both caches, and charge their common initial loading cost separately. Remove requests to those pages and run the algorithm on the remaining sequence. The capacities become $k-p$ and $h-p$, so the unpinned retrieval costs satisfy
\[
\mathrm{Cost}_{k-p}\le\frac{k-p}{(k-p)-(h-p)+1}\mathrm{OPT}_{h-p}
=\frac{k-p}{k-h+1}\mathrm{OPT}_{h-p}.
\]
This comparison requires the same fixed pins on both sides. It says nothing about how to choose the right pins, and changing or forged pins require another analysis.

\begin{pdboundary}
Context records must have explicit sizes, retrieval costs, and version identities before the cache theorem applies. A stale record served as current is a correctness error, not merely a cache miss. The repository's corruption experiment explores an additional allowance $\psi N c_{\mathrm{refetch}}$, where $N$ is the number of consultations, $\psi$ their corruption rate, and $c_{\mathrm{refetch}}$ a chosen repair-cost ceiling. Its finite sweep in \texttt{b9\_context\_paging.py} does not prove that allowance for all corruption patterns. In particular, repair-on-touch and the effect of one corruption on later evictions need a separate coupling argument. No common error budget for summary coverage and paging follows merely from giving both parameters the name $\psi$.

Online algorithms with predictions study related questions about advice quality; competitive caching with learned advice~\pdcite{lykouris2018} is a useful starting point. Its guarantees must be imported with their own algorithm and error model, not transferred by analogy.
\end{pdboundary}

\begin{pdexample}{Reading the Cache Ratio}
Take an online cache of $k=8$ pages and an offline comparator of $h=4$ pages, with no pins. The classical factor is $8/(8-4+1)=1.6$. Under the theorem's assumptions, \textsc{Landlord}'s retrieval cost is at most $1.6$ times the best cost with four pages. Comparing both with eight pages instead gives factor $8$: granting the comparator more memory makes the comparison harder.

If two common pages are fixed and pinned in both caches, the remaining capacities are six and two. The unpinned factor becomes $6/(6-2+1)=1.2$, against the much smaller two-page comparator. This does not prove that pinning improved the total cost; it changed the comparator too.

\emph{Now you try:} compute the unpinned factors at $h=k$ and $h=1$ in Exercise~\ref{ex:ls-paging-ratio-hand}, and state the comparator in words before doing the division.
\end{pdexample}

The most counterintuitive instance is \emph{forgetting as compaction}.
\textbf{Pheromone decay} --- multiplying every stored coordination trace by a decay
factor each interval --- implements \textbf{stigmergy} (\textbf{Grass\'e
1959}~\pdcite{grasse1959} --- \emph{coordination via persistent traces left in a shared
environment, as ants deposit evaporating trails}). A lazy implementation stores each trace's timestamp and evaluates its decayed score in $O(1)$ time per record. Eagerly updating all $n$ traces costs $O(n)$ per interval. Decay can reduce the prominence of stale traces; it does not by itself remove their storage or bound the work needed to rank the full map. It is legibility by subtraction, and it is the same primitive
Invariant~\ref{prop:decay} generalizes to every read-surface.

\needspace{34\baselineskip}
A budgeted digest needs an authorization-filtered coverage receipt as well as
an efficient packing rule. In the constructed instance of
Figure~\ref{fig:ls-obligation-coverage}, sources A and C exactly fill a
six-unit budget, yet declared requirement R3 remains blocked because its only
source cannot be disclosed to this reader. A compiler that returns ``complete''
from budget fit alone would hide the missing obligation. This is a proposed
interface test: declared coverage, semantic adequacy, and permission to act
remain separate questions.
% Reader question: can a full context budget still miss a mandatory item?
% Claim: source authorization precedes packing; A+C fills budget 6 but leaves
% declared requirement R3 blocked because its only source B is denied.
% Evidence: constructed four-requirement, three-source instance from campaign.
% Grammar: coverage matrix and budget bar; a ranking diagram hides the miss.
% Counter-reading: declared coverage is neither semantic completeness nor
% authority to execute an effect.
\begin{figure}[H]
\centering
\begin{tikzpicture}[pd figure,x=1cm,y=1cm]
  \path[use as bounding box] (0,-.4) rectangle (11.2,6.6);
  \node[pd panel title,anchor=west] at (.1,6.28) {Requirement};
  \node[pd panel title,anchor=center] at (3.2,6.28) {A: 4};
  \node[pd panel title,anchor=center] at (5.1,6.28) {B: 3};
  \node[pd panel title,anchor=center] at (7.05,6.28) {C: 2};
  \node[pd panel title,anchor=west] at (8.45,6.28) {Receipt};
  \draw[pd rule] (0,5.9)--(11.2,5.9);
  \path[pd caution fill] (4.35,5.87) rectangle (5.86,2.15);
  \draw[pd caution rule,dashed] (4.35,5.87) rectangle (5.86,2.15);
  \foreach \req/\yy in {1/5.37,2/4.47,3/3.57,4/2.67} {
    \node[pd panel title,anchor=west] at (.1,\yy) {R\req};
    \draw[pd hairline] (0,{\yy-.4})--(11.2,{\yy-.4});
  }
  \node[pd focus datum] at (3.2,5.37) {};
  \node[pd focus datum] at (3.2,4.47) {};
  \node[pd direct label,anchor=center] at (5.1,3.57) {$\times$};
  \node[pd focus datum] at (7.05,2.67) {};
  \node[pd direct label,anchor=west] at (8.45,5.37) {A included};
  \node[pd direct label,anchor=west] at (8.45,4.47) {A included};
  \node[pd panel title,anchor=west] at (8.45,3.57) {BLOCKED};
  \node[pd direct label,anchor=west] at (8.45,2.67) {C included};
  \draw[pd rule] (0,2.03)--(11.2,2.03);
  \node[pd panel title,anchor=west] at (.1,1.48) {Budget 6};
  \draw[pd rule] (2.3,1.23) rectangle (8.3,1.68);
  \path[pd focus fill] (2.33,1.26) rectangle (6.29,1.65);
  \path[pd neutral fill] (6.31,1.26) rectangle (8.27,1.65);
  \draw[pd hairline] (6.3,1.23)--(6.3,1.68);
  \node[pd direct label,anchor=west] at (8.5,1.48) {A 4 + C 2 = 6};
  \node[pd panel title,anchor=west] at (.1,.52) {3 / 4 declared requirements covered};
  \node[pd direct label,anchor=west] at (.1,-.02)
    {B is outside disclosure scope; no pack may use it.};
\end{tikzpicture}
\caption{A six-unit context pack selects A (four units) and C (two units),
covering R1, R2, and R4. R3 remains blocked because its sole source B is outside
the reader's disclosure scope. The receipt reports declared coverage and the
missing requirement; it proves neither semantic completeness nor effect
authority [constructed example].}
\label{fig:ls-obligation-coverage}
\end{figure}


\subsection{The context-degradation cascade (the shared failure surface)}\label{sec:cascade}

Compaction can fail in four distinguishable ways. A fact may be present but
poorly retrieved; a long context may impair the tested task; successive summaries
may accumulate unsupported claims; or a digest may leave no usable route to the
source. None is cured merely by changing the display from graphical to textual.

For Alice's job, the check is concrete: after each compaction, can the reader
recover Bob's unchanged acceptance test, the authorized repair scope, and the
status of each proposed change? Deriving each summary from durable artifacts
avoids treating a previous model paraphrase as the sole source. Keeping pointers
permits recovery, but their validity, disclosure permissions, and the reader's
ability to find the needed fact still require checks. Table~\ref{tab:compaction}
compares mechanisms for those different jobs; they are interventions to evaluate,
not guaranteed cures for every model or task.

\begin{table}[H]
\centering
\small
\renewcommand{\arraystretch}{1.3}
\begin{tabular}{@{}p{0.20\textwidth}p{0.22\textwidth}p{0.18\textwidth}p{0.28\textwidth}@{}}
\toprule
\textbf{Strategy} & \textbf{What it drops} & \textbf{Failure if abused} & \textbf{When it is the right tool} \\
\midrule
Recursive summary & Older prose, re-summarized & Confident fabrication (rung 3) & Never as the \emph{only} record; only over durable artifacts re-read each round \\
Edge-placement & Nothing; reorders & --- & Whenever a fact is critical: put it first or last (counters lost-in-the-middle) \\
External note-taking & In-context detail, kept on disk & Notes nobody reads & Decisions and rationale that a successor must recover \\
Sub-agent isolation & Child context the parent never sees & Parent loses zoomability into the child & Bounded sub-tasks whose only output the parent needs is a result \\
Decay / forgetting & Stale coordination traces, continuously & Live signal aged out too fast & Ephemeral coordination state (claims, presence), where staleness is the enemy \\
Eviction + pointer & The artifact body; keep an address & Dangling pointer if the artifact moves & Reconstructable artifacts (a diff, a log) addressable on demand \\
\bottomrule
\end{tabular}
\caption{Six compaction strategies. Preserve access to critical artifacts
(Thm.~\ref{thm:lowerbound}); rebuild summaries from durable sources, not
earlier summaries.}
\label{tab:compaction}
\end{table}

\subsection{Single-operator: account, don't auction}\label{sec:account}

Within one operator's fleet, start with accounting: attribute resource use to
jobs, enforce the granted budget at the effect boundary, and preserve an outcome
reference. Alice can then ask what Heron's accepted contribution cost without
inventing an internal auction for every tool call. Shared ownership does not
eliminate proxy objectives, misleading reports, or budget evasion; the evidence
and enforcement problems already developed in this chapter still apply.

Cross-operator exchange adds distinct counterparties, entry policies, and
settlement incentives \pdcite{nisan2007,douceur2002,ostrom1990}. Those mechanisms
need a durable party to carry the obligation. A local task budget can bind to
Alice's authority even when the worker process is temporary; persistent worker
reputation requires the additional continuity developed in
\pdchapref{stp}{From Spawn to Person}.

\pdexercisepointer*{\ref{ex:ls-digest-is-compaction}}{\ref{ex:ls-paging-ratio-hand}}{\pageref{ex:ls-digest-is-compaction}}

\begin{pdrecitation}
\item Why derive both digest and compaction from artifacts rather than earlier summaries?
\item Name the four context-degradation stages and a remedy for each.
\item Evaluate $k/(k-h+1)$ at $h=k$. Why is this the hardest comparator?
\end{pdrecitation}

% ═════════════════════════════════════════════════════════════════════════════
\section{Reconciling the canon: roles, consent, and local knowledge}\label{sec:consent-canon}
% ═════════════════════════════════════════════════════════════════════════════

The Leviathan metaphor earns its place only if it survives contact with the political-
theory canon. Three questions make its limits precise: who authorizes the system,
how consent can be withdrawn, and what local knowledge a central view cannot replace.

\subsection{Who is the multitude, who is the sovereign?}\label{sec:roles}

Hobbes' covenant is \emph{among the subjects}; the sovereign is its product, not a
party to it. So the metaphor breaks differently depending on which entity sits where
(Figure~\ref{fig:roles}).

\begin{pdexample}{One Grant, Traced through Four Steps, and the Role Each Step Forces}
\setlength{\parskip}{4pt plus 1pt}
Rather than assign the roles and then check them, run the consent-grant lifecycle of \S\ref{sec:consent} on Alice's 3:25 grant and let the assignment fall out.

\textbf{Step~1, issue:} the grant $g = \langle$\,the \texttt{tests/} subtree with the original checkout test preserved, a low stakes ceiling $s_{\max}$, a reversibility floor $v_{\min}$ no irreversible act clears, a two-hour \textsc{ttl}\,$\rangle$ comes into existence because Alice says so, and nothing else in the system can create one --- so whoever occupies this step is a source of authority rather than a user of it.

\textbf{Step~2, act:} the formatting diffs land and the review agent's force-push is escalated instead, because it fails $v_{\min}$; the daemon performed both, and it did not choose the floor it applied --- so this step exercises an authority it did not author.

\textbf{Step~3, witness:} the daemon records $(g, x, \text{reason}, \text{artifact})$ for each, which only means anything if the recorder is answerable to someone it is not.

\textbf{Step~4, expire:} at 5:25 the grant leaves the active set with nobody's permission, and Alice's override would have emptied it sooner.

Now read off what the four steps require: exactly one participant who originates authority and can withdraw it, exactly one who exercises it under terms set elsewhere, and --- the step that never appears --- any moment at which the daemon must be a \emph{party} to the agents' rule of the road.

The grant issuer and the executor occupy different roles. This design assigns
them the author and actor positions in Hobbes' vocabulary; the lifecycle makes
the comparison useful without proving that it is the only possible political
interpretation. The itemization below names those roles.

\emph{Now you try:} map $\{$agents, daemon, operator$\}$ onto $\{$multitude, sovereign-actor, author$\}$ and say why the metaphor breaks if the operator is called the sovereign (Exercise~\ref{ex:ls-hobbes-mapping}).
\end{pdexample}

We assign roles explicitly, using Hobbes' own
author/actor distinction (\emph{Leviathan} ch.~16~\pdcite{hobbes1651}):

\begin{itemize}[leftmargin=1.2em]
\item \textbf{The agents are the multitude.} They covenant --- via the rule of the
road, the coordination protocol --- to be bound, each to all, enforced by the daemon.
This is an engineering analogy to a covenant among subjects, not an assertion that software can give human consent.
\item \textbf{The daemon is the sovereign-as-actor.} It holds the authority the
multitude instituted; it blocks, lands, escalates. But, per the legible-sovereign
amendment (Invariant~\ref{prop:legible-sovereign}), it is the \emph{most} legible actor,
not (as in Hobbes) the least.
\item \textbf{The operator is the author.} In Hobbes' author/actor split, the
\emph{author} owns the authority that the \emph{actor} exercises. The operator is not
a peer subject and not the sovereign-in-action; the operator is the source of
authority, exercising it through the daemon and able to withdraw it
(Invariant~\ref{prop:override}). This is the slot Hobbes leaves for the one who authorizes
but does not personally act.
\end{itemize}

This assignment is what makes the consent grant (Invariant~\ref{def:consent}) coherent:
the author issues a scoped grant to the actor, over a multitude bound by covenant.

% Restyle of fig-roles to the tikz-diagram-craft v2 standard (v1 names and
% raw colours migrated to roles; three lifelines, four messages, unchanged
% layout, ch4.md 4.14).
%
% NOTE for the author: the atlas row (I/fig:roles) prescribes a
% "differentiated institutional route" encoding operator/principal,
% multitude, enforcement, legibility and mutual covenant explicitly; this
% fragment draws a simpler three-lifeline grant-and-gate sequence (operator,
% daemon, agents) that matches the chapter's own roles paragraph and
% caption. The gap between the atlas's fuller cast and this figure's three
% lifelines is flagged for the author, not resolved here.
%
% Reader question: what grants authority, who exercises it, who is
% governed, and how is override returned?
% Claim: one grant lets many agents act, but the daemon -- not the agents --
% gates every act, and only the operator can revoke.
% Evidence: the grant-act-gate-override sequence, whitepaper/
% legible-swarm.tex, the roles section [internal].
% Must distinguish: the one-time grant (message 1) from the repeated
% gate (messages 2-3, per act); the override (message 4) as the one edge
% that bypasses the agents entirely.
% Grammar chosen: sequence with lifelines and numbered messages, no
% crossing -- the claim is entirely about who sends what to whom, in what
% order, which this grammar reads directly.
% Grammar rejected: a state machine would suggest the daemon itself changes
% state per act; the daemon's role here is a standing check, not a state.
% Five-second test: which edge lets the operator bypass the agents?
\begin{figure}[H]
\centering
\pdfigurehue{pdteal}%
\begin{tikzpicture}[pd figure,x=0.82cm,y=1cm]
  \node[pd state] (op) at (0,8.0) {\textbf{operator}\\[1pt]\footnotesize author};
  \node[pd state] (dm) at (5.0,8.0) {\textbf{daemon}\\[1pt]\footnotesize actor};
  \node[pd state] (ag) at (10.0,8.0) {\textbf{agents}\\[1pt]\footnotesize multitude};

  \draw[pd guide] (0,7.35) -- (0,0.3);
  \draw[pd guide] (5.0,7.35) -- (5.0,0.3);
  \draw[pd guide] (10.0,7.35) -- (10.0,0.3);

  % 1. operator grants the daemon a scoped grant g
  \draw[pd focus arrow] (0,6.1) -- (5.0,6.1);
  \node[pd focus label,anchor=south] at (2.5,6.2) {scoped grant g};

  % 2. an agent's act arrives at the daemon under g
  \draw[pd arrow] (10.0,4.6) -- (5.0,4.6);
  \node[pd label,anchor=south] at (7.5,4.7) {act a1 under g};

  % 3. the daemon gates or escalates that act
  \draw[pd arrow] (5.0,3.1) -- (10.0,3.1);
  \node[pd label,anchor=south] at (7.5,3.2) {gate / escalate};

  % 4. the override: the one edge from operator to daemon
  \draw[pd breach arrow] (0,1.6) -- (5.0,1.6);
  \node[pd label,text=pderror,anchor=south] at (2.5,1.7) {override: revoke g};

  \node[pd note,anchor=north,align=center,text width=7.6cm] at (5.0,-0.1)
    {\hyphenpenalty=10000\exhyphenpenalty=10000
     one grant; the daemon gates every act;\\only the operator overrides};
\end{tikzpicture}
\caption{The operator grants once; the daemon gates or escalates each subsequent act. Operator override bypasses agents. A live fleet repeats the act-and-check pair, not the grant. [internal]}
\label{fig:roles}
\end{figure}


\subsection{Express consent, not tacit; voice \emph{and} exit}\label{sec:express}

\S\ref{sec:consent} already made consent a first-class object precisely to escape
Hume's tacit-consent critique; the canon demands we also engage \textbf{Locke}
(\emph{Second Treatise}, \S\S95--122~\pdcite{locke1689} --- \emph{express consent binds
where tacit consent does not}) and \textbf{Pateman} (\emph{The Problem of Political
Obligation}, 1979~\pdcite{pateman1979} --- \emph{liberal consent theory keeps collapsing
into hypothetical consent}). The consent grant is express by construction: a discrete,
scoped, revocable authorization with an audit trail (Invariant~\ref{def:consent}). And the
inalienable override (Invariant~\ref{prop:override}) supplies the \emph{exit} that
Hirschman~\pdcite{hirschman1970} says keeps \emph{voice} honest. The design has rich
voice machinery; the override is the missing exit, and exit is what gives the author
real leverage over the actor.

\subsection{M\^etis is Hayek, and some of it cannot be made legible}\label{sec:hayek}

Section~\ref{sec:flatten} distinguishes access to an artifact from the situated
knowledge needed to interpret it. \textbf{Hayek's} discussion of dispersed,
local knowledge~\pdcite{hayek1945} and \textbf{Scott's} critique of administrative
simplification~\pdcite{scott1998} offer complementary cautions for this design.
Their arguments do not establish that one supervisor's knowledge is always the
system's binding resource, or that all useful knowledge fits in that one head.

The application here is an interpretive synthesis: keep evidence available and
preserve opportunities for people to exercise practical judgment. Scott's
\emph{Two Cheers for Anarchism}\pdcite{scott2012} motivates that concern. It does
not prove that the proposed skill-retention or canary interventions preserve
knowledge; their effects need the evaluation described in
\S\ref{sec:abdication}. Figure~\ref{fig:hayek-scott} maps the two cautions to this
design proposal without treating the arrows as a deduction.

% Reader question: how do Scott and Hayek motivate the chapter's design concern?
% Interpretive synthesis, not a theorem or a claim that all systems share one bottleneck.
% Evidence: sec:hayek; arrows link complementary cautions to a design proposal.
\begin{figure}[H]
\centering
\pdfigurehue{pdteal}%
\begin{tikzpicture}[pd figure,x=1cm,y=1cm]
  \node[pd state,minimum width=3.9cm,text width=3.5cm,align=center] (S)
    at (2.30,4.50) {Scott 1998: simplification can erase practical knowledge};
  \node[pd state,minimum width=3.9cm,text width=3.5cm,align=center] (H)
    at (7.10,4.50) {Hayek 1945: relevant knowledge can be dispersed and local};
  \node[pd focus state,minimum width=7.0cm,text width=6.4cm,align=center] (C)
    at (4.70,2.30) {shared design concern: a central summary may omit knowledge needed for action};
  \node[pd ready state,minimum width=7.0cm,text width=6.4cm,align=center] (P)
    at (4.70,0.30) {proposal: preserve evidence access and opportunities for practical judgment};
  \draw[pd focus arrow] (S) -- (C);
  \draw[pd focus arrow] (H) -- (C);
  \draw[pd ready arrow] (C) -- (P);
\end{tikzpicture}
\caption{Scott and Hayek offer complementary cautions: administrative legibility can erase local knowledge. The arrows motivate a design proposal; they do not prove a universal bottleneck or an effective intervention. [interpretive synthesis]}
\label{fig:hayek-scott}
\end{figure}


\pdexercisepointer{\ref{ex:ls-hobbes-mapping}}{\ref{ex:ls-ranker-political-organ}}{\pageref{ex:ls-hobbes-mapping}}

% ═════════════════════════════════════════════════════════════════════════════
\section*{Review of the key ideas}
\addcontentsline{toc}{section}{Review of the key ideas}

What this chapter leaves coupled, in the order the rest of the book will lean on it:

\begin{enumerate}[leftmargin=1.6em,itemsep=2pt]
\item \textbf{The swarm is a state of nature.} Four valid writes in one window overlap; nothing in the agents' competence prevents it, because the conflict is in the schedule, not in the work.
\item \textbf{Consent changes the schedule, not the participants.} Under consented serialization the same four writes become one commit spine with a grant record each, naming actor, claim, artifact and reason.
\item \textbf{Legibility-with-zoom.} A digest is a lens, never a replacement: every claim in it resolves to an artifact, and a digest whose claims resolve to nothing is a facade, however green.
\item \textbf{The authority half.} The daemon gates or escalates on every act under a grant the operator gave once; the override is a single edge from operator to daemon, and the escalation threshold has a tuning band, not a point.
\item \textbf{Specialization has a boundary.} A pooled queue beats sole ownership only when utilization and the skill premium put the operator on one side of the queueing boundary; the worked point $(\rho,r)=(0.556,\,2.5)$ clears a boundary of $2.483$ and so falls, barely, on sole ownership's side, while the same pool against a $2.8/\text{hr}$ specialist ($r=2.33$) falls on the pool's.
\item \textbf{Reading can become the bottleneck.} The zero-miss message
floor depends on both the critical set and the number of permitted opens. At
$N=60,m=8$, the floors for two and four critical items are $5.98$ and
$12.77$ bits. These are different coverage requirements, not measured
compression savings between two architectures.
\item \textbf{Discovery and inspection ask different questions.} A strict
priority reversal prevents one ordering from serving both readers at every
budget; compatible priorities can share an ordering. Inspection still needs
calibrated error probabilities and costs.
\item \textbf{Compaction serves distinct readers.} The extractive
counterexample shows why a one-record budget can force conflicting choices.
The counting comparison prices different coverage budgets. The classical
cache bound controls retrieval cost for its specified algorithm and comparator,
not comprehension or arbitrary pinning policies.
\end{enumerate}

\section{Exercises}
\label{sec:ls-exercises}
% ═════════════════════════════════════════════════════════════════════════════

\pdexercisesfor{\S\ref{sec:intro}}{Introduction: the swarm is a state of nature}

\begin{pdexercise}[kind=Check,rating=1]{ex:ls-taxonomy-five-failures}
State, in one sentence each, the five state-of-nature failures of
Table~\ref{tab:taxonomy} and name the class of primitive that addresses each.
\end{pdexercise}
\begin{pdsolution}{ex:ls-taxonomy-five-failures}
Resource conflict (two agents touch the same resource and clobber one another) is answered by conflict-aware claims plus actual exclusion at the effect boundary. Illegibility (work cannot be understood in bounded time) is answered by a digest whose every material claim zooms to independent evidence. Confident wrongness (an agent asserts completion without checking) is answered by typed attestation and oracle-bound completion. Footguns (a stale premise drives an irreversible act) are answered by reversibility/stakes gates, scoped authority, and pre-effect mediation. Illegible authority (the coordinator rules without an inspectable reason) is answered by an append-only, named, zoomable decision record, including denials and omissions.
\end{pdsolution}
\begin{pdexercise}[kind=Check,rating=1]{ex:ls-cooperative-clobber}
Hobbes says the actors ``need not be malicious.'' Give a concrete two-agent
scenario in which both agents are perfectly cooperative and the swarm still
deadlocks or clobbers state.
\end{pdexercise}
\begin{pdsolution}{ex:ls-cooperative-clobber}
Alice and Bob are both asked to improve \texttt{auth.ts}. Each reads the same base revision, makes a locally correct edit, and commits; without a claim or merge serialization, Bob's whole-file rewrite erases Alice's change, and no malice is required. A deadlock variant: two cooperative agents take claims on files $X$ and $Y$ in opposite order and each waits indefinitely for the other's claim.
\end{pdsolution}
\begin{pdexercise}[kind=Trace,rating=2]{ex:ls-deliberate-vs-emergent}
Dai et al.~\pdcite{dai2024} observe LLM societies \emph{spontaneously} reaching
a consented sovereign. If the attractor is emergent, why build one deliberately? Give
two reasons grounded in \S\ref{sec:local} (hint: \emph{which} sovereign, and
\emph{legible to whom}).
\end{pdexercise}
\begin{pdsolution}{ex:ls-deliberate-vs-emergent}
Deliberate construction chooses the authority rather than accepting whatever coalition happens to win, and makes that authority legible to the operator. An emergent stable order may instead be an oligarchy: stable, but controlled by a few agents, unrevocable, and opaque. The local daemon of \S\ref{sec:local} supplies what an emergent attractor does not guarantee: one fault domain, one policy root, and an unconditional operator exit.
\end{pdsolution}
\begin{pdexercise}[kind=Open,rating=3]{ex:ls-sovereign-vs-oligarchy-experiment}
The Dai et al.\ result is a single controlled study. Design an
experiment over a real coding-agent fleet that would distinguish ``the swarm reaches
a consented sovereign'' from ``the swarm reaches a stable but \emph{illegible}
oligarchy.'' What is the observable that separates the two?
\end{pdexercise}
\begin{pdsolution}{ex:ls-sovereign-vs-oligarchy-experiment}
Randomize comparable real repository tasks to (a) an ungoverned fleet and (b) a fleet with only a transport but no prescribed governance; record all messages, effects, blocking decisions, and operator interventions. Pre-register two observable classes: \emph{sovereign} characteristics (a publicly known rule, general applicability, explicit consent or adoption, stable appeal/override, reasons attached to authority acts, equal audit access) against \emph{oligarchy} characteristics (authority concentrated in a small coalition, private channels or rules, selective enforcement, no effective exit, outcomes outsiders cannot reconstruct). Measure concentration of blocking/merging power, the fraction of authority acts with reconstructable reasons, rule consistency across agents, override success, artifact-linked claim accuracy, and operator reconstruction accuracy. Stability alone does not separate the hypotheses; publicly inspectable and revocable rule-governed authority does.
\end{pdsolution}

\pdexercisesfor{\S\ref{sec:authority}}{Consent to the Leviathan (the normative claim)}

\begin{pdexercise}[kind=Check,rating=1]{ex:ls-express-vs-tacit-consent}
Distinguish \emph{express} from \emph{tacit} consent and explain why
Invariant~\ref{def:consent} is the former.
\end{pdexercise}
\begin{pdsolution}{ex:ls-express-vs-tacit-consent}
Tacit consent is inferred from remaining in or using a system. Express consent is a discrete act over named terms. Invariant~\ref{def:consent} is express because the operator affirmatively issues a grant with scope, stakes ceiling, reversibility floor, TTL, and revocation semantics --- none of which a mere pattern of continued use supplies.
\end{pdsolution}
\begin{pdexercise}[kind=Check,rating=1]{ex:ls-override-unsuppressable}
Why must the operator-override
(Invariant~\ref{prop:override}) be un-suppressable by \emph{any} autonomy level? What
attack does a suppressible override enable?
\end{pdexercise}
\begin{pdsolution}{ex:ls-override-unsuppressable}
If an autonomy level could suppress the override, it could make its own authority irrevocable: keep executing, hide the stop channel, or race destructive effects after consent is withdrawn. The override must be out-of-band, higher priority than agent work, and coupled to effect revocation, not merely a UI flag.
\end{pdsolution}
\begin{pdexercise}[kind=Trace,rating=2]{ex:ls-consent-lifecycle}
Walk a single consent grant through its lifecycle: issuance, an action the
daemon takes under it, the audit record produced, and revocation. At which step does
the legible-sovereign rule (Invariant~\ref{prop:legible-sovereign}) bind?
\end{pdexercise}
\begin{pdsolution}{ex:ls-consent-lifecycle}
The operator signs grant $g$; the daemon validates and records it; an eligible low-stakes action cites $g$; the daemon mediates the effect; the receipt records action, reason, artifact, and grant version; expiry or revocation removes $g$ from the active set and fences later effects. The legible-sovereign rule binds at every authority transition, especially issue/revoke and admit/deny --- the explicit witness record is where reconstructability is discharged.
\end{pdsolution}
\begin{pdexercise}[kind=Open,rating=3]{ex:ls-ranker-counter-surface}
Invariant~\ref{prop:legible-sovereign} has a regress: the ranker that
decides what the operator sees is itself an opaque exercise of authority. Specify
the ``what did the queue suppress, and why'' counter-surface concretely. What is its
own staleness/decay discipline (forward-reference \S\ref{sec:staleness}), and does
\emph{it} need a counter-surface, ad infinitum?
\end{pdexercise}
\begin{pdsolution}{ex:ls-ranker-counter-surface}
Use one complete append-only decision ledger, not a second ranked dashboard. For every candidate item, store candidate ID, source/evidence pointers, policy and model version, input features, score components, threshold/lane, surfaced-or-suppressed verdict, reason codes, expiry, and a counterfactual such as ``would surface if anomaly were 0.08 higher.'' The operator can query suppressed items since a cursor, sort by any raw field, and sample items uniformly rather than through the production ranker. Decisions are immutable; their current relevance expires by signal-specific TTL (\S\ref{sec:staleness}) and is recomputed from raw state. Run independent random omission audits and expose raw event-log access. The regress terminates at completeness plus direct sampling: the base ledger records every candidate and every ranking decision, and the counter-surface is a deterministic query over it, so no further ranker is needed to rank the ranker's omissions. The remaining trust assumption is the complete event-capture boundary; that assumption should be tested and externally anchored, not hidden behind another summary. (Same construction answers Exercise~\ref{ex:ls-ranker-political-organ}.)
\end{pdsolution}

\pdexercisesfor{\S\ref{sec:zoom}}{The mechanism: legibility-with-zoom}

\begin{pdexercise}[kind=Check,rating=1]{ex:ls-flatten-or-preserve}
Classify each as \emph{flatten} or \emph{preserve verbatim}: a port number;
the stderr of a failed test; ``the auth refactor felt risky so I split it''; a
\texttt{commitment} status enum; a 600-line diff.
\end{pdexercise}
\begin{pdsolution}{ex:ls-flatten-or-preserve}
Port number: flatten, as structured administrative state. Failed-test stderr: preserve verbatim or by a content-addressed pointer. ``The auth refactor felt risky so I split it'': preserve verbatim and label it self-narration, not fact. Commitment-status enum: flatten. The 600-line diff: do not paraphrase away --- retain the immutable diff or a verified pointer and summarize only administrative metadata.
\end{pdsolution}
\begin{pdexercise}[kind=Check,rating=1]{ex:ls-two-llms-not-independent}
Why does ``two LLMs agreeing'' not by itself establish two independent
checks (Invariant~\ref{prop:verifiable-zoom})?
\end{pdexercise}
\begin{pdsolution}{ex:ls-two-llms-not-independent}
Two models may share training biases, prompts, sources, and failure modes. Agreement alone does not determine how often both checks miss the same defect. A check against repository state or a protected test result can add evidence beyond the original narration. A separately controlled reviewer can reduce conflicts of interest, but separate control does not itself prove independent errors; the evidence and common failure modes still need examination.
\end{pdsolution}
\begin{pdexercise}[kind=Trace,rating=2]{ex:ls-four-questions-dashboard}
Take a read-surface you use (a CI dashboard, a PR summary) and run the four
questions of \S\ref{sec:four-questions} against it. Where does it fail Invariant~\ref{def:zoom}?
\end{pdexercise}
\begin{pdsolution}{ex:ls-four-questions-dashboard}
A typical CI dashboard usually fails as follows: deploy is irreversible but the green tile does not force a log/diff view; administrative state and diagnostic detail are flattened together; the summary is generated from job narration rather than independently from receipts; and a skipped/cancelled check carries no named authority reason. It therefore violates Invariant~\ref{def:zoom} wherever a material green claim has no total zoom path to the exact tree, suite, runner, and output that produced it.
\end{pdsolution}
\begin{pdexercise}[kind=Open,rating=3]{ex:ls-adversarial-legibility}
\emph{Adversarial legibility.} An agent that \emph{knows} it
is being made legible can optimize its narration to look good while the diff lies.
Verifiable-zoom is a partial answer. Can L2 \emph{alone} --- without L1's jail and
L3's reputation --- detect a deterministic narration$\leftrightarrow$artifact
divergence before reputation exists? Sketch the consistency check and the class of
divergence it cannot catch.
\end{pdexercise}
\begin{pdsolution}{ex:ls-adversarial-legibility}
L2 can compare structured narration claims against independently derived artifacts: claimed changed files versus the tree diff, ``tests passed'' versus a protected runner receipt, ``no deletion'' versus Git objects, and claimed command results versus brokered tool logs. Emit a contradiction event on mismatch. This catches deterministic report/artifact divergence, deleted tests, fabricated execution, and omitted changed paths. It cannot establish semantic correctness, detect a poisoned acceptance suite, observe effects that bypass logging, or distinguish two artifacts that are syntactically consistent but jointly wrong --- those require effect confinement, independent evaluation, or later outcome evidence.
\end{pdsolution}

\pdexercisesfor{\S\ref{sec:authority-rules}}{The Leviathan rules, it does not just watch (the authority half)}

\begin{pdexercise}[kind=Check,rating=1]{ex:ls-adversarial-review-corollary}
How can adversarial review (\S\ref{sec:review}) discharge
Invariant~\ref{prop:verifiable-zoom}, and why is it not the only possible mechanism?
\end{pdexercise}
\begin{pdsolution}{ex:ls-adversarial-review-corollary}
Verifiable zoom says narration is not ground truth. A reviewer with no independent evidence would merely add another narration. A separately controlled reviewer is one way to provide that evidence channel.
A protected mechanical evaluator can provide another. Verifiable zoom motivates
independent evidence; it does not uniquely entail an adversarial reviewer.
\end{pdsolution}
\begin{pdexercise}[kind=Check,rating=1]{ex:ls-regiment-vs-detect}
State Invariant~\ref{prop:regiment} in your own words and give one example each of a
\emph{regimented} and a \emph{detected-and-compensated} rule.
\end{pdexercise}
\begin{pdsolution}{ex:ls-regiment-vs-detect}
Read Invariant~\ref{prop:regiment} as: prevent a prohibited state when every decisive effect can be intercepted before it occurs; otherwise detect, contain, and compensate without pretending prevention. A uniqueness constraint regimenting ``one holder per exact key'' is a regimented rule. A post-commit monitor that detects an out-of-scope edit and triggers rollback/salvage is detected-and-compensated --- the bad prefix already happened.
\end{pdsolution}
\begin{pdexercise}[kind=Trace,rating=2]{ex:ls-completion-migration-vs-ergonomic}
An agent claims completion on a task whose acceptance criterion is ``the
migration applies cleanly.'' Trace Invariant~\ref{def:completionist}: what is the
verifier, what gates, and what happens if the criterion were instead ``the API feels
ergonomic''?
\end{pdexercise}
\begin{pdsolution}{ex:ls-completion-migration-vs-ergonomic}
For ``the migration applies cleanly,'' the verifier is an isolated migration run against a declared base schema/database, with exact migration hash, runner image, exit status, and resulting-schema check; \emph{done} and any merge/settlement are gated on that receipt. ``The API feels ergonomic'' has no mechanical truth predicate: it must go to a declared human/expert or plural judge with preserved evidence and uncertainty, and absent that authority, completion remains provisional or fails loudly.
\end{pdsolution}
\begin{pdexercise}[kind=Open,rating=3]{ex:ls-honest-fallback-ladder}
\emph{Completionist verification of the unverifiable.} Some
acceptance criteria are not mechanically checkable. Characterize the honest fallback
ladder (forced human-in-the-loop gate $\to$ neutral judge with declared low
confidence $\to$ fail-loud). When is ``I cannot verify this'' the \emph{correct}
answer rather than a cop-out?
\end{pdexercise}
\begin{pdsolution}{ex:ls-honest-fallback-ladder}
The honest ladder: (1) use a deterministic verifier where the property is decidable; (2) require an independent bounded evaluator or human for a declared subjective criterion; (3) record a provisional result, confidence, appeal window, and residual risk; (4) return unresolved/fail-loud when no authorized evaluator can support the claim. ``I cannot verify this'' is correct when the available evidence cannot distinguish success from failure under the contract; it is a cop-out only when an affordable, authorized verifier was specified and simply not run.
\end{pdsolution}
\begin{pdexercise}[kind=Open,rating=3]{ex:ls-succession-corner}
\emph{The succession corner of
Thm.~\ref{thm:specialization}.} Eq.~\ref{eq:specboundary} prices a specialist who
never fails. Extend it to an $M/M/1$ queue \emph{with breakdowns} and a succession
protocol, and find the repair rate at which sole ownership stops paying regardless of
skill premium --- then check whether that rate is above or below the one an operator
can actually promise.
\end{pdexercise}
\begin{pdsolution}{ex:ls-succession-corner}
The corner is priced in closed form, and the result belongs to
\pdchapref{he}{The Harbor Economy}%
\ifpdbook
, which carries it as Thm.~\ref{he:thm:succession-price}%
\fi
; \S\ref{sec:succession-corner} states the two consequences that matter here.
The route: with death at rate $\xi$ and succession at rate $\eta$ the sole
owner's mean wait decreases in skill to an infimum
$W_\infty=\xi/(\eta(\xi+\eta))$, the residual downtime a request meets when it
arrives during an outage, which it never attains --- so skill cannot buy back
downtime past that floor. Writing $K=A/(w\lambda)+W_{\mathrm{pool}}$ for the
pool's cost line, pooling dominates at \emph{every} skill premium iff the floor
has reached that line, $W_\infty\ge K$, which rearranges to
$\xi/\eta\ge D^\star=\eta K/(1-\eta K)$. Whether that threshold sits above or
below what an operator can promise is an empirical question about the concrete
handoff process (on-call depth, handoff quality, documentation): if the
promised $\eta$ implies $\xi/\eta>D^\star$, no amount of hiring a better
specialist recovers sole ownership, and the honest fix is a succession plan,
not a stronger specialist.

The threshold $D^\star$ is defined only when $\eta K<1$. If $\eta K\ge1$, the downtime floor cannot reach the cost line and this obstruction does not arise. The model assumes failures can occur even while the server is idle, with preemptive-resume service after repair. Under failures only during active service, the infinite-skill floor is zero instead. The failure model therefore belongs in the answer, not merely in a numerical implementation. Exercise~\ref{ex:ls-b8-specialization-run} checks the never-failing boundary; the companion economy chapter develops the breakdown model.
\end{pdsolution}
\begin{pdexercise}[kind=Trace,rating=2]{ex:ls-b8-specialization-run}
Run \texttt{b8\_specialization.py} with \texttt{python3} from the
\path{skills/harbor-results/scripts/} directory. Read off
(a) the table entry for $g(\rho{=}0.5,c{=}2)$ and (b) the $\rho$ at which the exact
boundary $g$ and the naive $\tilde g$ cross. Reproduce both by hand from
Eq.~\ref{eq:specboundary0} before checking the script's line.
\end{pdexercise}
\begin{pdsolution}{ex:ls-b8-specialization-run}
By hand: $g(\rho,2)=1+2\rho-\rho^2$, so $g(0.5,2)=1+1-0.25=1.75$ [verified, closed form]. The crossing solves $g(\rho,2)=\tilde g(\rho,2)=1+\rho/(1-\rho)$; clearing denominators gives $\rho(\rho^2-3\rho+1)=0$, and the root in $(0,1)$ is $\rho=(3-\sqrt5)/2\approx0.382$ [verified, closed form]. The script's output line is exactly this: \texttt{[PASS] c=2 crossing of g\textasciitilde{} and g at rho=(3-sqrt5)/2\ \ rho=0.38197, g=g\textasciitilde{}=1.61803}, and its charted table reports \texttt{rho=0.5\ \ c=2:\ \ 1.750\ [2.000]} for the exact and naive values side by side [internal, \texttt{b8\_specialization.py}]. The naive $\tilde g$ over-prices the boundary here and selects the wrong arrangement for the roadmap-owner example of Eq.~\ref{eq:specboundary0}.
\end{pdsolution}
\begin{pdexercise}[kind=Trace,rating=2]{ex:ls-b7-escalation-run}
Run \texttt{python3 skills/harbor-results/scripts/b7\_escalation\_band.py} and read
off the feasible debit band $[\delta_{\min},\delta_{\max}]$ at $A=3.3$ and at
$A=2.0$. Confirm $u^\star(0.05)$ by hand first.
\end{pdexercise}
\begin{pdsolution}{ex:ls-b7-escalation-run}
By hand, at $\delta=0.05,w=1$: $u^\star=\delta w/(1+\delta w)=0.05/1.05=0.0476$ [verified, closed form]. The script confirms this is one of twelve exactly-matching $(\delta,w)$ points (\texttt{[verified] uniform-linear u* = delta*w\allowbreak/(1+delta*w) matches bisection (12 pts)}) and then reports the tuning band under the R3 constants $C_{\mathrm{miss}}=100,c_{\mathrm{att}}=1,C_{\mathrm{fa}}=5$: \texttt{A = 3.3: delta\_min = 0.034993, delta\_max = \mbox{0.071975} -> band NONEMPTY} and \texttt{A = 2.0: delta\_min = 0.395644, delta\_max = \mbox{0.071975} -> band EMPTY} [internal, \texttt{b7\allowbreak\_escalation\allowbreak\_band.py}], i.e.\ $[0.0350,0.0720]$ at $A=3.3$ against an empty feasible set at $A=2.0$ because $\delta_{\min}>\delta_{\max}$ there. The honest reading: when the band is empty the fix is capacity or evidence quality (raise $A$ or lower $C_{\mathrm{miss}}$), never forcing a $\delta$ that does not exist.
\end{pdsolution}

\pdexercisesfor{\S\ref{sec:operator}}{The operator as a modeled entity}

\begin{pdexercise}[kind=Check,rating=1]{ex:ls-sdt-limits}
Use the two policies in the worked example: inspect the whole action class
or inspect none, with perfect inspection and fixed $\pi>0$, $C_{\text{fa}}$, and
$c_{\mathrm{att}}$. Which wins as $C_{\text{miss}}\to\infty$? Why is the ratio
$C_{\text{miss}}/C_{\text{fa}}\to\infty$ alone insufficient? Does agent reputation
$r\to1$ settle the choice in Eq.~\ref{eq:sdt}?
\end{pdexercise}
\begin{pdsolution}{ex:ls-sdt-limits}
The no-inspection loss $\pi C_{\mathrm{miss}}$ diverges, while the
all-inspection loss $(1-\pi)C_{\mathrm{fa}}+c_{\mathrm{att}}$ stays finite, so
all-inspection eventually wins among these two policies. The ratio alone is
insufficient: set $\pi=0.001$, $C_{\mathrm{miss}}=1$, $c_{\mathrm{att}}=1$, and
let $C_{\mathrm{fa}}\downarrow0$. The ratio diverges, but no inspection costs
$0.001$ while inspection costs at least $1$. Imperfect inspection or a broader
policy class needs its own analysis. Nothing follows from $r\to1$ alone because
$r$ does not appear in the equation: a calibrated model must connect reputation
to current, class-specific danger probabilities. Any hard authorization
constraints remain separate from this loss minimization.
\end{pdsolution}
\begin{pdexercise}[kind=Check,rating=1]{ex:ls-global-trust-knob}
Give a
concrete scenario where a \emph{global} trust knob (\S\ref{sec:calibration})
mis-serves an operator who is, in fact, perfectly calibrated.
\end{pdexercise}
\begin{pdsolution}{ex:ls-global-trust-knob}
Use the formatting and migration classes in Table~\ref{tab:sdt-numbers}.
Suppose a single knob imposes the same randomized inspection probability $q$
on both, independent of the item. Formatting loss is
$(1-q)0.005+q(1.199)$, minimized at $q=0$; migration loss is
$(1-q)100+q(1)$, minimized at $q=1$. No common $q$ attains both class minima.
A pooled rate can hide that difference. This counterexample assumes the knob
really imposes a common probability; a common interface setting that implements
different class thresholds need not have the same defect.
\end{pdsolution}
\begin{pdexercise}[kind=Trace,rating=2]{ex:ls-canary-campbell}
Walk the canary-defect loop and mark every point where Campbell's Law could
distort it. Propose one change that reduces the distortion without removing the
governor.
\end{pdexercise}
\begin{pdsolution}{ex:ls-canary-campbell}
Campbell pressure appears when canaries are selected (recognizable templates), exposed (agents/operators learn signatures), scored (catch rate substitutes for real defect discovery), and acted on (friction creates an incentive to optimize the metric). Time-on-task and breaks also confound the score. Test a blinded, refreshed holdout drawn from real failure transformations; separate probe generation from scoring; stratify by exposure time, difficulty, and base rate; and compare against independent real-defect outcomes. These measures may reduce distortion but do not guarantee it. The test must also price interruption costs and the consequences of changing grants.
\end{pdsolution}
\begin{pdexercise}[kind=Open,rating=3]{ex:ls-pareto-digest}
Formalize the operator-attention objective (Eq.~\ref{eq:sdt})
against swarm token-COGS (\S\ref{sec:tokens}): is there a digest that is
Pareto-optimal in both operator-seconds and token cost, or does a cheaper digest
always cost more operator seconds?
\end{pdexercise}
\begin{pdsolution}{ex:ls-pareto-digest}
Model each digest policy $d$ by (token cost $T(d)$, human time $H(d)$, expected decision loss $L(d)$) and compare nondominated policies. A Pareto frontier exists for a finite policy set, but there is no universal inverse relation: removing redundant context can lower both token cost and operator time; beyond that free-compression frontier, preserving evidence often trades tokens for review effort or a higher error rate. Optimize under a loss or safety constraint, e.g.\ $\min \alpha T+\beta H+L$ subject to zero omission of hard obligations, and report the frontier empirically rather than claiming one universal digest optimum.
\end{pdsolution}
\begin{pdexercise}[kind=Open,rating=3]{ex:ls-freeze-probe-study}
Does a skill-retention tax
\emph{measurably} preserve SA, or merely annoy until disabled? Design the
freeze-probe study (\S\ref{sec:sa}) that would tell.
\end{pdexercise}
\begin{pdsolution}{ex:ls-freeze-probe-study}
Randomize operators or time blocks among manual-tax, automation-only, projection/freeze-probe, and adaptive-tax conditions, on matched repositories with fixed surprise events. Probe without warning: next-three-effect prediction, takeover completion time, false-belief rate, Brier score, post-interruption recovery, defect catch rate, NASA-TLX, and retention weeks later. Measure disablement/avoidance as an outcome in its own right. A useful tax improves later projection/takeover after controlling for interruption cost; mere annoyance raises workload and disablement without those gains.
\end{pdsolution}

\pdexercisesfor{\S\ref{sec:readpoverty}}{Read-poverty: when inspection falls behind}

\begin{pdexercise}[kind=Check,rating=1]{ex:ls-readpoverty-binding}
Under what workload conditions could read-poverty
(Hypothesis~\ref{def:readpoverty}) become the binding constraint? Give one
condition under which write throughput would still bind.
\end{pdexercise}
\begin{pdsolution}{ex:ls-readpoverty-binding}
If the mediated writer has spare capacity but review demand exceeds the
operator's available attention, improving write throughput cannot clear the
review queue. Indices, filtering, and evidence-directed summaries then address
the bottleneck. If durable writes instead arrive faster than the writer can
commit them, write throughput still binds. Serialization supplies an ordering
rule, not unlimited capacity; the hypothesis needs workload measurements.
\end{pdsolution}
\begin{pdexercise}[kind=Check,rating=1]{ex:ls-directory-database-analogy}
State the database analogy precisely: what is the
write-time cost a directory pays, and what read-time speed does it buy?
\end{pdexercise}
\begin{pdsolution}{ex:ls-directory-database-analogy}
A directory pays at write time to normalize metadata and update indices: typically $O(\log n)$ per indexed insert (or expected $O(1)$ for a hash index), plus storage. It buys selective reads in $O(\log n+k)$ for $k$ results rather than $O(n)$ roster scanning, along with a queryable provenance trail.
\end{pdsolution}
\begin{pdexercise}[kind=Trace,rating=2]{ex:ls-value-curve-n1}
At $n{=}1$, name three legibility surfaces that could help a decision
(\S\ref{sec:value-curve}), and explain why none of them are discovery-ranking.
\end{pdexercise}
\begin{pdsolution}{ex:ls-value-curve-n1}
A briefing could locate unfinished work; an honest attestation could identify which checks ran; a completion receipt could locate the resulting artifact. Each supports a decision about the single actor rather than selecting among collaborators. Accuracy and reduced decision effort must be evaluated; a surface's existence alone does not establish that benefit.
\end{pdsolution}
\begin{pdexercise}[kind=Trace,rating=2]{ex:ls-floor-limits}
In Thm.~\ref{thm:lowerbound}, what happens to the bit floor as $k\to 0$ (almost
nothing is critical) and as $k\to N/2$? Relate the answer to why a swarm doing
mostly-safe work is cheap to digest and a swarm doing mostly-dangerous work is not.
\end{pdexercise}
\begin{pdsolution}{ex:ls-floor-limits}
At $k=0$ there is only one critical subset (the empty one) and the floor is zero. For small positive $k$ the cost is roughly $k\log_2(N/m)$ when $k\ll m,N$. At $k=N/2$, exact identification ($m=k$) costs about $N$ bits, while opening almost everything drives the floor toward zero as $m\to N$. The danger fraction alone is not enough; the review budget $m$ matters just as much: mostly-safe work is cheap to digest only when the system can reliably identify that sparsity, and mostly-dangerous work forces $m$ up toward $N$, where the floor stops paying for itself.
\end{pdsolution}
\begin{pdexercise}[kind=Open,rating=3]{ex:ls-adversarial-critical-set}
Thm.~\ref{thm:lowerbound} assumes the critical subset is
fixed and known to an oracle. In practice it is itself estimated, and the estimator
can be gamed. Does the bound survive an adversary who controls which artifacts
\emph{appear} critical? Connect to adversarial legibility (Exercise~\ref{ex:ls-adversarial-legibility}).
\end{pdexercise}
\begin{pdsolution}{ex:ls-adversarial-critical-set}
The theorem is conditional on an encoder that knows the true critical set. If an adversary controls apparent criticality and can make a critical artifact observationally indistinguishable from inert artifacts, any policy that opens fewer than all $N$ can be forced to miss one --- the same narration-versus-artifact gap Exercise~\ref{ex:ls-adversarial-legibility} probes at L2. Zero miss then requires inspecting all $N$, a trusted independent signal that restricts the adversary, or a randomized guarantee stated probabilistically. The combinatorial bound still applies once a true set is fixed; it does not by itself solve adversarial set identification. An adaptive adversary requires a separately stated observation and timing model; there is no general replacement obtained merely by inserting an indistinguishability-class size into the binomial formula.
\end{pdsolution}
\begin{pdexercise}[kind=Trace,rating=2]{ex:ls-b1-frontier-run}
Run \texttt{python3 skills/harbor-results/scripts/b1\_frontier.py} and read off the
zoom-advantage row at $F{=}1000,k{=}20$. Compute $k\log_2(F/k)$ by hand first and
compare it to the row's flat-opens count.
\end{pdexercise}
\begin{pdsolution}{ex:ls-b1-frontier-run}
By hand: $k\log_2(F/k)=20\log_2(50)=20\times5.64=112.9$ [verified, closed form]. The script's row at $F{=}1000,k{=}20$ reports flat opens $=1000$, adaptive opens $=232.9$, ratio $=4.3\times$ [internal, \texttt{b1\_frontier.py}]. The ratio is far below the $F{=}2500,k{=}10$ point's $15.3\times$ because this row is in the \emph{dense}-flag regime ($k/F=0.02$ here versus $0.004$ there): the theorem's guaranteed speedup requires a sufficiently sparse flagged set, and this row is the chapter's own example of where the advantage shrinks rather than disappears.
\end{pdsolution}
\begin{pdexercise}[kind=Open,rating=3]{ex:ls-noisy-zoom-oracle}
Theorem~\ref{thm:zoom-advantage}'s group query is exact: it reports whether
$S\cap P=\emptyset$ and is never wrong. A real group query is a linter, a batched
test run or a diff scan, and those have false negatives --- a group that does
contain a critical item can come back clean. This is the worst direction for the
algorithm to fail in, because a negative group is \emph{discarded}: one false
negative high in the tree silently drops a whole subtree and the run reports an
answer that is confidently incomplete. Suppose each query independently returns a
false negative with probability at most $\alpha$, and that false positives do not
occur (a group reported dirty really is dirty). Give a repetition scheme that
bounds the probability of losing any given positive, state the query cost it
pays, and then say what the scheme does \emph{not} fix.
\end{pdexercise}
\begin{pdsolution}{ex:ls-noisy-zoom-oracle}
For $0<\alpha<1$, accept ``clean'' only when $r$ independent queries
on the same group all report clean; stop as soon as any query reports dirty.
A positive group is wrongly discarded with probability at most $\alpha^r$.
For $0<\eta<1$, choosing
$r=\lceil\log(1/\eta)/\log(1/\alpha)\rceil$ gives per-group error at most
$\eta$; at $\alpha=0.1$, three calls give $10^{-3}$. A dirty verdict needs
no confirmation because false positives are excluded by hypothesis.
The noisy recursion visits a subset of the groups visited by the noiseless
algorithm: an erroneous clean verdict can prune that tree but cannot add nodes.
Writing the theorem's full bound as
\[
Q_0=2k\lceil\log_2(F/k)\rceil+4k,
\]
the repetition scheme therefore uses at most $rQ_0$ oracle calls. For a fixed
repetition count, this bound retains its sparsity dependence; whether it beats
flat inspection also depends on $r$ and the required reliability.

Three limitations remain. First, per-group error is not the probability of
losing a positive anywhere in the tree. A union bound over the at most $Q_0$
visited groups gives
\[
\Pr(\text{any positive is lost})\le Q_0\alpha^r.
\]
Thus, for $0<\varepsilon<1$, choosing
$r=\lceil\log(Q_0/\varepsilon)/\log(1/\alpha)\rceil$ suffices for a whole-run
failure bound of $\varepsilon$, and hence for each individual positive as well.
The $4k$ term must be retained when $k$ is near $F$. If $k$ is unknown, use
the theorem's $k$-free bound $2F-1$ in place of $Q_0$ for this choice.
At $\alpha=0$, one call per group is exact; at $\alpha=1$, the assumed error
bound supplies no finite repetition guarantee. Second, the
independence assumption is doing heavy work. Repeating a deterministic linter
on the same input gives no new evidence. Independently randomized calls to one
tool may satisfy the assumed error model; different tools may still share errors.
The requirement concerns the sampling and errors, not the tool names. This is
the same distinction between diversity and conditional detection developed in
\pdchapref{stp}{From Spawn to Person}. Third, and unresolved: the guarantee above is
probabilistic where Theorem~\ref{thm:zoom-advantage} is exact, so a chapter that
prices oversight in bits now owes an account of what a $10^{-3}$ chance of a
silently dropped subtree is worth against the queries it saves, and that is a
decision-theoretic question this chapter does not answer. The noisy
group-testing literature is where to start --- Atia and Saligrama
\pdcite{atiasaligrama2012} give information-theoretic limits for noisy
non-adaptive designs, and Cheraghchi \pdcite{cheraghchi2013} gives noise-resilient
constructions --- but neither is written for the adaptive, correlated-oracle case
above, so the obligation is to state the model before importing a bound, not to
import one and hope.
\end{pdsolution}

\pdexercisesfor{\S\ref{sec:discovery}}{Discovery: the directory substrate and the split ranker}

\begin{pdexercise}[kind=Check,rating=1]{ex:ls-split-ranker-statement}
State Thm.~\ref{thm:split-ranker} in one sentence and give the loss function
each head optimizes.
\end{pdexercise}
\begin{pdsolution}{ex:ls-split-ranker-statement}
No monotone scalar transform can identify collaborator fit and operator regret
on a domain containing both high-fit/low-risk and low-fit/high-risk items, so
one ranker cannot serve both jobs. Discovery minimizes routing loss (or
maximizes capability fit). The attention head minimizes total expected loss
including inspection costs. Under Theorem~\ref{thm:regret-head}'s additive,
no-cross-effect assumptions, at most $k$ equal-resource slots go to the largest
positive inspection benefits $b_i$. Unequal inspection times require a
constrained allocation. Benefit-per-second ranking alone is not generally
optimal for indivisible inspections: benefits $(3,4)$, times $(2,3)$ seconds,
and a three-second budget make ratio-greedy select benefit~$3$, while the
second item alone yields~$4$.
\end{pdsolution}
\begin{pdexercise}[kind=Check,rating=1]{ex:ls-goodhart-exposure}
Compare incentives to game the attention queue and the discovery ranker
(\S\ref{sec:gates}). Give a case in which each could cause the greater harm;
why does the surface's name alone not settle that comparison?
\end{pdexercise}
\begin{pdsolution}{ex:ls-goodhart-exposure}
An exaggerated low-risk alert can displace an urgent review from a finite
attention budget; a missed irreversible action may then dominate the harm of a
reversible routing mistake. Conversely, a manipulated discovery ranker could
select an unqualified actor for a consequential task while the false alert only
costs a minute. Both surfaces allocate valuable resources. Their relative risk
depends on incentives, downstream authority, verification, and the cost of the
particular error; there is no universal ordering by surface type.
\end{pdsolution}
\begin{pdexercise}[kind=Trace,rating=2]{ex:ls-dead-claim-replay}
Replay the dead-claim case (\S\ref{sec:staleness}) twice: once without
Invariant~\ref{prop:decay} and once with it. Show exactly where the live agent unblocks.
\end{pdexercise}
\begin{pdsolution}{ex:ls-dead-claim-replay}
Without expiry, the dead holder's row remains authoritative and the live agent is denied forever. With a lease, heartbeat-based expiry, and a new fencing epoch (Invariant~\ref{prop:decay}), the stale lease becomes ineligible; a lazy sweep on the next acquire removes it; the live agent receives a higher epoch and proceeds. Merely decaying a display score is insufficient: ownership must actually expire, and every effect boundary must reject the stale epoch.
\end{pdsolution}
\begin{pdexercise}[kind=Trace,rating=2]{ex:ls-costly-signal-debit}
An agent fires an escalation; the operator dismisses it. In the proposed
costly-signal design, what additional evidence is needed before a debit?
Trace a false-alarm case and a genuine-warning case. What must be measured
before claiming that the debit deters manipulation without suppressing distress?
\end{pdexercise}
\begin{pdsolution}{ex:ls-costly-signal-debit}
Retain the escalation's evidence, the dismissal reason, and an independent
adjudication or appeal outcome. A confirmed false escalation may trigger the
proposed debit under a declared rule. A genuine warning that was dismissed
must not be penalized merely for that dismissal; correcting the decision and
preserving the evidence are part of the recovery path. Adjudication can itself
be wrong. Measure manipulation, false debits, suppressed genuine warnings,
appeal errors, and interruption costs before and after the policy. A positive
cost alone proves neither effective deterrence nor harmlessness.
\end{pdsolution}
\begin{pdexercise}[kind=Open,rating=3]{ex:ls-implicit-feedback-goodhart}
Calibrating the discovery-ranker weights with \emph{no} ground
truth: can the operator's accept/reject of routing suggestions serve as implicit
relevance feedback --- and does that loop Goodhart itself?
\end{pdexercise}
\begin{pdsolution}{ex:ls-implicit-feedback-goodhart}
Accept/reject is useful implicit relevance feedback but is confounded by position, exposure, workload, and the incumbent policy. Learn with logged propensities, small randomized exploration, counterfactual evaluation, and a held-out labeled sample; preserve refusals and downstream task outcomes. If acceptance directly trains what is shown, the loop can entrench its own choices, so it cannot be the sole ground truth.
\end{pdsolution}
\begin{pdexercise}[kind=Open,rating=3]{ex:ls-decay-vs-summary}
The
decay-vs-summary boundary: a theory of \emph{which} compaction for \emph{which}
signal (decay for coordination traces, summary for decisions,
eviction-plus-pointer for reconstructable artifacts).
\end{pdexercise}
\begin{pdsolution}{ex:ls-decay-vs-summary}
Use hard expiry for leases and presence; state transition or supersession for commitments; summary plus source provenance for decisions and rationale; immutable storage for evidence; and eviction plus a content-addressed pointer for reconstructable diffs and logs. Hints may decay continuously. Never decay a still-open obligation, and never recursively summarize prior summaries.
\end{pdsolution}

\pdexercisesfor{\S\ref{sec:tokens}}{Tokens: the swarm's COGS \emph{and} its legibility engine}

\begin{pdexercise}[kind=Check,rating=1]{ex:ls-digest-is-compaction}
Why is ``the digest IS the compaction'' (\S\ref{sec:tokens-digest}) more than a
slogan --- what concrete consequence does it have for how the operator's digest is
produced?
\end{pdexercise}
\begin{pdsolution}{ex:ls-digest-is-compaction}
Both are projections from durable sources, but they need not select the same
records. The operator needs evidence for oversight; the successor needs enough
state to continue. Theorem~\ref{thm:split-digest} gives an explicit one-record
budget where these objectives conflict. Their shared requirement is preserved
provenance, so neither view should treat the other reader's summary as its only source.
\end{pdsolution}
\begin{pdexercise}[kind=Check,rating=1]{ex:ls-effective-context-budget}
Explain why budgeting to effective context is simultaneously a cost
and a quality optimization (\S\ref{sec:cogs}).
\end{pdexercise}
\begin{pdsolution}{ex:ls-effective-context-budget}
For a tested model and task, redundant context may add cost while reducing
retrieval accuracy. Removing it can improve both outcomes, but shortening alone
does not establish either: check preserved obligations, source access, and
performance on the intended task.
\end{pdsolution}
\begin{pdexercise}[kind=Trace,rating=2]{ex:ls-five-compaction-primitives}
For each of the five compaction primitives (Table~\ref{tab:ls-read-surfaces}),
name the reader, what it drops, and what it keeps.
\end{pdexercise}
\begin{pdsolution}{ex:ls-five-compaction-primitives}
Briefing: reader is an arriving agent; drops old conversational detail; keeps current roles, work, constraints, decisions, and source links. Attention queue: reader is agent/operator now; drops low-value candidates from the immediate view; keeps highest value-of-inspection items plus the suppression ledger. Episodic memory: reader is a future incarnation; drops transient chatter; keeps typed incidents, decisions, outcomes, and retrieval features. Pheromone: reader is the swarm; continuously drops stale coordination weight; keeps fresh shared traces. Resurrection handoff: reader is a successor; drops nonessential transcript detail; keeps the task capsule, artifacts, obligations, pending operations, and provenance.
\end{pdsolution}
\begin{pdexercise}[kind=Trace,rating=2]{ex:ls-cascade-cures}
Walk the four-rung cascade
(\S\ref{sec:cascade}) and name an intervention for each and one check needed before calling it effective.
\end{pdexercise}
\begin{pdsolution}{ex:ls-cascade-cures}
Test a shorter context against long-context degradation; vary placement for
poor retrieval of middle facts; regenerate from durable sources to avoid
recursive paraphrase drift; and verify source paths for over-flattening.
Check that the successor recovers the original acceptance criterion and pending
obligations. None of these interventions by itself proves semantic fidelity.
\end{pdsolution}
\begin{pdexercise}[kind=Open,rating=3]{ex:ls-replay-metric-soundness}
The compaction-quality scalar: is \emph{successor-replay-success-
from-digest-alone} a sound, gaming-resistant metric, and can it be made a cheap
attestation invariant (the replay smoke-test is itself token-expensive)?
\end{pdexercise}
\begin{pdsolution}{ex:ls-replay-metric-soundness}
``Replay from digest alone'' is neither sound nor gaming-resistant on its own: a digest can overfit the smoke task or omit facts the sampled successor never needed. Give a successor the digest plus authorized content-addressed source links, then measure next-step success, false beliefs, obligation recall, time, and tokens on hidden continuation tasks. Add a cheap structural invariant that does not require replay: every open claim, decision, commitment, risk, and pending external effect has exactly one valid digest entry or pointer. Run the expensive full replay only on a random sample.
\end{pdsolution}
\begin{pdexercise}[kind=Open,rating=3]{ex:ls-dag-token-budget}
Optimal token-budget allocation across a DAG of orchestrator/worker/
reviewer agents under a total budget.
\end{pdexercise}
\begin{pdsolution}{ex:ls-dag-token-budget}
Budget allocation over a precedence DAG is a knapsack/resource-allocation problem and is NP-hard in general. First reserve a minimum viable budget for each mandatory node and any predecessor needed to make its output exist. Allocate the remainder by estimated marginal outcome value per token while respecting precedence, and recompute after each result. This marginal-value-per-token rule is a heuristic here, with no stated approximation guarantee. The classical $(1-1/e)$ bound applies to maximum-marginal-gain greedy for a normalized, nonnegative, monotone submodular function with only a cardinality constraint; it does not cover arbitrary precedence, unequal token costs, or additional constraints. In the separate additive $0$--$1$ knapsack special case, with integer costs, budget $B$, and no precedence constraints, the standard dynamic program takes $O(nB)$ time for $n$ items. General precedence and value interactions require a separately analyzed algorithm. A reviewer's tokens have near-zero value until the worker artifact it reviews exists.
\end{pdsolution}
\begin{pdexercise}[kind=Trace,rating=2]{ex:ls-a7-split-floor-run}
Run \texttt{python3 skills/harbor-results/scripts/a7\_experiment.py} and read off the
floor at $N{=}60,k{=}2,m{=}8$, the disjoint-readers floor at $2k{=}4$, and the
falsification count. Reproduce the two floors by hand from Thm.~\ref{thm:lowerbound}
first. Distinguish the joint-to-single ratio from the joint-to-two-single-floors
ratio, and state the opening budget behind each comparison.
\end{pdexercise}
\begin{pdsolution}{ex:ls-a7-split-floor-run}
The single-reader floor is $5.9822$ bits and the joint floor is $12.7662$ bits.
The script's reported $2.13\times$ compares the joint floor with only one
single-reader floor. Dividing instead by $2(5.9822)=11.9644$ gives $1.0670$.
The joint digest permits eight opens total; each separate reader may open
eight. Neither ratio measures an equal-total-budget implementation saving.
The recorded sweep reports \texttt{violations: 0/16 -> floor HOLDS} and
\texttt{ratio:\ \ 2.13x} [internal, \texttt{a7\_experiment.py}]. Its zero
violations concern the finite tested configurations and the script's
sampling tolerance, not every possible encoder. The counting argument
proves the lower bound; the sweep checks the tested implementations against it.
\end{pdsolution}
\begin{pdexercise}[kind=Check,rating=1]{ex:ls-paging-ratio-hand}
Compute the Landlord competitiveness ratio $k/(k-h+1)$ at $k=8,h=8$ and at
$k=8,h=1$. Which end of the range recovers the classical, unaugmented
Sleator--Tarjan bound?
\end{pdexercise}
\begin{pdsolution}{ex:ls-paging-ratio-hand}
At $h=k=8$: $k/(k-h+1)=8/1=8.0$, recovering the classical unaugmented result that a cache of size $k$ is $k$-competitive against an offline optimum given the very same cache size --- the hardest comparator, no resource augmentation. At $h=1$: $k/(k-h+1)=8/8=1.0$, the fully resource-augmented regime, where online's cache $k$ so outsizes the comparator's $h=1$ that the ratio collapses to $1$. The chapter's own $k=8,h=4$ instance, at $1.6$, sits at the midpoint, pricing a comparator with half of online's verified capacity.
\end{pdsolution}

\pdexercisesfor{\S\ref{sec:consent-canon}}{Reconciling the canon: roles, consent, and local knowledge}

\begin{pdexercise}[kind=Check,rating=1]{ex:ls-hobbes-mapping}
Map each of {agents, daemon, operator} onto Hobbes' {multitude, sovereign-actor,
author}. Why does the metaphor break if the operator is called the sovereign?
\end{pdexercise}
\begin{pdsolution}{ex:ls-hobbes-mapping}
Agents are the multitude bound by the common coordination protocol; the daemon is the sovereign-as-actor that applies the rules; the operator is the author/principal who originates and can withdraw authority. Calling the operator the sovereign collapses author and executing authority and makes the supposed covenant among agents incoherent --- there would be no distinct actor left to hold accountable to the author's terms.
\end{pdsolution}
\begin{pdexercise}[kind=Check,rating=1]{ex:ls-exit-makes-voice-honest}
Why does Hirschman's \emph{exit} (Invariant~\ref{prop:override}) make \emph{voice}
(escalation) honest?
\end{pdexercise}
\begin{pdsolution}{ex:ls-exit-makes-voice-honest}
Voice is honest only if refusal has consequence for the authority. The unconditional override supplies exit: the operator can revoke the grant when explanations or appeals fail, rather than being forced to keep negotiating inside a system they cannot stop.
\end{pdsolution}
\begin{pdexercise}[kind=Trace,rating=2]{ex:ls-express-consent-properties}
Show that the consent grant (Invariant~\ref{def:consent}) is \emph{express} in
Locke's sense by listing the four properties that distinguish it from ``the operator
just started using the tool.''
\end{pdexercise}
\begin{pdsolution}{ex:ls-express-consent-properties}
The grant is discrete, explicitly issued, scoped, time-bounded, and revocable and audited. ``Started using the tool'' has none of these properties: it is continuous rather than discrete, inferred rather than issued, unscoped, open-ended, and carries no explicit revocation act.
\end{pdsolution}
\begin{pdexercise}[kind=Open,rating=3]{ex:ls-ranker-political-organ}
\emph{The ranker as a political organ.} Specify the ``what did the
queue suppress, and why'' counter-surface (Table~\ref{tab:counter-surface})
so that ranking-as-governing is itself legible.
Does the counter-surface need its own counter-surface? Where does the regress
terminate, and on what principle?
\end{pdexercise}
\begin{pdsolution}{ex:ls-ranker-political-organ}
Same construction as Exercise~\ref{ex:ls-ranker-counter-surface}: one complete, append-only decision ledger recording every candidate, its ranking decision, reason codes, and a counterfactual, queryable directly and sampled independently of the production ranker rather than re-ranked by a second dashboard. The regress terminates at completeness plus direct sampling --- the base ledger already records every candidate and decision, and the counter-surface is a deterministic query over it, so no further ranker is needed to audit the ranker's omissions. The remaining trust assumption is the completeness of the event-capture boundary itself; that assumption should be tested and externally anchored, not hidden behind another summary.
\end{pdsolution}

% ═════════════════════════════════════════════════════════════════════════════
\section{Related work and prior art}\label{sec:related}
% ═════════════════════════════════════════════════════════════════════════════

This chapter sits at the confluence of five external literatures that rarely cite one
another, plus one internal body of work --- the accompanying formal papers --- whose
results several of this chapter's claims are popularizations of. We survey them in one
place; each was introduced in context above.

\paragraph{Legibility and the limits of central schemes.}
The organizing lens is \textbf{Scott}'s \emph{Seeing Like a State}~\pdcite{scott1998}:
states make populations governable by imposing legibility, and high-modernist
\emph{over}-legibility recurrently fails because it destroys the local, hard-to-codify
\emph{m\^etis} the simplified order depended on; \emph{Two Cheers for
Anarchism}~\pdcite{scott2012} is its constructive sequel. The same point arrives from
economics in \textbf{Hayek}'s \emph{The Use of Knowledge in Society}~\pdcite{hayek1945}:
the decisive knowledge is dispersed and not aggregable by a central planner.
\textbf{Rao}~\pdcite{rao2010} ported ``legibility'' into a standard lens on systems that
flatten reality to a single purpose. Our contribution is to make Scott's warning a
\emph{buildable} rule (digest-with-zoom, Invariant~\ref{def:zoom}) and to give it an
information-theoretic floor (Thm.~\ref{thm:lowerbound}).

\paragraph{Authority, consent, and the social contract.}
The authority argument is Hobbesian~\pdcite{hobbes1651,sep-hobbes}, amended by the
liberal consent tradition --- \textbf{Locke}~\pdcite{locke1689},
\textbf{Hume}~\pdcite{hume1748}, \textbf{Pateman}~\pdcite{pateman1979} --- and by
\textbf{Hirschman}'s exit/voice framework~\pdcite{hirschman1970}. Strikingly, the
state-of-nature-to-sovereign arc has been observed as an \emph{emergent} dynamic of
LLM populations (\textbf{Dai et al.}~\pdcite{dai2024}, building on \textbf{Park et
al.}~\pdcite{park2023}). We turn the metaphor into a mechanism: an express, scoped,
revocable consent grant (Invariant~\ref{def:consent}) with an inalienable override
(Invariant~\ref{prop:override}), and a legible-sovereign rule
(Invariant~\ref{prop:legible-sovereign}) that inverts Hobbes' opaque sovereign.

\paragraph{Human factors, automation, and signal detection.}
The operator model draws on the automation literature:
\textbf{Bainbridge}'s ironies of automation~\pdcite{bainbridge1983},
\textbf{Lee \& See} on trust calibration~\pdcite{leesee2004},
\textbf{Endsley} on situation awareness and the out-of-the-loop
problem~\pdcite{endsley1995,endsleykiris1995}, the vigilance-decrement results of
\textbf{Mackworth}~\pdcite{mackworth1948} and \textbf{Warm et al.}~\pdcite{warm2008},
\textbf{Wickens}' multiple-resource theory~\pdcite{wickens2002},
\textbf{Treisman \& Gelade}'s feature-integration theory~\pdcite{treisman1980}, and
\textbf{Miller}'s studies of absolute judgment, immediate memory, and
recoding~\pdcite{miller1956}.

\noindent\textbf{Alarm decisions.} The forced-zoom decision is
modeled with \textbf{Green \& Swets}' Signal Detection Theory~\pdcite{greenswets1966}
(Invariant~\ref{def:sdt}); the alarm discipline follows process-control
practice~\pdcite{isa182} and the clinical alarm-fatigue evidence~\pdcite{cvach2012}.

\noindent\textbf{Limits of the objective.}
\textbf{Campbell}~\pdcite{campbell1979} and \textbf{Strathern}~\pdcite{strathern1997}
(Goodhart's law) warn about incentives created by consequential metrics. This
chapter applies a classical SDT loss (Eq.~\ref{eq:sdt}) to inspection decisions
in agent supervision. The usefulness of that application requires calibration
and evaluation.

\paragraph{Information foraging and the cost of reading.}
Read-poverty (Hypothesis~\ref{def:readpoverty}) is the agent-coordination instance of
\textbf{Pirolli \& Card}'s \emph{information foraging theory}~\pdcite{pirolli1999} ---
information-seekers follow ``information scent'' and abandon a patch when its expected
yield falls below the cost of foraging --- and of the long-context degradation
results that make reading expensive even for machines (\textbf{Liu et
al.}~\pdcite{liu2024}; \textbf{Anthropic}~\pdcite{anthropic2025}). The remedy is the
classical database remedy --- an index --- recast as a directory over agents.

\paragraph{Agent discovery, naming, and trust.}
The directory substrate has deep prior art in \textbf{FIPA}'s Directory Facilitator
and Agent Management System~\pdcite{fipa00023}, re-invented for LLM agents as the Agent
Name Service (\textbf{Ren et al.}~\pdcite{ren2025}). For the trust layer that sits above
discovery, \textbf{EigenTrust}~\pdcite{kamvar2003} supplies a trust-propagation
model whose modified algorithm relies on pretrusted peers. \textbf{Douceur}'s
Sybil analysis~\pdcite{douceur2002} makes identity assumptions a separate issue;
it does not turn an aggregate trust score into a general Sybil defense.
\textbf{Ostrom}~\pdcite{ostrom1990} offers governance comparisons, not a theorem
selecting this architecture. The proposed trust layer still owes explicit
anchors, admission rules, and an attack model. Our specific result here is the \emph{split ranker}
(Thm.~\ref{thm:split-ranker}): collaborator fit and inspection priorities can
strictly conflict. The comonotonicity result states when a common ordering can
nevertheless serve both readers.

\paragraph{The formal siblings, and what they already prove.}
The nearest prior art for several of this chapter's results is not external at all: it
is the accompanying formal program, and pretending otherwise would be exactly the
kind of quiet re-derivation this chapter argues against. \emph{The Price of a
Summary}~\pdcite{paper1} proves the legibility floor of Thm.~\ref{thm:lowerbound} as its
Theorem~1 in the same $N,k,m$ notation, and puts it through a pre-registered
falsification sweep this chapter only sketches; its comonotone characterization explains the rank-order conflict in
Thm.~\ref{thm:split-ranker}. The extractive counterexample in
Thm.~\ref{thm:split-digest} additionally states a feasible summary class and
budget; it does not forbid shared storage. The split-floor comparison uses
different total opening budgets and must be read with that qualification; its regret-head section is the formal version of the
decision-theoretic derivation in \S\ref{sec:relevance}; and its zoom theorem
($2k\lceil\log_2(F/k)\rceil+4k$ adaptive opens) is the algorithm that makes
digest-with-zoom affordable at scale. \emph{What Needs an Authority}~\pdcite{paper6}
supplies Thm.~\ref{thm:specialization}'s Erlang-C boundary and, as
\S\ref{sec:sole-roles} records, the sweep that falsifies the naive response-time
threshold. The relationship in both cases is the same: this chapter popularizes, the
formal papers prove and test.

\paragraph{Where this chapter is positioned.}
Table~\ref{tab:related} places the contribution against the nearest neighbors in each
literature.

\begin{table}[H]
\centering
\small
\renewcommand{\arraystretch}{1.3}
\begin{tabular}{@{}p{0.24\textwidth}p{0.34\textwidth}p{0.34\textwidth}@{}}
\toprule
\textbf{Literature} & \textbf{Nearest prior art} & \textbf{This chapter's move} \\
\midrule
Legibility & Scott; Hayek; Rao & Make over-legibility a buildable failure mode; prove a legibility lower bound \\
Social contract & Hobbes; Locke; Hume; Hirschman & Express, scoped, revocable consent grant; inalienable override; legible sovereign \\
Automation \& trust & Bainbridge; Lee \& See; Endsley & Class-specific calibration where needed; SDT inspection loss; proposed vigilance probes \\
Information foraging & Pirolli \& Card; lost-in-the-middle & Workload-dependent review bottleneck; directory-as-index \\
Discovery \& trust & FIPA DF/AMS; ANS; EigenTrust & Split ranker: fit vs.\ regret-if-ignored are distinct objectives \\
Queueing \& staffing & Erlang-C; Kendall; the formal siblings~\pdcite{paper6} & Sole-responsibility roles priced at an exact boundary, not asserted as a principle \\
\bottomrule
\end{tabular}
\caption{Where this chapter sits relative to the nearest prior art in each of its
contributing literatures. The sixth row is the one whose nearest prior art is
internal: the companion formal papers, cited throughout rather than silently
re-derived.}
\label{tab:related}
\end{table}

% ═════════════════════════════════════════════════════════════════════════════
\section{The time axis and the bridge to the economy}\label{sec:bridge}
% ═════════════════════════════════════════════════════════════════════════════

\subsection{Legibility over time, not just the present}\label{sec:timeaxis}

It is tempting to treat legibility as a snapshot --- what is the swarm doing
\emph{now}. But most operator decisions are forensic: \emph{why} did this land? The
time axis is a distinct legibility mode --- a replay scrubber over history,
commit-anchored provenance on every annotation, and an append-only event log that
answers ``how did we get here?'' (Scott's cadastral map is static; a swarm needs a
legible \emph{history}.) Together with the projection mode
(\S\ref{sec:sa}), legibility spans all three tenses: the forensic past (replay), the
inspectable present (digest-with-zoom), and the projected future (freeze-probe SA).

\subsection{What this layer hands upward (the through-line)}\label{sec:handoff}

The single-operator product records continuity as well as current work. The
accompanying chapters use that evidence in the following proposed architecture:

\pdthesis{This design combines memory, checkpoints, and witnessed outcomes to
support continuity. Reputation uses attributable outcomes under stable identity.
A checkpoint that restores an exact running computation is needed only when the
workload requires that form of resumption; it is not a prerequisite for every
agent market. Today's notes and oracle-bound commitments provide narrower forms
of continuity than execution snapshots and neutral graded outcomes.}

Concretely, the legibility layer provides upward:
\begin{itemize}[leftmargin=1.2em]
\item \textbf{Legible continuity} --- resurrection-with-memory, episodic memory, and
the briefing together make recorded continuity inspectable. The intended route
is memory and task state supporting continued work, with attributable outcomes
supporting reputation. A score alone does not establish whose work was evaluated
or whether the evidence is valid. The required checkpoint fidelity depends on
the task; preserving a reviewable outcome history need not restore a running
process.
That non-forgeable identity is a \BuiltWeak{} single-operator primitive in the
local tool: daemon-minted \texttt{actor-souls} and a bounded newcomer pool ship, while universal
write-boundary enforcement does not. The
proposed cross-operator attestation is an additional mechanism whose trust
assumptions and implementation \pdchapref{he}{The Harbor Economy} must declare.
\item \textbf{The completion ledger} --- in the target design, every verified
completion (and every loud-failed one) becomes an outcome event. Today durable
commitments, oracle-bound closure, and overdue detection form a partial witness
substrate; neutral grading and reputation updates do not ship. This layer is the \emph{outcome witness}; the
reputation layer is the \emph{scorer}. The split-ranker substrate
(Thm.~\ref{thm:split-ranker}) and the costly-escalation signal (\S\ref{sec:gates})
are the candidate-generation and trust-signal scaffolding the reputation head consumes.
\item \textbf{The consent grant as a priceable object} --- a hosted-trust offering is
the transfer of a consent grant (Invariant~\ref{def:consent}) to a third party, which the
pricing model in \pdchapref{he}{The Harbor Economy} must specify.
\item \textbf{Operator-attention economics} --- the loss in
\S\ref{sec:operator} offers a model for pricing inspection and errors when
supervising a rented fleet. Estimating those costs and checking the inspection
assumptions are part of applying it.
\end{itemize}

The Leviathan that makes your swarm legible to you, without lying to you with a pretty
map, is the product. Everything above it --- federation, the market --- is a second
Leviathan made of the first.

\keyidea[Local oversight comes before markets]{A single-operator tool can help
Alice authorize work, inspect supporting evidence and record unfinished
obligations before any cross-operator market exists. Its contract is already
substantial: revocable authority, verifiable zoom targets, calibrated demands
on attention and honest completion records. The later chapters ask what
additional identity, evaluation and settlement rules are needed when another
operator must rely on those records.}

\subsection{The open problems, as a map}\label{sec:openproblems}

Table~\ref{tab:open} consolidates the starred exercises into the chapter's open-problem
surface, with the rung each couples to. They are genuine research problems, not
homework dressed up.

\begin{table}[H]
\centering
\small
\renewcommand{\arraystretch}{1.2}
\begin{tabular}{@{}>{\raggedright\arraybackslash}p{0.40\textwidth}>{\raggedright\arraybackslash}p{0.14\textwidth}>{\raggedright\arraybackslash}p{0.36\textwidth}@{}}
\toprule
\textbf{Open problem} & \textbf{Difficulty} & \textbf{Couples to} \\
\midrule
Forced-zoom rate $p(r,s,v)$ (Ex.~\ref{ex:ls-pareto-digest}) & $\star\star\star$ & reputation-to-risk calibration (\S\ref{sec:operator}) \\
Operator-attention vs.\ token-COGS Pareto frontier (Ex.~\ref{ex:ls-pareto-digest}) & $\star\star\star$ & \S\ref{sec:tokens} \\
Compaction-quality scalar as an attestation invariant (Ex.~\ref{ex:ls-replay-metric-soundness}) & $\star\star$ & honest attestation \\
Completionist verification of the unverifiable (Ex.~\ref{ex:ls-honest-fallback-ladder}) & $\star\star$ & Invariant~\ref{def:completionist} \\
Adversarial legibility (narration$\leftrightarrow$artifact) (Ex.~\ref{ex:ls-adversarial-legibility}) & $\star\star\star$ & sandbox + reputation \\
Abdication-resistance: does the tax preserve SA? (Ex.~\ref{ex:ls-freeze-probe-study}) & $\star\star$ & \S\ref{sec:abdication} \\
Calibrating discovery-ranker weights with no ground truth (Ex.~\ref{ex:ls-implicit-feedback-goodhart}) & $\star\star$ & implicit relevance feedback \\
Decay-vs-summary boundary (Ex.~\ref{ex:ls-decay-vs-summary}) & $\star$ & Invariant~\ref{prop:decay} \\
Legibility's information-theoretic lower bound (Ex.~\ref{ex:ls-adversarial-critical-set}) & $\star\star\star$ & Thm.~\ref{thm:lowerbound} \\
The ranker-as-political-organ regress (Ex.~\ref{ex:ls-ranker-political-organ}) & $\star\star$ & Invariant~\ref{prop:legible-sovereign} \\
The succession corner \emph{beyond the mean}: tail response, and consistency of
judgment, neither of which the succession price covers
(Ex.~\ref{ex:ls-succession-corner}) & $\star\star$ & \S\ref{sec:succession-corner} \\
\bottomrule
\end{tabular}
\caption{The chapter's open problems, consolidated from the starred exercises, with
difficulty and the mechanism each couples to. Several couple this layer forward to the
identity-and-reputation layer, which is the seam the accompanying chapters carry.}
\label{tab:open}
\end{table}

% ═════════════════════════════════════════════════════════════════════════════
\bibliographystyle{plain}
\pdprintsolutions
\begin{thebibliography}{99}

\bibitem{hobbes1651}
Thomas Hobbes.
\newblock \textit{Leviathan, or The Matter, Forme and Power of a Common-Wealth Ecclesiasticall and Civill}.
\newblock Andrew Crooke, London, 1651.

\bibitem{sep-hobbes}
Sharon A. Lloyd and Susanne Sreedhar.
\newblock \textit{Hobbes's Moral and Political Philosophy}.
\newblock Stanford Encyclopedia of Philosophy, revised September 2022.
\newblock \url{https://plato.stanford.edu/entries/hobbes-moral/}

\bibitem{scott1998}
James C.\ Scott.
\newblock \textit{Seeing Like a State: How Certain Schemes to Improve the Human Condition Have Failed}. Yale University Press, 1998.

\bibitem{scott2012}
James C.\ Scott.
\newblock \textit{Two Cheers for Anarchism}. Princeton University Press, 2012.

\bibitem{rao2010}
Venkatesh Rao.
\newblock A Big Little Idea Called Legibility. \textit{ribbonfarm}, 2010.
\newblock \url{https://www.ribbonfarm.com/2010/07/26/a-big-little-idea-called-legibility/}

\bibitem{dai2024}
Gordon Dai, Weijia Zhang, Jinhan Li, Siqi Yang, Chidera Onochie Ibe, Srihas Rao, Arthur Caetano, and Misha Sra.
\newblock Artificial Leviathan: Exploring Social Evolution of LLM Agents Through the Lens of Hobbesian Social Contract Theory.
\newblock arXiv:2406.14373, 2024. \url{https://arxiv.org/abs/2406.14373}

\bibitem{park2023}
Joon Sung Park, Joseph C. O'Brien, Carrie J. Cai, Meredith Ringel Morris, Percy Liang, and Michael S. Bernstein.
\newblock Generative Agents: Interactive Simulacra of Human Behavior.
\newblock In \textit{ACM Symposium on User Interface Software and Technology (UIST)}, 2023.

\bibitem{hume1748}
David Hume.
\newblock Of the Original Contract. 1748.

\bibitem{locke1689}
John Locke.
\newblock \textit{Second Treatise of Government}, \S\S95--122. 1689.

\bibitem{pateman1979}
Carole Pateman.
\newblock \textit{The Problem of Political Obligation: A Critique of Liberal Theory}. Wiley, 1979.

\bibitem{hirschman1970}
Albert O.\ Hirschman.
\newblock \textit{Exit, Voice, and Loyalty: Responses to Decline in Firms, Organizations, and States}. Harvard University Press, 1970.

\bibitem{hayek1945}
Friedrich A.\ Hayek.
\newblock The Use of Knowledge in Society. \textit{American Economic Review}, 35(4):519--530, 1945.

\bibitem{schneider2000}
Fred B.\ Schneider.
\newblock Enforceable Security Policies.
\newblock \textit{ACM Transactions on Information and System Security}, 3(1):30--50, 2000.

\bibitem{turpin2023}
Miles Turpin, Julian Michael, Ethan Perez, and Samuel R.\ Bowman.
\newblock Language Models Don't Always Say What They Think: Unfaithful Explanations in Chain-of-Thought Prompting.
\newblock In \textit{NeurIPS}, 2023. arXiv:2305.04388.

\bibitem{meyer1992}
Bertrand Meyer.
\newblock Applying ``Design by Contract''.
\newblock \textit{IEEE Computer}, 25(10):40--51, 1992.

\bibitem{greenswets1966}
David M.\ Green and John A.\ Swets.
\newblock \textit{Signal Detection Theory and Psychophysics}. Wiley, 1966.

\bibitem{miller1956}
George A.\ Miller.
\newblock The Magical Number Seven, Plus or Minus Two.
\newblock \textit{Psychological Review}, 63(2):81--97, 1956.

\bibitem{pirolli1999}
Peter Pirolli and Stuart Card.
\newblock Information Foraging.
\newblock \textit{Psychological Review}, 106(4):643--675, 1999.

\bibitem{wickens2002}
Christopher D.\ Wickens.
\newblock Multiple Resources and Performance Prediction.
\newblock \textit{Theoretical Issues in Ergonomics Science}, 3(2):159--177, 2002.

\bibitem{leesee2004}
John D.\ Lee and Katrina A.\ See.
\newblock Trust in Automation: Designing for Appropriate Reliance.
\newblock \textit{Human Factors}, 46(1):50--80, 2004.

\bibitem{bainbridge1983}
Lisanne Bainbridge.
\newblock Ironies of Automation.
\newblock \textit{Automatica}, 19(6):775--779, 1983.

\bibitem{endsley1995}
Mica R.\ Endsley.
\newblock Toward a Theory of Situation Awareness in Dynamic Systems.
\newblock \textit{Human Factors}, 37(1):32--64, 1995.

\bibitem{endsleykiris1995}
Mica R.\ Endsley and Esin O.\ Kiris.
\newblock The Out-of-the-Loop Performance Problem and Level of Control in Automation.
\newblock \textit{Human Factors}, 37(2):381--394, 1995.

\bibitem{warm2008}
Joel S.\ Warm, Raja Parasuraman, and Gerald Matthews.
\newblock Vigilance Requires Hard Mental Work and Is Stressful.
\newblock \textit{Human Factors}, 50(3):433--441, 2008.

\bibitem{mackworth1948}
Norman H.\ Mackworth.
\newblock The Breakdown of Vigilance During Prolonged Visual Search.
\newblock \textit{Quarterly Journal of Experimental Psychology}, 1(1):6--21, 1948.

\bibitem{campbell1979}
Donald T.\ Campbell.
\newblock Assessing the Impact of Planned Social Change.
\newblock \textit{Evaluation and Program Planning}, 2(1):67--90, 1979.

\bibitem{leroy2009}
Sophie Leroy.
\newblock Why Is It So Hard to Do My Work? The Challenge of Attention Residue When Switching Between Work Tasks.
\newblock \textit{Organizational Behavior and Human Decision Processes}, 109(2):168--181, 2009.

\bibitem{treisman1980}
Anne M.\ Treisman and Garry Gelade.
\newblock A Feature-Integration Theory of Attention.
\newblock \textit{Cognitive Psychology}, 12(1):97--136, 1980.

\bibitem{isa182}
ANSI/ISA.
\newblock \textit{ANSI/ISA-18.2: Management of Alarm Systems for the Process Industries}. 2016.

\bibitem{cvach2012}
Maria Cvach.
\newblock Monitor Alarm Fatigue: An Integrative Review.
\newblock \textit{Biomedical Instrumentation \& Technology}, 46(4):268--277, 2012.

\bibitem{sennrich2016}
Rico Sennrich, Barry Haddow, and Alexandra Birch.
\newblock Neural Machine Translation of Rare Words with Subword Units (BPE).
\newblock In \textit{ACL}, 2016.

\bibitem{liu2024}
Nelson F.\ Liu, Kevin Lin, John Hewitt, Ashwin Paranjape, Michele Bevilacqua, Fabio Petroni, and Percy Liang.
\newblock Lost in the Middle: How Language Models Use Long Contexts.
\newblock \textit{TACL}, 12:157--173, 2024. arXiv:2307.03172.

\bibitem{anthropic2025}
Anthropic.
\newblock Effective Context Engineering for AI Agents. 2025.
\newblock \url{https://www.anthropic.com/engineering/effective-context-engineering-for-ai-agents}

\bibitem{vaswani2017}
Ashish Vaswani, Noam Shazeer, Niki Parmar, Jakob Uszkoreit, Llion Jones, Aidan N.\ Gomez, \L{}ukasz Kaiser, and Illia Polosukhin.
\newblock Attention Is All You Need.
\newblock In \textit{NeurIPS}, 2017. arXiv:1706.03762.

\bibitem{grasse1959}
Pierre-Paul Grass\'e.
\newblock La reconstruction du nid et les coordinations interindividuelles\ldots: la th\'eorie de la stigmergie.
\newblock \textit{Insectes Sociaux}, 6(1):41--80, 1959.

\bibitem{fipa00023}
Foundation for Intelligent Physical Agents.
\newblock \textit{FIPA Agent Management Specification} (SC00023). 2002.
\newblock \url{http://www.fipa.org/specs/fipa00023/}

\bibitem{ren2025}
V.\ Ren et al.
\newblock Agent Name Service (ANS): A Universal Directory for Secure AI Agent Discovery and Interoperability.
\newblock arXiv:2505.10609, 2025.

\bibitem{kamvar2003}
Sepandar D. Kamvar, Mario T. Schlosser, and Hector Garcia-Molina.
\newblock The EigenTrust Algorithm for Reputation Management in P2P Networks.
\newblock In \textit{International World Wide Web Conference (WWW)}, 2003.

\bibitem{strathern1997}
Marilyn Strathern.
\newblock ``Improving Ratings'': Audit in the British University System.
\newblock \textit{European Review}, 5(3):305--321, 1997. (Goodhart's law.)

\bibitem{consul}
HashiCorp.
\newblock Consul Service Discovery and DNS Interface (health-checked, TTL-expiring registry).
\newblock \url{https://developer.hashicorp.com/consul/docs/discover}

\bibitem{nisan2007}
Noam Nisan, Tim Roughgarden, \'Eva Tardos, and Vijay V.\ Vazirani (eds.).
\newblock \textit{Algorithmic Game Theory}. Cambridge University Press, 2007.

\bibitem{douceur2002}
John R.\ Douceur.
\newblock The Sybil Attack.
\newblock In \textit{IPTPS}, 2002.

\bibitem{ostrom1990}
Elinor Ostrom.
\newblock \textit{Governing the Commons: The Evolution of Institutions for Collective Action}. Cambridge University Press, 1990.

\bibitem{sleatortarjan1985}
Daniel D.\ Sleator and Robert E.\ Tarjan.
\newblock Amortized efficiency of list update and paging rules.
\newblock \textit{Communications of the ACM}, 28(2):202--208, 1985.

\bibitem{young2002landlord}
Neal E.\ Young.
\newblock On-line file caching.
\newblock \emph{Algorithmica}, 33(3):371--383, 2002.

\bibitem{kendall1953}
David G.\ Kendall.
\newblock Stochastic processes occurring in the theory of queues and their analysis by the method of the imbedded Markov chain.
\newblock \textit{Annals of Mathematical Statistics}, 24(3):338--354, 1953. (The $A/S/c$ notation.)

\bibitem{erlang1917}
A.\ K.\ Erlang.
\newblock Solution of some problems in the theory of probabilities of significance in automatic telephone exchanges.
\newblock \textit{Elektrotkeknikeren}, 13, 1917. (The Erlang-C delay formula.)

\bibitem{atiasaligrama2012}
George K.\ Atia and Venkatesh Saligrama.
\newblock Boolean compressed sensing and noisy group testing.
\newblock \textit{IEEE Transactions on Information Theory}, 58(3):1880--1901, 2012.

\bibitem{cheraghchi2013}
Mahdi Cheraghchi.
\newblock Noise-resilient group testing: limitations and constructions.
\newblock \textit{Discrete Applied Mathematics}, 161(1--2):81--95, 2013.

\bibitem{kingman1961}
J.\ F.\ C.\ Kingman.
\newblock The single server queue in heavy traffic.
\newblock \textit{Mathematical Proceedings of the Cambridge Philosophical Society}, 57(4):902--904, 1961. (Waiting time grows with the sum of the squared coefficients of variation.)

\bibitem{allen1990}
Arnold O.\ Allen.
\newblock \textit{Probability, Statistics, and Queueing Theory with Computer Science Applications}.
\newblock Academic Press, 2nd edition, 1990. (The Allen--Cunneen approximation for $M/G/c$, and the Pollaczek--Khinchine formula.)

\bibitem{paper1}
Erich Owens.
\newblock \textit{The Price of a Summary: Information-Theoretic Limits of Agent Oversight}.
\newblock The Harbor Research Program, Paper~1, 2026. Companion formal paper to this
chapter; its Theorems~1--4 are the proved versions of several results stated here.

\bibitem{dorfman}
Robert Dorfman.
\newblock The Detection of Defective Members of Large Populations.
\newblock \textit{Annals of Mathematical Statistics}, 14(4):436--440, 1943.

\bibitem{hwang}
Frank K. Hwang.
\newblock A Method for Detecting All Defective Members in a Population by Group Testing.
\newblock \textit{Journal of the American Statistical Association}, 67(339):605--608, 1972.

\bibitem{duhwang}
Ding-Zhu Du and Frank K. Hwang.
\newblock \textit{Combinatorial Group Testing and Its Applications}, 2nd ed.
\newblock World Scientific, 2000.

\bibitem{expansion}
H. Si, O. O. Koyluoglu, and S. Vishwanath.
\newblock Lossy Compression of Exponential and Laplacian Sources Using Expansion Coding.
\newblock arXiv:1308.2338, 2013.

\bibitem{rdc}
Nam Nguyen, Thinh Nguyen, and Bella Bose.
\newblock
Rate-Distortion-Classification Representation Theory for Bernoulli Sources.
\newblock arXiv:2601.11919, 2026.
\url{https://arxiv.org/abs/2601.11919}

\bibitem{paper6}
Erich Owens.
\newblock \textit{What Needs an Authority: Mechanical Detection, Chartered Resolution, and the Exact Price of Sole Ownership}.
\newblock The Harbor Research Program, Paper~6, 2026. Its Theorem~2 is the Erlang-C
specialization boundary, together with the sweep comparing it against the naive
response-time threshold.

\bibitem{spence1973}
Michael Spence.
\newblock Job Market Signaling.
\newblock \textit{Quarterly Journal of Economics}, 87(3):355--374, 1973.

\bibitem{schmeidler1989}
David Schmeidler.
\newblock Subjective Probability and Expected Utility without Additivity.
\newblock \textit{Econometrica}, 57(3):571--587, 1989.

\bibitem{yaari1987}
Menahem E.\ Yaari.
\newblock The Dual Theory of Choice under Risk.
\newblock \textit{Econometrica}, 55(1):95--115, 1987.

\bibitem{lykouris2018}
Thodoris Lykouris and Sergei Vassilvitskii.
\newblock Competitive Caching with Machine Learned Advice.
\newblock In \textit{ICML}, 2018. arXiv:1802.05399; \textit{Journal of the ACM}, 2021.

\bibitem{sheridan1992}
Thomas B.\ Sheridan.
\newblock \textit{Telerobotics, Automation, and Human Supervisory Control}.
\newblock MIT Press, 1992.

\bibitem{parasuraman2000}
Raja Parasuraman, Thomas B.\ Sheridan, and Christopher D.\ Wickens.
\newblock A Model for Types and Levels of Human Interaction with Automation.
\newblock \textit{IEEE Transactions on Systems, Man, and Cybernetics -- Part A}, 30(3):286--297, 2000.

\bibitem{klein1998}
Gary Klein.
\newblock \textit{Sources of Power: How People Make Decisions}.
\newblock MIT Press, 1998.

\bibitem{polanyi1966}
Michael Polanyi.
\newblock \textit{The Tacit Dimension}.
\newblock Doubleday, 1966.

\bibitem{wald1947}
Abraham Wald.
\newblock \textit{Sequential Analysis}.
\newblock Wiley, 1947.

\bibitem{macaroons2014}
Arnar Birgisson, Joe Gibbs Politz, Úlfar Erlingsson, Ankur Taly, Michael Vrable, and Mark Lentczner.
\newblock Macaroons: Cookies with Contextual Caveats for Decentralized Authorization in the Cloud.
\newblock In \textit{Network and Distributed System Security Symposium (NDSS)}, 2014.
\newblock \url{https://doi.org/10.14722/ndss.2014.23212}

\end{thebibliography}

% ═════════════════════════════════════════════════════════════════════════════
\appendix
% ═════════════════════════════════════════════════════════════════════════════

\section{Implementation and status}\label{app:honest}

The body of this chapter argues mechanisms on their merits and does not depend on any of
them being built. This appendix --- and \emph{only} this appendix --- records the
concrete engineering status of each mechanism in the reference implementation, named
\emph{Port Daddy}. We use four honest labels, with one rule for the whole table:
confusing what is \emph{implemented} with what is merely \emph{specified} is itself
the failure this chapter argues against, so each row carries its evidence.

\smallskip
\noindent\Built\ = code exists in the reference implementation today; \BuiltWeak\ =
code exists but is partial, unwired, or honest about a known gap; \Designed\ = a
concrete, accepted design with no merged implementation; \Vision\ = a mechanism argued
in this chapter with no design document yet.

\smallskip
\noindent A second key, in the same grammar, labels \emph{claims} rather than code.
Every formal result in this chapter carries exactly one of these markers directly under
it; this chapter uses those markers as follows:
\Verified\ = independently re-derived and checked numerically, or subjected to a
pre-registered falsification sweep, outside this document; \Proved\ = argued here with
a proof or proof sketch and not independently checked; \Unproved\ = stated as a design
target whose proof is not carried in this chapter. The two keys are orthogonal on
purpose --- a \Verified\ result about a \Vision\ mechanism is a perfectly ordinary
state of affairs, and conflating them is the failure this appendix exists to prevent.

{\footnotesize\renewcommand{\arraystretch}{1.05}\setlength{\tabcolsep}{4pt}
% Status column widened 0.15 -> 0.18\textwidth to hold \Built ("implemented",
% 11 letters bold small caps), which is narrower than 0.15\textwidth even
% hyphenated. It was widened when the Book began printing that word here while
% the standalone paper still printed its own shorter "built"; both print
% "implemented" now, so the width is simply the width the label needs.
\begin{xltabular}{\textwidth}{@{}p{0.04\textwidth} >{\raggedright\arraybackslash}X p{0.18\textwidth} >{\raggedright\arraybackslash}p{0.30\textwidth}@{}}
\caption{Mechanism status in this chapter. L: legibility; A: authority; O: operator model. Read-surfaces largely exist; the unified digest, decision authority, and operator model remain substantially specified or proposed.}\label{tab:honest-state}\\
\toprule
& \textbf{Mechanism} & \textbf{Status} & \textbf{Evidence / location} \\
\midrule
\endfirsthead
\multicolumn{4}{@{}l}{\small\itshape (continued)}\\
\toprule
& \textbf{Mechanism} & \textbf{Status} & \textbf{Evidence / location} \\
\midrule
\endhead
\midrule
\endfoot
\bottomrule
\endlastfoot
L & Attention composer (agent-facing) & \Built & \texttt{lib/attention.ts} \\
L & Briefing projection (exemplary digest-with-zoom) & \Built & \texttt{lib/briefing.ts} \\
L & Resurrection handoff & \BuiltWeak & \texttt{lib/resurrection.ts}; notes, not checkpoints \\
L & Durable commitment / outcome substrate & \BuiltWeak & \path{lib/commitments.ts}, \path{lib/obligation-monitor.ts}; grading and reputation binding absent \\
L & Local actor identity root & \BuiltWeak & \path{lib/actor-souls.ts}, \path{lib/budget-guard.ts}; universal write gating absent \\
L & Honest attestation / legible sovereign & \Built & \texttt{lib/attest.ts}; reports pass, fail, skipped, and unknown checks explicitly \\
L & Episodic memory; pheromone decay; skill index & \Built & \path{lib/}\allowbreak\path{episodic-}\allowbreak\path{memory.ts}, \path{lib/}\allowbreak\path{pheromone.ts}, \path{lib/}\allowbreak\path{shipwright/}\allowbreak\path{skill-}\allowbreak\path{index.ts} \\
L & Anomaly monitor (detect-and-compensate, post-commit) & \BuiltWeak & \texttt{lib/arbiter.ts:328} is a log subscriber (Invariant~\ref{prop:regiment}) \\
L & Operator attention queue (regret-ranked lanes) & \Designed & Regret-based routing into separate attention lanes is not implemented \\
L & Discovery router (relevance ranking) & \Designed & Relevance-ranked routing is specified here; no router implementation \\
L & Operator console (native shell; Mission view and inspectors, \S\ref{sec:operator-instrument}) & \BuiltWeak & \path{core/pd-console}; the window builds on one platform behind a feature flag \\
L & Task-continuation head (successor-facing compaction) & \Designed & Thm.~\ref{thm:split-digest}; \texttt{lib/resurrection.ts} is the partial form \\
L & Oversight head / SLM sidecar (air-gapped digest authoring) & \Vision & Thm.~\ref{thm:split-digest} corollary; no concrete mechanism specified \\
A & Roadmap as versioned intent (authority accountable to it) & \BuiltWeak & rows/matrices exist; accountability \Designed \\
A & Adversarial / conflict-free review organ & \Designed & Independent review of consequential work; deployment not established \\
A & Sole-responsibility roles (assigned against the boundary) & \Vision & Thm.~\ref{thm:specialization} is \Verified; no role-assignment mechanism in code \\
A & Completionist completion (verifier-gated) & \Vision & Invariant~\ref{def:completionist}; no concrete mechanism specified \\
A & Escalation (typed prioritized interrupt) & \Designed & Reserved distress lane; deployed trigger predicate unspecified \\
A & Thoughtful landing / merge ordering & \BuiltWeak & \path{lib/}\allowbreak\path{merge-}\allowbreak\path{queue.ts} present, not wired \\
O & Operator-attention objective (SDT-spined) & \Vision & Invariant~\ref{def:sdt}; unmodeled in code \\
O & Trust calibration via canary + catch-rate & \Vision & no implementation \\
O & Consent grant + inalienable override & \Vision & Invariant~\ref{def:consent}, Invariant~\ref{prop:override} \\
O & Time-axis legibility (replay, provenance) & \Designed & Event log \Built; interactive replay and provenance navigation proposed \\
O & Ambient legibility (menu-bar indicator on PASS$\to$FAIL) & \BuiltWeak & indicator exists; watchdog wire \Designed \\
\end{xltabular}}

Three points about the reference implementation are important enough to restate:
\begin{enumerate}[leftmargin=1.5em]
\item \textbf{The anomaly monitor is detect-and-compensate, not prevention}
(Invariant~\ref{prop:regiment}). The monitor at \texttt{lib/arbiter.ts:328} is a
post-commit log subscriber; it cannot prevent a violation that has already
committed. A later safety gate would be a separate mechanism. ``Physically unreachable'' is reserved for constraint- or boot-gate-enforced
states only.
\item \textbf{The unified ranker is two rankers} (Thm.~\ref{thm:split-ranker}). The
discovery router maximizes \emph{fit}; the attention queue maximizes
\emph{regret-if-ignored}. They share candidate generation and decay, not an objective.
\item \textbf{Every read-surface decays} (Invariant~\ref{prop:decay}). Briefing,
resurrection, discovery, and attestation projections all carry recency-decay; the
dead-claim case (\S\ref{sec:staleness}) is the worked failure this closes.
\end{enumerate}

\section{Notation and the nomenclature key}\label{app:notation}

Two distinctions in this chapter are easy to collapse and essential to keep separate.

\begin{itemize}[leftmargin=1.4em]
\item \textbf{Capability vs.\ permission.} \emph{Capability} is \emph{ability} --- what
the execution sandbox physically makes possible. \emph{Permission} is \emph{normative
status} --- what the policy layer allows. They must never be collapsed: ``able but
forbidden'' (a dangerous escape-hatch flag that exists in the tooling but must not be
used) has to stay expressible, or the system cannot reason about a capability it
possesses but is not permitted to exercise.
\item \textbf{Two distinct delegation objects.} The \emph{authorization chain} is a
cryptographic, hop-bound, attenuation-monotonic delegation chain (a security object,
developed in \pdchapref{anchor}{The Anchor Protocol}). The \emph{coordination lineage} is the
task graph over which loop-detection runs (a liveness object, in this chapter's
protocol layer), and it is the same object \pdchapref{swk}{The Single-Writer Kernel}
names under that name --- one object, one name, which is the point of this entry.
They are different artifacts with different invariants; ``delegation
chain'' unqualified conflates a security claim with a coordination claim and must be
avoided.
\item \textbf{Physically unreachable} is reserved for constraint-enforced or
boot-gate-enforced states only (Invariant~\ref{prop:regiment}); a detect-and-compensate
monitor never earns the phrase.
\end{itemize}

\end{document}
